\documentclass[11pt]{article}
\usepackage[T1]{fontenc}
\usepackage[utf8]{inputenc}
\usepackage{lmodern}
\usepackage{amsmath,amssymb,amsthm,booktabs,longtable,array}
\usepackage{mathtools}
\usepackage{mathrsfs}
\usepackage[a4paper,margin=27mm]{geometry}
\usepackage[hidelinks]{hyperref}
\usepackage{url}
\setlength{\emergencystretch}{3em}
\newcommand{\GT}{\widehat{\mathrm{GT}}}
\newcommand{\GTz}{\widehat{\mathrm{GT}}_0}
\newcommand{\Z}{\mathbf Z}
\newcommand{\Q}{\mathbf Q}
\newcommand{\F}{\mathbf F}
\newcommand{\HH}{\mathcal H}
\DeclareMathOperator{\Gal}{Gal}
\DeclareMathOperator{\Aut}{Aut}
\DeclareMathOperator{\Out}{Out}
\newenvironment{statement}[3]{\paragraph{#1 (#2).}\textbf{#3}\par}{}
\title{DELTA23: two--three frame reciprocity,\quad
  order-five characters, and the remaining pentagon gate}
\author{Research run DELTA23\\GPT-6 Astra Pro, four cooperating workers}
\date{Working manuscript, 20--21 November 2026, Europe/Berlin}
\begin{document}
\maketitle
\begin{abstract}
The main question is whether the supplied continuous character
$\Delta:\GTz\to\F_2$ vanishes on full profinite $\GT$.
This working version does not assert a decision of that question.
It gives explicit Hilbert-symbol formulas and their limitations,
an actual integral order-five frame character, the full image of
the normalized adelic two-frame map, an exact remaining kernel
for relative quadratic characters, and a marked two-word
Atkin--Lehner test.  A new affine free-group evaluation detects
the cubic sign in a group of order $216$ and reduces the main
question to twelve joint finite source residues, six of mixed
parity, in a joint source of order $5184$. A checked six-dimensional
Krammer quotient retains this cubic detector. All its $216$
reduced pentagon assignments have common-word realizations, and
the harmonic $S_4$ phase remains independent at that word level.
We also give an explicit genus-zero quotient diagram for the
cubic detector and a general limitation of fixed-target
distribution tests. A canonical cubic descent identifies the
half-character $\Psi_3/2$ with a descended cubic frame sign.
Two exact orbital pentagons now identify the harmonic and cubic
residuals and their common periodic conjugacy; the missing parity
comparison is not inferred from these identities alone.
An actual full-GT Hurwitz quotient now gives twelve orthogonal
blocks. Its canonical cup form, simple factors and exact solvable
quotient lead to a complete finite list of normalized conjugators.
Uniqueness of the full first value and the literal raw pentagon
reduce that list from thirty-six conjugators to six, with three
mixed values each. Cubing and actual complex conjugation reduce
the existence question further to three explicit pairs. An odd product
with a complex conjugate reduces it to two. The same-associator
last-meridian test then leaves one existential normal form.
Every retained pair nevertheless passes the complete finite
common-word and mixed-cycle tests; none is asserted to lift to GT.
The corresponding raw product-fixed cover is also preserved.
Ongoing full-GT comparisons are separated
from these completed proofs. All dependence on supplied
lifting and arithmetic inputs is explicit.
\end{abstract}
\tableofcontents

\section{Conventions, status, and source inputs}
Write $G=\GTz$, $\mathcal H=\ker\lambda$, and
$\mathcal H_5=\GT\cap\mathcal H$.  In sections using $H$
for this kernel it means $\mathcal H$; the level-six section
explicitly defines its own subgroup $H\subset\widehat B$.
The modular group is $B=C_2\amalg C_3$, with generators
$s,r$ and $\tau=sr$.  We retain the positive pure frames
$\beta_h,A_h$ of LM7\_PUSH E.65.  With its signed correction
the two parities are $a_2=\epsilon_2(W)$ and
$a_3=\epsilon_3(T)$, and $\delta=a_2+a_3$ on $\mathcal H$.
The full relative extension is $\Delta$.

An equation in a finite matrix group is not a proof that a
candidate lifts to an actual gentle automorphism or full-GT
element.  Conversely, the E.58 and LI.8 elements are used as
actual elements supplied by the attached runs, subject to their
stated lifting input.  The earlier D.50 lifting proof is not
present in the attachment.  It has not been silently upgraded
to an independently proved result of this run.

The following are the current main states.  The machine-readable
questions and check events accompany this paper.
\begin{center}\small
\begin{tabular}{p{.12\linewidth}p{.78\linewidth}}
\toprule
Question&Current result or gate\\\midrule
Item 1&Explicit compressed Hilbert realization; global-product
failure refuted; specified matrix-local and Kubota interpretations
separated.\\
A.84&Open: D23.H2 normalizes any mixed actual state. The
coupled Hurwitz action, first-value uniqueness and raw pentagon
leave six conjugators and eighteen pairs. D23.C30--C31 and C33
reduce actual mixed existence to two pairs. D23.C34 excludes
the second and reduces existence to $(1,0)$. Its realization
or exclusion in full GT remains open.\\
A.24&Actual order-five target map and exact Galois scalar
comparison constructed. The further full-GT comparison is open.\\
A.86&Explicit marked involution and two necessary and sufficient
word equations. Every congruence test is trivial; the full
noncongruence comparison remains open.\\
B.40&All characters factoring through the full adelic frame pair
are $0,\Delta$. Remaining gate is the relative elementary-two
image of its kernel, D23.C8.\\\bottomrule
\end{tabular}\end{center}

\subsection{Exact supplied inputs}
The notations R2, R3, R4 in the attached DELTA\_ON\_GT run
refer respectively to LM7\_PUSH E.58, the finite covariance
calculation, and E.65.  The proofs below give their exact
equations whenever a sign or a normalization is used.
Other supplied inputs used here are DELTA\_ON\_GT A.12--A.14,
A.25, A.62, AD.1--AD.3, AF.4--AF.6, LI.8 (integrated A.94),
and FP.3, FP.6, FP.17--FP.21, FP.32--FP.39.  Their source
files and SHA-256 values are listed in the source manifest.
The ledger distinguishes use as background from a proof
actually reread or independently checked here.

\part{Findings with a completed whole-proof check}
\section{Local symbols and the product formula}
% DELTA23 root proof draft. All uses of E.58/E.65 are explicit supplied inputs.
\subsection{What a two--three Hilbert-symbol interpretation proves}

Write $\mathcal G=\widehat{GT}_0$, $\mathcal H=\ker\lambda$, and
$\mathcal H_5=\widehat{GT}\cap\mathcal H$. We retain the signed
frame coordinates of LM7\_PUSH E.65 and the CRT map of
DELTA\_ON\_GT A.6. Thus $a_2=\epsilon_2(W)$,
$a_3=\epsilon_3(T)$ and $\delta=a_2+a_3$ on $\mathcal H$.
The additive character $\Delta$ and the multiplicative Dirichlet
character $\chi_{12}$ have different codomains: the correctly typed
identity is $(-1)^{\delta(h)}=\chi_{12}(u(h))$.

\paragraph{D23.L1 (a well-defined Hilbert realization of the compressed frames).}
For $h\in\mathcal H$, choose any rational number $n$ which is a unit
at $2$ and $3$ and has residues
\[
n\equiv(-1)^{a_2(h)}\pmod4,\qquad
n\equiv(-1)^{a_3(h)}\pmod3.
\]
In particular one may choose the unique $n\in\{1,5,7,11\}$.
Define $L_2(h)=(n,3)_2$ and $L_3(h)=(n,3)_3$, using the quadratic
Hilbert symbols over $\mathbf Q_2$ and $\mathbf Q_3$. These are
independent of the indicated choice of $n$, are continuous characters,
and satisfy
\[
L_2=(-1)^{a_2},\qquad L_3=(-1)^{a_3},\qquad
L_2L_3=(-1)^\delta=\chi_{12}(u).
\tag{D23.L1.1}
\]
For the E.58 element $h_*$, in the supplied convention
$(a_3,a_2)=(1,0)$, one has $n=5$ and $(L_2,L_3)=(1,-1)$.
\begin{proof}
Use the standard local Hilbert formulas, with the precise locators
Voight, \emph{Quaternion algebras}, (12.4.10), 12.4.12(b), and
(12.4.14). For a $2$-adic unit $n$ they give
$(n,3)_2=(-1)^{(n-1)/2}$, where the exponent is read modulo two.
For a $3$-adic unit $n$ they give $(n,3)_3=(n/3)$.
These depend only on $n\bmod4$ and $n\bmod3$, respectively, and
give the first two identities. Their product gives the third by
the supplied definition of $\delta$. The four residues give
$\chi_{12}$ with values $1,-1,-1,1$, in the order $1,5,7,11$.
Continuity and multiplicativity follow from the already proved
frame characters, not from a purported multiplicative choice of
rational representatives. Finally E.58.23 and E.65.2--4 give
$T=-I$ at three and $W=I$, so the last evaluation follows.
\end{proof}

This is a positive answer for an explicitly defined pair of symbols
of the \emph{compressed} frame data. The definition has used the
two frame parities. It is not an independent arithmetic comparison
and does not prove that the two symbols have equal values.

\paragraph{D23.L2 (the proposed failure of global reciprocity is false).}
The rational pair $(n,3)$ in D23.L1 obeys Hilbert reciprocity for
every frame value. For the four positive representatives its complete
ramification table is
\[
\begin{array}{c|cc|c|c}
n&(n,3)_2&(n,3)_3&\text{other ramified finite primes}&(n,3)_\infty\\\hline
1&1&1&\varnothing&1\\
5&1&-1&5&1\\
7&-1&1&7&1\\
11&-1&-1&\varnothing&1
\end{array}
\tag{D23.L2.1}
\]
In particular E.58 does not violate a global product formula: its
extra factor is $(5,3)_5=-1$.
\begin{proof}
At an odd prime not dividing $3n$, both arguments are units and
the Hilbert symbol is one by Voight 12.4.12(a). At a prime dividing
one of $5,7,11$ the symbol is $(3/p)$. The square residues give
$(3/5)=(3/7)=-1$ and $(3/11)=1$ (since $5^2=3\pmod{11}$).
The two local values are D23.L1, and the real symbol is one
because $3>0$. These facts give the entire table. Each row has
an even number of minus signs, in agreement with Hilbert
reciprocity, Voight Proposition 14.2.1 (or Corollary 14.6.2).
\end{proof}

\paragraph{D23.L3 (the missing support hypothesis is exactly the old parity gate).}
For two bits $(b_2,b_3)$, the following are equivalent:
\begin{enumerate}
\item $b_2=b_3$;
\item there is $a\in\mathbf Q^\times$, a unit at $2$ and $3$, with
$(a,3)_2=(-1)^{b_2}$ and $(a,3)_3=(-1)^{b_3}$, and with
$(a,3)_v=1$ at every other place;
\item there is a rational quaternion algebra whose local invariants
at $2,3$ are $b_2/2,b_3/2$ and which is split at every other place.
\end{enumerate}
For the two equal-bit cases the algebras can be chosen as
$M_2(\mathbf Q)$ and the discriminant-six algebra $(2,3)$.
Thus the existence of such a two-place globalization for every
$h\in\mathcal H_5$ is \emph{equivalent} to A.84; it is not an
already proved arithmetic reason for A.84.
\begin{proof}
The implication from the third condition to the first is Hilbert
reciprocity (or the even-ramification part of Voight Main Theorem
14.6.1). The second condition implies the third by taking $(a,3)$.
For equal bits choose $a=1$ or $a=-1$. The local formulas give,
for $a=-1$, minus signs exactly at $2$ and $3$, and the real symbol
is again positive. This proves the first implication to the second.
Finally $(-1,3)$ is isomorphic to $(2,3)$. For an explicit change
of generators in $(-1,3)$, write $i^2=-1$, $j^2=3$, $ij=-ji$.
Then $i'=(1+j)i$ has square
$(1+j)(1-j)i^2=(1-3)(-1)=2$ and anticommutes with $j$.
Since $1+j$ is invertible, $i',j$ generate the original algebra.
This proves the asserted identification without any appeal merely
to coincident discriminants. Applying the equivalence to
$(b_2,b_3)=(a_2(h),a_3(h))$ gives the final assertion.
\end{proof}

\paragraph{D23.L4 (relative pro-three frames cannot recover the cubic bit).}
Let $U=F'F^3$ and let $\Gamma_3=B/O^3(U)$ be the specified
relative maximal pro-three quotient used in E.58. There is no
function of the induced automorphism of $\Gamma_3$ together with
the two retained pure frame values there which equals $a_3$ on
all of $\mathcal H$. In particular the cubic bit cannot be a
normalized local symbol depending only on that data.
\begin{proof}
The identity element has identity action and identity frames, with
$a_3=0$. The actual element supplied by E.58 has identity action
on every relative local factor. By E.58.20 its two frame values
in the factor at three are both one, while E.58.23 gives $a_3=1$.
The input data for the proposed function are equal and the
required outputs differ. This is a contradiction.
\end{proof}

The quotient $\Gamma_3$ in D23.L4 must not be confused with the
matrix target $\operatorname{SL}_2(\mathbf Z_3)$. The latter has
non-three-primary finite subgroups; reduction modulo three already
retains a quaternion group. D23.L4 refutes a precise pro-three
interpretation, not every possible arithmetic local comparison.

\paragraph{D23.L5 (full multipliers and the cyclotomic distinction).}
For $g\in\mathcal G$ choose $\sigma\in G_{\mathbf Q}$ with
$\chi(\sigma)=\lambda(g)$ and put $h=g\sigma^{-1}$.
The product $L_2(h)L_3(h)=(-1)^{\Delta(g)}$ is independent of this
choice, while the individual symbols need not be. Moreover neither
of the two individual bits is the restriction of a Dirichlet
character of the multiplier.
\begin{proof}
Choice independence of the product is E.65.5--6. The source
DELTA\_ON\_GT A.25 constructs $\sigma_0\in G_{\mathbf Q^{\rm ab}}$
which sends $r_0=(3/4)^{1/4}$ to $-r_0$, with
$a_2(\sigma_0)=a_3(\sigma_0)=1$ and $\lambda(\sigma_0)=1$.
For $g=1$, taking $\sigma=1$ gives the two symbols $(1,1)$,
whereas taking $\sigma=\sigma_0$ gives $(-1,-1)$. This proves
the individual ambiguity. A character of the multiplier is
trivial on every multiplier-one element, while each of these
frame bits takes value one at $\sigma_0$. This proves the last
assertion. The cited example concerns a nonabelian quartic-root
extension over $\mathbf Q$, quadratic only after adjoining
$\mathbf Q^{\rm ab}$; it is not the cyclotomic quadratic field
$\mathbf Q(\sqrt3)$.
\end{proof}

\subsection{A precise Kubota test which does not detect the discrepancy}

\paragraph{D23.L6 (the integral splitting-normalized test is trivial).}
Use the actual signed adelic frames $\mathcal T_h,\mathcal W_h$
of DELTA\_ON\_GT AD.2.  For a prime $p$, and a matrix
$g=\left(\begin{smallmatrix}a&b\\c&d\end{smallmatrix}\right)
\in\operatorname{SL}_2(\mathbf Q_p)$, put
\[
 x(g)=\begin{cases}c&c\ne0,\\d&c=0,\end{cases}\qquad
 a_p(g,h)=(x(g),x(h))_p
             (-x(g)^{-1}x(h),x(gh))_p.
\]
This is Kubota's factor set (Kubota 1967, equation (2), printed
p.~114).  Fix the splitting normalization
\[
 s_p(g)=\begin{cases}
 (c,d)_p&cd\ne0\text{ and }v_p(c)\text{ is odd},\\
 1&\text{otherwise},
 \end{cases}
 \qquad
 \kappa_p(g,h)=a_p(g,h)s_p(g)s_p(h)s_p(gh)^{-1}.
\tag{D23.L6.1}
\]
Then the concrete functions of the corrected frame pair
\[
 K_p(h)=\kappa_p(\mathcal T_{h,p},\mathcal W_{h,p}),
 \qquad p=2,3,
\]
are both identically $1$ on the entire gentle multiplier-one
group.  In particular $K_2K_3$ is not $(-1)^\delta$: the
stipulated actual E.58 element has $\delta=1$.

\begin{proof}
The standard local splitting theorem says that (D23.L6.1) is
one on $\Gamma_1(4,\mathbf Z_2)$ at two and on
$\operatorname{SL}_2(\mathbf Z_p)$ at every odd prime.
A primary source giving precisely this normalization is Hida,
\emph{Adjoint L-value as a period integral and the mass formula
of Siegel--Shimura}, version 20 September 2023, Section 1.3,
printed p.~7, the displayed definition of $s_p$ and paragraph
immediately before (1.5); it cites Gelbart, \emph{Weil's
Representation and the Spectrum of the Metaplectic Group},
Proposition 2.8.  We use this published splitting statement,
not a computation of a few residue classes.

AD.2 gives $\mathcal T_{h,2}\equiv\mathcal W_{h,2}\equiv I$
modulo four, so both matrices lie in the indicated group at
two.  At three both are integral matrices of determinant one.
The two conclusions follow directly from the respective
splittings.  E.58.23 and E.65 give $\delta(h_*)=1$, providing
the asserted failure of this proposed comparison.
\end{proof}

This refutes the specified test, not every conceivable local
construction.  A factor set changes by a coboundary when its
section is changed: replacing a section by $b(g)$ times that
section multiplies its factor set by
$b(g)b(h)b(gh)^{-1}$.  Thus an unqualified phrase ``the Kubota
symbol of two frames'' does not specify a function until the
matrix lifts and the section are fixed.  The explicit
normalization above fixes both and has the stated outcome.

\subsection{The two proposed quaternion algebras are different}
\paragraph{D23.L7.}
The rational quaternion algebras $(-1,-3)$ and $(2,3)$ are
not isomorphic.  Their ramification sets are, respectively,
$\{3,\infty\}$ and $\{2,3\}$.  Thus the former is definite
and has finite discriminant three; the latter is indefinite
and has discriminant six.  An identification between them cannot
serve as the missing two--three frame comparison.
\begin{proof}
For $(-1,-3)$ the odd-prime Hilbert formula is one away from
three, and at three it is the quadratic residue symbol
$(-1/3)=-1$.  At two the two odd units have residues $7,5$
modulo eight; the exponent in Voight (12.4.14) is
$(({-1}-1)/2)(({-3}-1)/2)=2$, which is even.
The real symbol is negative because both arguments are negative.
This gives its entire ramification set.  The ramification set
of $(2,3)$ was calculated in D23.L3.  A rational algebra
isomorphism would induce an isomorphism over $\mathbf Q_2$,
where the first algebra is split and the second is division,
which is impossible.
\end{proof}
This statement does not dispute a supersingular endomorphism
interpretation of $(-1,-3)$.  It identifies the extra comparison
which would be needed before that interpretation could bear on
the discriminant-six proposal.

\subsection{Direct local formulas for the actual matrix frames}
\paragraph{D23.L8.}
Write the three-components of the two signed corrected adelic
frames of $h\in\mathcal H$ as
\[
 \mathcal T_{h,3}=aI+bR_0,\quad a^2-ab+b^2=1,
 \qquad
 \mathcal W_{h,3}=cI+dS,\quad c^2+d^2=1.
\]
The two quantities $a+b$ and $c^2-d^2$ belong to
$\mathbf Z_3^\times$, and one has the direct matrix formulas
\[
 (a+b,3)_3=(-1)^{a_3(h)},\qquad
 (c^2-d^2,3)_3=(-1)^{a_2(h)}.
\tag{D23.L8.1}
\]
Their product is $(-1)^{\delta(h)}$.  These formulas use
actual local matrix coefficients, without first naming the bits,
but both are formulas at three.  They do not turn the harmonic
frame bit into a function of the two-adic frames; D23.C7 rules
out that latter factorization.
\begin{proof}
Modulo three the first norm equation is $(a+b)^2=1$.
Also $R_0$ acts as one on its one-dimensional fixed space,
so $a+b\bmod3$ is exactly the signed fixed-line eigenvalue
defining $a_3$.  For the second norm equation, the only pairs
$(c,d)\in\mathbb F_3^2$ with $c^2+d^2=1$ are
$(1,0),(-1,0),(0,1),(0,-1)$.  They represent respectively
$I,-I,S,-S$.  The first two have even $S$-exponent and the
last two have odd $S$-exponent.  The value $c^2-d^2$ is
respectively $1,1,-1,-1$.  This proves both the unit assertions
and the equality of the residue signs with the indicated frame
signs.  For a three-adic unit $u$, the odd-prime Hilbert formula
gives $(u,3)_3=(u/3)$, which in $\mathbb F_3^\times$ is
precisely this sign.  The product assertion follows.
For E.58, the two matrices are $-I,I$ and the values are
therefore $-1,1$, as in D23.L1.
\end{proof}


\section{Integral order-five construction}
% DELTA23: full proof draft, root, 2026-11-21 CET.
\subsection{Integral covariance of the five-strand symplectic representation}

\paragraph{D23.T1 (integral specialization of the rigidity method).}
Let $J$ be the alternating $4\times4$ integral matrix with
$J_{i,i+1}=1$, $J_{i+1,i}=-1$ and all other entries zero.
Put $t_v=I+vv^{\mathsf t}J$ and $t_i=t_{e_i}$.
The assignment $q(\sigma_i)=t_i$ defines a continuous homomorphism
$q:\widehat B_5\to\operatorname{Sp}(J,\widehat{\mathbf Z})$.
For every multiplier-one element $h\in\widehat{GT}$ there is a unique
$V_h\in\operatorname{Sp}(J,\widehat{\mathbf Z})$ such that
\[
V_he_1=e_1,\qquad q\Phi_h=\operatorname{Ad}(V_h)q.
\]
The map $h\mapsto V_h$ is a continuous homomorphism.
This statement concerns the genuine five-strand braid action; it does
not assert that this representation descends to the five-punctured
sphere group.

\begin{proof}
The determinant of $J$ is one. The transvections preserve $J$.
For $c=v^{\mathsf t}Jw$, direct multiplication gives
\[
t_vt_w-t_wt_v=c(vw^{\mathsf t}+wv^{\mathsf t})J,
\qquad
t_vt_wt_v-t_wt_vt_w=(1-c^2)(vv^{\mathsf t}-ww^{\mathsf t})J.
\tag{D23.T1.1}
\]
The entries of $J$ therefore give the braid and distant commutation
relations for the $t_i$. The discrete homomorphism extends to the
profinite completion, since reduction modulo each integer has finite
target.

Fix a prime $\ell$. The standard GT action has
$\Phi_h(\sigma_1)=\sigma_1$ and sends every other $\sigma_i$ to a
specified conjugate of itself. Hence
$q_\ell\Phi_h(\sigma_i)=t_{v_i}$ with $v_i$ primitive over
$\mathbf Z_\ell$ and $v_1=e_1$. The $v_i$ can be chosen continuously
using the specified conjugators in the GT braid formula.

We may first work over the field $\mathbf Q_\ell$. If $t_v,t_w$
commute, the first identity in (D23.T1.1) implies $c=0$.
Indeed, if $c\ne0$, it would give $vw^{\mathsf t}+wv^{\mathsf t}=0$.
Choose a linear functional $L$ with $L(v)=1$ and multiply this matrix
on the left by $L$; then $w=-L(w)v$, which implies $c=0$, a
contradiction. For a pair satisfying the braid relation, the second
identity shows that either $c^2=1$ or $t_v=t_w$.

An equality $t_{v_i}=t_{v_{i+1}}$ is impossible. A neighboring
generator $t_{v_{i+2}}$, when present, commutes with $t_{v_i}$ and
braids with $t_{v_{i+1}}$. Two commuting invertible elements which
braid are equal: $a^2b=b^2a$ implies $a=b$. Equality thus propagates
to the right, and the identical argument propagates it to the left.
All four generators would have equal image. But $\Phi_h$ is an
automorphism of $\widehat B_5$, so $q_\ell\Phi_h$ has exactly the
same image as $q_\ell$, and the latter is nonabelian, as witnessed by
$t_1,t_2$. This is a contradiction.

Consequently $c_i=v_i^{\mathsf t}Jv_{i+1}$ equals $1$ or $-1$,
and the pairings of distant vectors are zero. Define
$\epsilon_1=1$ and $\epsilon_{i+1}=\epsilon_i c_i$ and replace
$v_i$ by $w_i=\epsilon_i v_i$. The transvections do not change and
$w_i^{\mathsf t}Jw_{i+1}=1$. Thus the column matrix
$V_\ell=(w_1,w_2,w_3,w_4)$ satisfies
$V_\ell^{\mathsf t}JV_\ell=J$. Taking determinants shows
$\det(V_\ell)^2=1$, so $V_\ell$ is invertible over
$\mathbf Z_\ell$. It fixes $e_1$ and conjugates every $t_i$ to
$t_{v_i}$.

For uniqueness, the simultaneous centralizer of the $t_i$ in
$\operatorname{GL}_4(\mathbf Q_\ell)$ is the scalar group. To see
this, commuting with $t_i-I$ preserves its image $\mathbf Q_\ell e_i$,
so a centralizing matrix is diagonal. Its commutation with the
nonzero $(i,i+1)$ entries of $t_i-I$ equates all adjacent diagonal
entries. A scalar symplectic matrix is $\pm I$, and fixing $e_1$
chooses $I$.

The product $V_h=(V_\ell)_\ell$ lies in the indicated adelic
symplectic group. The displayed formulas using the original
continuous $v_i$ and their pairings prove its continuity, prime by
prime, including at $2$ and $5$. Both $V_{hk}$ and $V_hV_k$ fix
$e_1$ and intertwine $q\Phi_h\Phi_k$ with $q$. Uniqueness gives
$V_{hk}=V_hV_k$.
\end{proof}

\paragraph{D23.T2 (the exact integral order-five centralizer).}
Put $C=t_1t_2t_3t_4$. Then
\[
C=\begin{pmatrix}0&0&0&1\\-1&0&0&1\\0&-1&0&1\\0&0&-1&1\end{pmatrix},
\quad C^5=-I,\quad \det(XI-C)=X^4-X^3+X^2-X+1.
\]
For every commutative ring $R$, the endomorphism centralizer of
$C$ in $M_4(R)$ is $R[C]\simeq R[\zeta_5]$, with
$\zeta_5$ represented by $-C$. Its symplectic centralizer is
\[
\operatorname{Cent}_{\operatorname{Sp}(J,R)}(C)
=\{z\in R[\zeta_5]:z\overline z=1\},
\qquad \overline{\zeta_5}=\zeta_5^{-1}.
\]
\begin{proof}
Multiplication gives the matrix and polynomial identities. The four
vectors $e_1,Ce_1,C^2e_1,C^3e_1$ are respectively
$e_1,-e_2,e_3,-e_4$, so they form a basis over every $R$.
Any endomorphism commuting with $C$ is determined by its value on
$e_1$; writing that value in this basis shows that the endomorphism
is a unique polynomial of degree at most three in $C$. This proves
the asserted algebra centralizer without any restriction on the
residue characteristic.
For the symplectic adjoint $X^*=J^{-1}X^{\mathsf t}J$ one has
$C^*=C^{-1}$, since $C$ is symplectic. On the commutative algebra
$R[C]$ the adjoint is therefore exactly the displayed involution.
The symplectic condition is $z^*z=1$. It already implies that $z$
is a unit, and is precisely the norm-one equation.
\end{proof}

% DELTA23, agent /root/five_level_six.  New algebraic results.
% Model: GPT-6 Astra Pro. Effort inherited, not exposed in the runtime.
% Run start: 2026-11-20 22:59 UTC.  This fragment is not a final certification.

\subsection*{D23.F1. A rigid order-five frame with one distinguished pole}
Let $L_6=(\widehat\Gamma_{0,[6]})_6$ be the subgroup fixing the
sixth label, and use the standard identification
\[
 L_6=\widehat B_5/\overline{\langle z_5\rangle},\qquad
 z_5=(\sigma_1\sigma_2\sigma_3\sigma_4)^5.
\]
Put $a=\sigma_1\sigma_2\sigma_3\sigma_4$ and write $\bar a$
for its image in $L_6$.  Let $\mathcal H_5$ denote the
multiplier-one subgroup of full profinite $\widehat{GT}$.
There is a unique continuous map
\[
 A_5:\mathcal H_5\longrightarrow\widehat{PB}_5
\]
such that
\[
 \Phi_h(a)=A_5(h)aA_5(h)^{-1},\qquad
 \ell_{12}(A_5(h))=0,
 \tag{D23.F1.1}
\]
where $\ell_{12}$ is the integral linking coordinate for the first
two strands, normalized by $\ell_{12}(z_5)=1$.
Its composition rule is
\[
 A_5(hk)=\Phi_h(A_5(k))A_5(h).
 \tag{D23.F1.2}
\]
The image of $A_5(h)$ in $P_5$ after forgetting the distinguished
sixth point is the unique order-five frame of
DELTA\_ON\_GT, FP.6.

\begin{proof}
First $C_{P_6}(\bar a)=1$.  Forget the sixth point.  The image of a
pure element centralizing $\bar a$ lies in $C_{P_5}(a_5)=1$,
by DELTA\_ON\_GT FP.3.  Hence it lies in the free kernel
$\widehat F_4$ of the pure Birman sequence.  Geometrically the
remaining punctures are $\mu_5$, the sixth point is $\infty$, and
the action is rotation $z\mapsto\zeta_5 z$.  The quotient orbifold
of $\mathbb P^1\setminus\mu_5$ has one puncture and two cone
points of order five.  Thus
\[
 \widehat F_4\rtimes\langle\bar a\rangle
       \simeq C_5\amalg C_5,
\]
with the section at $\infty$ a factor.  The factor-centralizer
theorem (Herfort--Ribes 1985, Theorem B', Section 4, printed
p.~159, also used in FP.3) makes its centralizer meet
$\widehat F_4$ trivially.  This proves the assertion.  The
profinite Birman exactness can be justified without assuming
the five-point base group is free. After fixing three labels
at $0,1,\infty$, the discrete forgetful map is
$\operatorname{Conf}_3(\mathbb C\setminus\{0,1\})\to
\operatorname{Conf}_2(\mathbb C\setminus\{0,1\})$.
It has a continuous section choosing its extra real point
farther from the origin than the absolute values of all the
existing points plus two. Its fundamental-group extension
therefore splits and has finitely generated free kernel.
Cofinal characteristic finite-index subgroups of that kernel
are invariant under the base action, and the induced action
on each finite quotient has finite image. Taking a finite-index
subgroup of the base that kills that image proves that the
topology induced on the kernel is its full profinite topology.
The split finite quotients also separate the base coordinate,
so completion gives the asserted exact sequence. The same
finite-extension argument applies to the subgroup over the
order-five rotation.

Collas 2012, Theorem B, applies to the order-five element $\bar a$
in the full six-point group.  For $h\in\mathcal H_5$ it gives
$\Phi_h(\bar a)=g\bar a g^{-1}$.  The standard GT action is the
identity on the permutation quotient.  The permutation of $g$
therefore centralizes a five-cycle with one fixed label.  It
fixes that unique sixth label and is a power of the five-cycle.
Multiplying $g$ on the right by an appropriate power of $\bar a$
gives a pure conjugator $\alpha_5(h)\in P_6$.
The preceding centralizer calculation proves uniqueness.
The continuous injective orbit map from the compact group $P_6$
to the Hausdorff group $G_6$ is a homeomorphism onto its image;
therefore this frame is continuous.

The central extension of pure groups splits by $\ell_{12}$.
For clarity, on ordered configurations one has the algebraic
isomorphism
\[
 \operatorname{Conf}_5(\mathbb A^1)
 \simeq\mathbb A^1\times\mathbb G_m\times\mathcal M_{0,6},
\quad
 (z_i)\longmapsto
 \left(z_1,z_2-z_1,
       \frac{z_3-z_1}{z_2-z_1},\ldots,
       \frac{z_5-z_1}{z_2-z_1}\right).
\]
Here the last factor has the first two points normalized to $0,1$
and the distinguished sixth point to $\infty$.
The $\mathbb G_m$ loop is the full twist, and its winding number
is $\ell_{12}$.  Hence
\[
 \widehat{PB}_5\simeq
 \overline{\langle z_5\rangle}\times\ker\ell_{12},
 \qquad \ker\ell_{12}\xrightarrow{\sim}P_6.
\]
Lift $\alpha_5(h)$ through this last isomorphism; call the lift
$A_5(h)$.  The frame equation in $L_6$ lifts to
\[
 \Phi_h(a)=A_5(h)aA_5(h)^{-1}z_5^e
\]
for some $e\in\widehat{\mathbb Z}$.  The exponent-sum map on
$\widehat B_5$ sends $a$ to four and $z_5$ to twenty, and GT
acts on this map by its multiplier, which is one.  Thus $20e=0$.
Multiplication by twenty on $\widehat{\mathbb Z}$ is injective,
so $e=0$.  This proves (D23.F1.1).

At multiplier one, $\Phi_h$ is the identity on the abelianization
of the pure braid group.  Indeed, its standard generators are
conjugates of half-twist squares; their images are conjugates
of the same squares, with the same strand-pair permutation.
Their linking coordinates are consequently unchanged.
It follows that $\Phi_h(A_5(k))A_5(h)$ has linking coordinate
zero and is a conjugator for $\Phi_{hk}(a)$.  Uniqueness gives
(D23.F1.2).  Finally the forgetful map takes its frame equation
to the five-point frame equation, and pure uniqueness there
identifies the two frames.
\end{proof}

\subsection*{D23.F2. The actual order-five norm-one frame character}
Use the symplectic braid representation and normalized covariance
$q\Phi_h=\operatorname{Ad}(V_h)q$ of D23.T1--D23.T2.
Set $C=q(a)$ and
\[
 \mathcal C_5(h)=q(A_5(h))^{-1}V_h.
 \tag{D23.F2.1}
\]
Then $\mathcal C_5$ is a continuous homomorphism on
$\mathcal H_5$, taking values in
\[
 \operatorname{Cent}_{\operatorname{Sp}_4(\widehat{\mathbb Z})}(C)
  \simeq
 \{u\in(\widehat{\mathbb Z}[\zeta_5])^\times:
                    u\overline u=1\}.
 \tag{D23.F2.2}
\]
The torus is the order-five target of DELTA\_ON\_GT FP.19;
here it receives an actual frame character.  The source is the
six-point stabilizer, so this does not assert that the
symplectic representation descends to the five-point sphere
group refuted in FP.21.

\begin{proof}
Apply $q$ to (D23.F1.1) and use covariance:
$V_h C V_h^{-1}=q(A_5(h))Cq(A_5(h))^{-1}$.
This proves centralization.  The precise integral centralizer
is D23.T2: $C^5=-I$, $-C$ has polynomial $\Phi_5$, is cyclic
over every coefficient ring, and its symplectic adjoint is its
inverse.  Thus commuting symplectic matrices are exactly the
displayed norm-one units.

For composition, write $Q_h=q(A_5(h))$.  Equation (D23.F1.2)
gives $Q_{hk}=V_hQ_kV_h^{-1}Q_h$, while D23.T1 gives
$V_{hk}=V_hV_k$.  Hence
\[
 Q_{hk}^{-1}V_{hk}
 =Q_h^{-1}V_hQ_k^{-1}V_k
 =\mathcal C_5(h)\mathcal C_5(k).
\]
Continuity follows from the continuity of all its factors.
\end{proof}

\paragraph{Current arithmetic comparison gate.}
The construction of $\mathcal C_5$ is algebraic and unconditional
on full $\mathcal H_5$ once the stated published torsion and
standard GT-action inputs are granted.  Equality with the
Jacobian Tate action requires accounting for the normalization
$z_1=0,z_2=1$.  At the ordered CM point
$(1,\zeta_5,\ldots,\zeta_5^4)$ put $d=\zeta_5-1$ and
$t=(x-1)/d$.  The normalized monic hyperelliptic equation is
\[
 y_0^2=\prod_{j=0}^{4}
 \left(t-\frac{\zeta_5^j-1}{d}\right),
 \qquad y=d^{5/2}y_0.
\]
Thus this normalized curve is a quadratic twist of
$y^2=x^5-1$ by $d^5$, whose Kummer class is the class of $d$.
The exact scalar comparison is not to be suppressed.

% DELTA23 /root/five_level_six; proof draft for independent checking.

\subsection*{D23.F3. An original fifth-fibre Tate character on the gentle group}
Let $F=\widehat F_2=\langle x,y\rangle$ with $z=(xy)^{-1}$,
and let $\mathcal H=\ker\lambda\subset\widehat{GT}_0$ act by
$\alpha_h(x)=x$, $\alpha_h(y)=f_h^{-1}yf_h$.
Consider the cyclic quotient
\[
 F\twoheadrightarrow C_5\times C_2,
 \qquad x\mapsto(1,0),\quad y\mapsto(0,1),
\]
with kernel $K_{10}$, and let $N$ be the closed normal subgroup
of $F$ generated by $x^5,y^2,(xy)^{10}$.
Then $N\subset K_{10}$ and
\[
 \Pi_5=K_{10}/N
\]
is the completed fundamental group of the compact genus-two
curve $C_5:y^2=x^5-1$.  Every $h\in\mathcal H$ induces an
automorphism of $\Pi_5$ and a continuous character
\[
 \Omega_5:\mathcal H\longrightarrow
 \{u\in(\widehat{\mathbb Z}[\zeta_5])^\times:
                         u\overline u=1\}.
 \tag{D23.F3.1}
\]
On $G_{\mathbb Q^{\rm ab}}$ this is the full Tate representation
of $J_5=\operatorname{Jac}(C_5)$, in the fixed covering marking.
If $f_h\in F''$, then $\Omega_5(h)=1$.

\begin{proof}
The algebraic map
\[
 C_5\longrightarrow\mathbb P^1_t,
 \qquad t=x^5=y^2+1
\]
is Galois of degree ten over $\mathbb Q(\zeta_5)$, with deck
group $C_5\times C_2$.  Its ramification orders over $0,1,\infty$
are $5,2,10$.  Removing these fibres gives the covering of
$\mathbb P^1\setminus\{0,1,\infty\}$ specified by $K_{10}$;
filling them gives exactly the normal closure of the displayed
peripheral powers.  This identifies $\Pi_5$.  Riemann--Hurwitz
gives
\[
 2g-2=10\left(-2+(1-1/5)+(1-1/2)+(1-1/10)\right)=2,
\]
so $g=2$.

The gentle action is peripheral: each of $x,y,z$ is sent to a
conjugate of itself at multiplier one.  It is the identity on
the abelian quotient $C_5\times C_2$.  Hence it preserves
$K_{10}$ and $N$, with the same true statement for its inverse;
the resulting map on $\Pi_5$ is an automorphism.  It commutes
with the order-five deck action on $\Pi_5^{\rm ab}$, since it
fixes $x$ literally.

The induced map on the four-dimensional homology is symplectic.
Here is a direct boundary argument that specifies its multiplier.
The open cover has eight punctures: two over zero, five over one,
and one over infinity.  The two over zero are fixed because $x$
is fixed.  The five over one are fixed individually because the
conjugator $f_h$ lies in $F'\subset K_{10}$, so it has trivial
deck coordinate.  The puncture over infinity is unique.
All eight peripheral loops are therefore preserved up to conjugacy,
with exponent one.  Fill the first seven punctures but leave
the one over infinity.  The resulting completed free group has
rank four; in a geometric free basis its remaining boundary is
\[
 [a_1,b_1][a_2,b_2].
\]
Our induced automorphism preserves its conjugacy class.  In
the quotient of the lower central series of degrees two and
three, conjugation disappears and this boundary gives
$a_1\wedge b_1+a_2\wedge b_2$.  Preservation of that element
is precisely the symplectic equation.  This proof works after
every finite reduction, hence over $\widehat{\mathbb Z}$.

The order-five rotation on $C_5$ has eigenvalues
$\zeta_5,\zeta_5^2$ on $dx/y,x\,dx/y$, and the conjugate
eigenvalues on the antiholomorphic part.  Its characteristic
polynomial on integral first homology is $\Phi_5$.
Consequently this lattice is a torsion-free rank-one module over
$\mathbb Z[\zeta_5]$, a Dedekind ring.  Such a module is an
invertible ideal and its endomorphism ring is exactly
$\mathbb Z[\zeta_5]$; after completion the same statement
holds prime by prime.  The symplectic adjoint of rotation is
its inverse, so the adjoint on that endomorphism ring is complex
conjugation.  The commuting symplectic automorphisms are exactly
the norm-one units in (D23.F3.1).  Restriction and abelianization
respect composition, establishing the asserted character and
its continuity.

For the arithmetic identification use the tangential point
at $t=0$ directed toward $1$, lifted to $x=0,y=i$ with leading
coefficient one in the local parameter $x=t^{1/5}$.
All its constant choices lie in $\mathbb Q^{\rm ab}$.
The standard Galois inclusion in $\widehat{GT}_0$ is the action
on this marked base fundamental group.  Functoriality of the
specified algebraic cover, followed by filling its punctures,
therefore identifies its restricted homology action with the
Tate action of $J_5$.  The complex/etale comparison used here is
the Riemann existence theorem; for an exact modern locator see
Milne, \emph{Lectures on Etale Cohomology}, v2.21, Theorem 3.4.
No claim about a five-point braid frame enters this identification.

Finally, $F'\subset K_{10}$ implies $F''\subset K_{10}'$.
If $f_h\in F''$, the action on $F/K_{10}'$ fixes both $x$ and
$y$ and hence is the identity.  It is therefore the identity on
$K_{10}^{\rm ab}$ and on its quotient $\Pi_5^{\rm ab}$.
\end{proof}

\subsection*{D23.F4. A scalar cyclotomic-unit character for the normalization}
Put $K_5^{\rm cyc}=\ker(F\to C_5)$, with $x\mapsto1$,
$y\mapsto0$.  Its standard Schreier generators are
\[
 X=x^5,\qquad Y_j=x^j y x^{-j}\quad(0\leq j<5).
\]
For $h\in\mathcal H$ define
\[
 \kappa_5(h)=\operatorname{coeff}_{Y_1}
       (f_h\bmod (K_5^{\rm cyc})')\pmod2.
 \tag{D23.F4.1}
\]
This is a continuous homomorphism to $\mathbb F_2$.
With the cyclic covering $t=x^5$ and the usual radial path
from zero to one, its restriction to $G_{\mathbb Q^{\rm ab}}$
is the Kummer character of $\zeta_5-1$.

\begin{proof}
Every $f_h$ belongs to $F'\subset K_5^{\rm cyc}$.  The action
fixes $X$ and conjugates each $Y_j$ by an element of
$K_5^{\rm cyc}$.  It is therefore the identity on the
abelianization of that group.  The composition rule
$f_{hk}=f_h\alpha_h(f_k)$ now makes every integral coefficient
in that abelianization additive; reduction modulo two proves
the first assertion.

For the arithmetic coefficient, pull back the double cover of
the $x$-line given by
\[
 v^2=1-x/\zeta_5.
\]
Its sign character on the cyclic covering fundamental group
is one on $Y_1$, zero on the other $Y_j$, and zero on $X$.
The radial path from $x=0$ to $x=1$ lifts from $v=1$ to a
chosen square root of $1-\zeta_5^{-1}$.  The usual Galois
path cocycle is the lift of $f_h$, up to inversion according
to the convention; inversion has no effect on its sign.
Thus (D23.F4.1) is the quadratic Kummer character of
$1-\zeta_5^{-1}$.  Its ratio with $\zeta_5-1$ is $\zeta_5^{-1}$,
which has a square root in $\mathbb Q^{\rm ab}$.  The two
quadratic characters on $G_{\mathbb Q^{\rm ab}}$ coincide.
\end{proof}

\subsection*{D23.F5. A direct common fifth-fibre bit is impossible}
Let $k=\mathbb Q^{\rm ab}$ and retain the actual torsion fields
$L_3,L_4,L_5$ of DELTA\_ON\_GT AF.4--AF.6.  No continuous
quadratic character of $\operatorname{Gal}(L_5/k)$ pulls back
on $G_k$ to the common arithmetic bit
\[
 a_2|_{G_k}=a_3|_{G_k}
        =\kappa_{\sqrt[4]{3}}.
 \tag{D23.F5.1}
\]
Thus a proposed proof identifying both bits with one quadratic
character of the genuine $J_5$ Tate representation is false
already on the Galois subgroup.

\begin{proof}
The field of the nonzero character on the right is
$k(\sqrt[4]3)=L_3\cap L_4$, by E.65 and AF.6.
If it factored through $\operatorname{Gal}(L_5/k)$, Galois
correspondence would give $k(\sqrt[4]3)\subset L_5$.
But AF.6 states $L_3\cap L_4\cap L_5=k$, a contradiction.
This applies to the full Tate image, so passing to a finite
level cannot repair the contradiction.
\end{proof}

\paragraph{Scope.}
D23.F3 and D23.F4 construct actual homomorphisms on the gentle
group, with arithmetic interpretations.  D23.F1--D23.F2
construct a different, frame-induced character on full GT.
An equality between the two constructions is not part of
these statements.  The explicit quadratic twist above shows
which scalar must be included in a comparison at the chosen
normalization; it does not make such a comparison on all GT
automatic.

% Arithmetic comparison. Independent whole-proof reading by /root completed;
% explicit positive-rotation marking added after that reader's requested check.
% No all-GT equality is asserted.

\subsection*{D23.F9. The exact scalar on the Galois subgroup}
Fix the usual ordered tangential braid base point
\[
 b_5=(0,t^4,t^3,t^2,t)
\]
and the monic hyperelliptic family
$Y^2=\prod_{j=1}^5(X-v_j)$ over the ordered configuration
space.  Choose the geometric path to the ordered CM
configuration $(1,\zeta_5,\zeta_5^2,\zeta_5^3,\zeta_5^4)$
that identifies its positive order-five rotation with the
standard class of $a=\sigma_1\sigma_2\sigma_3\sigma_4$.
The marking on the original character $\Omega_5$ is transported
along this path, with the order-five rotation represented by
$-q(a)$.  Then for $\sigma\in G_{\mathbb Q^{\rm ab}}$,
\[
 \mathcal C_5(\sigma)
       =(-1)^{\kappa_5(\sigma)}\Omega_5(\sigma).
 \tag{D23.F9.1}
\]
Consequently the character
\[
 \mathcal E_5(h)=(-1)^{\kappa_5(h)}
                  \mathcal C_5(h)\Omega_5(h)^{-1}
 \quad(h\in\mathcal H_5)
\]
is trivial on $G_{\mathbb Q^{\rm ab}}$.  This is a relative
character on the multiplier-one full GT group; no claim of its
triviality away from Galois or extension to the full gentle
group is made.

\begin{proof}
Let $M_\sigma$ be the arithmetic Tate action of the displayed
family at $b_5$. The exact base point and the arithmetic
Drinfeld action for every number of strands are recorded in
Nakamura, \emph{On arithmetic monodromy representations of
Eisenstein type in fundamental groups of once punctured
elliptic curves}, RIMS-1691 (2010), Section 5.4,
equation (5.4.2), printed p.~28. Arithmetic/topological
functoriality gives
\[
 q\Phi_\sigma=\operatorname{Ad}(M_\sigma)q.
\]
For multiplier one this is a symplectic action.  In addition
$M_\sigma e_1=e_1$ in the chain basis used by D23.T1.
To check the sign in this last assertion, let the first two
roots coalesce while the others remain at $t^3,t^2,t$.
The unit multiplying the double root at zero has leading
coefficient $-t^6$ and square root $it^3$.  Thus the node
splits over $\mathbb Q^{\rm ab}((t))$, with branches individually
fixed by $G_{\mathbb Q^{\rm ab}}$; its oriented vanishing
cycle is the first chain vector $e_1$.  The local annulus
has inertia $\widehat{\mathbb Z}(1)$, on which Galois acts
by the cyclotomic character, which is one here.  Hence the
cycle is fixed with its sign.  The uniqueness statement of
D23.T1 therefore gives $M_\sigma=V_\sigma$.

We specify the periodic-braid sign before comparing fibres.
Rotate the five branch points counterclockwise through
$2\pi/5$. Along the angular parameter $0\leq\theta\leq2\pi/5$
the isomorphisms from the initial curve to the rotating
family are
\[
 (X,Y)\longmapsto(e^{i\theta}X,e^{5i\theta/2}Y).
\]
Their endpoint is $(X,Y)\mapsto(\zeta_5X,-Y)$.
The positive periodic braid of this isotopy has the standard
class of $a$; choose the comparison path to identify these
positive classes. In the disc braid group both lifts have
exponent sum four, so a possible difference $z_5^e$ has
$20e=0$ and is trivial. The homological action of the
positive periodic braid is therefore exactly $q(a)$.
The order-five CM deck map $(X,Y)\mapsto(\zeta_5X,Y)$
differs from it by the hyperelliptic involution, whose
action on homology is $-I$. Its action is consequently
$-q(a)$, as required. This argument fixes the sign and
the direction of the cyclotomic marking, rather than only
the projective action.

Normalize the first two configuration coordinates to $0,1$.
As shown in D23.F1, its fundamental group is exactly
$\ker\ell_{12}\subset\widehat{PB}_5$.
At the base point the normalizing difference is $t^4$;
its fifth power has the fixed square root $t^{10}$, so this
normalization introduces no scalar on
$G_{\mathbb Q^{\rm ab}}$.

Let $p$ be the normalized CM point and transport its arithmetic
action to this normalized tangential base point along the fixed
path.  Its order-five automorphism is defined over
$\mathbb Q(\zeta_5)$ and is fixed by $\sigma$.
Consequently the path cocycle is a pure conjugator of $\bar a$.
By the pure centralizer calculation in D23.F1 it is precisely
$\alpha_5(\sigma)$.  Since the path lies in the normalized
configuration space, its lift in the pure braid group has
$\ell_{12}=0$ and is exactly $A_5(\sigma)$.
The standard base-point change formula for a local system now
identifies the Tate action at $p$, transported to the fixed
homology lattice, with
\[
 q(A_5(\sigma))^{-1}M_\sigma
       =\mathcal C_5(\sigma).
\]

Finally set $d=\zeta_5-1$.  The coordinate change
$X=1+dt$ in the original CM equation gives
\[
 X^5-1=d^5\prod_{j=0}^4
             \left(t-\frac{\zeta_5^j-1}{d}\right).
\]
The normalized monic family and the original curve are
therefore related by $Y=d^{5/2}Y_0$.
The only nontrivial scalar in their Tate actions is
$\sigma(\sqrt d)/\sqrt d$, because $d^2$ is already in
$\mathbb Q(\zeta_5)$.  D23.F4 identifies that scalar as
$(-1)^{\kappa_5(\sigma)}$.  The original Tate action is
$\Omega_5(\sigma)$ by D23.F3, giving (D23.F9.1).
All three factors are homomorphisms into the commutative
norm-one torus, so their product defines the stated relative
character.
\end{proof}

\paragraph{Scope of the comparison.}
The formula retains the specified tangential point, ordered
path, vanishing-cycle sign and base-point-change convention.
In particular the order-five deck rotation in $\Omega_5$
is identified with $-q(a)$. No equality between the three
characters on all of the full GT group is inferred from
their arithmetic equality.


\section{Level-six marking and exact comparison}
% DELTA23 /root/five_level_six. Exact marking and decision criterion.

\subsection*{D23.F6. An explicit marked Atkin--Lehner involution}
Use the literal modular matrices
\[
 s=\begin{pmatrix}0&1\\-1&0\end{pmatrix},\quad
 r=\begin{pmatrix}0&-1\\1&-1\end{pmatrix},\quad
 \tau=sr=\begin{pmatrix}1&-1\\0&1\end{pmatrix}
\]
in $\operatorname{PSL}_2(\mathbb Z)$, and put
$H=\overline{\Gamma_0(6)/\{\pm I\}}\subset\widehat B$.
Then $H$ is the free profinite group on
\[
 a=\tau,\qquad b=r\tau^5r,\qquad
 c=r\tau^2r\tau r\tau^2r.
 \tag{D23.F6.1}
\]
Their projective matrix representatives are respectively
\[
 A=\begin{pmatrix}1&-1\\0&1\end{pmatrix},\quad
 B=\begin{pmatrix}1&-1\\6&-5\end{pmatrix},\quad
 C=\begin{pmatrix}5&-3\\12&-7\end{pmatrix}.
\]
The Atkin--Lehner matrix
\[
 W_3=\begin{pmatrix}3&-2\\6&-3\end{pmatrix},\qquad
 \det W_3=3,\quad W_3^2=-3I
\]
induces the literal involution
\[
 w(a)=c^{-1},\qquad w(b)=b^{-1},\qquad w(c)=a^{-1}
 \tag{D23.F6.2}
\]
on this marked free group. Its outer class is the $w_3$ used
in DELTA\_ON\_GT FP.17.

\begin{proof}
Here is an explicit Bass--Serre/Schreier certificate for the
free-basis assertion.  The twelve right cosets of $H$ are
the projective bottom rows modulo two and three.  Numbering
them starting with $H$ as zero, the right actions of the
generators are
\[
 s=(0\ 1)(2\ 3)(4\ 5)(6\ 8)(7\ 9)(10\ 11),
\]
\[
 r=(0\ 2\ 1)(3\ 4\ 6)(5\ 7\ 10)(8\ 11\ 9).
\]
These permutations can be checked by multiplying rows by
the displayed two matrices in $\mathbb P^1(\mathbb F_2)$
and $\mathbb P^1(\mathbb F_3)$; their stabilizer at zero
is exactly the condition that the lower-left entry vanish
modulo six.  Use two consecutive edges in each of the four
displayed $r$-cycles, starting at $0,3,5,8$, and the three
$s$-edges $(2,3),(4,5),(6,8)$ as a tree.  After collapsing
that tree and the four $r^3$ relations, precisely the three
$s$-edges $(0,1),(7,9),(10,11)$ remain.  Their oriented
Schreier loops are exactly (D23.F6.1).  The $s^2$ relations
only reverse the chosen orientation of each such edge.
Thus these loops form a discrete free basis.  Completion of
this finite-index subgroup has the induced topology from
$B$, and gives the asserted profinite free basis.

Direct multiplication yields the three displayed matrices.
It also gives, in projective matrices,
\[
 W_3 A W_3^{-1}=C^{-1},\qquad
 W_3 B W_3^{-1}=B^{-1},\qquad
 W_3 C W_3^{-1}=A^{-1}.
\]
The free-basis result now proves (D23.F6.2) as full group
equalities, not only in a finite quotient.  In particular
$W_3$ normalizes the modular subgroup.  Its shape
$\left(\begin{smallmatrix}3u&v\\6w&3z\end{smallmatrix}\right)$
and determinant three give the usual level-six Atkin--Lehner
coset for the exact divisor three.  Replacing a representative
of that coset by a member of $H$ changes the induced map only
by an inner automorphism, so the outer class is the prescribed
one.  The supplied exact program \texttt{level6\_words.py}
independently performs the row enumeration, Schreier rewriting,
matrix conjugations, and square check.
\end{proof}

\subsection*{D23.F7. Two exact word tests for the level-six comparison}
For $h\in\mathcal H$ let $\Phi_h$ be the normalized modular
automorphism fixing $\tau$, and let $\rho_h$ be E.65's
integral harmonic coordinate.  Set
\[
 \psi_h=\operatorname{Ad}(a^{8\rho_h})\Phi_h\big|_H.
 \tag{D23.F7.1}
\]
This represents $\rho_6(h)$ of FP.17 and fixes $a$.
With $y=s\tau^{-2}s$ put
\[
 k_h=\psi_h(y^{-1})y\in H,\qquad
 d_h=\operatorname{coeff}_a(k_h^{\rm ab})-
                     \operatorname{coeff}_c(k_h^{\rm ab}),
 \qquad v_h=k_hc^{d_h}.
 \tag{D23.F7.2}
\]
In this formula $\psi_h$ on $y^{-1}$ means the same conjugated
automorphism of all of $\widehat B$ before restriction.
Then the exact criterion for A.86 is
\[
 AL_{6,3}(h)=1\quad\Longleftrightarrow\quad
 \begin{cases}
 v_h w(v_h)=1,\\
 \psi_h(b)^{-1}=v_h w(\psi_h(b))v_h^{-1}.
 \end{cases}
 \tag{D23.F7.3}
\]
The balanced frame $v_h$ does not depend on which conjugator
$k_h\in H$ is used for $\psi_h(c)=k_hck_h^{-1}$.
There is no remaining unspecified inner conjugator in this test.

\begin{proof}
E.65.9--10 proves
$q_{12}\Phi_h=\operatorname{Ad}(U^{8\rho_h})q_{12}$
on the actual modular quotient, with $U=q(\tau^{-1})$.
Thus (D23.F7.1), defined first on $\widehat B$, is the identity
on its level-six quotient and preserves $H$.  It fixes $a$
since both factors do.  It represents the restriction outer
action of FP.17 because its correction is inner in
$\widehat B$.

In the literal matrices one has
\[
 g=y^{-1}=\begin{pmatrix}1&0\\2&1\end{pmatrix},
 \qquad c=ga^{-3}g^{-1}.
\]
Consequently $k_h=\psi_h(g)g^{-1}$ conjugates $c$ to
$\psi_h(c)$.  Because $\psi_h$ is the identity modulo six,
$k_h$ belongs to the principal level-six subgroup, hence to
$H$.  In terms of the original parameter this is the explicit word
\[
 k_h=a^{8\rho_h}f_h^{-1}y^{-1}f_h a^{-8\rho_h}y,
 \quad
 \psi_h(b)=a^{8\rho_h}f_h^{-1}r a^5 f_h^{-1}r a^{-8\rho_h}.
 \tag{D23.F7.4}
\]

For completeness, the following elementary free-group argument
eliminates the inner ambiguity.  Write $\psi=\psi_h$, $k=k_h$.
The condition $[\psi,w]=1$ in $\operatorname{Out}(H)$ says
that for some $v\in H$,
$\psi w=\operatorname{Ad}(v)w\psi$.
Apply this equality to $a$.  It gives
$k c^{-1}k^{-1}=v c^{-1}v^{-1}$.  Since a free generator
has centralizer its procyclic subgroup, $v=kc^e$ for a
unique $e\in\widehat{\mathbb Z}$.
Applying it to $c$ gives
\[
 k c^e w(k)\in\overline{\langle a\rangle}.
\]
If $k^{\rm ab}=p a+q b+t c$, then (D23.F6.2) shows that
the abelianization on the left is
$(p-t)a+(t+e-p)c$.
Thus necessarily $e=p-t=d_h$, and the displayed element
must be exactly $a^{d_h}$.  This is equivalent to
$v_hw(v_h)=1$.  The equality on the final free generator
$b$ is exactly the second line of (D23.F7.3).
Conversely those two equalities give the required identity
on all three free generators and hence the equality of
continuous automorphisms.  This proves the criterion.

Changing $k$ to $kc^j$ changes $d_h$ to $d_h-j$, so
$kc^j c^{d_h-j}=kc^{d_h}$.  This proves the asserted
independence.  The centralizer fact used here is the
free-factor centralizer theorem for
$H=\overline{\langle c\rangle}\amalg\widehat F_2$
(and similarly for $a$); it is the same Herfort--Ribes
Theorem B' used in D23.F1.
\end{proof}

\paragraph{What the criterion does and does not decide.}
For the element $h_*$ of E.58 one has $\rho_{h_*}=0$, so
\[
 k_{h_*}=f_{h_*}^{-1}y^{-1}f_{h_*}y,
 \qquad \psi_{h_*}(b)=f_{h_*}^{-1}r a^5 f_{h_*}^{-1}r.
\]
E.58 specifies a relative-pronilpotent lift and finite frame
coordinates, rather than these two full free-group values.
No equality or inequality in (D23.F7.3) on that element is
claimed here.  The computation supplies a marked basis and
an exact normalization for the test that was absent in FP.17;
the asserted equivalence with $\Delta=0$ remains a separate
decision.

\subsection*{D23.F8. The Atkin--Lehner matrices do not generate the division algebra of discriminant six}
An algebra generated over $\mathbb Q$ by rational
$2\times2$ Atkin--Lehner matrices is a subalgebra of
$M_2(\mathbb Q)$.  It cannot be isomorphic to the division
quaternion algebra $(2,3)$ of discriminant six.  In particular
the literal matrices used to mark A.86 do not provide an
identification with that quaternion algebra.

\begin{proof}
The Hilbert-symbol calculation $(2,3)_2=(2,3)_3=-1$,
with all other local symbols positive, is the standard
local calculation already made in D23.L3.  It is therefore
a division quaternion algebra over $\mathbb Q$.
An injective homomorphism from that four-dimensional algebra
to $M_2(\mathbb Q)$ would be an isomorphism by dimension.
But $M_2(\mathbb Q)$ has nonzero noninvertible elements and
a division algebra does not, a contradiction.  Any nonzero
unital homomorphism from a division algebra is injective,
so even that weaker possibility is excluded.
This obstruction concerns a literal matrix-algebra
identification; it does not exclude a separately constructed
Brauer class or a Shimura-curve comparison.
\end{proof}

\subsection*{D23.F10. The extra-CM element passes every pro-two level-six test}
Let $h_*$ be the actual element of LM7\_PUSH E.58, and let
$H^{(2)}$ be the maximal pro-two quotient of the marked
level-six subgroup $H$.  The automorphism $\rho_6(h_*)$
is the identity on $H^{(2)}$. Consequently each of
$AL_{6,2}(h_*),AL_{6,3}(h_*),AL_{6,6}(h_*)$ is the identity
on that entire quotient, although $\Delta(h_*)=1$.
No test factoring through the induced action on
$H^{(2)}$ can establish $AL_{6,3}=1\Longleftrightarrow\Delta=0$.

\begin{proof}
Put $U=F'F^3$, as in E.58.  That theorem supplies
$\Phi_{h_*}=1$ on $\widehat B/R_{\rm nil}(U)$ and
$\rho_{h_*}=0$.  The latter makes the correction in
(D23.F7.1) trivial.
The finite source calculation E.58.11--13, reduced modulo
six, gives
\[
 q_6(U)=\{1\}\times V_4
 \quad\hbox{in }S_3\times A_4.
\]
On the other hand $q_6(H)=C_2\times C_3$ by FP.17.
Their intersection is trivial. Thus
$V=U\cap H=\ker(q_6|_U)$ is normal of index four in $U$,
with quotient $V_4$.

For a normal subgroup $V$ of a profinite group $U$ with
finite two-group quotient, one has $O^2(V)=O^2(U)$.
Indeed, if $N$ is an open normal subgroup of $V$ of
two-power index, intersect its finitely many $U$-conjugates.
The resulting $N_0$ is normal in $U$; the quotient
$V/N_0$ embeds into a finite product of two-groups, and
$U/N_0$ is an extension of that two-group by $U/V$.
Hence $O^2(U)\subset N_0\subset N$. Intersecting over $N$
gives $O^2(U)\subset O^2(V)$.  Conversely restriction of
every finite two-group quotient of $U$ to $V$ gives
$O^2(V)\subset O^2(U)$.
Since every finite nilpotent residual lies in every
pro-prime residual, and every homomorphism from $H$ to a
finite two-group restricts to such a homomorphism on $V$,
we obtain
\[
 R_{\rm nil}(U)\subset O^2(U)=O^2(V)\subset O^2(H).
\]
For each $h\in H$ the difference
$\Phi_{h_*}(h)h^{-1}$ belongs to $R_{\rm nil}(U)$,
so its image in $H^{(2)}$ is one.  Thus the induced
automorphism is the identity.  The characteristic subgroup
$O^2(H)$ is preserved by each $w_e$, so their crossed
commutators with this automorphism are identity too.
The sign $\Delta(h_*)=1$ is E.58.23 and E.65.4.
\end{proof}

\subsection*{D23.F12. One actual mixed-sign element is invisible to every pronilpotent level-six test}
Assume the supplied relative peripheral lifting theorem
LM7\_PUSH E.58, in its stated version for every open normal
subgroup.  There is an actual element $h_6\in\mathcal H$
with
\[
 \Delta(h_6)=1,\qquad\rho_{h_6}=0,
 \qquad \rho_6(h_6)|_{H/R_{\rm nil}(H)}=1.
 \tag{D23.F12.1}
\]
Consequently all three level-six Atkin--Lehner defects of
this one element are the identity on every finite nilpotent
quotient of $H$, and on its entire maximal pronilpotent
quotient.  This is an actual gentle element, rather than
a collection of unrelated finite shadows.

\begin{proof}
Let
\[
 U=N_6=\ker(\widehat B\longrightarrow
              Q_6=\operatorname{SL}_2(\mathbb Z/6)/\{\pm I\}).
\]
Then $U$ is open normal and $U\subset H\cap F$.
Let $q=q_{12}$ be the exact projective level-twelve quotient
of E.58. Its image $q(U)=\ker(Q_{12}\to Q_6)$ is an
elementary abelian two-group of order eight.  Indeed every
class in this kernel has a unique representative congruent
to $I$ modulo six, of the form $I+6M$ with
$M\in M_2(\mathbb F_2)$ and $\operatorname{tr}M=0$.
Multiplication is addition of $M$, proving the assertion.
The central element $z_0=[5I]=[7I]$ belongs to this kernel
and corresponds to $M=I$.

Thus $q$ factors through
$\Gamma_2=\widehat B/O^2(U)$, and the kernel of
$\Gamma_2\twoheadrightarrow Q_{12}$ is pro-two.
The element $r$ has order three in both groups.
The coprime centralizer-lifting argument in the proof of
E.58.19 gives a lift
\[
 c_2\in C_{\Gamma_2}(r)\cap U^{(2)},
 \qquad q(c_2)=z_0.
\]
For specificity, that argument lifts the centralizer
element $z_0$, then conjugates its lifted conjugate of
$\langle r\rangle$ back by the pro-two kernel using
finite Schur--Zassenhaus complement conjugacy and compactness.
Since $z_0$ is identity in the retained base $Q_6$, the
corrected lift lies in $U^{(2)}$.

Use E.58.2 to identify the relative pronilpotent quotient
with the fibre product of the relative pro-prime quotients
over $Q_6$.  Give its identity automorphism the specified
peripheral conjugators
\[
 b_s=1,\qquad b_\tau=1,\qquad
 b_r=(c_2\text{ at two},\ 1\text{ at every other prime}).
\]
Their base coordinates agree, and $b_r$ really centralizes
$r$ in that quotient.  Apply the supplied relative lifting
theorem, and normalize the full lift to fix $\tau$ exactly.
This produces an actual automorphism $\Phi$ of
$\widehat B$, identity on
$\widehat B/R_{\rm nil}(U)$, with unique pure frames
\[
 \beta\in R_{\rm nil}(U),\qquad A\in U,
 \qquad q(\beta)=1,\quad q(A)=z_0.
\]
The literal parameter $f=s\Phi(s)$ lies in
$R_{\rm nil}(U)\subset U'\subset F'$.
The torsion and product calculation E.58.21, applied to
this normalized lift, gives the two gentle equations and
the actual free-group automorphism.  Hence this is an
element $h_6\in\mathcal H$.  Because $\beta\in U'\subset F'$,
its harmonic coordinate is zero.  The same signed
calculation as E.58.23 gives $(a_3,a_2)=(1,0)$ and
$\Delta(h_6)=1$.

Every continuous map from $H$ to a finite nilpotent group
restricts to such a map on $U$, since $U\subset H$.
Therefore
\[
 R_{\rm nil}(U)\subset R_{\rm nil}(H).
\]
The identity of $\Phi$ modulo the former residual proves
the last assertion in (D23.F12.1).  Each Atkin--Lehner
automorphism preserves the characteristic subgroup
$R_{\rm nil}(H)$, so its crossed commutator with this
identity action is trivial on the whole quotient.
All finite nilpotent tests follow at once from this one
full quotient identity.
\end{proof}

\paragraph{The remaining level-six gate after D23.F12.}
To distinguish this witness from the kernel of $\Delta$
using $AL_{6,3}$, a test must retain a genuinely
nonnilpotent quotient of the free level-six group.
No claim that $AL_{6,3}(h_6)$ itself is trivial is made.

\subsection{Every congruence Atkin--Lehner test is trivial}

\paragraph{D23.F18.}
Let $H$ be the marked profinite level-six subgroup of
D23.F6 and let $H^{\rm cong}$ be its full congruence
completion, namely the closure of
$\Gamma_0(6)/\{\pm I\}$ in
$\operatorname{SL}_2(\widehat{\mathbf Z})/\{\pm I\}$.
For every multiplier-one gentle element $h$, its
induced outer action on $H^{\rm cong}$ is trivial.
The three Atkin--Lehner involutions extend to this
completion. Consequently every outer defect
$AL_{6,e}(h)$, $e=2,3,6$, is trivial on the entire
congruence completion and on every finite quotient
of it to which the Atkin--Lehner action descends.

In particular the congruence version of
$AL_{6,3}=1\Longleftrightarrow\Delta=0$ is false:
the supplied E.58 witness has $\Delta=1$ and still
has trivial congruence defect. Any decision of the
full A.86 comparison must use noncongruence
information, not just a finer congruence level.

\begin{proof}
Write $q$ for the adelic congruence representation,
projectivized by the global sign when its source
is $\widehat B$.  The signed covariance theorem
DELTA\_ON\_GT AD.1 gives a uniquely defined
$t_h\in\widehat{\mathbf Z}$ with
\[
 q\Phi_h=\operatorname{Ad}(U^{t_h})q,
 \qquad U=\begin{pmatrix}1&1\\0&1\end{pmatrix}.
\]
The marked cusp generator of $H$ is $a=\tau$,
with $q(a)=U^{-1}$.  Thus $U^t$ belongs to
$H^{\rm cong}$ for every $t\in\widehat{\mathbf Z}$.
This can also be seen directly from its lower-left
entry, which is zero.

For the representative used in D23.F7,
$\psi_h=\operatorname{Inn}(a^{8\rho_h})\Phi_h$,
the congruence covariance is
\[
 q\psi_h
  =\operatorname{Ad}(U^{t_h-8\rho_h})q.
\tag{D23.F18.1}
\]
The conjugating element is itself in $H^{\rm cong}$.
Hence the induced automorphism of that completion
is inner. In particular the congruence kernel of
$H$ is preserved: equation (D23.F18.1) proves this,
so no kernel invariance is assumed in forming
the quotient action.

For precision about bad primes, the closure in
question is exactly the projective image of
\[
 \{g\in\operatorname{SL}_2(\widehat{\mathbf Z}):
                         g_{21}\in6\widehat{\mathbf Z}\}.
\]
Indeed reductions of $\operatorname{SL}_2(\mathbf Z)$
surject onto every finite congruence quotient,
and lifting at a modulus divisible by six retains
the displayed lower-left condition. Thus it
is the inverse image of the usual level-six
subgroup and not an unspecified larger adelic
group.

Every rational Atkin--Lehner matrix normalizing
$\Gamma_0(6)$ acts continuously by conjugation
on its local closures, with a continuous inverse,
and therefore on the full congruence completion.
For the specified $W_3$ this is completely
explicit. If
\[
 W_3=\begin{pmatrix}3&-2\\6&-3\end{pmatrix},
 \qquad g=\begin{pmatrix}a&b\\6c&d\end{pmatrix},
\]
then
\[
 W_3gW_3^{-1}=
 \begin{pmatrix}
 -3a+12c-6b+4d&2a-8c+3b-2d\\
 6(-a+3c-2b+d)&4a-12c+6b-3d
 \end{pmatrix}.
\]
These integral formulas make sense over
$\widehat{\mathbf Z}$, preserve determinant one
and the level-six condition, and their square
is the identity because $W_3^2=-3I$.
The same closure argument applies to $W_2,W_6$;
their classes square to inner automorphisms as
usual for Atkin--Lehner involutions.

In the outer automorphism group of
$H^{\rm cong}$ the class of (D23.F18.1) is
therefore identity. Its crossed commutator with
any Atkin--Lehner class is identity as well.
If a chosen finite normal congruence kernel is
not initially Atkin--Lehner stable, intersect
its finitely many translates under the
Atkin--Lehner group; their squares act inner
and hence fix every normal subgroup. The
result is a finite-index invariant congruence
kernel, and the same inner-action argument
applies there. This explains the precise
finite-quotient formulation of the statement.

Finally E.58.23 and E.65 supply an actual gentle
element with $\Delta=1$. The just-proved
congruence conclusion applies to it, proving
the stated refutation of the congruence version.
No conclusion that its full noncongruence
Atkin--Lehner defect is trivial is drawn.
\end{proof}


\section{Relative characters and the full adelic image}
% DELTA23 characters workstream. Author: GPT-6 Astra Pro (/root/characters).
% First proof reading completed on 2026-11-21 Europe/Berlin.
% This file uses the input notation G=GT_0-hat, H=ker(lambda).

\subsection{A full relative character from the harmonic descent defect}

\paragraph{D23.C1 (proved from the cited inputs).}
Let $G=\widehat{\mathrm{GT}}_0$, $H=\ker\lambda$, and let
$D=D_P:G\to G$ be the endomorphism of
DELTA\_ON\_GT A.14, in its literal current marking.  Write
$a_2=\epsilon_2W$ and $a_3=\epsilon_3T$ on $H$.
The character
\[
 d_H(h)=a_2(h)+a_2(Dh)\quad(h\in H)
\]
has a unique continuous extension
\[
 d:G\longrightarrow\mathbb F_2,
 \qquad d|_{G_{\mathbb Q}}=0.
\]
For $g\in G$, its choice-independent formula is
\[
 d(g)=d_H(g\sigma^{-1}),\qquad
 \sigma\in G_{\mathbb Q},\quad\chi(\sigma)=\lambda(g).
\]
On $H$, this is the mixed duplication defect $d_2$ of
DELTA\_ON\_GT A.103 and A.107.  Its nonvanishing is not asserted.

\emph{Proof.}
The two individual characters $a_2,a_3$ are invariant under
conjugation by all of $G$: this is the separate normalizer
calculation in LM7\_PUSH E.65, equations (E.65.16)--(E.65.17),
not just invariance of their sum.  Since $D$ is a homomorphism
preserving $\lambda$, it maps $H$ to $H$, and
\[
 a_2(D(ghg^{-1}))=a_2(Dg\,Dh\,(Dg)^{-1})=a_2(Dh).
\]
Thus $d_H$ is a continuous, $G$-invariant character of $H$.
The arithmetic marking comparison DELTA\_ON\_GT A.62 says
$D(\tau)=c\tau c$ for $\tau\in G_{\mathbb Q}$, where
$c=(-1,1)$.  If also $\tau\in H$, conjugation invariance gives
$a_2(D\tau)=a_2(\tau)$; hence $d_H(\tau)=0$.

Every element of $G$ has the form $h\sigma$ with $h\in H$ and
$\sigma\in G_{\mathbb Q}$, because the cyclotomic character is
surjective.  If $h\sigma=h'\sigma'$, then
$h'^{-1}h=\sigma'\sigma^{-1}\in H\cap G_{\mathbb Q}$,
on which $d_H$ is zero.  Therefore $d(h\sigma)=d_H(h)$ is
well defined.  Multiplication in this form gives
\[
 (h\sigma)(k\tau)=h(\sigma k\sigma^{-1})(\sigma\tau),
\]
so conjugation invariance proves additivity.  The map
$H\times G_{\mathbb Q}\to G$, $(h,\sigma)\mapsto h\sigma$,
is a continuous surjection from a compact space to a Hausdorff
space.  The continuous function $(h,\sigma)\mapsto d_H(h)$ is
constant on its fibres, so the resulting $d$ is continuous.
The same decomposition proves uniqueness.  A.107, equation
(MD.8b), identifies its restriction with $d_2$.  No lifting
argument or pentagon relation enters this proof.\hfill$\square$

\paragraph{D23.C2 (proved; exact restricted classification).}
For an integer $n\geq0$, define the finite parity map
\[
 p_n:H\longrightarrow(\mathbb F_2^2)^{n+1},\qquad
 p_n(h)=\bigl((a_2(D^jh),a_3(D^jh))\bigr)_{0\leq j\leq n}.
\]
Among the continuous characters $G\to\mathbb F_2$ trivial on
$G_{\mathbb Q}$ whose restriction to $H$ factors through $p_n$,
all are precisely the span of
\[
 \Delta\circ D^j\quad(0\leq j\leq n),\qquad
 d\circ D^j\quad(0\leq j<n).
\]
This gives generators, not a claim that they are independent.
In particular, allowing finite strings of descendant frame parities
gives the $\mathbb F_2[D^*]$-module generated by $\Delta$ and $d$,
where $D^*\eta=\eta\circ D$.  It does not classify characters
without this factorization hypothesis.

\emph{Proof.}
Let $S_n=p_n(H)$, an $\mathbb F_2$-subspace of the displayed
finite vector space.  By DELTA\_ON\_GT B.14 the pair $(a_2,a_3)$
on $G_{\mathbb Q}\cap H$ has image $\{(0,0),(1,1)\}$.
The arithmetic formula $D\tau=c\tau c$, together with separate
conjugation invariance of the two parities, shows that the image
of $G_{\mathbb Q}\cap H$ under $p_n$ is exactly the diagonal line
spanned by the vector with all $2n+2$ entries equal to one.

A character factoring through $p_n$ is a linear functional on
$S_n$.  Extend it to the ambient vector space by extending a
basis of $S_n$ to a basis of that space.  Since the diagonal
vector lies in $S_n$ and its functional value is zero, the
extended functional has an even sum of coordinate coefficients.
The even-sum coefficient vectors are spanned by the following
$2n+1$ edge vectors: the sums of the two coordinates at a fixed
level $j$, and the sums of the harmonic coordinates at consecutive
levels $j,j+1$.  To check this without any rank assumption on
$S_n$, view the $2n+2$ coordinate positions as vertices: these
edges form a tree.  Removing a leaf and its incident edge eliminates
one specified coefficient; induction leaves the zero coefficient
at the final vertex because the total sum is even.  The corresponding
functions on $H$ are respectively
\[
 (a_2+a_3)\circ D^j=\Delta|_H\circ D^j,
 \qquad (a_2+a_2\circ D)\circ D^j=d_H\circ D^j.
\]
They extend to the global relative characters in the statement by
E.65 and D23.C1.  Two global characters trivial on $G_{\mathbb Q}$
that agree on $H$ agree on $G=HG_{\mathbb Q}$.  This proves both
containments and the asserted exact restricted classification.
\hfill$\square$

\paragraph{D23.C3 (proved; a separation and the harmonic lifting obstruction).}
Assume the actual element $k$ of DELTA\_ON\_GT LI.8
(integrated statement A.94) with
$Dk=1$, $\rho_k=0$, $q_{12}(\beta_k)=1$, and $\Delta(k)=1$.
Then
\[
 \Delta(k)=1,\qquad
 (\Delta\circ D^j)(k)=0\ (j\geq1),\qquad
 (d\circ D^j)(k)=0\ (j\geq0).
\]
In particular, if $d$ is nonzero it is independent of $\Delta$.
More generally $\Delta$ is outside the span of every positive
iterate of $\Delta$ and every iterate of $d$.

The elementary coprime-centralizer construction used in E.58
cannot, with identity relative action and a seed lying in the
relative kernel, be turned into a witness with
$q_{12}(\beta)=W=\pm S$ at three by replacing its prime two
component with a prime three component.

\emph{Proof.}
The retained frame and zero shear in LI.8 give $a_2(k)=0$.
Since $Dk=1$, the definition in D23.C1 yields $d(k)=0$ and
all positive iterates vanish.  Evaluating any proposed linear
relation at $k$ proves the separation claims.

For the last assertion, the relative lift in E.58 first chooses an
open normal subgroup $U$ and uses a seed whose finite frame value
belongs to $q(U)$, with $q(U)$ a $p$-group.  A harmonic value
$\pm S\in\mathrm{SL}_2(\mathbb F_3)$ has order four, so such
a $q(U)$ must be a two-group: it cannot be a three-group or any
odd-primary group.  But the torsion element being centralized is
$s$, of order two.  The hypothesis that its order be prime to the
pro-$p$ kernel, required for the coprime complement-conjugacy
argument, fails at this forced prime $p=2$.  Therefore that specific
construction does not establish the required centralizer lift.
This proves an obstruction to the proposed proof, not a theorem
that no other harmonic lift exists.\hfill$\square$

% DELTA23 characters workstream, 2026-11-21 Europe/Berlin.
% D23.C4 is an arithmetic derivation. D23.C5 was expanded after a root
% review to separate ordinary Legendre paths from the projective orbifold.
% Root independently read and passed C4--C7 on 2026-11-21, 00:29 CET;
% see the root check register for its exact event time.

\subsection{The arithmetic size of the normalized adelic frame target}

Put $O_3=\mathbb Z[\zeta_3]$, $O_4=\mathbb Z[i]$,
$K_j=O_j\otimes\mathbb Q$, and write a bar for the nontrivial
automorphism of $K_j$.  Set
\[
 N_j=\{u\in(O_j\otimes\widehat{\mathbb Z})^\times:u\bar u=1\},
 \qquad N_j(4)=\{u\in N_j:u\equiv1\pmod4\}.
\]
Here the congruence modulo four concerns the component over two.
The integral matrix embeddings are the ones in DELTA\_ON\_GT
AD.2 (integrated A.72), so $N_3(4)$ and $N_4(4)$ are precisely
its two normalized targets.

\paragraph{D23.C4 (arithmetic derivation; first proof reading passed).}
For $E_3:y^2=x^3+1$ and $E_4:y^2=x^3-x$, choose CM-compatible
bases of their full Tate modules.  Let $\rho_j$ be their adelic
representations and set $G^{\mathrm{cyc}}=G_{\mathbb Q^{\rm ab}}$.
Then
\[
 \rho_4(G^{\mathrm{cyc}})=N_4(4).
\]
For $r_0=(3/4)^{1/4}$ and
$\kappa(\sigma)=\sigma(r_0)/r_0\in\{\pm1\}$ on
$G^{\mathrm{cyc}}$, the normalized representation
\[
 \rho_3^\circ(\sigma)=\kappa(\sigma)\rho_3(\sigma)
\]
has image $N_3(4)$ and has the same kernel as $\rho_3$ on
$G^{\mathrm{cyc}}$.  Both full division fields are
$\mathbb Q(E_j[\mathrm{tors}])=K_j^{\rm ab}$.

\emph{Proof.}
We use the main theorem of complex multiplication in the torsion-point
form: Milne, \emph{Modular Functions and Modular Forms}, version 1.10,
Theorem 12.27, printed p.124
(\url{https://www.jmilne.org/math/CourseNotes/MF110.pdf}).
For a class-number-one imaginary quadratic ring $O$, it states that
the image of the full torsion representation after quotient by the
geometric automorphism group $O^\times$ is all of
$\widehat O^\times/O^\times$, with the usual ray-class interpretation.
It also implies that adjoining the Weber values of all torsion points
generates $K^{\rm ab}$.  All torsion coordinates are abelian over $K$
because the CM action is defined there and its Tate-module centralizer
is the commutative group $\widehat O^\times$.  Hence the full division
field over $K$ is $K^{\rm ab}$.

Let $\pi_4=1+i$ and $\pi_3=\sqrt{-3}$.  The entire subgroup
$E_4[\pi_4^3]$ is $K_4$-rational: besides the four two-torsion
points, the four points with abscissae $i,-i$ have ordinates
$\pm(1-i),\pm(1+i)$, respectively, and are the remaining points
of this kernel.  This last kernel identification can be checked
using the automorphism $[i](x,y)=(-x,iy)$ and the chord-and-tangent
addition law: applying $1+i$ to any of those four points gives
one of the nonzero two-torsion points in $E_4[\pi_4^2]$ which
maps to $(0,0)$ under $1+i$.
Thus $\rho_4(G_{K_4})$ is contained in
\[
 C_4=\{u\in\widehat O_4^\times:u\equiv1\pmod{\pi_4^3}\}.
\]
The reduction $\mu_4\to(O_4/\pi_4^3)^\times$ is an isomorphism:
the target has four elements, and the four roots of unity are distinct
modulo $\pi_4^3$.  Consequently $C_4\to\widehat O_4^\times/\mu_4$
is an isomorphism.  The main theorem therefore gives equality
$\rho_4(G_{K_4})=C_4$.

For $E_3$, all two-torsion points have abscissae
$-1,-\zeta_3,-\zeta_3^2$, and the kernel of $1-\zeta_3$ is
$\{O,(0,1),(0,-1)\}$.  The coprime kernels show that the whole
$E_3[2\pi_3]$ is $K_3$-rational.  The reduction
\[
 \mu_6\longrightarrow (O_3/(2\pi_3))^\times
       \simeq\mathbb F_4^\times\times\mathbb F_3^\times
\]
is an isomorphism: the cubic roots give the first factor and the
sign gives the second.  The same argument proves
\[
 \rho_3(G_{K_3})=C_3
   =\{u\in\widehat O_3^\times:u\equiv1\pmod{2\pi_3}\}.
\]
The CM fields themselves occur in the torsion fields (use the
three roots of the two-division polynomial for $E_3$, and the
Weil pairing on four-torsion for $E_4$), proving the asserted
identities over $\mathbb Q$ as well.

Complex conjugation acts on each Cartan by $u\mapsto\bar u$.
Since $C_j$ is abelian and the quotient of the full Galois image
by $C_j$ has order two, its closed derived subgroup is
$\{u/\bar u:u\in C_j\}$.  This image is compact, so no further
closure changes it.  Kronecker--Weber identifies
$G^{\mathrm{cyc}}$ with the closed derived subgroup of
$G_{\mathbb Q}$; taking the image of a closed derived subgroup
under a compact surjection gives the closed derived subgroup of
the target.  Therefore
\[
 \rho_j(G^{\mathrm{cyc}})=\{u/\bar u:u\in C_j\}.
\]

At every unramified finite place, the map $u\mapsto u/\bar u$
from local units is onto the norm-one local units.  For a split
place this is the explicit map $(a,b)\mapsto(a/b,b/a)$.
For an unramified quadratic local field extension, Hilbert 90
provides a preimage, and multiplying by a base-field uniformizer
power makes it a unit because the valuations have the same value
group.  At the inert prime two of $K_3$, a norm-one unit equal to
one modulo two has a Hilbert-90 unit preimage whose residue is
fixed by the residue-field involution.  That residue is in
$\mathbb F_2^\times=\{1\}$, so the preimage already lies in
$1+2O_{3,2}$.  Conversely such a preimage has a ratio equal to
one modulo two.

At the ramified prime three of $K_3$, a norm-one unit $v$ has
residue $\pm1$ modulo $\pi_3$.  If $v\equiv1$, a Hilbert-90
preimage has even $\pi_3$-valuation: a preimage of odd valuation
has ratio with residue $\pi_3/\bar\pi_3=-1$.  Multiplying by
a rational power of three makes the preimage a unit; multiplying
by $1$ or $-1$ makes its residue one.  Conversely a unit preimage
has ratio congruent to one.  Hence
\[
 \rho_3(G^{\mathrm{cyc}})
 =\{u\in N_3:u_2\equiv1\pmod2,
                     \ u_3\equiv1\pmod{\pi_3}\}=:M_3.
\]

At the ramified prime two of $K_4$, a unit
$u=a+bi\equiv1\pmod{\pi_4^3}$ has even $b$, and therefore
$u/\bar u\equiv1\pmod4$.  Conversely let $v\in N_{4,2}$
be one modulo four.  Put $b_0=(1+v)/2$.  It is a unit, is one
modulo two, and $b_0/\bar b_0=v$.  Among the four residue classes
of units modulo $\pi_4^3$, the two equal to one modulo two are
$1$ and $-1$.  Thus one of $b_0,-b_0$ belongs to
$1+\pi_4^3O_{4,2}$ and has the same ratio.  This proves
$\rho_4(G^{\mathrm{cyc}})=N_4(4)$.

The norm-one subgroup of $(O_3/4)^\times$ which is one modulo
two is $\{1,-1\}$.  For example write a representative as
$1+2a$; its norm is one modulo four exactly when the trace of
$a\in\mathbb F_4$ is zero, i.e. when $a=0$ or $1$.
By DELTA\_ON\_GT A.13, or its explicit four-division formulas,
the sign of $\rho_3(\sigma)$ modulo four is exactly
$\kappa(\sigma)$.  Thus $\rho_3^\circ$ takes values in $N_3(4)$.

The map $M_3\to N_3(4)$ which multiplies by this modulo-four
sign is an isomorphism.  Its inverse sends $v\in N_3(4)$ to
$\varepsilon(v_3)v$, where $\varepsilon(v_3)\in\{\pm1\}$ is
the residue of its three-component modulo $\pi_3$.  Indeed this
product has three-component equal to one modulo $\pi_3$, and
two-component equal to that sign modulo four.  Both compositions
are identity.  This proves surjectivity and the assertion about
the kernel of the normalized representation.\hfill$\square$

\paragraph{D23.C5 (proved; uniform signed integral path comparison).}
In the positive normalization of LM7\_PUSH E.65 and
DELTA\_ON\_GT AD.2, for every
$\sigma\in G_{\mathbb Q^{\rm ab}}$ one has, in CM-compatible
transported bases,
\[
 \mathcal W_\sigma=\rho_4(\sigma),\qquad
 \mathcal T_\sigma=\kappa(\sigma)\rho_3(\sigma).
\]
These are full adelic identities, not merely identities modulo twelve.

\emph{Proof, with the path marking and signed normalization made explicit.}
Use the Legendre family
$\mathscr E_t:Y^2=X(X-1)(X-t)$ on
$Y=\mathbb P^1\setminus\{0,1,\infty\}$, with the fixed positive
peripheral generators and the rational tangent $\overrightarrow{01}$.
The embedded Galois action on $F$ is this normalized tangential
action.  Its geometric monodromy in a symplectic
basis is $x\mapsto U^2$, $y\mapsto L^2$: a positive loop about
either collision of two branch points is the square of the half
twist, whose lift to the double cover is the corresponding Dehn
twist, and the two vanishing cycles have intersection one.  This
gives the signed free-subgroup representation used in AD.2.

At the tangential cusp over $t=0$, the Tate degeneration has a
primitive rank-one vanishing submodule and a rank-one quotient.
The special cubic is $Y^2=X^2(X-1)$; its two branches at the
node have slopes $\pm i$.  Thus its only splitting character is
the quadratic character of $\mathbb Q(i)$, and this is trivial on
$G_{\mathbb Q^{\rm ab}}$.  The vanishing line has cyclotomic
action and the quotient has trivial action.  Both therefore have
action one on that subgroup, including integrally at two.
The actual arithmetic matrix $V^{\rm geom}_\sigma$ at the cusp
is consequently upper unipotent with lower-right entry one.
Functoriality gives
$q\Phi_\sigma=\operatorname{Ad}(V^{\rm geom}_\sigma)q$ on $F$.
This already identifies it with $\mathcal V_\sigma$: an upper
unipotent matrix commuting with $U^2,L^2$ is the identity over
every $\mathbb Z_\ell$, including $\ell=2$, because those rings
are domains.  The normalized intertwiner of AD.1 has the same
covariance on $F$, so $V^{\rm geom}_\sigma=\mathcal V_\sigma$.

Here is the independent orbifold identification of the paths.
The quotient stack $[Y/S_3]$ has geometric fundamental group
$B=C_2\amalg C_3$ and pure subgroup $F$.  The transposition
interchanging $0$ and $1$ fixes $p=1/2$; its isotropy is the
class of $s$.  Choose an ordinary geometric path $\xi$ in $Y$
from the tangential base point to $p$ so that the transported
orbifold isotropy is the named $s$.  Such a choice exists: all
order-two factor subgroups are conjugate in $B$, and conjugators
preserving the prescribed transposition in $S_3$ lie in
$F\langle s\rangle$, so the conjugation may be made pure.
For the order-three point one uses one of
$p=(1\pm\sqrt{-3})/2$ and its positive local isotropy, and makes
the same choice with the named $r$.  These two points account
for the two possible pure conjugacy classes; select the one
corresponding to $r$.  Both points and their coordinate-permutation
isotropies are fixed by $G_{\mathbb Q^{\rm ab}}$.

Paths are composed as maps of fibre functors: $\xi$ goes from
the base point to $p$, so
\[
 c_\sigma=(\sigma\xi)^{-1}\xi\in F
\]
is a loop at the base point.  If $s_p$ denotes the isotropy at
$p$ and $s=\xi^{-1}s_p\xi$, then
\[
 \Phi_\sigma(s)=(\sigma\xi)^{-1}s_p(\sigma\xi)
                  =c_\sigma s c_\sigma^{-1}.
\]
The same formula holds with $r$.  The actual arithmetic action
on this orbifold is the source's normalized extension: it agrees
on $F$ with the given tangential action, fixes $\tau$ when the
multiplier is one, and its extension is unique since the
centralizer of $F$ in $B$ is trivial.  Uniqueness of the pure
torsion frames therefore gives $c_\sigma=\beta_\sigma$ or
$c_\sigma=A_\sigma$, respectively.

For the remaining linear calculation use the Tate sheaf only on
$Y$, where the Legendre family is an actual family; no descent of
that sheaf through the projective modular stack is asserted.
If $\rho_p(\sigma)$ is its arithmetic matrix at the ordinary
point $p$, functoriality on the path gives
\[
 q(\sigma\xi)=\rho_p(\sigma)q(\xi)\mathcal V_\sigma^{-1}.
\]
Define $M_p=q(\xi)^{-1}\rho_p(\sigma)q(\xi)$ in the
transported integral basis.  Taking the inverse and multiplying
by $q(\xi)$ shows, with the order fixed,
\[
 q(c_\sigma)=\mathcal V_\sigma M_p^{-1}.
\]
Consequently
\[
 q(\beta_\sigma)M_p(\sigma)=\mathcal V_\sigma
 \quad\hbox{or}\quad
 q(A_\sigma)M_p(\sigma)=\mathcal V_\sigma.
\]
The matrices $M_p$ are exactly $\mathcal W_\sigma$ and
$\mathcal T_\sigma$.  The path cocycle is an actual element of
$F$, so its signed value is the already specified representation
$x\mapsto U^2,y\mapsto L^2$.  There is no independent scalar
lift to choose.  The calculation holds on every finite torsion
sheaf with compatible transition maps, hence on the full integral
Tate module at every prime.

The harmonic fixed point for the interchange is $t=1/2$.
Putting $X=X'+1/2$ in its equation gives
$Y^2=X'^3-X'/4$, the quadratic twist of $E_4$ by $1/2$.
Its twisting square root belongs to $\mathbb Q^{\rm ab}$, so
its representation on $G_{\mathbb Q^{\rm ab}}$ is $\rho_4$.
For the equianharmonic point $t=(1+\sqrt{-3})/2$, since
$t^2-t+1=0$, the substitution
$X=X'+(1+t)/3$ gives
\[
 Y^2=X'^3+\frac{\sqrt{-3}}9
      =X'^3+\left(\frac1{\sqrt{-3}}\right)^3.
\]
This is the quadratic twist of $E_3$ by $d=1/\sqrt{-3}$.
If a square root of $d$ is chosen, its quotient by $r_0$ lies
in $\mathbb Q^{\rm ab}$ (its square is $-2i/3$, and one may
use $\sqrt{2/3}\,\zeta_8^{-1}$).  Thus its twisting character
on $G_{\mathbb Q^{\rm ab}}$ is $\kappa$, as claimed.
The conjugate equianharmonic point changes the twist by $-1$,
whose square root lies in $\mathbb Q^{\rm ab}$, and yields the
same conclusion.  Choose the transported CM identifications
so that their centralizers are those of $S$ and $R_0$, in the
signed convention of AD.2.  This gives the asserted exact
integral matrix identities in those bases.\hfill$\square$

\paragraph{D23.C6 (proved from D23.C4--C5 and the stipulated inputs).}
The signed adelic frame pair has exact images
\[
 (\mathcal T,\mathcal W)(H)=N_3(4)\times N_4(4),
\]
\[
 (\mathcal T,\mathcal W)(G_{\mathbb Q^{\rm ab}})
 =\{(t,w):\epsilon_3(t\bmod3)=\epsilon_2(w\bmod3)\}.
\]
Every continuous relative abelian character of $G$ whose
restriction to $H$ factors through the full adelic frame pair
factors through $\Delta$.  In particular its quadratic possibilities
are exactly $0$ and $\Delta$.

\emph{Proof.}
D23.C4 and D23.C5 show that the arithmetic pair projects onto
both normalized tori, and that their kernels on
$G_{\mathbb Q^{\rm ab}}$ cut out $K_3^{\rm ab}$ and
$K_4^{\rm ab}$.  DELTA\_ON\_GT A.12 gives their exact
intersection
\[
 K_3^{\rm ab}\cap K_4^{\rm ab}
       =\mathbb Q^{\rm ab}(3^{1/4}),
 \qquad [\mathbb Q^{\rm ab}(3^{1/4}):\mathbb Q^{\rm ab}]=2.
\]
The compact Galois image is therefore the fibre product over a
quotient of order two, and has index two in the torus product.
E.65's already checked comparison modulo twelve puts it in the
equal-parity subgroup displayed above.  That subgroup itself
has index two, so they are equal.

AD.2 puts the full gentle frame image inside the normalized torus
product.  It contains the arithmetic index-two subgroup and the
E.58 element, whose parity pair is $(a_3,a_2)=(1,0)$.
A subgroup containing an index-two subgroup and one element
outside it is the entire group.  This proves the first equality.
The quotient of the full frame image by the arithmetic image is
thus $C_2$, with quotient character $a_2+a_3$.  Any continuous
abelian character satisfying the factorization and arithmetic
vanishing hypotheses descends to that quotient.  Finally
$G=HG_{\mathbb Q}$ makes its full extension unique, and E.65
identifies the nonzero quadratic one with $\Delta$.
\hfill$\square$

\paragraph{D23.C7 (proved; precise failure of a place-two frame factorization).}
There is an actual $\sigma\in G_{\mathbb Q^{\rm ab}}$ with
\[
 a_2(\sigma)=a_3(\sigma)=1,
 \qquad
 q_2(A_\sigma)=q_2(\beta_\sigma)
     =\mathcal V_{\sigma,2}
     =\mathcal T_{\sigma,2}
     =\mathcal W_{\sigma,2}=I.
\]
Here $q_2$ means the full $2$-adic representation, not reduction
modulo two.  Thus neither parity is a function solely of the pair
of raw signed $2$-adic frame matrices, or of the corrected pair,
even if the $2$-adic intertwiner is also supplied.  This is a
statement about those precise congruence data; it makes no claim
about every possible meaning of a ``local'' geometric invariant.

There is also an actual $h\in H$
with $\Delta(h)=1$, $a_3(h)=0$, and $a_2(h)=1$, whose entire
corrected adelic pair has
\[
 \mathcal T_h=I,\qquad
 \mathcal W_{h,3}=S,\qquad
 \mathcal W_{h,\ell}=I\quad(\ell\ne3).
\]
In particular its corrected $2$-adic pair is the identity pair.

\emph{Proof.}
Fix an embedding into a decomposition group at three and write
$I_3$ for inertia.  Its wild subgroup is pro-three and its tame
quotient is $\prod_{\ell\ne3}\mathbb Z_\ell(1)$.  A pro-two
Sylow subgroup $P$ therefore maps isomorphically to the
$\mathbb Z_2(1)$ factor: its intersection with the wild subgroup
is trivial, and the image of a Sylow subgroup under a compact
surjection is a Sylow subgroup.  These assertions can also be
obtained by taking compatible Sylow subgroups in the finite
inertia quotients.  Choose a topological generator $\tau$ of $P$.
The extension generated by $r_0$ is tamely ramified of degree
four after adjoining $\mu_4$, which is unramified at three;
indeed $(2r_0)^4=12$ is Eisenstein at three.  It follows that
$\tau(r_0)=i^{\pm1}r_0$.  Put $\sigma=\tau^2$.  Then
$\sigma(r_0)=-r_0$.

For every $\ell\ne3$, the $\ell$-adic cyclotomic character is
unramified at three.  The image of the pro-two group $P$ under
$\chi_3$ is contained in $\{1,-1\}\subset\mathbb Z_3^\times$,
so its square is one.  Consequently $\chi(\sigma)=1$ and
$\sigma\in G_{\mathbb Q^{\rm ab}}$.  R4's exact finite arithmetic
comparison identifies both parities with the sign on $r_0$;
therefore $a_2(\sigma)=a_3(\sigma)=1$.

The two Legendre fibres used in the path proof of D23.C5 have
good reduction at three.  The harmonic fibre at $t=1/2$ has
discriminant $16t^2(1-t)^2=1$.  For either equianharmonic fibre,
$t(1-t)=1$, so the discriminant is $16$, a unit at three.
The latter fibre is defined over $K_3$, which $\sigma$ fixes;
$\sigma\in I_3\cap G_{K_3}$ is in its local inertia group.
The good-reduction implication of Serre--Tate,
\emph{Good reduction of abelian varieties}, Theorem 1,
\emph{Annals of Mathematics} 88 (1968), p.493, makes both full
$2$-adic Tate actions of $\sigma$ the identity.  The integral
path identity in D23.C5 therefore gives
$\mathcal T_{\sigma,2}=\mathcal W_{\sigma,2}=I$.

The actual tangential action on the free pro-two group is
unramified at three.  One explicit source for exactly this
automorphism statement is Hain--Matsumoto,
\emph{Weighted completion of Galois groups and Galois actions on
the fundamental group of $\mathbb P^1-\{0,1,\infty\}$},
\emph{Compositio Mathematica} 139 (2003), Section 8.5,
the paragraph before Theorem 8.7 (take $N=1$ and $\ell=2$).
Thus $\Phi_\sigma$ acts identically on $F^{(2)}$.
The image $q_2(F)$ lies in the principal congruence subgroup
modulo two of $\mathrm{SL}_2(\mathbb Z_2)$, which is pro-two.
Hence $q_2\Phi_\sigma=q_2$.  The covariance says that the
upper unipotent $\mathcal V_{\sigma,2}$ commutes with $U^2,L^2$,
forcing it to be $I$.  The corrected-frame definitions then give
$q_2(A_\sigma)=q_2(\beta_\sigma)=I$ as well.

The identity element has all those same matrices and parity zero,
which proves the asserted failure of factorization.  For the final
assertion, the displayed adelic pair belongs to
$N_3(4)\times N_4(4)$: its component over two is the identity,
and $S$ has norm one at three.  D23.C6 supplies an actual $h$
with that value.  Its parities and discrepancy have the stated
values by their definitions.\hfill$\square$

\paragraph{D23.C8 (proved; an explicit smaller gate for B.40).}
Let
\[
 N=\ker\bigl((\mathcal T,\mathcal W):H\longrightarrow
                   N_3(4)\times N_4(4)\bigr),
 \qquad
 Q=G/\overline{[G,G]G^2G_{\mathbb Q}}.
\]
Then, as a set product,
\[
 \ker\Delta=N G_{\mathbb Q}.
\]
If $\bar\Delta:Q\to\mathbb F_2$ is the induced nonzero
character, its kernel is exactly the image of $N$.  Consequently
the assertion that all relative continuous quadratic characters
are $0,\Delta$ is equivalent to the concrete vanishing condition
\[
 N\subseteq\overline{[G,G]G^2G_{\mathbb Q}}.
\]
This computes the entire remaining kernel in terms of elements
whose two full adelic frames are identity.  It does not assert
that this subgroup has trivial image in $Q$.

For the two particular characters $d$ and $\Delta\circ D$,
the known actual kernel witness of D23.C3 gives the further
equivalences
\[
 \begin{split}
 d\ne0&\quad\Longleftrightarrow\quad
      \exists n\in N:\ a_2(Dn)=1,\\
 \Delta\circ D\ne0&\quad\Longleftrightarrow\quad
      \exists n\in N:\ \Delta(Dn)=1.
 \end{split}
\]
In particular both characters vanish if and only if
$(a_2,a_3)(D(N))=\{(0,0)\}$.

\emph{Proof.}
Write $g=h\sigma$ with $h\in H$ and $\sigma\in G_{\mathbb Q}$.
If $g\in\ker\Delta$, then $h\in\ker\delta$, so its two
adelic frame values have equal parity.  D23.C6 gives an actual
$\tau\in G_{\mathbb Q^{\rm ab}}$ with exactly the same entire
frame pair as $h$.  The two frame maps are homomorphisms on
$H$, hence $n=h\tau^{-1}\in N$ and $g=n(\tau\sigma)$.
The converse follows from $\Delta|_N=\Delta|_{G_{\mathbb Q}}=0$.
No normality assertion about $N$ in all of $G$ is needed.

The closed subgroup in the definition of $Q$ is the inverse
image of the image of $G_{\mathbb Q}$ in the maximal elementary
abelian pro-two quotient, so the quotient is well defined.
The preceding set equality shows
$\ker\bar\Delta=\operatorname{im}(N\to Q)$.  Continuous
quadratic characters separate points of an elementary abelian
pro-two group.  Thus they are exactly $0,\bar\Delta$ if and
only if this kernel is trivial.  Pullback gives the asserted
equivalence on $G$.  In the integrated relative-abelianization
statement P.12, $Q=K_G/2K_G$.

On $N$, one has $a_2=a_3=0$, so
$d(n)=a_2(Dn)$.  Any relative character vanishing on $N$
vanishes on $\ker\Delta=N G_{\mathbb Q}$ and is therefore
$0$ or $\Delta$.  The actual $k$ of D23.C3 has
$d(k)=(\Delta\circ D)(k)=0$ but $\Delta(k)=1$;
it excludes the second possibility for both characters.
This proves the two equivalences.  Their right sides both
fail exactly when $a_2(Dn)=0$ and
$a_2(Dn)+a_3(Dn)=0$ for every $n\in N$, equivalently when
both descended parities vanish there.\hfill$\square$


\section{Finite comparison gates and failed shortcuts}
% DELTA23 agent /root/pentagon. New finite-quotient obstruction, not A.84 decision.
\subsection{Independent central flips in all four-point forgetting quotients}

\paragraph{D23.P1 (independent central coordinates).}
Let $P_n=\widehat\Gamma_{0,n}$, $n\geq4$. For each four-element
subset $I$ of the labelled punctures, let $\pi_I:P_n\to P_4$ be
forgetting the other punctures, with any orientation-preserving marking
of the target. Use the signed modular quotient of LM7\_PUSH E.58.11,
\[
 q:P_4\cong\widehat F_2\longrightarrow
 V_4\times\operatorname{SL}_2(\mathbb F_3),
\]
whose level-twelve projective notation sends $(1,-I)$ to $[5I]$.
Write $z=(1,-I)$. Then the simultaneous map has exact image
\[
 (q\pi_I)_I(P_n'')=\prod_{|I|=4}\langle z\rangle.
 \tag{D23.P1.1}
\]
Here the primes denote closed derived subgroups. In particular the
central signs in distinct four-point quotients can be prescribed
independently by one actual element of $P_n''$.

\begin{proof}
The second derived subgroup of $V_4\times\operatorname{SL}_2(\mathbb F_3)$
is $\{1\}\times\{I,-I\}$. The relevant explicit matrices are
\[
 x=\begin{pmatrix}1&2\\0&1\end{pmatrix},\qquad
 y=\begin{pmatrix}1&0\\1&1\end{pmatrix}.
\]
They generate $\operatorname{SL}_2(\mathbb F_3)$; its commutator
subgroup is the quaternion group of order eight and that group's
commutator subgroup is $\{I,-I\}$. These finite assertions, as well
as the matrix calculation below, are enumerated in the certificate
\path{agents/pentagon/check_four_point_central_flips.py}.
This gives the inclusion from left to right in (D23.P1.1).

Fix $I=\{i_1,i_2,i_3,i_4\}$. Choose a disk neighbourhood of a
chain joining these four punctures, with no other punctures in it.
Let $a,b,c\in P_n$ be the positive twists about the consecutive
pairs $(i_1,i_2),(i_2,i_3),(i_3,i_4)$ in that disk, and set
\[
 w_I=[[a,b],a[b,c]a^{-1}],\qquad [u,v]=uvu^{-1}v^{-1}.
 \tag{D23.P1.2}
\]
This lies in $P_n''$. If any puncture of $I$ is forgotten, one of
the two inner commutators is trivial: forgetting $i_1$ kills $a$,
forgetting $i_4$ kills $c$, and forgetting $i_2$ or $i_3$ kills
both twists in at least one inner commutator. A twist about a disk
containing at most one retained puncture is trivial in the punctured
sphere mapping class group. Therefore $\pi_J(w_I)=1$ for every
four-element subset $J\ne I$.

After forgetting all punctures outside $I$, the curves defining $a$
and $c$ are complementary and isotopic in the four-punctured sphere.
Thus $\pi_I(c)=\pi_I(a)$, and in the chosen standard marking their
signed modulo-three images are $q_3(a)=q_3(c)=x$, $q_3(b)=y$.
Direct multiplication gives
\[
 [x,y]=\begin{pmatrix}1&2\\2&2\end{pmatrix},\quad
 x[y,x]x^{-1}=\begin{pmatrix}1&1\\1&2\end{pmatrix},\quad
 [[x,y],x[y,x]x^{-1}]=-I.
\]
The modulo-four coordinate is trivial because $w_I\in P_n''$ and
the modulo-four target is abelian. Hence $q\pi_I(w_I)=z$.
Any other orientation-preserving marking is obtained by a mapping
class on the four-punctured sphere. Its conjugation action preserves
the modular quotient and its central element $z$; equivalently one
uses the quotient $\Gamma_{0,[4]}\to C_2\amalg C_3$ and its normal
level-twelve kernel. Thus the value $z$ does not depend on the
chosen chain marking. The images of the $w_I$ are the coordinate
basis of the product in (D23.P1.1). Their finite products give every
prescribed tuple. This proves the reverse inclusion.
\end{proof}

\paragraph{D23.P2 (finite simultaneous frames do not remove the cubic flip).}
Let a finite subgroup $H$ of $\widehat\Gamma_{0,[n]}$ act on the
four-element subsets by its label permutation, and retain the finite
quotient obtained from all the preceding four-point modular maps
and the label permutation map. For every $H$-orbit $\mathcal O$ of
four-element subsets there is an element of the image of $P_n''$
that centralizes the image of $H$ and whose four-point coordinates
are $z$ on $\mathcal O$ and $1$ elsewhere.
Consequently, whenever a finite simultaneous frame for $H$ is used to
read an equianharmonic frame sign at a subset $I\in\mathcal O$,
this centralizer changes that cubic sign by one. At any four-point
harmonic coordinate it leaves the harmonic parity unchanged. It also
leaves the abelianization of the simultaneous pure frame unchanged.

\begin{proof}
Conjugation by an element of the full mapping class group transports
forgetting maps according to its permutation of the labels. On the
second-derived central coordinates its effect is exactly that
permutation: $z$ is the distinguished nonidentity element of a group
of order two and is fixed by changes of target marking. By D23.P1,
the orbit-indicator tuple is the image of a product of the $w_I$.
It is fixed by conjugation by every element of $H$, hence centralizes
its image in the stated finite quotient. Multiplying any chosen
simultaneous frame by this element preserves all its finite
conjugator equations.

On the cubic coordinate, $z$ is $-I$ in
$C_{\operatorname{SL}_2(\mathbb F_3)}(R_0)\cong C_6$; its sign
on the rational fixed line is $-1$, so $a_3$ changes by one.
On the harmonic coordinate, $z=-I=S^2$ in
$\langle S\rangle\cong C_4$; the parity of the exponent of $S$
does not change. Trivial coordinates do not change that parity
either. Since the multiplying element lies in $P_n''$, the pure
frame's abelianization is also unchanged.
\end{proof}

\paragraph{Exact scope.}
D23.P2 is a statement about the specified finite quotient and finite
conjugator equations. It does not construct a different actual
unique frame in the full profinite group, nor a different gentle
or GT element. In particular, the additional affine five-point
quotient used in DELTA\_ON\_GT FP.32 records more than these
four-point projections and can exclude this ambiguity by its
relation $a_3=\kappa_d\bmod2$. No claim is made that D23.P2
survives that additional quotient. This computation therefore
blocks one simultaneous-frame shortcut but does not decide A.84.

% DELTA23: /root/pentagon, 2026-11-21 CET. Full proof draft.
\subsection{A small affine evaluation detecting the cubic frame sign}

\paragraph{D23.P3 (an order-1944 free-group detector).}
Put $k=\mathbb F_9=\mathbb F_3[i]$, $i^2=-1$, and use affine
multiplication $(v,M)(w,N)=(v+Mw,MN)$ on $k^2\rtimes\operatorname{GL}_2(k)$.
Retain the modular generators $s,r$, $x=(sr)^{-2}$, $y=sxs$ of
LM7\_PUSH E.58 and E.65. Set
\[
 S=\begin{pmatrix}0&1\\-1&0\end{pmatrix},\quad
 R=\begin{pmatrix}0&-1\\1&-1\end{pmatrix},\quad
 U=\begin{pmatrix}1&1\\0&1\end{pmatrix},\quad v=e_1.
\]
There is a continuous homomorphism
\[
 \eta:\widehat B=C_2\amalg C_3\longrightarrow E,
 \qquad \eta(s)=(0,iS),\qquad \eta(r)=(v,R),
 \tag{D23.P3.1}
\]
onto the finite group
\[
 E=k^2\rtimes H,\qquad H=\langle iI\rangle\operatorname{SL}_2(\mathbb F_3),
 \quad |E|=3888.
\]
Its restriction to the pure free group $F=\widehat F(x,y)$ has
image $E_F=k^2\rtimes\operatorname{SL}_2(\mathbb F_3)$, of order
$1944$. The image of $C_E(\eta(r))$ in $H$ is exactly $\langle R\rangle$,
and its order is $27$.

For an actual multiplier-one gentle element $h=(1,f)$, write
\[
 \eta(f)=(d,M_\rho),\quad d=(d_1,d_2)^{\mathsf t},\quad
 M_\rho=\begin{pmatrix}1&\rho\\\rho&1+\rho^2\end{pmatrix},
 \qquad\rho\in\mathbb F_3.
\]
Here $\rho$ is the reduction of the frame parameter in E.65; it is
also read directly as the lower-left entry of the linear part.
Then
\[
 \boxed{\quad (-1)^{a_3(h)}
       =(1-\rho)(1-d_1)-d_2\quad\text{in }k.\quad}
 \tag{D23.P3.2}
\]
The right-hand side belongs to $\{1,-1\}\subset\mathbb F_3$ for
every actual gentle element. Moreover, if $\epsilon=(-1)^{a_3(h)}$,
its full affine evaluation is exactly
\[
 d_1=\frac{1-\rho-\epsilon}{1+i},\qquad
 d_2=(\rho+i)d_1.
 \tag{D23.P3.3}
\]
Thus every such $f$-evaluation lies in an explicitly specified set of six elements of $E_F$.
For the E.58 witness one obtains
\[
 \eta(f_{58})=((1-i,1+i)^{\mathsf t},I),
 \tag{D23.P3.4}
\]
not merely the value of its sign.

\begin{proof}
The matrix $iS$ squares to $I$. In characteristic three,
$(R-I)^2=0$, so $I+R+R^2=0$. These identities prove the two
torsion relations in (D23.P3.1). Hence the homomorphism exists
by the defining free-product presentation, and continuity follows
because its target is finite.

The linear image is $H$: indeed $iSR=iU^{-1}$ and
$(iU^{-1})^3=-iI$, so it contains $iI,S,R$, which generate $H$.
The intersection $\langle iI\rangle\cap\operatorname{SL}_2(\mathbb F_3)$
is $\{I,-I\}$; hence $|H|=4\cdot24/2=48$.
The translation subgroup $N$ of the image is an $H$-invariant
$\mathbb F_3$-subspace of $k^2$. Scalar multiplication by $iI$
makes it a $k$-subspace. The natural representation of
$\operatorname{SL}_2(\mathbb F_3)$ over $k$ is irreducible:
a common invariant line for the two nonidentity unipotent
matrices $U$ and $L=\left(\begin{smallmatrix}1&0\\-1&1\end{smallmatrix}\right)$
would have to be both $ke_1$ and $ke_2$. Thus $N$ is zero or $k^2$.

If $N=0$, the affine image is the graph of a cocycle $c:H\to k^2$.
The central element $z=-I$ acts as $-1$. Comparing $c(gz)$ and
$c(zg)$ gives $2c(g)=(1-g)c(z)$, so
$c(g)=(1-g)t$ with $t=c(z)/2$. In particular this would imply
$v=c(R)\in(1-R)k^2$. But $(1-R)k^2=k(1,-1)^{\mathsf t}$
and $v=e_1$ is outside that line. This contradiction proves
$N=k^2$ and the asserted full image.

The linear images of $x,y$ are $-U^2,-L^2$. Cubing either
produces $-I$, so they generate $\operatorname{SL}_2(\mathbb F_3)$.
The pure image is normal in $E$, since $F$ is normal in $B$.
Its intersection with $k^2$ is again zero or all of $k^2$.
The zero alternative would make the quotient of $E$ by the pure
image have order $3888/24=162$, whereas it is a quotient of
$B/F\cong S_3$. It follows that the pure image is exactly $E_F$.

An element of the centralizer of $R$ in $H$ has the form
$i^jR^m$, $j\in\mathbb Z/4$, $m\in\mathbb Z/3$.
Indeed $C_{\operatorname{SL}_2(\mathbb F_3)}(R)
=\langle-I\rangle\times\langle R\rangle$, and the scalar
extension adds $iI$. For an affine element $(t,i^jR^m)$ to
centralize $(v,R)$ one must have
$(1-R)t=v-i^jR^m v$. In the one-dimensional coinvariant quotient
$k^2/(1-R)k^2$, the nonzero class of $v$ is fixed by $R$.
Consequently $i^j=1$. Conversely all powers of $(v,R)$ centralize
it, so the linear image of its centralizer is exactly
$\langle R\rangle$. Each of its three linear values has
$|\ker(1-R)|=9$ translation lifts. This proves the order $27$.

We next justify the signed comparison between the linear part of
$\eta$ and the supplied $q_{12}$, rather than identifying the
representations without a sign convention. There is a character
$\chi_4:\operatorname{SL}_2(\mathbb Z/4)\to\langle i\rangle$
with $\chi_4(U)=\chi_4(L)=-i$, $\chi_4(S)=i$,
$\chi_4(R)=1$, and $\chi_4(-I)=-1$. Its finite definition is
verified by enumerating all $48$ matrices and propagating the
values along multiplication by $U,L$; the certificate checks
consistency on every such edge. Thus
\[
 [(g_4,g_3)]\longmapsto\chi_4(g_4)g_3
\]
is a well-defined homomorphism from
$(\operatorname{SL}_2(\mathbb Z/4)\times
\operatorname{SL}_2(\mathbb F_3))/\langle(-I,-I)\rangle$ to $H$.
It sends the modular $s,r$ to $iS,R$, precisely the linear part
of $\eta$. For the actual pure frame $A_h$, E.65 proves
$q_4(A_h)=I$, so its linear image under $\eta$ is exactly
$q_3(A_h)$. For $f\in F'$ the modulo-four coordinate is also
trivial. E.65.9 therefore gives its displayed linear part
$M_\rho$, and $V=U^{2\rho}$ is the reduction of E.65's shear.

Write $\eta(A_h)=(t,A_0)$ and $T=A_0^{-1}V$. By E.65,
$T\in\langle-I\rangle\times\langle R\rangle$, and its sign
on the fixed line is $\epsilon=(-1)^{a_3(h)}$. The exact frame
identity $h(r)=A_hrA_h^{-1}=f^{-1}r$ gives, on translations,
\[
 M_\rho^{-1}(v-d)
       =A_0v+(1-A_0RA_0^{-1})t.
\]
Multiplying by $V^{-1}$ and passing to $R$-coinvariants makes
the last term zero. The row vector $\ell=(1,1)$ realizes this
coinvariant quotient; it satisfies $\ell R=\ell$, $\ell v=1$,
and $\ell T^{-1}=\epsilon\ell$. We obtain
\[
 \ell V^{-1}M_\rho^{-1}(v-d)=\epsilon.
\]
Direct multiplication gives
$V^{-1}M_\rho^{-1}=\left(\begin{smallmatrix}1&0\\-\rho&1\end{smallmatrix}\right)$,
which proves (D23.P3.2).

Finally $h(s)=sf$ has order two. The translation part of
$\eta(sf)^2=1$, with the displayed matrices, is equivalent to
$d_2=(\rho+i)d_1$. Substitution in (D23.P3.2) gives
$\epsilon=1-\rho-(1+i)d_1$, proving (D23.P3.3).
E.58 has $\rho=0$ and $a_3=1$, which gives (D23.P3.4).
The exact finite program \path{agents/pentagon/affine_cubic_detector.py}
independently enumerates both images, the centralizer, and all
six affine evaluations, and checks the coinvariant calculation
for all $486$ combinations of sign, cubic power, and translation.
\end{proof}

\paragraph{D23.P4 (an explicit finite evaluation gate for A.84).}
Retain $p=(123)$ and $t=(1234)$ in $S_4$. The simultaneous map
\[
 F\longrightarrow E_F\times S_4,\qquad
 x\longmapsto(\eta(x),p),\quad
 y\longmapsto(\eta(y),t)
 \tag{D23.P4.1}
\]
is onto a group of order $46656$. For an actual gentle element
$h=(1,f)$, write its $S_4$-evaluation as
$f(p,t)=t^{2\nu_4}p^{a_f}$, as in DELTA\_ON\_GT FP.37.
For actual full GT elements, the existing pentagon consequences
FP.33 and FP.37 give $a_f=\rho$ and $a_2=\nu_4$.
Consequently A.84 is equivalent to the following concrete finite
comparison on actual full-GT source residues:
\[
 \boxed{\quad
 (1-\rho)(1-d_1)-d_2=(-1)^{\nu_4}.
 \quad} \tag{D23.P4.2}
\]
After the two previously proved pentagon consequences, there are
only twelve possible joint evaluations: $\rho\in\mathbb F_3$
and $(a_3,\nu_4)\in\mathbb F_2^2$. Six obey (D23.P4.2) and
six fail it. The remaining question is whether any of the latter
six is the image of an actual full-GT element. This is not an
assertion that any particular finite candidate lifts to GT.

\begin{proof}
Since $-I$ acts as $-1$ on $k^2$ and $2$ is invertible, commutators
with $-I$ generate the translation subgroup. Hence
$E_F^{\mathrm{ab}}\cong\operatorname{SL}_2(\mathbb F_3)^{\mathrm{ab}}
\cong C_3$. On the other hand $S_4^{\mathrm{ab}}\cong C_2$.
Both groups are finite solvable. Any nontrivial common quotient
would be solvable and would have a nontrivial abelian quotient;
its abelianization would be a quotient of both $C_3$ and $C_2$,
which is impossible. Both projections of (D23.P4.1) are onto,
so Goursat's lemma proves that its image is the full product.
Its order is $1944\cdot24=46656$.

D23.P3 computes the cubic bit from its first coordinate, while
FP.37 computes the harmonic bit on full GT from its second
coordinate. Thus their equality is exactly (D23.P4.2).
The six possible first coordinates are (D23.P3.3); the second
coordinates have the six values $t^{2\nu_4}p^{a_f}$.
Imposing $a_f=\rho$ leaves the stated twelve candidates and
partitions them according to the displayed sign equality.
The reduction from arbitrary multiplier to multiplier one is
LM7\_PUSH E.65.6, using Galois vanishing, exactly as in
DELTA\_ON\_GT AO.1. This proves the equivalence with A.84.
\end{proof}

\subsection{The cubic bit in an order-216 marked affine quotient}

\paragraph{D23.P10 (a nine-fold smaller exact main gate).}
Let $V=\mathbb F_3^2$, $E_0=V\rtimes\operatorname{SL}_2(\mathbb F_3)$,
and use affine multiplication $(v,M)(w,N)=(v+Mw,MN)$.
There is a continuous surjection from the same marked pure free
group $F=\widehat F(x,y)$ as in E.65,
\[
 \xi:F\twoheadrightarrow E_0,\qquad
 \xi(x)=\left(0,\begin{pmatrix}-1&1\\0&-1\end{pmatrix}\right),\qquad
 \xi(y)=\left(e_1,\begin{pmatrix}-1&0\\-1&-1\end{pmatrix}\right).
\tag{D23.P10.1}
\]
The two linear matrices are $-U^2,-L^2$ in the source's positive
convention. For every actual multiplier-one gentle element
$h=(1,f)$ one has the exact formula
\[
 \boxed{\quad
 \xi(f)=\left(
   \begin{pmatrix}e\\\rho+(\rho+1)e\end{pmatrix},
   \begin{pmatrix}1&\rho\\\rho&1+\rho^2\end{pmatrix}
  \right),\qquad e=a_3(h)\in\{0,1\}\subset\mathbb F_3.
 \quad}
\tag{D23.P10.2}
\]
Thus the cubic bit is literally the first translation coordinate
of this order-$216$ finite evaluation of $f$.

The joint marked map
\[
 F\longrightarrow E_0\times S_4,\qquad
 x\longmapsto(\xi(x),(123)),\qquad
 y\longmapsto(\xi(y),(1234))
\tag{D23.P10.3}
\]
is onto, of order $5184$. Its image on $F'$ has order $864$ and
is exactly
\[
 (V\rtimes Q_8)\times A_4.
\tag{D23.P10.4}
\]
For actual full-GT elements, the supplied pentagon consequences
$a_f=\rho$ and $a_2=\nu_4$ leave precisely the same twelve
candidate states as D23.P4, now in a nine-fold smaller source
quotient. In these coordinates A.84 is exactly the assertion
that none of the six states with
\[
                  e\ne\nu_4
\tag{D23.P10.5}
\]
is attained by an actual full-GT element. The number of possible
mixed states has not decreased; the finite group realizing the
exact comparison has.

Among quotients of $E_0$, no proper quotient can retain this
detector on the supplied E.58 witness: every nontrivial normal
subgroup of $E_0$ contains $V$, whereas
\[
       \xi(f_{58})=((1,1)^{\mathsf T},I)\ne1.
\tag{D23.P10.6}
\]
This is a minimality statement within the quotient family of
$E_0$, not a claim of global minimality among all finite detectors.

\begin{proof}
First construct the map directly from the free presentation.
Write $X=\xi(x)$ and $Y=\xi(y)$. Their linear images generate
$\operatorname{SL}_2(\mathbb F_3)$: cubing either gives $-I$,
after which $U^2,L^2$ and hence $U,L$ are available. Moreover
\[
              [X^3,Y]=(e_1,I).
\]
Indeed $X^3=(0,-I)$, and conjugation by it negates the
translation of $Y$ while leaving its linear part unchanged;
the resulting commutator translation is $-2e_1=e_1$ in
characteristic three. Conjugating this translation by $Y$
gives the translation $(-1,-1)^{\mathsf T}$. Together these
span $V$. Thus the image is all of $E_0$, of order $9\cdot24=216$.

We next compare the marked map with D23.P3; this step fixes
the actual cubic sign rather than guessing from an isomorphic
finite group. On $E_F=\mathbb F_9^2\rtimes\operatorname{SL}_2(\mathbb F_3)$,
real-part projection of the translation is a group homomorphism,
because every linear matrix has entries in $\mathbb F_3$.
D23.P3's two pure generators have real translation $(1,1)$
and linear parts $-U^2,-L^2$. Change affine origin by
$t=(0,2)^{\mathsf T}$, which changes a translation $d$ of
linear part $M$ into $d-(I-M)t$. The two resulting translations
are $0,e_1$, precisely (D23.P10.1).

For an actual gentle $f$, D23.P3 gives
\[
 d_1=\frac{1-\rho-\epsilon}{1+i},\qquad
 d_2=(\rho+i)d_1,\qquad \epsilon=(-1)^{a_3(h)}.
\]
Since $(1+i)^{-1}=2+i$, writing
$c=1-\rho-\epsilon\in\mathbb F_3$ gives
\[
 \operatorname{Re}d_1=2c,\qquad
 \operatorname{Re}d_2=(2\rho-1)c.
\]
The origin correction is $(I-M_\rho)t=(\rho,\rho^2)^{\mathsf T}$.
If $(e,b)^{\mathsf T}$ denotes the corrected translation, then
\[
 e=2(1-\rho-\epsilon)-\rho,\qquad \epsilon=1+e.
\]
Thus $e=0$ when $a_3=0$ and $e=1$ when $a_3=1$.
Substituting $c=-\rho-e$ into the second coordinate gives
\[
 b=(2\rho-1)(-\rho-e)-\rho^2=\rho+(\rho+1)e.
\]
This proves (D23.P10.2), and E.58's $\rho=0,a_3=1$ gives
(D23.P10.6). The use of E.58 here is the supplied actual
witness; its absent infinite lifting input is not independently
reproved by this finite calculation.

Because $-I$ acts as $-1$ on $V$, commutators with it generate
all translations. Therefore $E_0'=V\rtimes Q_8$ and
$E_0^{\mathrm{ab}}\simeq C_3$. Since
$S_4^{\mathrm{ab}}\simeq C_2$ and both groups are solvable,
they have no nontrivial common quotient: such a quotient would
have a nontrivial abelian quotient common to $C_3,C_2$.
Goursat's lemma proves surjectivity of (D23.P10.3). Taking
derived subgroups proves (D23.P10.4). The further derived orders,
also verified in the exact certificate, are
\[
 |E_0^{(j)}|:216,72,18,9,1;\qquad
 |(E_0\times S_4)^{(j)}|:5184,864,72,9,1.
\]
On full GT, supplied FP.33 and FP.37 identify the marked
$S_4$ value as $t^{2\nu_4}p^\rho$ and identify $a_2=\nu_4$.
Together with (D23.P10.2), their equality criterion is exactly
(D23.P10.5). Reduction from arbitrary multiplier to multiplier
one is the supplied E.65.6 normalization, as in D23.P4.

Finally let $N\lhd E_0$ be nontrivial. If $N\cap V\ne1$,
the irreducibility of the natural $\operatorname{SL}_2(\mathbb F_3)$
module implies $V\subset N$. If $N\cap V=1$, then
$[N,V]\subset N\cap V$ forces $N$ to centralize $V$.
The centralizer of all translations in $E_0$ is $V$ itself,
since the linear action is faithful. Hence $N\subset V$,
contradicting $N\cap V=1$ and $N\ne1$. Every proper quotient
therefore kills the nonidentity translation in (D23.P10.6),
as claimed.

The exact program \path{agents/pentagon/cubic_detector_216.py}
checks the marked commutator, enumerates the affine image,
checks all six projected values from D23.P3, and independently
computes the order and derived series of the joint action on
nine affine points plus four permutation points.
\end{proof}

\subsection{The cubic detector does not extend through one natural rank-four module}
\paragraph{D23.P5 (a checked obstruction to the twisted symplectic affine route).}
Let $k=\mathbb F_9=\mathbb F_3[i]$, and let $t_1,t_2,t_3,t_4$ be
D23.T1's rank-four symplectic transvections. On the standard subgroup
$B_4\subset B_5$, put $A_j=-it_j$, $1\leq j\leq3$. Every
continuous affine lift of this linear representation has a common
fixed point; equivalently every continuous cocycle
$B_4\to k^4$ is a coboundary. In particular the cubic affine detector
D23.P3, on the nondegenerate span $W=\langle e_1,e_2\rangle$ of
the first two braid generators, does not extend through this
rank-four module. The analogous extension to $B_5$ also fails.

\begin{proof}
Write $z_j$ for the translation of an affine lift of $A_j$.
The matrix $I-A_1$ is invertible, because $t_1$ is unipotent and
$1+i\ne0$. A change of affine origin therefore makes $z_1=0$.
The distant commutation of the first and third generators now gives
$(I-A_1)z_3=0$, hence $z_3=0$. The two remaining braid equations
on translations give
\[
 (A_1-I-A_2A_1)z_2=0,\qquad
 (I+A_2A_3-A_3)z_2=0.
\]
The matrices are, respectively,
\[
 K_1=\begin{pmatrix}
 -i&1-i&0&0\\
 -1&-1-i&1&0\\
 0&0&-i&0\\
 0&0&0&-i
 \end{pmatrix},\qquad
 K_2=\begin{pmatrix}
 i&0&0&0\\
 1&1+i&-1&-1\\
 0&1-i&i&-1+i\\
 0&0&0&i
 \end{pmatrix}.
\]
Rows three and four of $K_1$ force the third and fourth coordinates
of $z_2$ to vanish; row one of $K_2$ then forces its first coordinate
to vanish, and row one of $K_1$ forces the second to vanish.
All translations are zero in the new origin, which proves the
coboundary assertion.

The restriction of $k^4$ to $B_3=\langle\sigma_1,\sigma_2\rangle$
splits as $W\oplus W^\perp$, where
$W^\perp=\langle e_1+e_3,e_4\rangle$. Projection of a coboundary
to $W$ is again a coboundary. D23.P3's cocycle is not a coboundary
on $W$: its value on $r$ is $e_1\notin(1-R)W$.
This excludes the proposed extension. The exact solver in
\path{agents/pentagon/twisted_affine_extension.py} checks the
matrices and the full affine braid equations independently.

For $B_5$, its central full twist acts on this twisted module by
$(-i)^{20}C^5=-I$, using D23.T2. For any cocycle $c$, comparison
of its value on $gz$ and $zg$, with this central $z$, gives
$2c(g)=(1-g)c(z)$, so it is again a coboundary. This gives an
independent conceptual obstruction in the five-strand case.
\end{proof}

\paragraph{Scope.}
This rules out an affine lift through the specified rank-four
linear representation. It does not rule out higher-dimensional,
non-affine, relative, or Brunnian representations, and it is not
a counterexample to the pentagon or to the desired comparison.

% DELTA23 /root/five_level_six. An attempted main-gate detector that collapses.

\subsection*{D23.F11. The new symplectic covariance has only the old shear}
Let $t_h\in\widehat{\mathbb Z}$ be the full signed
three-strand shear of DELTA\_ON\_GT AD.1 for
$h\in\mathcal H=\ker\lambda\subset\widehat{GT}_0$.
Use the chain form and the representation $q$ of D23.T1, and
write $w=e_1+e_3$. Define
\[
 V(t)=\begin{pmatrix}
 1&t&0&t\\0&1&0&0\\0&0&1&t\\0&0&0&1
 \end{pmatrix}.
 \tag{D23.F11.1}
\]
Even without assuming the pentagon, the four Drinfeld
generator formulas, evaluated in this symplectic representation,
are exactly conjugation by $V(t_h)$.
For an actual full-GT element, this is the $V_h$ of D23.T1.
In particular
\[
 V_he_1=e_1,\quad V_he_3=e_3,\quad V_hw=w.
\]
The putative sign on the line $\widehat{\mathbb Z}w$ is
identically positive; it cannot detect $\Delta$.

\begin{proof}
The integral lattice is the orthogonal sum of the symplectic
planes
\[
 E=\langle e_1,e_2\rangle,\qquad
 E^\perp=\langle w,e_4\rangle,
\]
with each displayed ordered pair having pairing one.  The
matrix $V(t)$ is $\operatorname{diag}(U^t,U^t)$ in this
integral symplectic basis, so it preserves the chain form.

Put $t_i=q(\sigma_i)$ and
$z_i=(\sigma_1\cdots\sigma_{i-1})^i$.
The exact Drinfeld formula can be written
\[
 \sigma_1\mapsto\sigma_1,\qquad
 \sigma_i\mapsto f_h(z_i,\sigma_i^2)^{-1}
                   \sigma_i f_h(z_i,\sigma_i^2)\quad(i\geq2).
 \tag{D23.F11.2}
\]
See Nakamura, RIMS-1691, Section 5.4, equation (5.4.2),
printed p.~28.  This is equivalent to the usual formula
with $y_i=\sigma_{i-1}\cdots\sigma_1^2\cdots\sigma_{i-1}$:
$z_i=y_i z_{i-1}$, and $z_{i-1}$ commutes with both $y_i$
and $\sigma_i^2$.  Multiplying the first argument of a
commutator word by such a commuting central factor has no
effect, because its exponent sum in that argument is zero.

For $i=1$, $V(t)$ fixes $e_1$ and hence its transvection.
For $i=2$, the two arguments act on $E$ as $U^2,L^2$ and
fix $E^\perp$ pointwise.  AD.1 says precisely that the
resulting image of $t_2$ is its conjugate by $U^{t_h}$ on
$E$.  Thus it is conjugation by $V(t_h)$ on the full lattice.

For $i=3$, direct multiplication gives
\[
 q(z_3)=
 \begin{pmatrix}-1&0&2&0\\0&-1&0&0\\0&0&1&0\\0&0&0&1\end{pmatrix}.
\]
It preserves the plane $\langle e_1,w\rangle$, while
$t_3^2$ is identity on that plane.  Their restricted actions
commute.  Therefore the image of every commutator word,
in particular $f_h$, is identity on that plane and fixes
$e_3=w-e_1$.  The evaluated conjugate of $t_3$ is exactly
$t_3$, which is also its conjugate by $V(t_h)$.

For $i=4$, the integral identity
\[
 q(z_4)=(t_1t_2t_3)^4=t_w^2
\]
holds, where $t_w=I+ww^{\mathsf t}J$.
Both $t_w^2$ and $t_4^2$ act on $E^\perp$ as exactly
$U^2,L^2$, in the basis $(w,e_4)$, and fix $E$ pointwise.
The identical AD.1 calculation therefore gives the shear
$U^{t_h}$ on this second plane.  This is again conjugation
by $V(t_h)$.
All four formulas have now been checked.  For a full-GT
element they give the actual braid action, so uniqueness
in D23.T1 identifies $V_h$ with (D23.F11.1).  Its action
on $w$ is visibly identity.
\end{proof}

\paragraph{Scope of the failed detector.}
All the evaluated braid relations and this line sign are
automatic for gentle elements, including the extra-CM
element.  This does not assert that every separate lift
of the five-cycle word has trivial symplectic evaluation.
The representation does not descend to the five-point
sphere group, and any such lifted-pentagon assertion
requires its own calculation with its central factors retained.

\subsection{A five-point Krammer quotient retaining the cubic detector}

\paragraph{D23.P6 (a concrete extension outside the low-derived tests).}
Let $k=\mathbb F_3[i]$, $i^2=-1$. In the ordered basis
$(F_{12},F_{13},F_{14},F_{23},F_{24},F_{34})$, put
\[
A=\begin{pmatrix}
-i&0&0&-1-i&-1-i&0\\
0&0&0&i&0&0\\
0&0&0&0&i&0\\
0&1&0&1-i&0&0\\
0&0&1&0&1-i&0\\
0&0&0&0&0&1
\end{pmatrix},\quad
B=\begin{pmatrix}
0&i&0&0&0&0\\
1&1-i&0&0&0&0\\
0&0&1&0&0&0\\
0&1-i&0&-i&0&-1-i\\
0&0&0&0&0&i\\
0&0&0&0&1&1-i
\end{pmatrix},
\]
\[
C=\begin{pmatrix}
1&0&0&0&0&0\\
0&0&i&0&0&0\\
0&1&1-i&0&0&0\\
0&0&0&0&i&0\\
0&0&0&1&1-i&0\\
0&0&1-i&0&1-i&-i
\end{pmatrix}.
\tag{D23.P6.1}
\]
The assignment $(\sigma_1,\sigma_2,\sigma_3)\mapsto(A,B,C)$
is a continuous finite representation of $\widehat B_4$ and
$(ABC)^4=-I$. It therefore defines an actual projective quotient
of $L=B_4/\langle z_4\rangle$, the label-five stabilizer in
$\Gamma_{0,[5]}$. Induction from that subgroup gives a finite
quotient of the full five-point group if desired.

The subspace $W=\langle F_{12},F_{13},F_{23}\rangle$ is
invariant under $A,B$. Its representation is exactly the Frobenius
contragredient of the affine detector D23.P3. More precisely, with
\[
P=\begin{pmatrix}
1-i&1-i&-1\\
-1-i&1-i&i\\
-1&-1-i&1
\end{pmatrix},
\quad
P^{-1}=\begin{pmatrix}
0&-1&i\\
-i&-1&0\\
1-i&1-i&1+i
\end{pmatrix},
\tag{D23.P6.2}
\]
and $\overline{i}=-i$, one has
\[
 A|_W\,P=P\,\overline{\eta(\sigma_1)}^{-\mathsf T},\qquad
 B|_W\,P=P\,\overline{\eta(\sigma_2)}^{-\mathsf T}.
\tag{D23.P6.3}
\]
Here affine matrices are written homogeneously in dimension three,
$\sigma_1=(sr)^{-1}$ and $\sigma_2=s\sigma_1s$.
Consequently $f(A^2,B^2)|_W$ determines the cubic sign $a_3(h)$
for every actual gentle $h=(1,f)$ by D23.P3. This is an actual
five-point quotient retaining that sign; it is not merely a
four-point forgetful product.
The same recovery works from the projective evaluation: the
homogeneous affine image $E_F$ contains no nonidentity scalar
matrix, so its affine normalization selects the unique lift on $W$.

The generated first-pair image
$P_{12}=\langle A^2,B^2\rangle$ has order $46656$, with derived
series orders
\[
46656,\quad3888,\quad486,\quad243,\quad3,\quad1.
\tag{D23.P6.4}
\]
In particular $P_{12}''$ is not nilpotent. Its quotient on $W$
is D23.P3's $E_F$, whose second derived group is
$k^2\rtimes\langle-I\rangle$, a nonnilpotent group of order $162$.
Thus the nilpotent-second-derived exclusion criterion of supplied
FP.36 does not dispose of this representation.

\begin{proof}
The matrices are the Krammer matrices at $q=i,t=-i$ in the
corrected six-case formula of Bigelow, \emph{The Lawrence--Krammer
representation}, arXiv:math/0204057v1, section4.2, printed page13.
No assertion identifying this specialized lattice with the
homological Lawrence lattice is used. Alternatively, the three
braid/distant relations and $(ABC)^4=-I$ follow by direct
multiplication of (D23.P6.1); the exact script
\path{agents/pentagon/krammer_f9_probe.py} verifies them.
All matrices are invertible, as follows either from the formula
or from their explicitly invertible diagonal blocks after the
corresponding basis permutation. Hence they define the asserted
finite representation, which extends continuously to the profinite
completion. Killing the scalar full twist is precisely the
projective passage needed for $L$; no spherical relation is
assumed for an unsupported direct $B_5$ representation.

Inspection of the first, second and fourth columns shows the
invariance of $W$. Its two matrices are
\[
K_1=\begin{pmatrix}-i&0&-1-i\\0&0&i\\0&1&1-i\end{pmatrix},
\qquad
K_2=\begin{pmatrix}0&i&0\\1&1-i&0\\0&1-i&-i\end{pmatrix}.
\]
Direct multiplication verifies both $PP^{-1}=I$ and
(D23.P6.3). In encoded $\mathbb F_9$ arithmetic these are
checked independently in
\path{agents/pentagon/krammer_affine_conjugacy.py}.
Since Frobenius followed by contragredience is an invertible
operation on matrix groups, these identities identify the entire
marked $B_3$ representation on $W$, rather than merely its order
or composition factors. Thus its pure image is exactly the
Frobenius contragredient of $E_F$, and its value on $f$ recovers
$\eta(f)$ by
\[
 \eta(f)=\overline{\bigl(P^{-1}f(A^2,B^2)|_W P\bigr)^{-\mathsf T}}.
\]
The cubic-sign formula now follows from D23.P3.

For completeness, the finite-order statements are certified in two
ways. The script \path{agents/pentagon/krammer_f9_probe.py}
uses a faithful permutation action on a union of vector orbits
containing a basis, and computes its stabilizer chain and derived
series. The independent matrix enumeration
\path{agents/pentagon/krammer_first_pair.py} enumerates exactly
$46656$ distinct matrices closed under the two generators and
exactly $3888$ matrices generated by the normal closure of their
commutator. The latter commutator has $486$ conjugates and three
of them generate the derived subgroup. The associated finite data
are retained in the JSON/NPZ certificates. These checks agree.

Finally, $E_F'=k^2\rtimes Q_8$ and
$E_F''=k^2\rtimes C_2$, since $-I$ acts as $-1$ on $k^2$.
The latter group is nonnilpotent: its Sylow two subgroup does not
commute with its normal Sylow three subgroup. A quotient of a
nilpotent finite group is nilpotent. The surjection
$P_{12}''\twoheadrightarrow E_F''$ therefore proves the final
assertion independently of the complete derived-series counts.
\end{proof}

\paragraph{Current scope.}
D23.P6 constructs and checks a new finite five-point test. It does
not prove the pentagon in this quotient for every gentle element,
and it does not prove $\Delta|_{\widehat{GT}}=0$ or exhibit a
full-GT counterexample. The exact first-pair evaluations are further
reduced to six matrices in D23.F13. The correlations with the other
two pentagon pair types remain the substantive calculation.

\subsection{The Krammer quotient and the disappearing central phase}

\paragraph{D23.F13 (six necessary first-pair values).}
Let $k=\mathbf F_3[i]$, $i^2=-1$, and use the specialized
Lawrence--Krammer matrices $K_1,K_2,K_3$ of D23.P6, in the
ordered basis $(12,13,14,23,24,34)$, with $q=i,t=-i$.
Set $X=K_1^2$, $Y=K_2^2$, $G=\langle X,Y\rangle$, and
\[
 W=\langle12,13,23\rangle,\qquad Q=k^6/W.
\]
The following statements hold.
\begin{enumerate}
\item The image of $G$ on $Q$ is dihedral of order eight.
For every actual gentle multiplier-one pair $(1,f)$,
$f(X,Y)|_Q=I$.
\item The map from $G$ to its two diagonal blocks has image
$E_F\times D_8$ of order $15552$ and kernel of order three.
Here $E_F$ is the cubic affine image of D23.P3, transported
to $W$ by the Frobenius contragredient comparison in D23.P6.
The kernel is central even in $\langle K_1,K_2\rangle$ and is
generated by
\[
 n=((K_1K_2)^3)^4=I+N,
 \qquad
 N_{W,Q}=\begin{pmatrix}
 -1+i&-1+i&-1+i\\
 1+i&1+i&1+i\\
 1-i&1-i&1-i
 \end{pmatrix},
\tag{D23.F13.1}
\]
all other blocks of $N$ being zero.
\item The actual first-pair matrix $f(X,Y)$ is uniquely
determined by $(\rho_h,a_3(h))\in\mathbf F_3\times\mathbf F_2$.
There are six explicit necessary matrices, one for each such
pair, in \texttt{krammer\_quotient\_extension.json} and
\texttt{krammer\_six\_first\_pair.npz}.  This assertion does
not assert a profinite lift for an arbitrary finite candidate.
In particular neither the quotient block nor the central
order-three factor supplies an additional independent state
of the first pair.
\end{enumerate}

\begin{proof}
The quotient matrices on $(14,24,34)$ are
\[
 A=\begin{pmatrix}0&i&0\\1&1-i&0\\0&0&1\end{pmatrix},
 \qquad
 B=\begin{pmatrix}1&0&0\\0&0&i\\0&1&1-i\end{pmatrix}.
\]
Direct multiplication gives $(A^2)^2=(B^2)^2=I$ and
$A^2B^2$ of order four.  Two involutions whose product has
order four generate the dihedral group of order eight:
each word reduces to $(A^2B^2)^j$ or $(A^2B^2)^jA^2$,
and these eight words are distinct.  Its commutator subgroup
is central of order two.

For completeness, the consequence of the gentle relation
uses no congruence-covariance assumption.  In a two-generator
group of nilpotency class at most two, a word from the closed
commutator subgroup of the free profinite group evaluates
as $[x,y]^c$ for its commutator coefficient
$c\in\widehat{\mathbf Z}$ modulo the order of $[x,y]$.
Writing $z=(xy)^{-1}$, the class-two identities give
$[y,z]=[z,x]=[x,y]$.  The multiplier-one three-cycle relation
therefore gives $[x,y]^{3c}=1$.  In $D_8$ the commutator has
order two, and hence $[x,y]^c=1$.  This proves the first
assertion for every actual gentle $f$.

By D23.P3 and the explicit Frobenius contragredient
intertwiner, the image on $W$ is
$E_F=k^2\rtimes\operatorname{SL}_2(\mathbf F_3)$, of order
$1944$.  Its abelianization is $C_3$: the central element
$-I$ makes the translation subgroup a commutator subgroup,
and $\operatorname{SL}_2(\mathbf F_3)^{\rm ab}=C_3$.
A nontrivial finite $2$-group quotient of $E_F$ would have
a further quotient $C_2$, which this abelianization excludes.
Thus $E_F$ and $D_8$ have no common nontrivial quotient.
Goursat's lemma applied to the subdirect diagonal image
gives the full product, of order $1944\cdot8=15552$.

Here is the exact finite certificate used for the remaining
kernel count; its scope is finite matrix arithmetic, not a
formal proof-assistant verification.  The program
\texttt{krammer\_quotient\_extension.py} starts from the two
literal six-dimensional matrices over $k$, represents
$a+bi$ by the block
$\left(\begin{smallmatrix}a&-b\\b&a\end{smallmatrix}\right)$
over $\mathbf F_3$, and enumerates $\langle X,Y\rangle$.
It starts with the identity and repeatedly adjoins right
products by $X,Y$ until no new matrix appears; full exact
matrix byte strings are used as keys, with equality tested
by equality of their entries.  Every enumerated element is
a word in $X,Y$, and at termination the finite set is closed
under right multiplication by these generators.  Since
the generators are invertible, the set is exactly $G$.
This exhaustive enumeration has $46656$ distinct matrices.
It is rebuilt independently of the matrix array supplied
by D23.P6, and the two sets are compared for equality.
The program checks the diagonal block image has $15552$
distinct values, and that exactly three matrices have both
blocks the identity.  Thus the kernel has order three.
The explicit product $((K_1K_2)^3)^4$ is exactly (D23.F13.1).
It is nonidentity, $N^2=0$, and $(I+N)^3=I$; it consequently
generates the kernel.  The braid relation makes
$(K_1K_2)^3$ central in $\langle K_1,K_2\rangle$, proving
the asserted centrality without relying on the enumeration.

It remains to remove the apparent choice of a central phase.
Put $S=K_1K_2K_1$.  The braid relation gives
$SXS^{-1}=Y$ and $SYS^{-1}=X$.  The two-cycle relation on
an actual $f$ is therefore the literal matrix equation
\[
 S f(X,Y)S^{-1}=f(X,Y)^{-1}.
\tag{D23.F13.2}
\]
If two possible matrices $F,F'$ have the same $W$ block
and trivial $Q$ block, then $F'=Fn^j$, $j\in\mathbf F_3$.
Since $n$ is central, substituting into (D23.F13.2) gives
$n^j=n^{-j}$.  As $n$ has order three, this forces $j=0$.
Thus the $W$ block uniquely determines an actual value.

For an explicit description, D23.P3 gives the six affine
matrices
\[
 \eta(f)=\begin{pmatrix}
 1&\rho&d_1\\
 \rho&1+\rho^2&d_2\\
 0&0&1
 \end{pmatrix},\quad
 d_1=\frac{1-\rho-\epsilon}{1+i},\quad
 d_2=(\rho+i)d_1,\quad \epsilon=(-1)^{a_3}.
\]
Their $W$ matrices are
$P\,\overline{\eta(f)}^{-T}P^{-1}$, for
\[
 P=\begin{pmatrix}1-i&1-i&-1\\-1-i&1-i&i\\-1&-1-i&1\end{pmatrix}.
\]
The certificate constructs these six blocks directly from
the displayed formulas, searches the enumerated group for
matrices with each block and trivial quotient block, and
finds three in each fiber.  It then tests (D23.F13.2) and
finds precisely one in each fiber.  These are the six
matrices in the stated files.  All restrictions used were
necessary identities of actual gentle elements.  No finite
matrix solution is being promoted to an actual profinite
pair, and no conclusion about $a_2=a_3$ follows from this
first-pair computation alone.
\end{proof}

\paragraph{Source and marking.}
The polynomial matrix convention is Stephen Bigelow,
\emph{The Lawrence--Krammer representation}, arXiv:math/0204057,
Section 4.2, as used and checked in D23.P6.  The passage
from the $W$ block to the affine cubic detector is the
explicit intertwiner in
\texttt{krammer\_affine\_conjugacy.py}; it is Frobenius
contragredient, not an ordinary conjugacy of the affine
homogeneous matrices.

% DELTA23 independent mixed Krammer analysis, /root/characters.
% Matrix certificates use F9=F3[i], i^2=-1, integer encoding a+3b.

\subsection{The two mixed Krammer pairs and their exact gentle correlation}

Use the six-dimensional Krammer representation of $B_4$ at
$q=i,t=-i$ over $k=\mathbb F_9$.  Its three literal matrices
$K_1,K_2,K_3$ are constructed in the companion pentagon workstream
from Bigelow's displayed Krammer formula.  Set
\[
 X_1=K_1^2,\quad X_2=K_2^2,\quad X_3=K_3^2,
 \quad X_4=(K_1K_2K_1)^2,\quad X_5=(K_3K_2K_3)^2.
\]
The exact matrices, all the bases below, and the full six-dimensional
intertwiners are recorded in
\path{agents/characters/records/krammer_mixed_correlation.json}; the executable
\path{agents/characters/code/check_krammer_mixed_correlation.py} verifies every
displayed matrix identity in the literal original basis.

In the matrix arrays in this subsection, an integer from zero to
eight means $a+bi\in\mathbb F_9$ when its integer value is $a+3b$,
with $0\le a,b\le2$.  Thus $3=i$, $6=-i$, and $2=-1$.

\paragraph{D23.C9 (proved by exact finite linear algebra and group certificates).}
Each mixed pair $(X_3,X_4)$ and $(X_4,X_5)$ preserves a direct
sum of two three-dimensional spaces.  On the first its image is
a solvable group $H$ of order $108$, with derived orders
\[
 108,27,3,1.
\]
On the second its image is a full unitary group $U(3,3)$ of order
$24192$, with perfect derived group $SU(3,3)$ of order $6048$.
The orders and successive derived orders of the full mixed images are
\[
\begin{array}{c|rrrr}
 \text{pair}&G&G'&G''&G'''=G''''\\\hline
 (X_3,X_4)&1306368&163296&18144&6048\\
 (X_4,X_5)&653184&163296&18144&6048.
\end{array}
\]
Their full linear images inject into the projective images obtained
by quotienting by the scalar $-I$.

There are explicit $C,Q,D\in\mathrm{GL}_6(k)$ such that, writing
$z=(X_4X_5)^{-1}$,
\[
 X_3C=CDz,\qquad X_4C=CX_4,\qquad
 [D,X_4]=[D,X_5]=1,
\]
\[
 zQ=Qz,\qquad X_4Q=QX_5.
\]
Consequently, for every actual multiplier-one gentle word $f$, if
\[
 R=f(X_3,X_4),\qquad r=C^{-1}RC,
\]
then its two mixed values obey the literal six-dimensional identity
\[
 \boxed{f(X_4,X_5)=Q^{-1}rQr^{-1}.}
\]
In particular the second mixed value is not an independent variable
in a finite necessary pentagon test.  The first mixed value belongs
to a group of order $163296$.

\emph{Proof.}
Let $B_{34}$ and $B_{45}$ be the splitting matrices from the
certificate, with the solvable summand first.  Their columns are
the indicated bases of the two invariant spaces, and have rank six.
Writing the corresponding three-dimensional generator pairs as
$(x_H,y_H),(x_U,y_U)$ and $(u_H,v_H),(u_U,v_U)$, the certificate
checks
\[
 B_{34}^{-1}(X_3,X_4)B_{34}
    =((x_H\oplus x_U),(y_H\oplus y_U)),
\]
\[
 B_{45}^{-1}(X_4,X_5)B_{45}
    =((u_H\oplus u_U),(v_H\oplus v_U)).
\]
The invariant-space discovery enumerated all $101$ projective
eigenvectors of one generator in each pair.  Discovery is not
needed as an assumption: the displayed bases and intertwining
identities independently certify the direct sums.

The two solvable block groups were computed as permutation groups
on unions of vector orbits containing a full $\mathbb F_3$-basis
of $k^3$.  Such actions are faithful, since a matrix fixing those
vectors is the identity.  Exact stabilizer chains give orders
$108,27,3,1$ in the derived series.  The first derived group has
exponent three.  The determinant maps the whole group onto
$\mu_4\subset k^\times$, so its kernel, of order $27$, is exactly
the derived group.

For the unitary blocks, the two respective invariant Hermitian
forms are
\[
 J_{34}=\begin{pmatrix}0&1&5\\1&0&3\\8&6&0\end{pmatrix},
 \qquad
 J_{45}=\begin{pmatrix}0&3&6\\6&0&4\\3&7&0\end{pmatrix}.
\]
They are Hermitian for $i\mapsto-i$, have nonzero determinants
$2$ and $1$, and the exact identities $g^*Jg=J$ hold for both
generators.  A faithful union of vector orbits has $224$ vectors;
the stabilizer-chain calculation gives order $24192$ and derived
order $6048$, with the derived group perfect.  The full unitary
group has order
\[
 3^3(3+1)(3^2-1)(3^3+1)=24192.
\]
For a direct count at this small field, there are $252$ vectors of
norm one in dimension three, $24$ in their orthogonal complement
of dimension two, and $4$ in dimension one, giving the same product.
Thus containment and the computed order identify the full group.
The determinant image has order four and the derived group lies
in its kernel; the kernel has order $6048$, so it is exactly the
special unitary group.  This identification does not need a
classification or simplicity theorem for finite groups.

For completeness, the full group orders do not require enumerating
the large six-dimensional image.  Each full mixed image is a
subdirect product of its solvable block $H$ and unitary block $U$.
Any common quotient of $H$ and $U$ is solvable, hence kills the
perfect group $U'=SU(3,3)$ and is abelian.  It therefore also kills
$H'$.  Goursat's lemma gives $H'\times U'$ inside the mixed image;
the image is the full preimage of its determinant-pair image.
For pair $34$, the generator determinant pairs are
\[
 (-1,1),\quad(-i,-i),
\]
generating the equal-parity subgroup of $\mu_4^2$ of order eight.
For pair $45$ they are both $(-i,-i)$ and generate the diagonal
subgroup of order four.  Multiplication by
$|H'|\,|U'|=27\cdot6048$ gives the two stated group orders.
Applying the same solvable/perfect Goursat argument to the derived
subgroup gives $G'=H'\times U'$, and then the remaining orders
follow by taking derived subgroups of a direct product.  Each
solvable block has a common fixed line.  Thus the full mixed image
does not contain scalar $-I$, proving the assertion about the
projective quotient.

Here are small matrices constructing the required full intertwiners:
\[
 C_H=\begin{pmatrix}4&1&0\\3&0&0\\0&0&1\end{pmatrix},
 \quad C_U=\begin{pmatrix}1&3&2\\3&5&0\\0&1&0\end{pmatrix},
\]
\[
 Q_H=\begin{pmatrix}4&0&6\\2&1&0\\1&0&0\end{pmatrix},
 \quad Q_U=\begin{pmatrix}2&0&2\\2&5&0\\6&2&1\end{pmatrix}.
\]
Put $z_H=(u_Hv_H)^{-1}$ and $z_U=(u_Uv_U)^{-1}$.  Direct
multiplication gives
\[
 x_HC_H=C_Hz_H^3,\quad y_HC_H=C_Hu_H,
 \quad x_UC_U=C_U(-z_U),\quad y_UC_U=C_Uu_U,
\]
and $z_H^2$ commutes with $u_H,v_H$.  For each block one also has
$z_jQ_j=Q_jz_j$ and $u_jQ_j=Q_jv_j$.  Define
\[
 \begin{split}
 C&=B_{34}(C_H\oplus C_U)B_{45}^{-1},\\
 Q&=B_{45}(Q_H\oplus Q_U)B_{45}^{-1},\\
 D&=B_{45}(z_H^2\oplus(-I_3))B_{45}^{-1}.
 \end{split}
\]
All three matrices are invertible; the stated literal
six-dimensional identities follow by conjugating the block
identities.  The companion computation also checks them directly
in the original basis, independently of this deduction.

If a profinite word $f$ belongs to $F'$, multiplying an input by
an element centralizing both inputs does not change its value:
the extra factor is raised to the corresponding abelian exponent,
which is zero.  Consequently
\[
 C^{-1}f(X_3,X_4)C=f(Dz,X_4)=f(z,X_4)=r.
\]
The other intertwiner gives
$f(z,X_5)=Q^{-1}rQ$.  The multiplier-one two-cycle equation says
$f(X_5,z)=f(z,X_5)^{-1}=Q^{-1}r^{-1}Q$.  In the source's exact
order, its three-cycle equation is
\[
 f(z,X_4)f(X_5,z)f(X_4,X_5)=I
\]
(see DELTA\_ON\_GT P.9, the displayed two- and three-cycle
equations in the normalized lift proof).  Substitution and inversion
give $f(X_4,X_5)=Q^{-1}rQr^{-1}$, as asserted.  Finally
$f\in F'$ puts its first mixed value in $G_{34}'$, whose order
has already been computed.\hfill$\square$

\paragraph{D23.C10 (proved; the nilpotent residual does not erase this block).}
Under either mixed-pair evaluation, the image of $F''$ is
$H''\times SU(3,3)$, of order $3\cdot6048$.  The image of
$R_{\rm nil}(F'')$ is exactly
\[
 \{1\}\times SU(3,3).
\]
These assertions concern the images of all words in the indicated
subgroups.  They do not assert that actual gentle words fill those
images or satisfy the pentagon there.

\emph{Proof.}
Surjectivity of the free evaluation and D23.C9 give the claimed
image of $F''$ as the second derived subgroup.  The $H''$ factor
has order three, so it is a finite nilpotent quotient and is
annihilated by $R_{\rm nil}(F'')$.  Let $N$ be the image of that
residual in the perfect $SU(3,3)$ factor.  The finite group
$SU(3,3)/N$ is a quotient of the pronilpotent group
$F''/R_{\rm nil}(F'')$, hence is nilpotent.  It is also perfect,
being a quotient of a perfect group.  A finite perfect nilpotent
group is trivial, so $N=SU(3,3)$.\hfill$\square$

\subsection{Exact reduced pentagon assignments in the new quotient}

\paragraph{D23.P7 (a finite calculation with a precise scope).}
Use D23.P6's matrices and put
\[
 X_1=A^2,\quad X_2=B^2,\quad X_3=C^2,\quad
 X_4=(ABA)^2,\quad X_5=(CBC)^2,\qquad J=ABCABA.
\]
Let $F_{\rho,e}$, $\rho\in\mathbb F_3,e\in\mathbb F_2$, be
the six exact first-pair matrices of D23.F13. Let $C_0,Q_0$ be
the full six-dimensional intertwiners of D23.C9. In the finite
derived subgroup $G_{34}'=\langle X_3,X_4\rangle'$, of order
$163296$, consider the equation
\[
 R(JR^{-1}J^{-1})(JF_{\rho,e}^{-1}J^{-1})
 \bigl(Q_0^{-1}rQ_0r^{-1}\bigr)F_{\rho,e}=I,
 \qquad r=C_0^{-1}RC_0.
\tag{D23.P7.1}
\]
For each of the six pairs $(\rho,e)$, exactly $36$ matrices $R$
satisfy this equation. There are no solutions with the right side
replaced by $-I$. Thus the reduced matrix equation admits both
cubic-sign values at every $\rho$.

Every actual multiplier-one full-GT element gives one of these
$216$ assignments, because D23.F13 fixes its first-pair value,
the two-cycle gives the two reversing-half-twist factors, and
D23.C9 gives the displayed fourth factor. This is a necessary
finite assignment theorem. The count by itself does not assert
that every listed assignment comes from a common word, an actual
gentle element, or an actual full-GT element.

\begin{proof}
The reversing half twist satisfies $J^2=-I$ and interchanges
$X_1,X_3$ and $X_4,X_5$ by conjugation while fixing $X_2$.
These are direct matrix identities. Therefore the two-cycle
expresses the five-cycle exactly as in (D23.P7.1), using the
actual inverse $J^{-1}$ rather than silently treating the lift
$J$ as an involution.

The group decomposition D23.C10 gives
\[
 G_{34}'=B_{34}\bigl(H'\times U'\bigr)B_{34}^{-1},
 \qquad |H'|=27,\quad |U'|=6048,
\]
where the two three-dimensional blocks and the splitting matrix
$B_{34}$ are literal matrices in D23.C9's certificate. The
enumeration program reconstructs each of $H',U'$ as the normal
closure of the commutator of its two displayed generators. It
then runs through all $27\cdot6048$ block-diagonal pairs, without
a sampling assumption.

For efficiency the program uses the following equivalent form of
(D23.P7.1). Put
\[
 A_R=RJR^{-1},\qquad
 B_R=rQ_0^{-1}r^{-1}Q_0J.
\]
Multiplying (D23.P7.1) on the indicated sides yields
\[
 F_{\rho,e}A_R=B_RF_{\rho,e}.
\tag{D23.P7.2}
\]
The program tests equality of all entries for all six $F_{\rho,e}$
and every enumerated $R$. It obtains $36$ each; testing the
negative equality gives zero each. Every retained assignment was
then checked again by multiplying the original five factors in
(D23.P7.1), rather than by reusing (D23.P7.2).

All arithmetic is exact over $\mathbb F_3$, after the faithful
restriction-of-scalars embedding of $\mathbb F_9$ matrices. The
reproducible certificate is
\path{agents/pentagon/krammer_pentagon_search.py}; its counts are
in the matching JSON file, and all $216$ matrices with their
$(\rho,e)$ labels are retained in
\path{agents/pentagon/krammer_pentagon_solutions.npz}.

The actual full-GT pentagon lies in the label-five stabilizer
$B_4/\langle z_4\rangle$. The displayed word is a product of
evaluations of $f\in\widehat F_2'$, so its braid exponent sum is
zero. If its lift to $\widehat B_4$ is $z_4^a$, abelianization
gives $12a=0$ in $\widehat{\mathbb Z}$, hence $a=0$. Thus the
literal identity in (D23.P7.1), not merely an arbitrary scalar,
is necessary. In this finite calculation the only additional
possible scalar lift would in any event be $-I$, and it yields
no solutions.
\end{proof}

\paragraph{Unresolved comparison.}
D23.P7 is stronger than discarding the new representation by a
nilpotence criterion: its first-pair cubic sign enters a genuinely
new pentagon calculation. Its outcome nevertheless leaves both
signs. The common-word correlations and the distinction between
finite assignments and actual profinite lifts must still be
respected. In particular no inference $e=\nu_4$ is contained in
the count $36$.

\subsection{Common-word realization on the three independent pair types}

\paragraph{D23.P8 (the three-pair commutator image is a direct product).}
For the matrices of D23.P6, the simultaneous evaluation map
\[
 \widehat F'\longrightarrow
 P_{12}'\times G_{34}'\times G_{45}'
\tag{D23.P8.1}
\]
is surjective. Here the factors are the derived groups of the
three corresponding pair images. In particular each of the $216$
reduced assignments in D23.P7 has a common commutator-word
preimage for these three pair values.

This statement does not yet prescribe the reversed mixed-pair
value, and it does not assert that the common word satisfies
the profinite gentle identities. These qualifications matter:
the reductions in D23.P7 used those identities for actual gentle
elements, whereas surjectivity in (D23.P8.1) concerns unrestricted
commutator words.

\begin{proof}
Write $G_{34}'=H_{34}'\times S_{34}$ and
$G_{45}'=H_{45}'\times S_{45}$ in the exact splitting bases of
D23.C9, where $|H_{34}'|=|H_{45}'|=27$ and
$S_{34}\simeq S_{45}\simeq\operatorname{SU}_3(3)$.

First, the two solvable full pair blocks $H_{34},H_{45}$ have
abelianization $C_4$ via determinant. With $-i$ as generator,
their two marked determinant phase pairs are $(2,1)$ for $x$
and $(1,1)$ for $y$. In any common abelian quotient of the two
blocks, if $a,b$ denote the images of their respective cyclic
abelianization generators, the common marking would impose
$2a=b$ and $a=b$. Hence $a=b=0$. If their simultaneous image
were a proper subdirect product, Goursat's lemma would give a
nontrivial common solvable quotient, whose abelianization would
be nontrivial. The displayed calculation excludes that.
Thus the full simultaneous image is $H_{34}\times H_{45}$,
and its derived image is $H_{34}'\times H_{45}'$.

Next the two simple factors have different kernels even on
$\widehat F'$. The following ordinary commutator word is a
small exact witness:
\[
 w=y^4x^2yx^{-1}y^{-1}x^{-1}y^{-4}xy^2x^2yxy^{-1}x^{-3}y^{-2}x^{-1}.
\tag{D23.P8.2}
\]
Its exponent sums are both zero. In the unitary block matrices
of D23.C9 it has values
\[
 w_{34}=I,\qquad
 w_{45}=\begin{pmatrix}5&7&4\\8&8&6\\5&7&5\end{pmatrix}\ne I,
\tag{D23.P8.3}
\]
using the same encoding $a+3b$ for $a+bi$ as the certificates.
The script \path{agents/pentagon/krammer_commonword_witness.py}
verifies the two values by literal word multiplication.
Since each factor is nonabelian simple, a subdirect product
of the two is either the whole product or the graph of an
isomorphism. The latter would give equal kernels, contrary
to (D23.P8.3). Thus the perfect joint image is
$S_{34}\times S_{45}$.

The standard simplicity of $\operatorname{SU}_3(3)$ is the
primary-source input used here. Its centre is trivial since
$\gcd(3,4)=1$; the GAP Character Table Library's
\emph{Simple Groups by Group Order}, row $U_3(3)$, lists its
order $6048$, and the corresponding group record identifies
it with $\operatorname{PSU}(3,3)$. This is the same cited
input used in D23.F14, not an inference from the order alone.

Finally $P_{12}'$ is solvable, and its abelianization has
order $3888/486=8$, by D23.P6. The group
$H_{34}'\times H_{45}'$ is a three-group. Any nontrivial
common quotient would be solvable and would have a nontrivial
abelian quotient, simultaneously a quotient of a two-group
and a three-group. This is impossible. Hence Goursat gives
the full product of these three solvable factors. It has no
nontrivial common quotient with the perfect product
$S_{34}\times S_{45}$, because a finite group both solvable
and perfect is trivial. A final application of Goursat proves
(D23.P8.1).

The image of the ordinary free commutator subgroup in any
finite target equals the image of its profinite closure.
Therefore ordinary commutator-word preimages can be chosen.
This proves the stated finite-word realization and its scope.
\end{proof}

\paragraph{Primary-source locators.}
\url{https://www.math.rwth-aachen.de/homes/Thomas.Breuer/ctbllib/ctbltoc/views/simplebyorder.html},
row $U_3(3)$;
\url{https://www.math.rwth-aachen.de/homes/sam/ctbllib/ctbltoc/data/U3(3).html},
``Group constructions in GAP'' entry $\operatorname{PSU}(3,3)$.

\subsection{The full common-word audit for the Krammer test}

\paragraph{D23.P9 (all reduced pentagon assignments have finite-word realizations).}
Retain D23.P6--P8's notation, and write
$G_{43}=\langle X_4,X_3\rangle$ with the indicated reversed
marking. The simultaneous map
\[
 \widehat F'\longrightarrow
 P_{12}'\times G_{34}'\times G_{45}'\times G_{43}'
\tag{D23.P9.1}
\]
is onto. Every one of the $216$ assignments of D23.P7 can
therefore be realized by an ordinary free commutator word whose
five actual Krammer evaluations multiply to $I$ in the displayed
pentagon order. No unrelated-word substitution is involved.

This is a statement about ordinary finite word images. It does
not prove that one of these words is a gentle parameter. In
particular it does not construct an actual full-GT element with
nonzero $\Delta$.

\begin{proof}
Use D23.C9's faithful decompositions of each of the three mixed
pair groups into a solvable block of order $108$ and a unitary
block of order $24192$. On the reversed pair use the same basis
as pair $34$ and exchange its two generators. Their derived
groups are respectively $H_j'$ of order $27$ and simple
$S_j=\operatorname{SU}_3(3)$ of order $6048$.

The exact finite calculation
\path{agents/pentagon/krammer_allfive_commonword.py} forms the
three block-diagonal solvable generator pairs and computes the
normal closure of their joint commutator. It contains exactly
$19683=27^3$ matrices. The conjugacy orbit of that commutator
has $432$ elements and six of them generate the entire closure.
The image is contained in the product of the three order-$27$
derived groups, so this count proves equality with that full
product. This is a complete finite enumeration, not a random
test. Normal closure of the commutator computes the derived
subgroup of a two-generator group because the resulting
quotient makes the two generators commute.

The simple factors are independent as well. D23.P8's word
$w_1$ of length $28$ has identity pattern
\[
  (w_1(S_{34}),w_1(S_{45}),w_1(S_{43})):
       \quad(I,\ne I,\ne I).
\]
A second explicit word is
\[
 w_2=yxy^2xyx^{-1}y^{-3}x^{-1}y^{-1}
       xyx^2yxy^{-1}x^{-3}y^{-1}x^{-1}.
\tag{D23.P9.2}
\]
It has exponent sums zero and satisfies
$w_2(S_{45})=I$, $w_2(S_{43})\ne I$.
The program verifies these literal word values in the stored
three-dimensional matrices after restriction of scalars to
$\mathbb F_3$. Hence the kernels on $F'$ of the three simple
projections are pairwise distinct. A subdirect product of
nonabelian simple groups is a product of diagonal factors;
distinct kernels exclude every identification. The joint image
is therefore $S_{34}\times S_{45}\times S_{43}$.
The standard simplicity input and primary locators are those
recorded explicitly in D23.P8 and D23.F14.

By D23.P6, the abelianization of the solvable group $P_{12}'$
has order eight. The joint solvable mixed image is a three-group.
They have no nontrivial common quotient, because any nontrivial
finite solvable quotient has a nontrivial abelian quotient.
Goursat therefore separates $P_{12}'$ from the three solvable
mixed factors. A further application of Goursat separates this
solvable product from the perfect product of the three simple
factors. This proves (D23.P9.1).

Now take a listed solution $(F,R)$ of D23.P7 and put
\[
 H=Q_0^{-1}(C_0^{-1}RC_0)Q_0(C_0^{-1}R^{-1}C_0).
\]
Surjectivity provides a common commutator word $f$ with
\[
 f(X_1,X_2)=F,\quad f(X_3,X_4)=R,\quad
 f(X_4,X_5)=H,\quad f(X_4,X_3)=R^{-1}.
\]
The first-pair half twist $S=ABA$ exchanges $X_1,X_2$.
Thus, for every word, not merely for gentle words,
$f(X_2,X_1)=Sf(X_1,X_2)S^{-1}$. The selected six first-pair
matrices satisfy $SFS^{-1}=F^{-1}$ by D23.F13. The reversing
half twist $J=ABCABA$ gives
\[
 f(X_2,X_3)=JF^{-1}J^{-1},\qquad
 f(X_5,X_1)=JR^{-1}J^{-1}.
\]
The actual five values of this one word are consequently the
five factors in D23.P7.1. Their product is $I$ by the checked
assignment. The finite image of the ordinary free commutator
subgroup equals that of its profinite closure, so $f$ can be
chosen as an ordinary word.

The argument prescribes only finitely many evaluations and their
displayed relations. It does not establish the identities in
the full profinite free group required for a gentle pair, which
is the stated limitation.
\end{proof}

\subsection{Complete outer three-puncture closure on the mixed perfect blocks}

\paragraph{D23.C15 (the mixed perfect closure retains all 216 candidates).}
In the unitary block of the pair $(X_4,X_5)$ from D23.C9 put
$a=X_{4,U}$, $b=X_{5,U}$, $c=(ab)^{-1}$. For a pair of these
matrices, its kernel below means the kernel of evaluation restricted
to $\widehat F'$. The six oriented pairs have exactly three kernel
classes:
\[
 \{(a,b),(b,a)\},\qquad
 \{(b,c),(a,c)\},\qquad
 \{(c,a),(c,b)\}.
\tag{D23.C15.1}
\]
Every coincidence is certified by an explicit linear intertwiner,
and every noncoincidence by a literal commutator word whose value
is identity in one pair and nonidentity in the other. The six pairs
are the full orbit modulo simultaneous inner conjugation under
permuting the three peripheral conjugacy classes.

All 216 necessary Krammer pentagon solutions of D23.P7 satisfy
all value correlations arising from (D23.C15.1). They also satisfy
the opposite-orientation three-cycle equation in the full original
six-dimensional representation. Thus this complete closure on the
mixed perfect blocks removes none of D23.P7's 36 candidates for
each $(\rho,a_3)$, including the cases $(0,1)$.
This is a finite necessary-equation result, not an assertion that
any candidate lifts to an actual gentle or full-GT element.

\begin{proof}
For $g\in U(3,3)$ put $g^\circ=(\det g)g$. The determinant is in
$\mu_4$, so $\det(g^\circ)=(\det g)^4=1$. Central multipliers of
inputs do not change evaluations of words in $F'$. We can therefore
use the pairs of determinant-one matrices without changing the
kernels or values under discussion. Their images are $SU(3,3)$:
each pair generates $U(3,3)$ before normalization, and its derived
image is already $SU(3,3)$ by D23.C9.

In the field encoding $a+3b\leftrightarrow a+bi$ used in D23.C9,
intertwiners for the three displayed coincidences, in the same
order, are
\[
 J_0=\begin{pmatrix}2&7&2\\0&3&0\\2&6&1\end{pmatrix},\quad
 J_1=\begin{pmatrix}2&8&2\\8&6&8\\5&4&1\end{pmatrix},\quad
 J_2=\begin{pmatrix}2&0&2\\2&5&0\\6&2&1\end{pmatrix}.
\tag{D23.C15.2}
\]
The certificate verifies $p_iJ=Jq_i$ for both generators of each
normalized source pair $(p_1,p_2)$ and target pair $(q_1,q_2)$.
Hence $f(q_1,q_2)=J^{-1}f(p_1,p_2)J$ for every $f\in F'$.
No semilinear automorphism or classification of automorphisms of
$SU(3,3)$ is used.

For all twelve pairs of distinct classes the certificate records a
word in the alphabet $\{1,2,-1,-2\}$, with zero exponent sum in
each generator, and checks its two evaluations exactly. These
words were obtained by a Cayley-graph collision and then taking
a commutator; the verification needs only the recorded words and
the displayed matrices. Thus they prove noncoincidence on $F'$
itself, rather than merely a failed search for linear intertwiners.
Together these checks prove (D23.C15.1).

Let $R=f(X_3,X_4)$ and $r=C^{-1}RC$ in the full six-dimensional
notation of D23.C9. The three oriented values on $(a,b,c)$ are the
unitary blocks of
\[
 A=Q^{-1}rQr^{-1},\qquad B=Q^{-1}r^{-1}Q,\qquad C_0=r.
\tag{D23.C15.3}
\]
They satisfy $C_0BA=1$ by construction. Two-cycle gives the three
reversed values $A^{-1},B^{-1},C_0^{-1}$. Restricting each of
D23.P7's exact matrix solutions to the pair-$45$ splitting basis,
the certificate checks all three intertwiner relations (D23.C15.2)
on these six values. Every solution passes. In particular the
self-duality of $A$ under the first reflection, which is not just
the original displayed cycle, is retained by every candidate.

For completeness the three-cycle at the reversed input $(b,a)$
has third element $c'=(ba)^{-1}=a^{-1}ca=bcb^{-1}$. Its exact
form after two-cycle and simultaneous conjugation is
\[
              bB^{-1}b^{-1}a^{-1}C_0^{-1}aA^{-1}=1.
\tag{D23.C15.4}
\]
The separate certificate checks this identity directly in the full
six-dimensional matrices for every one of the 216 candidates,
with no projective scalar discarded. Cyclic rotation supplies the
other cycles of this orientation; the original orientation is
already (D23.C15.3). The six representatives suffice because
swapping the first two generators changes the third to an inner
conjugate, as the displayed formulas for $c'$ show.

D23.C9 identifies the pair-$34$ perfect evaluation with $(c,a)$,
up to a central input scalar; its reverse is $(a,c)$. Hence the
same closure covers both mixed perfect blocks, not just pair $45$.
No additional statement about the solvable blocks is needed or
claimed here.

The executable records are
\path{agents/characters/code/krammer_nielsen_perfect.py} and
\path{agents/characters/code/check_krammer_reverse_cycle.py};
their JSON outputs record the three intertwiners, all twelve
commutator witnesses, all 216 value tests, and the six counts
$(0\text{ failures},36\text{ passes})$.
\end{proof}

\subsection{The full solvable three-puncture closure of the Krammer test}

\paragraph{D23.F24.}
Retain the literal Krammer matrices, the six first-pair values, and
the $216$ pentagon assignments of D23.F13, D23.C9 and D23.P7.
Close each mixed solvable pair under all permutations of its three
peripheral conjugacy classes.  On this simultaneous finite word image,
impose the two-cycle equation and both three-cycle equations obtained
from multipliers $1$ and $-1$.  None of the $216$ assignments is removed:
there remain $36$ assignments for every $(\rho,a_3)\in
\mathbb F_3\times\mathbb F_2$.

More precisely, the simultaneous derived image of the twelve marked
solvable evaluations has order $19683=27^3$.  Exactly twelve elements
of that image satisfy all the stated equations, and their projections
to the original $34$ and $45$ pairs are distinct.  Every listed
pentagon assignment projects to one of these twelve pairs.  All six
first-pair values satisfy the two corresponding three-cycle equations.
These first-pair and mixed-solvable data can be realized simultaneously
by a single ordinary word in $F'$.  This assertion concerns the stated
finite evaluations; it does not give a lift to an actual gentle pair.

\begin{proof}
Write $(A,B)$ for either of the two faithful solvable blocks of
D23.C9, and put $Z=(AB)^{-1}$.  Representatives for the six outer
three-puncture markings are
\[
 (A,B),\quad (B,Z),\quad (Z,A),\quad
 (B,A),\quad (Z,B),\quad (A,Z).
\tag{D23.F24.1}
\]
They represent the full orbit modulo simultaneous conjugation.
Indeed the cyclic operation sends $(A,B)$ to $(B,Z)$ and has order
three; interchanging $A,B$ replaces $Z$ by
\[
 Z'=(BA)^{-1}=A^{-1}ZA=BZB^{-1}.
\]
Consequently its remaining cyclic representatives are conjugates of
the last two pairs in (D23.F24.1), with exactly the corrections used
below.  No uncorrected literal $S_3$ action on the ordered free basis
is being assumed.

Let $v_0,\ldots,v_5$ be the six evaluations of a common word, in the
displayed order.  The two-cycle and the positive three-cycle require
\[
 v_{j+3}=v_j^{-1}\quad(0\leq j\leq2),\qquad
                    v_2v_1v_0=1.
\tag{D23.F24.2}
\]
For an actual multiplier-one gentle pair, right composition with
$(-1,1)$ gives a multiplier-minus-one pair with the same parameter
$f$.  Its three-cycle equation is therefore the further necessary
identity
\[
                  v_2Z^{-1}v_1B^{-1}v_0A^{-1}=1.
\tag{D23.F24.3}
\]
At the reversed base pair $(B,A)$ the three cyclic values, in their
correct order, are
\[
 h_0=v_3,\qquad h_1=A^{-1}v_5A,\qquad
 h_2=Bv_4B^{-1}.
\]
The two reversed equations checked by the certificate are
\[
 h_2h_1h_0=1,\qquad
 h_2(Z')^{-1}h_1A^{-1}h_0B^{-1}=1.
\tag{D23.F24.4}
\]
Cyclic rotations of each identity supply the same equations at the
other markings, so (D23.F24.2)--(D23.F24.4) cover the full orbit.

Here is a reproducible exhaustive computation of the simultaneous
word image and its solution set.  The program
\path{agents/five_level_six/krammer_solvable_s3.py} reads the two
literal $3\times3$ matrices over $\mathbb F_9$ for each block from
\path{agents/characters/records/krammer_mixed_correlation.json}.
It restricts scalars to $\mathbb F_3$, enumerates each group of order
$108$, and constructs its complete multiplication and inversion
tables.  Thus all subsequent operations are exact table operations.
Form the two simultaneous generators $\mathbf X,\mathbf Y$ in the
product of the twelve marked groups.  The orbit of
$[\mathbf X,\mathbf Y]$ under conjugation by
$\mathbf X^{\pm1},\mathbf Y^{\pm1}$ has $432$ elements.
Adjoining these orbit elements successively and closing under
multiplication yields a group of order $19683$, generated by six
of them.  This procedure computes the complete derived image:
the normal closure of the commutator in a two-generator group is
its derived group, since its quotient makes the two generators
commute.  It is not a bounded-length word search.

Testing every one of these $19683$ elements against the displayed
equations for both base pairs gives exactly twelve solutions.  Their
$(v_{34,0},v_{45,0})$ projections are twelve distinct pairs.  The
JSON record stores these counts, and the accompanying NPZ record
stores every one of the twelve complete tuples.

To compare against the original six-dimensional candidates, let
$R=f(X_3,X_4)$ and use D23.C9's literal matrices $C,Q$ to put
\[
 r=C^{-1}RC,\qquad R_{45}=Q^{-1}rQr^{-1}.
\tag{D23.F24.5}
\]
For every one of D23.P7's $216$ matrices $R$, the program conjugates
$R$ and $R_{45}$ into their two faithful splitting bases and takes
their solvable blocks.  The resulting pair belongs to the twelve
allowed pairs in every case.  Counts by the six labels are
\[
\begin{array}{c|rrrrrr}
 (\rho,a_3)&(0,0)&(0,1)&(1,0)&(1,1)&(2,0)&(2,1)\\ \hline
 \text{candidates}&36&36&36&36&36&36\\
 \text{passed}&36&36&36&36&36&36.
\end{array}
\tag{D23.F24.6}
\]
In particular no independent-word substitutions enter this test.

The first-pair orbit can be checked directly without another group
enumeration.  Set $X=K_1^2$, $Y=K_2^2$,
$D=K_1K_2$, and $Z_3=D^3$.  Literal matrix multiplication gives
\[
 DXD^{-1}=Y,\qquad DYD^{-1}=Z_3(XY)^{-1}.
\tag{D23.F24.7}
\]
The matrix $Z_3$ centralizes this braid-group restriction.
Multiplying an argument of a commutator word by such a central
factor does not change its value, because its exponent sums vanish.
Hence, if $F=f(X,Y)$, its two cyclic values must be
$F_1=DFD^{-1}$ and $F_2=D^2FD^{-2}$.
For each of the six actual candidate matrices of D23.F13 the
certificate checks, with $Z=(XY)^{-1}$,
\[
       F_2F_1F=I,\qquad
       F_2Z^{-1}F_1Y^{-1}FX^{-1}=I.
\tag{D23.F24.8}
\]
Their two-cycle is already the defining exact test of D23.F13,
using the half twist $K_1K_2K_1$ which exchanges $X,Y$.
The reversed cycles follow from the same identities and the
explicit simultaneous-conjugation corrections above.

Finally the first-pair derived group has abelianization of order
eight, by D23.P6, whereas the joint mixed-solvable image just
enumerated is a three-group.  They have no common nontrivial
quotient: any nontrivial quotient of the latter is a nontrivial
three-group, so has a quotient $C_3$, while the former has no such
abelian quotient.  Goursat's lemma therefore makes their joint
$F'$ image their full product.  Every first-pair value and allowed
twelve-tuple can be realized simultaneously.  The ordinary free
commutator subgroup has the same image as its profinite closure
in any finite group, giving the stated ordinary-word realization.

The program, its exact JSON counts, and its NPZ solution list are
named \path{krammer_solvable_s3.py},
\path{krammer_solvable_s3.json}, and
\path{krammer_solvable_s3.npz} in
\path{agents/five_level_six/}.  D23.C15 separately handles the
mixed perfect blocks.  Neither result asserts the full identities
in $\widehat F$ or the automorphism condition required of a gentle
parameter; thus the nonempty finite solution set is not a
counterexample to the main full-GT comparison.
\end{proof}

\subsection{The two solvable mixed blocks have independent word images}

\paragraph{D23.F15.}
Let $(x_H,y_H)$ and $(u_H,v_H)$ be the solvable
three-dimensional block pairs of D23.C9 for pairs $34$
and $45$, respectively.  The simultaneous free-group
evaluation has image
\[
 \operatorname{im}\bigl(\widehat F\longrightarrow
       H_{34}\times H_{45}\bigr)=H_{34}\times H_{45},
\]
and hence
\[
 \operatorname{im}\bigl(\widehat F'\longrightarrow
       H_{34}'\times H_{45}'\bigr)=H_{34}'\times H_{45}',
 \qquad |H_{34}'\times H_{45}'|=729.
\tag{D23.F15.1}
\]
In particular the cyclic Nielsen automorphism
$\phi(x)=(xy)^{-1}$, $\phi(y)=x$ does not preserve
the kernel of the $45$-block evaluation on $\widehat F'$.
There is no internal conjugacy, even allowing central
factors in the two arguments, that makes this Nielsen
operation an automorphism of this marked derived image.

\begin{proof}
Both block groups are solvable of order $108$, and
their determinant maps identify their abelianizations
with $C_4$, by D23.C9.  Use $-i$ as generator of $C_4$.
The determinant exponents of the two simultaneous
generators are
\[
 x\longmapsto(2,1),\qquad y\longmapsto(1,1).
\]
Let $J\subset H_{34}\times H_{45}$ be the simultaneous
image. It is subdirect. By Goursat's lemma, a proper
such subgroup would give a nontrivial common quotient
$T$ of the two groups in which these two marked
generators agree.  The group $T$ is solvable and
nontrivial, so $T^{\rm ab}$ is nontrivial.  If $a,b$
are the images in this abelianization of the two
determinant generators, marked agreement gives
$2a=b$ and $a=b$.  Subtraction yields $a=b=0$.
Both maps onto $T^{\rm ab}$ would consequently be
trivial, a contradiction.  Thus $J$ is the full
direct product.  Images commute with taking the
closed derived subgroup under a finite surjection,
giving (D23.F15.1).

There is also an explicit short witness to failure
of kernel invariance.  The word
\[
 w=[x^2,y]\in\widehat F'
\]
is identity in the $34$ block because $x_H^2=I$.
In the $45$ block direct multiplication of the
matrices in D23.C9 gives
\[
 [u_H^2,v_H]=
 \begin{pmatrix}
 -i&0&-1+i\\
 0&1&1\\
 -1+i&0&-1+i
 \end{pmatrix}\ne I.
\tag{D23.F15.2}
\]
The central correction and the intertwiner $C_H$
from D23.C9 give, for every commutator word $f$,
\[
 C_H^{-1}f(x_H,y_H)C_H
       =f((u_Hv_H)^{-1},u_H).
\]
Consequently the $34$-block kernel on $\widehat F'$
is the inverse image under $\phi$ of the $45$-block
kernel.  The two kernels differ by (D23.F15.2), so
the latter is not $\phi$-invariant.  An internal
conjugacy of the marked pair, up to factors
centralizing both inputs, would imply exactly that
invariance on commutator words, which is impossible.
This does not discard the actual gentle correlation
of D23.C9: that correlation uses the two-cycle and
three-cycle identities, which arbitrary commutator
words need not satisfy.
\end{proof}

\subsection{The simultaneous Krammer word image cannot recover the harmonic phase}

\paragraph{D23.F14 (no $S_4$ quotient of the joint Krammer image).}
Use the same specialized matrices as D23.F13 and put
\[
 (X_1,X_2,X_3,X_4,X_5)
 =(K_1^2,K_2^2,K_3^2,(K_1K_2K_1)^2,(K_3K_2K_3)^2).
\]
Let $\Gamma$ be the image of the free profinite group
$\widehat F(x,y)$ under the simultaneous five evaluations
at the ordered pairs
\[
 (X_1,X_2),\quad(X_2,X_3),\quad(X_3,X_4),\quad
 (X_4,X_5),\quad(X_5,X_1).
\tag{D23.F14.1}
\]
Then $S_4$ is not a quotient of $\Gamma$.
Adjoin the marked evaluation $x\mapsto p=(123)$,
$y\mapsto t=(1234)$, and let
$\widetilde\Gamma\subset\Gamma\times S_4$ be the joint
image.  One has
\[
 \{1\}\times V_4\ \subset\ \widetilde\Gamma',
\tag{D23.F14.2}
\]
where $V_4$ is the group of double transpositions.
Consequently, for every commutator word $f$ whose marked
$S_4$ value is $t^{2\nu}p^\rho$, there is another
commutator word $f_*$ with exactly the same five Krammer
values and marked value $t^{2(\nu+1)}p^\rho$.
This is an exact statement about finite word images.
It does not assert that the altered word is an actual
gentle parameter or that it satisfies all profinite
two-cycle and three-cycle identities.

\begin{proof}
We first record exactly which previously computed matrix
facts enter the argument.  The first two pair images are
conjugate, up to exchanging generators, to the group $G$
of D23.F13.  The three remaining pair images have faithful
three-plus-three block decompositions into a solvable
group $H$ of order $108$, with $|H'|=27$, and the full
three-dimensional unitary group $U=\operatorname U_3(3)$
of order $24192$.  In each case the projection onto $U$
is surjective.  These are the literal block decompositions
and exact matrix group computations of D23.C9, in
\texttt{characters/records/krammer\_mixed\_correlation.json}
and the associated mixed-pair certificate. The proof
below needs only these stated facts, not an unverified
identification of a matrix order with a group structure.
The third mixed pair, not separately listed in D23.C9,
is accounted for by the reversing half twist
\[
 D_4=K_1K_2K_3K_1K_2K_1.
\]
The braid identities give
$D_4K_1D_4^{-1}=K_3$, $D_4K_2D_4^{-1}=K_2$ and
$D_4K_3D_4^{-1}=K_1$. Hence
$D_4X_3D_4^{-1}=X_1$ and
$D_4X_4D_4^{-1}=X_5$. Thus pair $(X_5,X_1)$ is
the conjugate by $D_4$ of the reversed pair $(X_4,X_3)$,
which has the same image group as pair $(X_3,X_4)$.
These literal matrix identities have also been checked
directly from the three stored Krammer matrices.
The same half twist proves the asserted relationship
between the first two pair groups.

By D23.F13, $G$ has a normal subgroup $N$ of order
$3^5=243$ whose quotient is
\[
 G/N\simeq\operatorname{SL}_2(\mathbf F_3)\times D_8.
\]
Indeed $N$ is the inverse image of the translation
subgroup of $E_F$ under the diagonal-block map, with
the $D_8$ coordinate constrained to be identity.
It is an extension of a group of order $3^4$ by the
central kernel of order three and is therefore a
three-group.  The group
$\operatorname{SL}_2(\mathbf F_3)$ has normal quaternion
Sylow two-subgroup $Q_8$ and quotient $C_3$.
Taking the inverse image of $Q_8\times D_8$ gives a
normal series for $G$ whose successive factors are,
in order, a three-group, a two-group and a three-group.
Each solvable block $H$ has the shorter series
$1\leq H'\leq H$, with factors a three-group and a
two-group; adjoining a central $C_4$ leaves this property
unchanged.

The determinant map on $U$ takes values in the cyclic
group of norm-one scalars $C_4$.  On scalar matrices
its value is $c^3$, which is an automorphism of $C_4$.
Thus scalar multiplication splits determinant and gives
the direct product
\[
 U=C_4\times S,\qquad S=\operatorname{SU}_3(3).
\]
Here $S$ is the nonabelian simple group of order $6048$;
its scalar centre is trivial because $\gcd(3,4)=1$.
For the standard simplicity fact, see W.~L.~Edge,
\emph{The simple group of order 6048}, Proc. Cambridge
Philos. Soc. \textbf{56} (1960), 189--204, introduction
on p.~189; see also the primary GAP Character Table
Library's \emph{Simple Groups by Group Order}, row
$U_3(3)$ of order $6048$, and its $\operatorname{PSU}(3,3)$
identification in the corresponding character-table record.

After these splittings, the simultaneous image embeds
in $A\times S^3$, where $A$ is a product of two copies
of $G$ and three copies of $H\times C_4$.
Write $A_0$ and $D$ for its projections onto $A$ and
$S^3$.  The first projection is solvable, and $D$ is
subdirect in the three nonabelian simple factors.
Such a subdirect product is a product of diagonal
copies of $S$, and in particular is perfect.  Here is
the short argument in the precise form used.  Induct
on the number of factors.  Goursat's lemma for a
subdirect subgroup of $S^r\times S$ either adds the
last factor independently or identifies it with a
simple quotient of $S^r$.  Every surjection from
$S^r$ to $S$ is projection onto one factor followed
by an automorphism: the factor images are commuting
normal subgroups of the simple centreless group $S$,
so exactly one image is nontrivial.  This proves the
claim by induction.

Apply Goursat again to the subdirect embedding
$\Gamma\subset A_0\times D$.  Any common quotient
would be both solvable, as a quotient of $A_0$, and
perfect, as a quotient of $D$.  A finite group with
both properties is trivial.  Therefore
\[
 \Gamma=A_0\times D.
\]
Every solvable quotient of $\Gamma$ consequently
factors through $A_0$.

Groups admitting a normal series with factors a
three-group, a two-group and a three-group are closed
under finite direct products, subgroups and quotients:
take respectively products, intersections or images
of the three terms of the series.  Hence $A_0$ and
all its quotients have such a series.  But $S_4$ does
not.  Its largest normal three-subgroup is trivial,
and its largest normal two-subgroup is $V_4$; the
quotient $S_4/V_4\simeq S_3$ is not a three-group.
No smaller normal two-subgroup can have three-power
index either.  It follows that $S_4$ is not a quotient
of $A_0$, proving the first assertion.

The map to $S_4$ in the statement is surjective, since
the displayed $p,t$ generate $S_4$.  Goursat's lemma
for $\widetilde\Gamma\subset\Gamma\times S_4$ says
that the quotient of $S_4$ by
\[
 N_4=\{v\in S_4:(1,v)\in\widetilde\Gamma\}
\]
is a quotient of $\Gamma$.  The already proved result
excludes $N_4=1$.  Every nontrivial normal subgroup
of $S_4$ contains $V_4$, so
$\{1\}\times V_4\subset\widetilde\Gamma$.
Moreover $[S_4,V_4]=V_4$: the action permutes the
three nonidentity elements transitively and has no
nontrivial quotient on which it acts trivially.
Taking commutators with lifts of $S_4$ proves
(D23.F14.2).

The image of $\widehat F'$ under a finite surjection
is exactly the derived subgroup of its image.
By (D23.F14.2) choose $u\in\widehat F'$ whose five
Krammer images are identity and whose $S_4$ image
is $t^2$.  Then $f_*=uf$ has unchanged Krammer
values and $S_4$ value
$t^2t^{2\nu}p^\rho=t^{2(\nu+1)}p^\rho$, as claimed.
All images are finite, so a word in the ordinary free
commutator subgroup can equally well be chosen.
The argument never supplies the other identities
needed for an actual gentle pair, which explains the
scope stated above.
\end{proof}

\paragraph{Primary-source URLs checked in DELTA23.}
Edge's original paper record:
\url{https://doi.org/10.1017/S0305004100034472}.
GAP Character Table Library, simple-group table:
\url{https://www.math.rwth-aachen.de/homes/Thomas.Breuer/ctbllib/ctbltoc/views/simplebyorder.html}.
The corresponding primary group record is
\url{https://www.math.rwth-aachen.de/homes/sam/ctbllib/ctbltoc/data/U3(3).html}.
The proof distinguishes the notation $U$ for the full
unitary group of order $24192$ from the ATLAS notation
$U_3(3)$ for its simple projective special factor of
order $6048$.

\subsection{The whole alternating phase is independent of the Krammer word data}

\paragraph{D23.F17 (strengthening D23.F14).}
Let $\Gamma$ be the simultaneous five-pair word image
of D23.F14. Then $S_3$ is not a quotient of $\Gamma$.
For the joint image $\widetilde\Gamma\subset
\Gamma\times S_4$ with the marked map
$x\mapsto(123)$, $y\mapsto(1234)$, one has
\[
 \{1\}\times A_4\subset\widetilde\Gamma'.
\tag{D23.F17.1}
\]
Thus all twelve alternating $S_4$ values can be
prescribed independently of any fixed value of a
commutator word in the five Krammer evaluations.
In particular both of the finite $S_4$ coordinates
occurring in FP.37 can vary at the word level.
No claim of independent variation on actual gentle
or full-GT elements is made.

\begin{proof}
We use the full simultaneous derived-image theorem
D23.P9.  After the reversing-half-twist identifications,
it identifies the image of $F'$ with
\[
 P'\times H_{34}'\times H_{45}'\times H_{43}'
       \times S^3,
\tag{D23.F17.2}
\]
where $P$ is the first-pair group, each $H_j$ is the
solvable order-$108$ block, and $S=SU_3(3)$ is
perfect.  This is a statement about the joint image,
not merely its separate surjective projections.

The abelianization of $P'$ has order $8$, by D23.P6.
Each $H_j'$ has order $27$, exponent three, and
commutator subgroup of order three.  Its abelianization
is therefore the two-dimensional vector space over
$\mathbf F_3$. The explicit conjugation actions of
the marked generators $x,y$ on these three spaces are
\[
 \begin{array}{c|cc}
 &x&y\\\hline
 H_{34}'/H_{34}''&-I&J\\
 H_{45}'/H_{45}''&J&J\\
 H_{43}'/H_{43}''&J&-I
 \end{array},\qquad
 J=\begin{pmatrix}0&-1\\1&0\end{pmatrix}.
\tag{D23.F17.3}
\]
Here $43$ means the reversed $34$ pair, so its row
is obtained by interchanging the marked generators.
To give a complete certificate for the two independent
rows, put $A=[x_H,y_H]$, $B=y_HAy_H^{-1}$ and
$Z=[A,B]$ in the literal block matrices of D23.C9.
The $27$ products $A^aB^bZ^c$, $a,b,c\in\mathbf F_3$,
are distinct and form the entire derived block.
The program \texttt{krammer\_solvable\_action.py}
checks this by exact matrix multiplication, checks
that $Z$ is central of order three, and expresses
the conjugates of $A,B$ by each generator in those
coordinates.  The resulting first two coordinates
are exactly the matrices in (D23.F17.3); all
underlying three-dimensional matrices are retained
in \texttt{krammer\_solvable\_action.json}.
In particular $J^2=-I$ and its characteristic
polynomial $T^2+1$ has no root in $\mathbf F_3$.

Suppose that $\pi:\Gamma\twoheadrightarrow S_3$
were surjective. Then its restriction to $\Gamma'$
would surject onto $S_3'=C_3$, and would be
$\Gamma$-equivariant for conjugation.  The factors
$P'$ and $S^3$ of (D23.F17.2) must be killed: the
former has two-primary abelianization, while the
latter is perfect. On each of the three remaining
factors the induced map is a linear functional
$\ell:\mathbf F_3^2\to\mathbf F_3$.
Conjugation by any element of $S_3$ on $C_3$ is
either $+1$ or $-1$.  Each row of (D23.F17.3) has
a generator acting by $J$. Equivariance would
therefore require $\ell J=\ell$ or
$\ell J=-\ell$. Since $J$ has no eigenvalue $1$
or $-1$ over $\mathbf F_3$, this forces $\ell=0$.
All factor restrictions vanish, contradicting
surjectivity onto $C_3$. Hence no such quotient
$S_3$ exists, even without retaining a marking.

In the Goursat argument of D23.F14, the normal
subgroup $N_4\triangleleft S_4$ now cannot be
$1$ or $V_4$: the corresponding common quotients
would be $S_4$ or $S_3$, both excluded. Thus
$N_4$ contains $A_4$. Furthermore
$[S_4,A_4]=A_4$. One can see this from
$A_4'=V_4$ and the fact that an odd permutation
acts by inversion on $A_4/V_4=C_3$.
Taking commutators with lifts of $S_4$ proves
(D23.F17.1).  Finally the image of $F'$ is exactly
the derived joint image, so multiplication by
a preimage of any $(1,a)$, $a\in A_4$, changes
the $S_4$ value arbitrarily while retaining all
five Krammer values.  This constructs finite
word images only and does not supply the full
gentle identities for the modified word.
\end{proof}

% DELTA23, /root/characters.  Source overlap is stated explicitly below.

\subsection{An induced harmonic quotient and the exact finite ambiguity}

\paragraph{D23.C11 (proved; an audit of the proposed induced $S_4$ quotient).}
Induce the marked permutation representation
\[
 F\longrightarrow S_4,\qquad x\longmapsto(123),\quad
 y\longmapsto(1234),
\]
from the index-six pure subgroup to $B=C_2\amalg C_3$.
The induced representation has degree $24$.  Its $B$ image has
order $41472$ and its $F$ image has derived orders
\[
 6912,1728,64,1.
\]
In particular the abelianization of the first derived image is
a three-group of order $27$.

Join this representation to the signed projective congruence quotient
$Q_{12}$, represented faithfully on its $48$ primitive row vectors
modulo global sign.  The joint $B$ image has order $995328$ and
its pure $F$ image has order $165888$.  The centralizer of the
image of $s$ within that pure group has order $192$ and contains
an element $b$ with positive signed residue
\[
 q_{12}(b)=\begin{pmatrix}9&4\\8&9\end{pmatrix}.
\]
This residue is $I$ modulo four and $S$ modulo three, so its
harmonic parity is odd.  The same $b$ belongs to the first derived
pure image, of order $13824$, whose $s$-centralizer has order $96$.
Thus the finite conjugacy equation for a pure $s$-frame in this
quotient has an odd-parity ambiguity even with zero abelianized
frame.  This is a finite ambiguity assertion, not an assertion
that the ambiguous frame choices lift to actual gentle elements.

\emph{Proof.}
The companion script \path{agents/characters/code/induced_s4_harmonic.py}
uses the literal words
$x=(sr)^{-2}$ and $y=sxs$.  Its six coset representatives are
$1,s,r,r^2,sr,sr^2$.  Every Schreier cocycle is reduced in the
free product and rewritten exactly as a word in $x,y$; the full
rewriting table and the two induced permutations are recorded in
\path{agents/characters/records/induced_s4_harmonic.json}.
The checks $s^2=r^3=1$ and the restrictions
$x|_{\text{identity block}}=(123)$,
$y|_{\text{identity block}}=(1234)$ verify the source marking.
Exact permutation-group stabilizer chains give the displayed orders
and derived subgroups.  The joint image uses the same right-action
convention on primitive row vectors, so its generators are the
same literal $s,r$ in both factors.
For comparison with FP.37's left-action convention, inversion of
the two four-point permutations gives $(p^{-1},t^{-1})$.
This pair is simultaneously conjugate to $(p,t)$ by $(13)$;
thus it has exactly the same marked free-group kernel.  This
convention conversion does not alter the induced quotient used here.

The displayed centralizer witness has induced permutation
\[
 (0\ 1)(2\ 3)(4\ 5)(6\ 7)(12\ 15)(13\ 14)
             (16\ 18)(17\ 19),
\]
with the other points fixed, using zero-based labels.
Its full $72$-point permutation is in the certificate.  Exact
membership in the derived pure group and commutation with the
joint image of $s$ are checked.  The matrix calculation modulo
four and three gives the asserted parity.\hfill$\square$

The Ferrari extension of this quotient contributes no additional
pentagon equation by itself.  Under $a=c=u,b=v$ one has
$(X_1,X_2,X_3,X_4,X_5)=(x,y,x,1,1)$; the two surviving terms
are $f(y,x)f(x,y)=1$.  This is precisely the general forgetful
obstruction already proved in source FP.30, not a new result.

\subsection{A graph proof of the existing relative homology detector}

\paragraph{D23.C12 (proved; independent graph reproof and exact implementation).}
Let $p$ be prime, $q\ge2$, and
$\Gamma=C_p\amalg C_q=\langle a,b\rangle$ be the profinite
free product.  Let $\pi:\Gamma\twoheadrightarrow Q$ be a finite
quotient injective on both factors, and set $U=\ker\pi$.
Then
\[
 \operatorname{pr}_Q C_{\Gamma/([U,U]U^p)}(a)
       =\langle\pi(a)\rangle.
\]
For the literal quotient $Q_{12}$ of $B=C_2\amalg C_3$, the
kernel $U$ is free profinite of rank $97$.  The modulus-two
relative homology extension separates the harmonic frame residues,
and the modulus-three extension separates the cubic ones.  Their
combined quotient
\[
 \Gamma_6=B/([U,U]U^6),\qquad |\Gamma_6|=576\,6^{97},
\]
therefore determines both pure $Q_{12}$ frame residues from the
actual multiplier-one action, and hence determines $a_2,a_3$.

\emph{Proof.}
Use the Bass--Serre tree of the free product, with edges indexed
by $\Gamma$, and vertices indexed by the cosets of its two
factors.  The torsion-free kernel $U$ acts freely.  Its quotient
is the finite connected graph with edges $Q$ and vertices
$Q/\langle\pi(a)\rangle$ and $Q/\langle\pi(b)\rangle$.
Its cycle space over $\mathbb F_p$ identifies with
$U/[U,U]U^p$.  All statements may first be made for the discrete
free product and its finite-index free kernel, then completed;
the finite quotient and its finite homology are unchanged.

Choose as base vertex $o$ the identity coset of the first factor.
An element with $Q$-image $g$ has a path-chain coordinate $v$
with boundary $g o-o$.  Changing the lift by an element of $U$
adds an arbitrary cycle.  The generator $a$ fixes the base vertex
and has zero path-chain coordinate.  Consequently a lift of
$g\in C_Q(\pi(a))$ centralizes $a$ in the relative extension
only if some path chain $v$ with that boundary is fixed by $a$.
This follows directly from the path cocycle rule
$v_{gh}=v_g+g v_h$; the commutator condition is
$(a-1)v=0$.

The action of $a$ on the edge set $Q$ is free, with all orbits
of length $p$.  Every invariant edge chain is a linear combination
of their orbit sums.  At a fixed vertex, the $p$ contributions
of such an orbit have the same incidence sign, and sum to zero
over $\mathbb F_p$.  Hence the boundary of an invariant chain
has coefficient zero at every $a$-fixed vertex.

Both $o$ and $g o$ are $a$-fixed when $g$ centralizes $a$.
If they are distinct, the prescribed boundary $g o-o$ has a
nonzero coefficient at both, a contradiction.  Thus $g o=o$,
or $g\in\langle\pi(a)\rangle$.  Powers of $a$ themselves
centralize $a$ and project onto that factor, proving equality.

For $Q_{12}$, the quotient graph has $576$ edges and
$288+192$ vertices, hence rank $576-288-192+1=97$.
Apply the first assertion with $a=s,p=2$ and, after interchanging
the factors, with $a=r,p=3$.  The finite images of the two
factors meet the pure subgroup $q(F)$ trivially.  Therefore two
pure conjugators of $s$ with the same finite conjugate have the
same $Q_{12}$ residue; the analogous statement holds for $r$.
Actual frames supply existence.  Their positive signed residues
are selected by their identity reduction modulo four.  R3 then
recovers the shear from the frame abelianization modulo three,
and the corrected frame parities follow.\hfill$\square$

The independent computation
\path{agents/characters/code/check_relative_frame_homology.py}
constructs the actual graph, a spanning tree, and all $97$ cycle
coordinates.  It solves every lifting equation
$(a-1)m=-(a-1)v$ over the relevant field.  At two the original
centralizer has order $16$, its liftable image has order $2$,
and $\operatorname{rank}(s-1\mid H_1)=45$.
At three the corresponding numbers are $18,3,63$.
It also checks the fixed-vertex boundary obstruction directly
on every edge-orbit sum.  Full data are in
\path{agents/characters/records/relative_frame_homology.json}.

\paragraph{Exact overlap and remaining scope.}
The source already proves the same frame-separation mechanism:
FP.23, in \path{sources/DELTA_ON_GT/agent_five/five_point_results.tex},
uses $\mathbb F_2[Q/H_2]\oplus\mathbb F_3[Q/H_3]$ and fixed
cosets; RH.2, in \path{sources/DELTA_ON_GT/sections/A_homology.tex},
explicitly says that detector factors through the same $\Gamma_6$.
D23.C12 is a direct graph proof, its general $C_p*C_q$ formulation,
and an exact rank-$97$ implementation of the lifting obstruction.
It is not claimed as a new detector or a new pentagon constraint.
The rank-$129$ and rank-$4097$ Ferrari-kernel tests concern different
cover homologies and retain five-point information.  As FP.30 proves,
simply applying Ferrari and inducing any of the present four-point
detectors leaves the pentagon word identically one for every
two-cycle word.  A comparison with a representation that retains
the relevant five-point kernel remains necessary.

\subsection{A smaller quotient when only the two parities are needed}

\paragraph{D23.C13 (proved; strict compression of the supplied detector).}
There is an explicit quotient of the FP.23 source image with order
\[
 576\cdot2^{34}\cdot3^{52}
\]
such that the value of an actual multiplier-one gentle word $f$
in this quotient determines both corrected frame parities
$a_2,a_3$.  Its ambient permutation modules have dimensions
$72$ over $\mathbb F_2$ and $64$ over $\mathbb F_3$.
This is a quotient for reading actual normalized action values,
or equivalently the word value $f$; invariance of its kernel
under all gentle automorphisms is not asserted.  The stable
quotient $\Gamma_6$ of D23.C12 maps onto it, so an actual action
on $\Gamma_6$ supplies the required readout.

The original FP.23 source image has exact order
$576\cdot2^{97}\cdot3^{96}$.  Thus the present quotient is
strictly smaller; the proved readout here concerns the two
parities, with no full frame-residue claim.  It still gives no additional
pentagon equation after Ferrari extension alone.

\emph{Proof.}
Write $C_s=C_Q(q(s))$ and $C_r=C_Q(q(r))$, of orders $16$ and
$18$.  A signed lift of an element of $C_s$ commutes with $S$
at both primes: a possible simultaneous anticommutation is
excluded modulo four, since it would require a sum of two
squares to be $-1$ modulo four.  At three, such a lift is
$S^j$; at four it is $aI+bS$ with $a^2+b^2=1$.
Define
\[
 \eta_2(g)=j+b\pmod2.
\]
This is independent of the global sign of the lift: changing that
sign adds two to $j$ and leaves $b$ unchanged modulo two.
Both summands are homomorphisms, the second being reduction to
the two-element centralizer of $S$ over $\mathbb F_2$.
It is a nonzero character and $\eta_2(q(s))=1+1=0$.
On the positive pure centralizers with residue $I$ modulo four,
it is exactly the harmonic exponent parity.

For an element of $C_r$, choose a signed lift.  At both moduli
four and three its centralizer has the form
$\{\pm R_0^j: j\in\mathbb Z/3\}$.
The lift commutes with $R_0$ genuinely: anticommutation would
negate its trace $-1$, already impossible modulo three.
Let $e_4,e_3\in\mathbb F_2$ be the respective signs and put
\[
 \eta_3(g)=e_4+e_3.
\]
Changing the global lift changes both signs and preserves their
sum.  This is again a nonzero character, trivial on $q(r)$;
on a positive pure centralizer it is exactly the cubic sign.

Put $H_2^+=\ker\eta_2\subset C_s$ and
$H_3^+=\ker\eta_3\subset C_r$.  Their orders are $8,9$ and
they contain $q(s),q(r)$, respectively.  Form
\[
 V_2=\mathbb F_2[Q/H_2^+],\qquad
 V_3=\mathbb F_3[Q/H_3^+],
\]
of dimensions $72$ and $64$, and let $e_2,e_3$ denote the
identity-coset basis vectors.  As in FP.23 the elements
\[
 S^+=(e_2,0,q(s)),\qquad R^+=(0,e_3,q(r))
\]
in $(V_2\oplus V_3)\rtimes Q$ have orders two and three.
They define a finite homomorphism $\theta^+:B\to E^+$.
Since $\langle q(s)\rangle\subset H_2^+$ and
$\langle q(r)\rangle\subset H_3^+$, the two coset projections
give a surjective module map from FP.23's ambient modules to
these modules, compatible with both generators.  In particular
$E^+$ is a quotient of the actual FP.23 image.

If $(m,g)$ centralizes $S^+$, its projection centralizes $q(s)$
and satisfies $(1-q(s))m_2=(1-g)e_2$.  Passing to coinvariants
shows that the two $q(s)$-fixed cosets $H_2^+$ and $gH_2^+$
must coincide; hence $g\in H_2^+$.  Thus every centralizer
ambiguity has $\eta_2(g)=0$.  The same argument with $R^+$
gives $g\in H_3^+$ and $\eta_3(g)=0$.

Given the value $\theta^+(f)$ for an actual gentle element,
the normalized formulas give
\[
 \theta^+\Phi(s)=S^+\theta^+(f),\qquad
 \theta^+\Phi(r)=\theta^+(f)^{-1}R^+.
\]
Solve the finite conjugacy equations for frames in $\theta^+(F)$
whose positive signed $Q$-residues are $I$ modulo four.
Actual frames guarantee solutions.  The ratio of two harmonic
solutions centralizes $S^+$ and has positive residue $I$
modulo four, so its $\eta_2$ is precisely its possible change
of harmonic parity; that change is zero.  The corresponding
argument with $\eta_3$ proves uniqueness of the cubic parity.
The shear is also known: the positive signed value of $f$
has $I$ residue modulo four, and its upper-right entry modulo
three is $\rho$ by R3.  Hence $V=U^{8\rho}$ modulo twelve
is determined before the corrected parities are read.
This proves the asserted readout without assuming every finite
admissible conjugacy solution lifts to an actual gentle element.

For the exact size, the implementation uses $576$ Schreier
representatives for $U=\ker q$.  If $v_g$ is the affine
translation of a chosen representative and $t_a$ that of a
generator $a\in\{s,r\}$, the Schreier kernel generator has
translation
\[
 v_g+g t_a-v_{ga}.
\]
These $1152$ vectors span exactly the image of $U$ in each
module.  Exact row reduction gives ranks $34$ and $52$.
The joint image of $U$ is the full product of these elementary
abelian groups, because their orders are coprime and the joint
image surjects onto both.  Surjectivity onto $Q$ gives the
stated total order.  The same calculation in FP.23's original
modules gives ranks $97$ and $96$, proving strictness.
Every character equation, every coset obstruction, the Schreier
spans, and the image ranks are checked by
\path{agents/characters/code/check_compressed_parity_detector.py},
with exact records in
\path{agents/characters/records/compressed_parity_detector.json}.
Finally FP.30 applies to every quotient factoring through the
forgetful map, so compression does not change its automatic
pentagon conclusion.\hfill$\square$

\subsection{Inducing the marked $S_4$ quotient does not remove the harmonic frame ambiguity}

\paragraph{D23.F16.}
Combine the first-pair Krammer representation of D23.F13
with the induced $24$-point representation of
$\widehat B=C_2\amalg C_3$ from D23.C11, and retain
the signed projective level-twelve quotient as well.
This combination is formed as a representation of
$\widehat B_3$, pulling the latter two factors back
along $\widehat B_3\to\widehat B$. In its Krammer
coordinate, ``centralizing $s$'' means centralizing
the half-twist lift $s_P=K_1K_2K_1$, whose square
is central; the six-dimensional representation is
not asserted to factor through $\widehat B$.
The finite centralizer of $s$ in the image of $F'$
still contains an element $\beta$ with
\[
 q_4(\beta)=I,\qquad q_3(\beta)=S
 =\begin{pmatrix}0&1\\-1&0\end{pmatrix}.
\tag{D23.F16.1}
\]
Its induced $24$-point coordinate can be the identity.
Thus this combined quotient does not determine the
harmonic parity from the finite harmonic frame equation,
even after imposing $\beta\in\operatorname{im}(F')$.
This is a finite ambiguity, not an assertion that the
nonidentity finite conjugator is an actual pure frame.

\begin{proof}
Let $P$ be the Krammer first-pair group.  The exact
derived orders of D23.P6 are
$46656,3888,486,243,3,1$. Thus $P'$ is solvable
and its abelianization has order $8$.
Let $E$ be the restriction to $F$ of the induced
representation in D23.C11. Its exact derived orders
are $6912,1728,64,1$, so $E'$ is solvable and its
abelianization has order $27$.
The simultaneous image of $F'$ in $P'\times E'$
is subdirect. A nontrivial common quotient of these
solvable groups would have a nontrivial abelianization,
which would be a quotient of both a two-group and
a three-group. This is impossible. Goursat's lemma
therefore gives the full product $P'\times E'$.

Under the diagonal-block projection of D23.F13,
\[
 P'\longrightarrow E_F'\times D_8'
\]
is surjective, with central kernel $\langle n\rangle$
of order three.  In the affine coordinates of D23.P3
choose the first coordinate
$\gamma=(0,S)\in k^2\rtimes Q_8=E_F'$,
and choose the second coordinate to be identity.
It follows that there is $\beta_P\in P'$ with
these diagonal values.

Put $s_P=K_1K_2K_1$. Its affine $W$ coordinate is
$(0,iS)$, which commutes with $\gamma$; its quotient
coordinate commutes with identity. Thus
$[s_P,\beta_P]\in\langle n\rangle$.
This kernel is central in $\langle K_1,K_2\rangle$,
and $s_P^2=(K_1K_2)^3$ is also central there.
Writing $s_P\beta_Ps_P^{-1}=\beta_Pn^j$ and
conjugating once more gives $n^{2j}=1$.
Since $n$ has order three, $j=0$. Every such lift
already centralizes $s_P$.

By the full-product statement, the pair
$(\beta_P,1)$ is the image of an element of $F'$.
It centralizes $s$ in both representations.  The
signed level-twelve coordinate of an element of
$F'$ has identity component modulo four: the two
signed generators $U^2,L^2$ commute modulo four,
so their commutator subgroup is killed there.
On $F'$ the linear part of the affine detector is
exactly $q_3$, since the modulus-four character
occurring in D23.P3 is trivial. Our choice of
$\gamma$ therefore gives $q_3(\beta_P)=S$.
This also shows that adjoining the level-twelve
coordinate imposes no further restriction on the
chosen pair: on $F'$ it is already determined by
the Krammer linear coordinate.  The resulting
matrix is
\[
 q_{12}(\beta)=\begin{pmatrix}9&4\\8&9\end{pmatrix},
\]
whose two reductions are those of (D23.F16.1).

For a literal finite certificate, the program
\texttt{krammer\_induced\_frame\_ambiguity.py}
finds the three lifts of $(\gamma,1)$ inside the
enumerated $P'$, checks that all commute with
$s_P$, and chooses the unique lift of order four.
Its complete six-dimensional matrix is in the
companion JSON. The existence proof above does
not depend on a guessed lift or on an assumed
congruence-kernel invariance.

At the identity action the harmonic frame equation
is simply that the frame conjugator centralizes $s$.
Both the identity and the constructed finite
conjugator solve this equation and satisfy the
positive modulus-four normalization. The latter
has odd harmonic residue modulo three, whereas
the former has even residue. Thus the finite
equation does not recover the actual parity.
The actual pure centralizer in the profinite
source is trivial; our assertion concerns only
the specified finite quotient.
\end{proof}

\subsection{A two-primary harmonic cover still separates from the cubic Krammer state}

\paragraph{D23.F19.}
Let $\Gamma_K$ be the common-word Krammer image of
D23.F14 and D23.P9. Its maximal two-group quotient
on the derived subgroup is
\[
 \Gamma_K'\twoheadrightarrow Q_8\times C_2.
\tag{D23.F19.1}
\]
On the six actual-gentle necessary first-pair values
of D23.F13, this quotient takes the value
\[
 \left(M_\rho,1\right),\qquad
 M_\rho=\begin{pmatrix}1&\rho\\\rho&1+\rho^2\end{pmatrix},
\tag{D23.F19.2}
\]
independently of the cubic bit $a_3$.

Let $E_2$ be any finite quotient of $\widehat B$
mapping onto $Q_{12}$ with two-group kernel; this
includes the relative mod-two harmonic graph
cover of D23.C12. Put $T_2=\operatorname{im}(F'\to E_2)$.
In the simultaneous image of $F'$ in
$\Gamma_K'\times T_2$, every Krammer value in the
kernel of (D23.F19.1) can be paired with the
identity of $T_2$.
Consequently all the reduced Krammer pentagon
solutions with $\rho=0,a_3=1$ can be combined
with the identity harmonic-cover word value.
By D23.P9 these combinations have a common
commutator-word realization. This refutes a
finite certificate obtained merely by adjoining
this harmonic cover by product and the Ferrari
forgetful map. It does not produce an actual
mixed-sign gentle solution of the full pentagon.

\begin{proof}
For the first-pair group $P$, D23.F13 gives
a central order-three kernel over
$E_F\times D_8$. On derived groups the quotient
is $E_F'\times D_8'$, where
$E_F'=k^2\rtimes Q_8$ and $D_8'=C_2$.
The inverse image of $k^2\times1$ is a normal
three-subgroup $N$ of order $3^5=243$ in $P'$,
and
\[
 P'/N\simeq Q_8\times C_2.
\]
Every homomorphism from $P'$ to a two-group
kills $N$, since the order of an image of a
three-group must be both a power of three
and a power of two. Conversely the displayed
quotient is itself a two-group. It is therefore
the maximal two-group quotient of $P'$.

D23.P9 identifies the remaining factors of
$\Gamma_K'$ as three groups $H_j'$ of order
$27$ and three perfect groups $SU_3(3)$.
Every homomorphism to a two-group kills the
former by order and the latter because a
finite two-group is solvable and cannot
contain a nontrivial perfect image. Thus
the maximal two-group quotient of the
entire common-word image is precisely
(D23.F19.1).

In the affine $W$ coordinate, the six values
have linear part $M_\rho$, by D23.P3; their
translation part is where $a_3$ occurs.
Their $D_8$ quotient coordinate is identity,
by D23.F13. Killing $N$ removes exactly
these translations and the central
order-three lift, which proves (D23.F19.2).
In particular both values with $\rho=0$
are in the kernel of the maximal two-group
quotient.

The signed level-twelve image of $F'$
is $Q_8$: modulo four it is identity,
and modulo three it is
$\operatorname{SL}_2(\mathbf F_3)'=Q_8$.
The first assertion follows directly
because $U^2,L^2$ commute modulo four;
the exact image modulo three is the
standard pure-pair image. Thus $T_2$ is
an extension of a subgroup of a two-group
by a subgroup of the assumed two-group
kernel of $E_2\to Q_{12}$, and is itself
a two-group.

Let $J\subset\Gamma_K'\times T_2$ be
the joint image. It is subdirect. The
common quotient in Goursat's lemma is
a two-group because it is a quotient
of $T_2$. The corresponding map from
$\Gamma_K'$ factors through its maximal
two-group quotient. Hence if $g$ belongs
to the kernel of (D23.F19.1), its value
in that common quotient is identity;
the fibre-product description gives
$(g,1)\in J$.

D23.P7 supplies $36$ reduced solutions
for each pair $(\rho,a_3)$, and D23.P9
supplies common commutator words for
their full five-factor assignments.
For $\rho=0,a_3=1$ their first value
has trivial maximal two-group image,
and none of the mixed factors contributes
to that quotient. The just-proved
Goursat assertion therefore supplies
common words realizing those same
Krammer assignments and the identity
in $T_2$.

Under the Ferrari forgetful map the
five cyclic inputs are exactly
$(x,y,x,1,1)$, as in FP.30. Consequently
a commutator word that is identity in
the harmonic cover has identity in
each of its five evaluation factors:
the reversed pair is conjugate by
the half twist, and factors with an
identity argument are identity.
The combined finite pentagon is
therefore satisfied. If the harmonic
frame is separated by the cover of
D23.C12, the identity normalized
action has the even harmonic residue.
Thus the added cover cannot discard
the unequal-bit finite assignments.
The common words need not satisfy
the full profinite gentle relations;
no claim of such a lift is used.
\end{proof}

\paragraph{D23.F20 (the whole relative pro-two Ferrari family).}
Let $L=\widehat B_4/\langle z_4\rangle$, let
$q_L:L\to\widehat B$ be the Ferrari forgetful map,
and put $N=\ker(q_{12}q_L)$.
The actual mixed-sign gentle witness $k$ of
DELTA\_ON\_GT FP.36, supplied by the strengthened
relative lifting input LI.8, satisfies every
five-cycle factor separately in
\[
 L/O^2(N).
\]
This holds for the entire relative pro-two
completion, not just its first homology layer.

\begin{proof}
The kernel of the map from this quotient
to $Q_{12}$ is pro-two. For each of
the five adjacent evaluation pairs,
its image in $Q_{12}$ is the ordinary
pure pair or a cyclic group with
one generator trivial, by FP.30.1--2
and the reversing half twist.
The pure pair has derived series
of orders $96,8,2,1$; in particular
its second derived group is $C_2$.
The other base pairs have trivial
second derived group. Hence the
second derived subgroup of every
pair in $L/O^2(N)$ is pro-two:
it is an extension of a pro-two
kernel by a finite two-group.

FP.36 gives one actual gentle
$k=(1,f_k)$ with $\Delta(k)=1$
and $f_k\in R_{\rm nil}(F'')$.
Every continuous homomorphism
from $F''$ into a pro-two group
kills this nilpotent residual,
because all its finite quotients
are nilpotent. Apply this to
each evaluation $F''\to P_j''$.
Each of the five values of $f_k$
is identity, proving the claim.
This uses the stated actual
lifting input; it does not
independently reconstruct LI.8.
\end{proof}

\subsection{The joint cubic--quartic cyclic cover}

\begin{statement}{D23.F21}{Proved}{The degree-twelve cyclic cover is blind at every two-adic and three-adic homology level}
\end{statement}
Let $C$ be the connected cyclic cover of the five-punctured sphere with
deck group $C_{12}$ and local monodromies
\[
 (9,9,10,4,4)\in(\mathbb Z/12)^5.
\]
Fill all punctures upstairs, and let $T$ denote the positive deck
transformation on the resulting integral homology. The compact curve has
genus $11$. Its degree-four and degree-three quotient covers have branch
data $(1,1,2,0,0)$ and $(0,0,1,1,1)$, respectively, and each has genus one.
Thus this single cover is a natural geometric attempt to retain the
quartic and cubic elliptic structures together.

Use the five pure braid elements
\[
 X_1=\sigma_1^2,\quad X_2=\sigma_2^2,\quad
 X_3=\sigma_3^2,\quad X_4=(\sigma_1\sigma_2\sigma_1)^2,
 \quad X_5=(\sigma_3\sigma_2\sigma_3)^2,
\]
and descend their homology action from $PB_4$ to the pure five-point
mapping class group by the linking-number normalization specified below.
For $p=2$ and for $p=3$, the closed subgroup generated by any adjacent
pair $(X_j,X_{j+1})$, with indices modulo five, has pro-$p$ second
derived subgroup in its action on $H_1(C,\mathbb Z_p)$.

Consequently the supplied actual gentle element $k$ of
\textup{FP.36}, with
\[
 f_k\in R_{\rm nil}(\widehat F(x,y)'')\qquad\hbox{and}\qquad
 \Delta(k)=1,
\]
has all five evaluations $f_k(X_j,X_{j+1})$ equal to the identity on
both $H_1(C,\mathbb Z_2)$ and $H_1(C,\mathbb Z_3)$, hence on every
finite quotient of these lattices. This assertion uses the supplied
existence statement for $k$; it does not reconstruct that lifting input.
It proves no claim about other primes or about nonabelian covering
invariants.

\begin{proof}
The sum of the branch weights is zero modulo $12$, and the weights
generate $\mathbb Z/12$, proving that the cover is well-defined and
connected. A puncture of weight $w$ has $\gcd(12,w)$ preimages, each
with ramification index $12/\gcd(12,w)$. The five numbers of preimages
are $3,3,2,4,4$. Removing the five branch values gives Euler
characteristic $12(2-5)=-36$; adding the $16$ points upstairs gives
$\chi(C)=-20$, so its genus is $11$. The same count at degrees four
and three gives Euler characteristic zero for both stated quotient
covers.

Here is a complete finite cellular model, including the marking used
by the certificate. Remove the fifth point and represent the other four
inertias by the free generators $x_1,x_2,x_3,x_4$. On the degree-twelve
cover of the resulting four-petal rose put
\[
 C_0=\mathbb Z^{12},\qquad C_1=\mathbb Z^{48},\qquad
 \partial e_{j,g}=v_{g+w_j}-v_g,
 \quad(w_1,w_2,w_3,w_4)=(9,9,10,4).
\]
Here $g\in\mathbb Z/12$, and $T(e_{j,g})=e_{j,g+1}$.
The graph is connected, so $\ker\partial$ has rank $37$.
The compact surface is obtained by filling all lifted peripheral
circles: for the first four punctures these are the translates of
$x_j^{12/\gcd(12,w_j)}$, and for the fifth they are the translates of
$(x_1x_2x_3x_4)^3$. Choose one translate for every coset of
$\langle w_j\rangle$ or $\langle8\rangle$, respectively.
There are sixteen such circles, with one relation, so their image has
rank fifteen and the quotient homology has rank $22$.
Equivalently, filling the boundary circles of a connected orientable
surface gives precisely their sole total-boundary relation. This also
shows that the quotient lattice is torsion-free.

The positive Artin automorphism is
\[
 \sigma_i(x_i)=x_ix_{i+1}x_i^{-1},\qquad
 \sigma_i(x_{i+1})=x_i,
\]
and fixes the other generators. The usual braid relations follow by
substitution in this free group; the certificate also checks all three
relations relevant to $B_4$ as literal reduced words. Each pure braid
preserves the weight map, acts on the displayed graph chains, preserves
the peripheral submodule, and commutes with $T$.

Write $P=x_1x_2x_3x_4$. The full twist
$z_4=(\sigma_1\sigma_2\sigma_3)^4$ acts by conjugation by $P$.
On the homology of the kernel of the weight map, this is the deck
action $T^8$. This identity follows because $P$ has weight $8$ and
changing a lift by an element of the kernel changes the action only
by an inner automorphism, which is trivial on its abelianization.
Let $\ell_{12}:PB_4\to\mathbb Z$ be the first-two-strands linking
number. Since $\ell_{12}(z_4)=1$, the corrected representation
\[
 \gamma\longmapsto T^{-8\ell_{12}(\gamma)}q(\gamma)
\]
is trivial on $z_4$ and descends to $PB_4/\langle z_4\rangle$.
For the five displayed elements the linking numbers are
$(1,0,0,1,0)$, so the correction multiplies precisely $X_1$ and
$X_4$ by $T^4$. Because the deck action is central, this correction
does not change the evaluation of any word $f\in\widehat F'$ at a
pair of elements.

Reduce the homology modulo $p$. The semisimple submodule, or socle,
for the deck action is
\[
 S_p=\bigoplus_h\ker h(T),
\]
where $h$ runs through the distinct monic irreducible factors of
$T^{12}-1$ in $\mathbb F_p[T]$. For $p=2$ these are
$T-1,T^2+T+1$; for $p=3$ they are $T-1,T+1,T^2+1$.
The chain calculation just described gives the following table. Each
entry is the sequence of orders of the successive derived groups of
the indicated adjacent-pair image on the indicated socle summand;
the last group is trivial.
\[
\begin{array}{c|c|c|ccccc}
 p&h&\dim_{\mathbb F_p}\ker h(T)&12&23&34&45&51\\\hline
2&T-1&1&1&1&1&1&1\\
2&T^2+T+1&6&(9,1)&(36,4,1)&(36,4,1)&(36,4,1)&(9,1)\\
3&T-1&1&1&1&1&1&1\\
3&T+1&3&(2,1)&(2,1)&1&(2,1)&(2,1)\\
3&T^2+1&6&(36,9,1)&(72,9,1)&(2,1)&(4,1)&(72,9,1)
\end{array}
\]
In particular every adjacent-pair image on $S_p$ is metabelian.

For reproducibility, \texttt{cyclic12\_cover.py} builds the boundary
matrix and all sixteen peripheral columns over the integers, computes
the compact quotient separately over $\mathbb F_2$ and $\mathbb F_3$,
and derives each pure braid matrix by tracing the literal Artin words
through the covering graph. It checks that every matrix preserves the
relation subspace, that the compact dimension is $22$, that $T^{12}=1$,
that the pure matrices commute with $T$, and that the full twist acts
as $T^8$. It then restricts to every irreducible deck socle summand.
For each pair it closes the orbit of a basis under the pair's matrices;
the resulting permutation action is faithful, since a linear map
fixing that orbit fixes the basis. Exact permutation derived-group
calculations produce the table. The outputs are
\texttt{cyclic12\_cover.json} and the two \texttt{.npz} matrix records.
Thus the computation starts with the specified covering and free-word
action, rather than assuming guessed braid matrices or a quotient
identification.

We next justify why this finite table controls all $p$-adic depths.
For a finite-dimensional $\mathbb F_p[T]$-module $V$, the kernel of
\[
 \operatorname{Aut}_{\mathbb F_p[T]}(V)
 \longrightarrow\operatorname{Aut}_{\mathbb F_p[T]}(\operatorname{soc}V)
\]
is a $p$-group. Indeed, if the kernel had an element of prime order
$q\ne p$, extend scalars to an algebraic closure. That element is
semisimple and, if nonidentity, has a nonzero eigenspace with eigenvalue
different from one. The eigenspace is $T$-stable because the two actions
commute. A nonzero $T$-stable space contains a $T$-eigenvector, which
lies in the scalar extension of the socle. This contradicts trivial
action on the socle. Cauchy's theorem now proves the claim.

Every adjacent-pair second derived group acts trivially on $S_p$ by
the table, so its image on $H_1(C,\mathbb F_p)$ is a $p$-group.
The kernel of reduction from $\operatorname{GL}_{22}(\mathbb Z_p)$
to $\operatorname{GL}_{22}(\mathbb F_p)$ is pro-$p$, as is seen from
the successive congruence kernels. Therefore the closed second
derived group of each $p$-adic pair image is pro-$p$.

Finally the restriction of each evaluation homomorphism
$\widehat F\to\operatorname{GL}_{22}(\mathbb Z_p)$ to
$\widehat F''$ has pro-$p$ image. Every finite continuous quotient
of that image is a finite $p$-group and hence nilpotent. It follows
directly from the definition of $R_{\rm nil}(\widehat F'')$ as the
intersection of kernels of finite nilpotent quotients that the
restriction kills $R_{\rm nil}(\widehat F'')$. Apply this to the
supplied $f_k$ for each of the five pairs and for $p=2,3$.
\end{proof}

\paragraph{Scope and source locators.}
The five pure braid elements and their pentagon order are those in
\texttt{DELTA\_ON\_GT/agent\_five/five\_point\_results.tex},
\textup{FP.30}. The sole actual-element existence input is
\textup{FP.36} in the same file, equivalently the strengthened
relative lifting witness used in \textup{A.94} and \textup{A.109}.
The cyclic covering, its chain action, the socle calculation, and the
all-depth implication above are new calculations in this run.

\subsection{Two noncentral Hurwitz tests and an explicit finite escape}

\paragraph{D23.P11 (genuine five-point Hurwitz quotients).}
For a finite group $G$ and $p\in G$, the Hurwitz moves
\[
 (g_i,g_{i+1})\longmapsto(g_ig_{i+1}g_i^{-1},g_i)
\tag{D23.P11.1}
\]
act on four-tuples of product $p$, modulo simultaneous conjugation
by $C_G(p)$. Every orbit is a permutation representation of
$B_4/\langle z_4\rangle$, where $z_4=(\sigma_1\sigma_2\sigma_3)^4$.
In the following two explicit cases the product is nonidentity:

\begin{enumerate}
\item Let $G=\mathbb F_3^2\rtimes\operatorname{GL}_2(\mathbb F_3)$,
of order $432$, and let all four entries lie in the conjugacy
class of $(0,\operatorname{diag}(1,-1))$, of size $36$.
Take $p=((0,1)^{\mathsf T},U)$. Its centralizer has order $9$.
The quotient set of such tuples has $891$ elements and four
$B_4$-orbits of sizes $9,72,162,648$. Use the orbit of size
$648$, containing
\[
 \left(
  (0,\begin{psmallmatrix}0&1\\1&0\end{psmallmatrix}),
  (0,\begin{psmallmatrix}0&1\\1&0\end{psmallmatrix}),
  (0,\begin{psmallmatrix}1&1\\0&2\end{psmallmatrix}),
  ((1,2)^{\mathsf T},\begin{psmallmatrix}1&2\\0&2\end{psmallmatrix})
 \right).
\]
Call this representation $q_{648}$.
\item Let $G=E=\mathbb F_9^2\rtimes H_{48}$ be D23.P3,
with $s=(0,iS)$ and $r=(e_1,R_0)$. The tuple $(s,s,r,r)$
has product $p=r^2\ne1$, whose centralizer has order $27$.
Its Hurwitz orbit modulo this centralizer has size $20736$.
Call this representation $q_{20736}$.
\end{enumerate}

Both tests contain mixed evaluation pairs with nonnilpotent
second derived subgroups; thus neither of these selected tests
is dismissed by the supplied FP.36 nilpotence criterion.
For example, a degree-$144$ component of
$\langle X_3,X_4\rangle$ in $q_{648}$ has derived orders
\[
408146688,\quad102036672,\quad34012224,\quad531441,\quad1.
\tag{D23.P11.2}
\]
It admits the literal marked quotient
$X_3\mapsto(123),X_4\mapsto(1234)$ onto $S_4$.
A degree-$96$ mixed component of $q_{20736}$ has derived orders
\[
3265173504,\quad272097792,\quad4251528,\quad531441,\quad1.
\tag{D23.P11.3}
\]
Here, as throughout this branch,
\[
 X_1=A^2,\quad X_2=B^2,\quad X_3=C^2,\quad
 X_4=(ABA)^2,\quad X_5=(CBC)^2.
\]

\begin{proof}
The elementary Hurwitz moves satisfy the braid and distant
relations, preserve the product, and commute with simultaneous
conjugation. The full twist acts by simultaneous conjugation
by the product (or its inverse if the opposite action convention
is chosen); this is trivial after quotienting by $C_G(p)$.
Thus these are quotients of the label-five stabilizer used by
the actual five-cycle, not the product-one forgetful family.

All displayed orbit sizes and component orders are finite exact
computations. The program
\path{agents/pentagon/hurwitz_affine_reflections.py}
enumerates the first set using the first three entries of a tuple,
the fourth being determined by its product. It quotients by the
explicit centralizer and checks all braid and full-twist relations.
The program \path{agents/pentagon/hurwitz_cubic_affine.py}
builds $E$ from its literal affine generators, enumerates the
single orbit of $(s,s,r,r)$, and checks the same relations.
There is no random orbit sampling in these constructions.

For (D23.P11.2)--(D23.P11.3), the exact permutation groups are
computed on a component of the specified evaluation pair, and
their derived subgroups are computed successively. The order-$531441$
third derived subgroup is elementary abelian of exponent three;
the preceding derived subgroup is checked to be nonnilpotent.
Adjoining the marked four-point action to the degree-$144$
component leaves its group order unchanged, proving the stated
marked $S_4$ quotient. See
\path{hurwitz_affine_affine_unipotent_3_components.json},
\path{hurwitz_cubic_ssrr_components.json}, and
\path{hurwitz_affine_mixed_marked_quotients.json}.
The first of these analyzes all components of the selected
$648$-point orbit; the second leaves larger components unanalyzed,
and no assertion about those components is made here.
\end{proof}

\paragraph{D23.P12 (one ordinary word defeats both finite tests).}
Use $[a,b]=aba^{-1}b^{-1}$ and let $\iota$ interchange $x,y$.
Define ordinary words
\[
\begin{aligned}
 c&=[x,y],&d&=[c,xcx^{-1}],&u&=[d,xdx^{-1}],\\
 v&=u(c^4uc^{-4}),&w&=[c^6,v]^{-1},\\
 b&=(xy)w^{-1}(xy)^{-1},& f_*&=\iota(b)b^{-1}.
\end{aligned}
\tag{D23.P12.1}
\]
Then $f_*\in F'''$ and $\iota(f_*)=f_*^{-1}$ as an exact
identity of ordinary free words. Its expanded length is $904$
(freely reduced length $588$). In D23.P3 and D23.P10 it has
the exact values
\[
 \eta(f_*)=((1-i,1+i)^{\mathsf T},I),\qquad
 \xi(f_*)=((1,1)^{\mathsf T},I),
\tag{D23.P12.2}
\]
while its marked $S_4$ value is $1$. In both $q_{648}$ and
$q_{20736}$, all ten orientations of the five pentagon pairs
have identity $f_*$-value. In particular the entire literal
five-cycle equals $1$ in both representations.
The same identity values hold on the entire Nielsen $S_3$
closure of these pairs, modulo inner conjugacy. Thus all gentle
two- and three-cycle identities in that finite closure hold.

Moreover, its complete first-pair Krammer value is the matrix
$F_{0,1}$ of D23.F13, not merely the projection of that matrix
to the cubic detector. Thus these finite tests, even taken
together with the exact first-pair Krammer state and the marked
$S_4$ state, do not discard the mixed state
$(\rho,a_3,\nu_4)=(0,1,0)$.

This statement concerns an ordinary word and specified finite
images. It does not claim the free three-cycle for $f_*$,
gentle admissibility, the Krammer pentagon, or a full-GT lift.

\begin{proof}
The derived-depth assertion follows directly from the nested
commutators: $c\in F'$, $d\in F''$, $u\in F'''$, and
normality of $F'''$ puts $v,w,b,f_*$ in $F'''$. The identity
\[
 \iota(f_*)=b\iota(b)^{-1}=f_*^{-1}
\]
is formal. The $S_4$ value is therefore trivial already because
$S_4'''=1$; this part does not use an unproved gentle identity.

The remaining assertions are exact evaluations of the explicitly
given ordinary word. To make the computation reproducible,
\path{agents/pentagon/hurwitz_escape_word.py} constructs the
expanded word from (D23.P12.1), freely checks the two-cycle,
checks the saved Hurwitz generators' braid and full-twist
relations again, and evaluates the word on the five ordered pairs
\[
 (X_3,X_4),(X_5,X_1),(X_2,X_3),(X_4,X_5),(X_1,X_2)
\tag{D23.P12.3}
\]
and on their reversals. All twenty resulting permutations are
identity. In fact the intermediate $448$-letter word $w$ has
identity values on these same twenty pairs.
For each ordered pair $(a,b)$ and its reversal the certificate
also sets $z=(ab)^{-1}$ and evaluates $f_*$ on
$(a,b),(b,z),(z,a)$. All sixty values across the two families
are identity. These representatives suffice for the entire
Nielsen $S_3$ closure: the transformations
$\tau:(x,y)\mapsto(y,(xy)^{-1})$ and
$\iota:(x,y)\mapsto(y,x)$ satisfy
$\tau^3=\iota^2=1$ and $(\iota\tau)^2=\operatorname{Inn}(x^{-1})$.
Inner conjugation merely conjugates a word value. In particular
each product $f_*(z,a)f_*(b,z)f_*(a,b)$ is identity in both tests.

The original affine evaluation, with the encoding $a+3b$ for
$a+bi\in\mathbb F_9$, is
\[
 \eta(w)=((1,1-i)^{\mathsf T},I),\qquad
 \eta(b)=\eta(f_*)=((1-i,1+i)^{\mathsf T},I).
\]
The change from $w$ to $b$ is literal affine conjugation; the
change from $b$ to $f_*$ is also evaluated, rather than inferred
from an assumed three-cycle. The program independently checks
the compressed $3\times3$ homogeneous matrix
\[
 \xi(f_*)=\begin{pmatrix}1&0&1\\0&1&1\\0&0&1\end{pmatrix}.
\]
Finally it evaluates $f_*$ at the actual $12\times12$ first-pair
Krammer matrices and compares all entries with the saved
D23.F13 matrix labelled $(0,1)$. This equality also follows
from D23.F13: the quotient $D_8$ is killed by derived depth,
the affine block has the required value, and the exact two-cycle
fixes the remaining central phase uniquely.

The exact JSON certificate stores the entire word, all zero
support counts, both affine values, the Krammer label, and hashes
of the two finite Hurwitz actions. No existence of an unspecified
word or a profinite lift is used in these finite evaluations.
\end{proof}

\paragraph{Scope of the failure.}
The two tested Hurwitz representations are genuinely new mixed
tests: one sees the marked harmonic $S_4$ quotient, and both
have a nonnilpotent second derived mixed pair. Their failure is
therefore established by an explicit word, not by misapplying
the earlier nilpotence criterion. The full Nielsen closure was
checked as a further test and is also annihilated by this one
word. D23.C15 and D23.F24 subsequently verify the full Krammer
Nielsen cycle constraints. D23.P14 combines the two families
with Krammer by one common finite word; no actual gentle or
full-GT lift is supplied.

\subsection{A common finite word for the combined Hurwitz--Krammer tests}

\paragraph{D23.P14 (combined finite equations still admit a mixed state).}
There is an ordinary word $f\in F'''$ whose values simultaneously
satisfy the following finite conditions:
\begin{enumerate}
\item the complete Krammer pentagon of D23.P7, with first-pair
state $F_{0,1}$ and identity solvable blocks in the mixed pairs;
\item every two-cycle, three-cycle, reversed three-cycle and
$m=-1$ three-cycle in the full Nielsen closure of those Krammer
pairs, as checked in D23.C15 and D23.F24;
\item identity values on the entire Nielsen closure of all five
pairs in both Hurwitz representations of D23.P11;
\item the cubic detector value $\xi(f)=((1,1)^{\mathsf T},I)$
and the marked $S_4$ value $1$.
\end{enumerate}
Thus these combined finite equations admit a common-word mixed
state $(\rho,a_3,\nu_4)=(0,1,0)$.
The values also pass the mixed-pair generation test described
below, and the reflection conjugacy test in D23.C15's supplement.
No actual gentle automorphism or full-GT element is asserted.
In particular a simultaneous full-source lift of all peripheral
frames is not supplied by this proposition.

\begin{proof}
Let $S=\operatorname{SU}_3(3)$, the nonabelian simple group
of order $6048$. D23.P9 gives the exact image of $F'$ in the
first Krammer pair and the three marked mixed factors:
\[
 P_{12}'\times(H'\times S)^3.
\]
Here $|P_{12}|=46656$, $|H|=108$, and $H'''=1$.
Taking two further derived subgroups proves that the image
of $F'''$ in the three perfect coordinates is all of $S^3$,
while its image in every mixed $H$ coordinate is identity.
The same assertion holds for their full Nielsen closure,
because D23.C15 identifies exactly the same three perfect
kernel classes and D23.F24 identifies the corresponding
three solvable classes. The first-pair cyclic Nielsen values
are conjugates by $K_1K_2$, up to a central argument that
cancels on a commutator word.

Every evaluation-pair group in both Hurwitz representations
has order divisible only by $2$ and $3$. For the $648$-point
action this follows from the complete $66$-component order
inventory. For the $20736$-point action, the earlier inventory
computes all $372$ components of degree at most $216$;
\path{hurwitz_large_pair_orders.py} completes the remaining
$58$ exact orders. Every prime-to-six part is one.
No timed-out computation is included in this conclusion.
The three basic pair types suffice: reversal and the braid
half twist cover all five pentagon pairs, and a Nielsen
change of two generators does not change the group they
generate. Thus the full joint Hurwitz image, including all
these Nielsen markings, also has only the primes $2,3$.

Let $L$ be the joint image of $F'''$ in all these Hurwitz
coordinates, the first Krammer pair, all mixed solvable
coordinates, and the marked $S_4$ coordinate. Its order has
only the primes $2,3$. The common-source image in $L\times S^3$
projects onto both factors. There is no nontrivial common
quotient: every nontrivial quotient of $S^3$ is a product
of one or more copies of $S$, and hence has order divisible
by $7$, whereas every quotient of $L$ has order prime to $7$.
Goursat's lemma therefore gives the full product
\[
                 \operatorname{im}(F'''\to L\times S^3)
                    =L\times S^3.
\tag{D23.P14.1}
\]

The explicit word $f_*$ of D23.P12 supplies a particular
$L$-value: first Krammer state $F_{0,1}$, identity in every
Hurwitz Nielsen coordinate, identity in every mixed solvable
coordinate, and identity marked $S_4$ value. Among the P7
solutions labelled $(0,1)$ there are exactly three with
identity mixed solvable block. Choose any one of them and
take its three perfect coordinates. Equation (D23.P14.1)
gives one ordinary word $f\in F'''$ with this prescribed
$L$-value and these three perfect coordinates.

Its complete Krammer values are therefore exactly that P7
assignment. D23.C15 and D23.F24 verify all the stated Nielsen
cycle equations on every one of the $216$ assignments,
including these three. Its Hurwitz values are unchanged from
$f_*$ and so are identity on the entire finite closure.
The cubic value follows from the exact first-pair projection
of D23.P6 and the identification of $F_{0,1}$ in D23.F13;
the marked $S_4$ value was retained explicitly in $L$.
All assertions are consequently properties of one common
ordinary word, rather than independent assignments of the
same parameter in unrelated factors.
\end{proof}

\paragraph{A necessary generation test which needs no kernel rigidity.}
If $(1,f)$ is actual gentle, its normalized automorphism of
$F$ is $x\mapsto x$, $y\mapsto f^{-1}yf$. For every finite
surjection $q:F\twoheadrightarrow G=\langle X,Y\rangle$,
the composition $q\Phi_f$ is onto. Therefore necessarily
\[
                    \langle X,f(X,Y)^{-1}Yf(X,Y)\rangle=G.
\tag{D23.P14.2}
\]
This assertion does not require $\ker q$ to be invariant,
and is weaker than a claim that $\Phi_f$ induces an automorphism
of the particular quotient.

The exact program \path{krammer_gentle_generation.py} tests
(D23.P14.2) for both mixed Krammer pairs on all $216$ P7
assignments. In every case the solvable projection still has
order $108$ and the unitary projection still has order $24192$.
There are $100$ distinct block calculations after caching
identical finite matrices. The common abelian determinant
image is unchanged by conjugation, so the Goursat description
in D23.C10 shows that the entire original mixed pair group
is generated. Reversing a pair gives a conjugate generating
pair by the two-cycle, and the remaining Nielsen markings
are the already identified marked classes. All $36$ states
for each label $(\rho,a_3)$ pass this necessary test.

For the reflection supplement, the swap-stable perfect pair
has a literal intertwiner $J_0$ implementing $x\leftrightarrow y$,
with $J_0^2$ scalar. After projectivization it extends the
quotient to $F\rtimes\langle s\rangle$, $s^2=1$. The actual
identity $\beta s\beta^{-1}=sf$ consequently requires
$J_0f$ to be conjugate to $J_0$ by the pair image. The scalar
center of $U_3(3)$ does not affect conjugation, so an
$\operatorname{SU}_3(3)$ conjugator suffices. A projective
equality is also literal: comparing determinants forces a
scalar $\zeta$ with $\zeta^3=1$ in $\mathbb F_9^\times$, hence
$\zeta=1$. The checked conjugacy class has size $63$, and
all $216$ values pass. This is a necessary actual-source
frame test, not a GT-rigidity assertion for the Krammer kernel.

The result closes the finite equations and pointwise necessary
conditions listed above. It does not settle a global peripheral
frame lifting problem or decide A.84.


\section{The cubic dessin and distribution comparisons}
\subsection{An explicit rational quotient of the cubic detector dessin}

\paragraph{D23.C14 (the marked genus-zero quotient and its cyclic pullback).}
Let $\xi:F\twoheadrightarrow\mathbb F_3^2\rtimes\operatorname{SL}_2(\mathbb F_3)$
be the marked cubic detector of D23.P10. Its degree-nine action has a
canonical extension to
\[
 \Gamma_+=\ker(C_2*C_3\longrightarrow C_2)
       =\langle u=r,v=srs\mid u^3=v^3=1\rangle.
\]
In the same affine origin as D23.P10, the images are
\[
 u=\left((2,2)^{\mathsf T},
       \begin{pmatrix}0&2\\1&2\end{pmatrix}\right),\qquad
 v=\left((1,1)^{\mathsf T},
       \begin{pmatrix}2&2\\1&0\end{pmatrix}\right).
\tag{D23.C14.1}
\]
Here $x=(vu)^{-1}$, and the branch-cycle types of $(x,v,u)$ are
\[
                   (621),\qquad(33111),\qquad(333).
\tag{D23.C14.2}
\]
The corresponding degree-nine rational Belyi map is
\[
 \Psi(z)=\frac{432z^6(z-1)(4z-1)^2}
                   {(19z^3-9z^2-3z+1)^3}.
\tag{D23.C14.3}
\]
Put
\[
 \begin{split}
 C(z)&=53z^3+15z^2-3z-1,\\
 Q(z)&=z^2-4z+1,\qquad H(z)=19z^3-9z^2-3z+1.
 \end{split}
\]
Its cyclic cubic pullback has the rational model
\[
                  w^3=C(z),
\tag{D23.C14.4}
\]
which is explicitly birational over $\mathbb Q$ to
$E:Y^2=X^3-4$. The formulas in both directions are
\[
 \begin{split}
 G&=-X^2-6X-3Y+10,\\
 z&=(2X+Y-2)/G,\qquad w=(X^2-2X+2Y+4)/G,\\
 X&=(4+16z-2w)/(1+z+w),\\
 Y&=6(w^2-19z^2-8z-1)/(1+z+w)^2.
 \end{split}
\tag{D23.C14.5}
\]
Let $n=(1+\sqrt{-3})/2$, $\bar n=1-n$, and let $\phi$ be the
normalized first-orbit degree-nine map in the retrieved primary BelyiDB data,
with $x_{\rm old}=9X/1600$ and $y_{\rm old}=27Y/64000$.
Then its compact expression and the precise composition identity are
\[
 \phi=\bar n\,\frac{H+wQ}{H-nwQ},\qquad
 \boxed{\ \Psi(z)=1-\left(\frac{\phi-\bar n}{\phi-n}\right)^3.\ }
\tag{D23.C14.6}
\]
Thus the square formed by $E\xrightarrow{\phi}\mathbb P^1_\lambda$,
$E\xrightarrow{z}\mathbb P^1_z$, $\mathbb P^1_z\xrightarrow{\Psi}\mathbb P^1_b$,
and $b=1-((\lambda-\bar n)/(\lambda-n))^3$ commutes with the stated
marking, not just with the stated passports.

The map $E\to\mathbb P^1_z$ is cyclic of degree three over
$\mathbb Q(\sqrt{-3})$ and has three ramification points, the points
$w=0$. It is therefore the quotient by a geometric order-three
symmetry with fixed points. In particular, the fact that $u$ has no
fixed sheet in (D23.C14.2) does \emph{not} make this symmetry a
fixed-point-free translation: the fixed sheets of $v$ account for
its three fixed points.
More explicitly, $w\mapsto\zeta^2w$, with $\zeta=n^2$, is
\[
       \alpha(Q)=\zeta^2Q+(1-\sqrt{-3},2\sqrt{-3})
\]
in the group law on $E$. It cycles the sixfold ramification points
$O$, $(1-\sqrt{-3},2\sqrt{-3})$, and
$(1+\sqrt{-3},-2\sqrt{-3})$.

\begin{proof}
The subgroup presentation is the index-two Reidemeister--Schreier
presentation of $C_2*C_3$. Restrict D23.P3's affine representation to
$\Gamma_+$. Its linear matrices then lie in
$\operatorname{SL}_2(\mathbb F_3)$, so taking the real part of the
translation is a homomorphism. The origin shift $(0,2)^{\mathsf T}$
from D23.P10 gives (D23.C14.1), and direct multiplication gives its
marked $x,y$. Their action on the nine affine points has the cycles
(D23.C14.2). In the labeling $(a,b)\mapsto3a+b$, the simultaneous
conjugator to the BelyiDB rational branch pair $(\sigma_0,\sigma_1)$ is
\[
                     [4,0,5,2,1,6,3,8,7].
\]
The certificate checks the two literal conjugacy equations. The
rational database passport has size one, and its map, after the
source coordinate change $x_{\rm DB}=7z$, factors as (D23.C14.3).
Equivalently, the formula itself has the required branch data since
\[
                 \Psi(z)-1=\frac{Q(z)^3C(z)}{H(z)^3},
\tag{D23.C14.7}
\]
all indicated factors are separable and coprime, and the partitions
in (D23.C14.2) have total ramification $16=2\cdot9-2$.

For the orbifold topology, let $N$ be the inverse image of a point
stabilizer in $\Gamma_+$, and put $U=N\cap F$. There are three
cusps, from the cycles of $x$, no elliptic points above $u$, and
three order-three elliptic points above $v$. Since
$\chi(\Gamma_+)=-1/3$ and $[\Gamma_+:N]=9$, the signature equation is
\[
 -3=2-2g(N)-3-3(1-1/3),
\]
so the compactification of $N$ has genus zero. A conjugate of $v$
fixing a sheet maps nontrivially to $\Gamma_+/F\simeq C_3$.
Consequently $[N:U]=3$ and $U\lhd N$. The three elliptic points are
exactly the ramification points of the compact cyclic cover.
The cusps split into nine cusps because $x\in F$. Therefore
$\chi(U)=-9=2-2g(U)-9$, and its compactification has genus one.
This also proves the fixed-point assertion without choosing a
potentially non-normalizing lift of $r$.

For the exact algebraic map, write
$L=2X+Y-2$ and $N_0=X^2-2X+2Y+4$. Modulo $Y^2-X^3+4$, one has
\[
 N_0^3=53L^3+15L^2G-3LG^2-G^3.
\]
This proves (D23.C14.4) under (D23.C14.5). Substitution verifies both
inverse formulas and the equation of $E$; hence these are
birational maps of smooth projective curves. Since $C$ is squarefree,
with discriminant $-11664$, $w^3=C(z)$ gives precisely three branch
points and no ramification at infinity.
The automorphism $w\mapsto\zeta^2w$ acts by $\zeta^2$ on the regular
differential $dz/w^2$, and its image of $O$ is
$(1-\sqrt{-3},2\sqrt{-3})$ by the inverse formulas. These two facts
give the displayed group-law expression for $\alpha$ and its orbit.

The normalized database numerator has the form $P(X)Y+R(X)$ and
denominator $D(X)$, recorded in the root's exact normalization file.
With
\[
 H_h=19L^3-9L^2G-3LG^2+G^3,\qquad
 Q_h=L^2-4LG+G^2,
\]
the following polynomial identity modulo $Y^2-X^3+4$ fixes the sign
and all scalar factors:
\[
 (PY+R)(H_h-nN_0Q_h)
       =\bar n D(H_h+N_0Q_h).
\]
It proves the first expression in (D23.C14.6). Solving for the
cross-ratio gives
\[
              \frac{\phi-\bar n}{\phi-n}
                       =n\,\frac{wQ}{H}.
\]
Because $n^3=-1$ and $w^3=C$, its cube is $-CQ^3/H^3$.
Equation (D23.C14.7) proves the boxed composition identity, with its
orientation as written. An earlier direct cross-multiplication of
the degree-27 composition checked the same identity and rejected
its opposite orientation; the displayed lower-degree identities
provide a shorter exact replay certificate.

For a normalization check, let $P_0=(2,2)$ on $E$. Then
$2P_0=(5,-11)$ and
\[
 (\phi)_0=6(O)+2(2P_0)+(-P_0),\qquad
 z(O)=0,\quad z(2P_0)=1/4,\quad z(-P_0)=1.
\]
The point with coordinates $(5,11)$ is not a zero of $\phi$.
These signs follow directly from the normalized numerator, not
from an abstract identification of the elliptic curves.

The exact scripts are
\path{agents/characters/code/check_cubic_orbifold_topology.py} and
\path{agents/characters/code/check_belyi_rational_quotient.py};
the corresponding JSON certificates are in
\path{agents/characters/records/}.
The primary database records are Michael Musty's BelyiDB,
\path{belyi_db/9/9T23-[6,3,3]-621-33111-333-g0.m} and
\path{belyi_db/9/9T23-[6,6,6]-621-621-621-g1.m},
fields \texttt{BelyiDBPermutationTriple}, \texttt{BelyiDBCurve},
and \texttt{BelyiDBBelyiMap}, first Galois orbit for the genus-one
record. The repository is
\url{https://github.com/michaelmusty/BelyiDB}.
The retrieved genus-zero record and its MIT license are retained
in the certificate directory. The normalized genus-one coefficients
are in \path{records/belyi_9T23_normalized.txt}; the marked
permutation comparison is in \path{records/belyi_9T23_match.json}.
\end{proof}

\paragraph{D23.C14a (the precise five-point comparison left by this diagram).}
In the coordinate
\[
                 t=(\lambda-\bar n)/(\lambda-n),\qquad b=1-t^3,
\]
the source of the cyclic cover is
$\mathbb P^1\setminus(\{0,\infty\}\cup\mu_3)$.
Writing $\zeta=n^2$, the points
$t=\zeta^2,\zeta,1$ are respectively $\lambda=0,1,\infty$,
and $t=0,\infty$ are $\lambda=\bar n,n$.
Thus recovering the original marked $\lambda$ sphere fills
$t=0,\infty$. The standard Enriquez distribution map instead
retains $t=0,1,\infty$ and fills the other two roots of unity.
The required relation between these two forgetting maps is an
additional marked comparison to be established; it is not being
asserted as a consequence of the existence of the diagram.

Also, $\Psi$ is indecomposable as a rational cover. Indeed its
monodromy is the affine group $\mathbb F_3^2\rtimes
\operatorname{SL}_2(\mathbb F_3)$ in its primitive degree-nine
action: a block containing zero is a translation subgroup stable
under the irreducible linear stabilizer. There is no proper such
subgroup. Consequently $\Psi$ cannot itself be a nontrivial
composition of degree-three rational Kummer covers. This rules out
that direct shortcut to applying distribution, while leaving the
five-point comparison available.

\subsection{A general limitation of fixed-target distribution tests}

\paragraph{D23.D1 (simultaneous distribution tests in a normally generated target).}
Let $F=\langle x,y\rangle$ be the free group, let $G$ be a finite
group generated by $X=q(x),Y=q(y)$, and assume that the normal
closure of $Y$ in $G$ is all of $G$. For $N\ge2$ let
\[
 K_N=\ker(F\longrightarrow\mathbb Z/N),\qquad
 x\longmapsto1,\quad y\longmapsto0.
\]
On the free basis $x^N,y_j=x^jyx^{-j}$, $0\le j<N$, define
\[
 \delta_N(x^N)=x,\quad\delta_N(y_0)=y,\quad
 \delta_N(y_j)=1\ (j>0),\qquad
 c_N(w)=\operatorname{exp}_y(\delta_N(w)).
\]
Given a finite set $\mathcal S\subset\{2,3,\ldots\}$, elements
$g_0,g_N\in G'$, and arbitrary integers $b_N$, there exists an
ordinary commutator word $f\in F'$ such that
\[
 q(f)=g_0,\qquad c_N(f)=b_N,\qquad
 q\delta_N(f)=Y^{b_N}g_N\quad(N\in\mathcal S).
\tag{D23.D1.1}
\]
For an infinite set $\mathcal S$ and arbitrary $b_N\in\widehat{\mathbb Z}$,
the same assertion holds with a profinite commutator word.

In particular, in one fixed such target the entire family of
Enriquez distribution equations can hold for a word with any
prescribed $G'$ value. This does not assert the gentle identities
for that word, or the distribution equations in other targets.

\begin{proof}
Write $L=\operatorname{lcm}(\mathcal S)$ and adjoin $L$ to
$\mathcal S$ if necessary, prescribing arbitrary extra data there.
The case of an empty set is immediate. For every positive $m$,
$q(K_m)=G$: this image is normal in $G$ and contains $Y$.
For $N\mid L$, the restriction $\delta_N(K_L)$ is exactly
$K_{L/N}$. This follows on the free basis: $x^L$ maps to
$x^{L/N}$, and $y_j$ maps to $x^{j/N}yx^{-j/N}$ if $N\mid j$
and to $1$ otherwise. Consequently
\[
 q\delta_N(K_L)=G.
\tag{D23.D1.2}
\]

First construct $b\in F'\cap K_L$ with $c_N(b)=b_N$ for all
the prescribed indices. Use the basis $x^L,y_0,\ldots,y_{L-1}$
of $K_L$. The coefficient of $y_0$ sets $c_L$ and affects every
$c_N$ by the same integer. In descending order of $N<L$ in
$\mathcal S$, the coefficient of $y_N$ sets the remaining
value of $c_N$ without altering a value already set at an index
$M>N$: indeed $c_M(y_N)=0$ for $M>N$, while $c_N(y_N)=1$.
Finally the coefficient of $y_1$ corrects the total $y$ exponent
to zero without altering any $c_N$. All these words have zero
$x$ exponent. The resulting word therefore belongs to $F'$.

Now put $V_L=K_L$. For $N<L$ put
\[
 V_N=\langle\!\langle y_N\rangle\!\rangle_{K_L}.
\]
By (D23.D1.2), the image of $V_N$ under $q\delta_N$ is the
normal closure of $XYX^{-1}$ in $G$, hence all of $G$.
For $V_L$ this follows directly from (D23.D1.2). Thus
\[
 q\delta_N(V_N')=G'.
\tag{D23.D1.3}
\]
If $M>N$ is another prescribed index, then $\delta_M(y_N)=1$
and hence $\delta_M(V_N)=1$. Every $V_N'$ lies in $F'$ and
has zero value under every homomorphism $c_M$.

Starting with $b$, proceed again in descending order of $N$.
The two elements $q\delta_N(b)$ and $Y^{b_N}g_N$ have the
same abelianization in $G$, since $\delta_N(b)$ has zero
$x$ exponent and $y$ exponent $b_N$. Their quotient therefore
belongs to $G'$. Equation (D23.D1.3) lets us correct it by
right multiplication by an element of $V_N'$. This leaves
all larger prescribed indices and all $c_M$ unchanged.
After finitely many corrections the distribution coordinates
have exactly the required values.

To set the original coordinate, put
$V_0=\langle\!\langle y_1\rangle\!\rangle_{K_L}$.
Its image under $q$ is all of $G$, because $q(K_L)=G$ and
$q(y_1)=XYX^{-1}$ normally generates $G$. Thus $q(V_0')=G'$.
But $\delta_N(V_0)=1$ for every $N\ge2$ dividing $L$.
One final factor from $V_0'$ adjusts $q(b)$ to $g_0$ without
changing any previously fixed coordinate. All constructions
are ordinary group operations with integer powers; finite
images of normal closures and their derived groups give
ordinary word representatives. This proves the finite assertion.

The same proof works in the profinite free group for finitely
many prescribed profinite integers, using their powers in the
profinite cyclic subgroups. Alternatively one can impose these
integers modulo successively larger moduli. All conditions in
(D23.D1.1) are closed conditions in the compact group
$\widehat F'$. Every finite subcollection has just been shown
to be satisfiable. Compactness therefore proves the infinite
assertion. It produces a single profinite word in this fixed
target, and does not promote that word to a gentle automorphism.
\end{proof}

\paragraph{D23.D2 (application to the compressed two-sign test).}
Let $G=E_0\times S_4$ be D23.P10's marked group of order
$5184$. Every one of its twelve necessary full-GT values
\[
 g_{\rho,e,\nu}=
 \left(\left(\binom{e}{\rho+(\rho+1)e},M_\rho\right),
               t^{2\nu}p^\rho\right),
 \quad\rho\in\mathbb F_3,\quad e,\nu\in\mathbb F_2,
\]
can be the value of one profinite word satisfying all the
Enriquez distribution equations in $G$, simultaneously for
all positive $N$, with arbitrarily prescribed $c_N$ for $N\ge2$.
In particular the fixed-target distribution equations do not
discard any of the six mixed-sign states. No actual gentle or
full-GT lift of these words is asserted.

\begin{proof}
The group is solvable. Its abelianization is $C_3\times C_2$,
and the image of $Y$ generates it: the affine linear part
$-L^2$ generates $\operatorname{SL}_2(\mathbb F_3)^{\rm ab}=C_3$,
and the four-cycle $t$ is odd. Hence the quotient of $G$ by
the normal closure of $Y$ is both solvable and perfect, so
is trivial. Every displayed state belongs to
$G'=(\mathbb F_3^2\rtimes Q_8)\times A_4$; this is also
checked in D23.P10's six-row certificate. Apply D23.D1 with
$g_0=g_N=g_{\rho,e,\nu}$ for every $N\ge2$. The case $N=1$
is the identity equation with $c_1=0$.
\end{proof}

\paragraph{Exact source and computational checks.}
The source theorem is B. Enriquez, \emph{Quasi-reflection algebras
and cyclotomic associators}, Selecta Math. 13 (2007),
Section 8, Theorem 8.1 and equation (50); arXiv:math/0408035v2,
printed pp.~30--32. It states these distribution equations
on the full profinite GT group. Remark 8.6 gives the precise
identity $c_{MN}=c_M+c_N$ under multiplication of the integer indices; that requirement
can be imposed on the arbitrary prescribed values above.
The supplied LM7\_PUSH E.6 already used this theorem in a
different finite target. The general normal-generation argument
here is a new proof of the stated fixed-target limitation.
As an independent finite check, \path{code/check_distribution_216.py}
constructs the induced wreath actions and verifies every one of
the twelve states for $N=2,3,4,6$. At each index the joint
projected commutator image has order $8957952=864^2\cdot12$;
all twelve choices of $c_N\bmod12$ occur. The proof above,
not those four finite checks, establishes the simultaneous
all-index statement.

\subsection{A cubic descent and the exact remaining cubic defect}

\begin{statement}{D23.F22}{Proved}{The index-three half-character is a descended cubic frame sign}
\end{statement}
There is a canonical continuous endomorphism
\[
 \mathcal D_E:\mathcal H\longrightarrow\mathcal H
\]
of the multiplier-one actual gentle group with the following explicit
description. In $B=\langle s,r\mid s^2=r^3=1\rangle^{\wedge}$ put
$\tau=sr$, $x=\tau^{-2}$, $y=sxs$, and
\[
 E=\langle r,srs\rangle\cong C_3\amalg C_3,
 \qquad C_f=f(x,r),\qquad t_f=\frac{3\Psi_3(f)}2.
\]
For $h=(1,f)\in\mathcal H$ first form the automorphism
\[
 \psi_f(r)=r,\qquad
 \psi_f(s)=f(r,srs)^{-1}s.
\]
Then
\[
 \boxed{\quad
 \mathcal D_E(h)=\operatorname{Inn}(x^{-t_f}C_f^{-1})\psi_f,
 \qquad A_{\mathcal D_E(h)}=x^{-t_f}C_f^{-1}.
 \quad}                                                     \tag{D23.F22.1}
\]
In particular
\[
 \boxed{\quad
 a_3(\mathcal D_E(h))=\Psi_3(h)/2\pmod2,\qquad
 \rho(\mathcal D_E(h))=\rho(h)\pmod3.
 \quad}                                                     \tag{D23.F22.2}
\]
Thus the cubic comparison defect of \textup{A.103}, on the actual
full-pentagon subgroup, has the intrinsic form
\[
 e_3=a_3+a_3\circ\mathcal D_E.
                                                               \tag{D23.F22.3}
\]
The fixed-point condition has the exact primitive formulation
\[
 \mathcal D_E(h)=h
 \quad\Longleftrightarrow\quad
 A_h^{-1}=f(x,r)x^{3\Psi_3(h)/2}.                              \tag{D23.F22.4}
\]
This is the equianharmonic primitive form with a specifically
normalized coefficient. Neither the fixed-point condition nor the
weaker equality $a_3=a_3\circ\mathcal D_E$ is asserted on all full
profinite GT here.

\begin{proof}
The kernel of the map $B\to C_2$ sending $s$ to the nonidentity
element and $r$ to one is $E=C_3\amalg C_3$, with the displayed
generators. This follows directly by the two-coset Reidemeister
rewriting, and the presentation holds profinitely as well. The
standard pure subgroup $F$ has index three in $E$; $r$ splits
$E\to C_3$, and conjugation by $r$ cyclically permutes
$x,y,(xy)^{-1}$.

The input automorphism $X\mapsto X$, $Y\mapsto f(X,Y)^{-1}Yf(X,Y)$
of the free profinite group preserves the closed normal subgroup
generated by $X^3,Y^3$, as does its inverse. Hence it induces an
automorphism of $E$, with $X=r,Y=srs$. Put $g=f(r,srs)$.
The two-cycle identity gives $sgs=g^{-1}$, so $(g^{-1}s)^2=1$.
Consequently the displayed formulas define an endomorphism $\psi_f$
of $B$. Its restriction to $E$ is exactly the induced automorphism,
and it induces the identity on $B/E$. It is therefore an
automorphism: surjectivity and injectivity follow successively from
the quotient and the restriction. The pair-composition rule
$f_{hk}=f_h\Phi_h(f_k)$ proves
$\psi_{f_{hk}}=\psi_{f_h}\psi_{f_k}$ by evaluating it in $E$.
In particular $h\mapsto\psi_f$ is a continuous homomorphism.

Both $g$ and $C_f$ belong to $F$: in $E/F=C_3$ the image of any
commutator word is one. Hence $\psi_f$ induces the identity on
$B/F=S_3$. The peripheral identity of \textup{A.100}, applied to
the pair $(r,srs)$, gives
\[
 \psi_f(y)=f(y,r)y f(y,r)^{-1},
 \qquad y=(rs)^{-2}.
\]
As $r^{-1}yr=x$ and $\psi_f(r)=r$, this gives
$\psi_f(x)=C_fxC_f^{-1}$.
It also gives $\psi_f(\tau)=C_f\tau C_f^{-1}$, with no choice
of square root. Indeed $v=C_f^{-1}\psi_f(\tau)C_f$ satisfies
$v^2=\tau^2$ and has the same image as $\tau$ in $S_3$.
The element $v\tau^{-1}\in F$ centralizes $x=\tau^{-2}$.
The pure cusp centralizer is $C_F(x)=\overline{\langle x\rangle}$;
thus $v=\tau x^b$ for some $b\in\widehat{\mathbb Z}$.
The equality of squares gives $2b=0$, whence $b=0$.

We compute the required integral normalization rather than choosing
it arbitrarily. Identify the $r$-module $F^{\rm ab}$ with
$\widehat{\mathbb Z}[\omega]$, where $[x]=1$, $[y]=\omega$,
$\omega^2+\omega+1=0$, and $r$ acts by multiplication by $\omega$.
In $F^{\rm ab}\rtimes\langle r\rangle$ the elements $r,x,srs$
are respectively
\[
 (0,\omega),\qquad(1,1),\qquad(-1,\omega^{-1}).
\]
The argument of \textup{A.90} works over the entire coefficient
ring $\widehat{\mathbb Z}[\omega]$: $1-\omega$ is a non-zero-divisor
because its norm is $3$, which is a non-zero-divisor in
$\widehat{\mathbb Z}$. The peripheral identity therefore gives
$A_f(\omega,\omega^{-1})=0$. The two-cycle and three-cycle then
give that the translation of
$f((0,\omega),(1,1))$ is a scalar $t\in\widehat{\mathbb Z}$.
Explicitly the two nontrivial translations in the three-cycle are
$-\bar t$ and $t$, so $t=\bar t$; the fixed ring is
$\widehat{\mathbb Z}$, again because multiplication by
$\omega-\omega^{-1}$ has nonzero norm $3$.

If $m_0,m_1,m_2$ are the three index-three Schreier coefficients,
the translation is $t=m_0+m_1\omega+m_2\omega^2$, with
$m_0+m_1+m_2=0$. Taking its trace gives
$2t=3m_0=3\Psi_3$. Since \textup{A.90} proves that $\Psi_3$ is
even, $t=t_f=3\Psi_3/2$.
The two-cycle now gives
\[
 C_f^{\rm ab}=-t_f[x],\qquad g^{\rm ab}=0.                    \tag{D23.F22.5}
\]
For the second assertion, evaluate the already vanishing expression
at $(0,\omega),(-1,\omega^{-1})$; changing the translation by a
scalar factor changes a commutator-word translation by the same
factor.

Let $\theta_f=\operatorname{Inn}(C_f^{-1})\psi_f$. It fixes
$\tau$, and its normalized pair coordinate is
\[
 f_0=s\theta_f(s)=(sC_f^{-1}s)gC_f.
\]
Conjugation by $s$ exchanges $[x]$ and $[y]$, so (D23.F22.5)
gives $f_0^{\rm ab}=(-t_f,t_f)$.
Conjugating once more by $x^{-t_f}$ changes its pair coordinate to
\[
 f_E=y^{-t_f}f_0x^{t_f}\in F'.
\]
The resulting automorphism fixes $\tau$, induces the identity on
$S_3$, and has commutator pair coordinate; equivalently it is an
actual multiplier-one gentle automorphism. It sends $r$ to its
conjugate by $x^{-t_f}C_f^{-1}\in F$, so uniqueness of the pure
cubic frame proves (D23.F22.1).

There is only one such normalized gentle representative in this
pure inner class. If two differ by $\operatorname{Inn}(u)$ with
$u\in F$, then fixing $\tau$ forces $u\in C_F(x)=\langle x\rangle$.
Writing $u=x^n$ changes a commutator pair coordinate by the
abelian vector $(-n,n)$, so $n=0$. The homomorphism law for
$\psi_f$ and this uniqueness prove the homomorphism law for
$\mathcal D_E$. Its formula and the continuity of $\Psi_3$ give
continuity.

The sign computation must retain the signed lift. Put
\[
 U=\begin{pmatrix}1&1\\0&1\end{pmatrix},\quad
 L=\begin{pmatrix}1&0\\-1&1\end{pmatrix},\quad
 S=ULU,\quad R=-UL
 \qquad\hbox{over }\mathbb F_3.
\]
Define $q_E:E\to\operatorname{SL}_2(\mathbb F_3)$ by
$r\mapsto R$ and $srs\mapsto SRS^{-1}$. Both images have order
three, so this is well-defined. Direct multiplication gives
\[
 q_E(x)=-U^2,\qquad q_E(y)=-L^2.
\]
Consequently $q_E(w)=(-I)^{\ell(w)}q_3(w)$ on $F$, where
$\ell(w)$ is the sum modulo two of its two abelian coordinates,
and $q_3$ is the positive signed free-subgroup lift of \textup{AD.2}.
In particular $q_E=q_3$ on the pure cubic frame in (D23.F22.1),
whose abelianization is zero by (D23.F22.5).

Write $\rho=\rho(h)\bmod3$ and
$M_\rho=\left(\begin{smallmatrix}1&\rho\\\rho&1+\rho^2\end{smallmatrix}\right)$.
The supplied signed covariance calculation gives
$f(U^2,L^2)=M_\rho$. Since $f$ has zero exponent sums, a central
sign in either input does not change its value. Using
$R=U^2L^2U^{-2}$ gives
\[
 C_E:=q_E(C_f)=f(-U^2,R)=U^2M_\rho U^{-2}.                   \tag{D23.F22.6}
\]
Put $b=\Psi_3(h)/2$, so $t_f=3b$. Because $U^3=I$,
\[
 q_3(A_{\mathcal D_E(h)})=q_E(A_{\mathcal D_E(h)})
       =(-I)^bC_E^{-1}.
\]
For $\rho_D=\rho(\mathcal D_E(h))\bmod3$, the corrected
cubic frame is therefore
\[
 T_{\mathcal D_E(h)}=(-I)^b C_E U^{2\rho_D}.
\]
It centralizes $R$ by \textup{AD.2}. Exact multiplication for
$\rho=0,1,2$ gives
$C_EU^{2\rho}=R^\rho$. Since no nonidentity power of $U$
centralizes $R$, this forces $\rho_D=\rho$, and hence
\[
 T_{\mathcal D_E(h)}=(-I)^bR^\rho.
\]
The sign character of $\langle-I\rangle\times\langle R\rangle$
is precisely $a_3$, proving (D23.F22.2). All signed identities,
all three shear residues, and both values of $b\bmod2$ are
independently checkable in \texttt{cubic\_descent.py}; its record
is \texttt{cubic\_descent.json}.

On full GT, \textup{A.76} identifies $a_3$ with
$\kappa_d\bmod2$, and \textup{A.103} defines
$e_3=\kappa_d+\Psi_3/2$. This proves (D23.F22.3).
Finally two normalized gentle automorphisms fixing $\tau$ and
having the same pure $r$-frame coincide: their images of $r$
coincide and $s=\tau r^{-1}$. Applying (D23.F22.1) proves
(D23.F22.4).
\end{proof}

\paragraph{Source boundary.}
The normalizations, actual gentle group, pure-frame uniqueness, and
the positive signed covariance are the supplied \textup{R3--R4},
\textup{LM7\_PUSH E.65}, and \textup{DELTA\_ON\_GT AD.1--AD.2}.
The character definition, evenness and peripheral identity are
\textup{A.88--A.90} and \textup{A.100}.
The new assertion here is the cubic endomorphism and its exact
parity formula. \textup{A.65}, equation \textup{(M8)}, records the
equianharmonic primitive form with an existential even coefficient;
the computation above identifies the coefficient for this descent.
The primary geometric account of the two fillings is
Lochak--Nakamura--Schneps, \emph{Eigenloci of 5 point configurations},
\S6, equations (6.3)--(6.4), author PDF pp.30--31, and its
Theorem 2, equation (HF1), author PDF p.4. That theorem has Galois
scope. Enriquez, \emph{Quasi-reflection algebras and cyclotomic
associators}, arXiv:math/0408035v2, Theorem 8.1 and equation (50),
pp.30--32, supplies the full-GT cyclic distribution relation but
is not silently substituted for the missing fixed-point or parity
comparison in (D23.F22.4).

\subsection{The harmonic sign is unchanged by cubic descent}

\paragraph{D23.P13 (an actual-source relation between the two descents).}
For every actual multiplier-one gentle element $h=(1,f)$,
the cubic descent of D23.F22 satisfies
\[
 \boxed{\quad a_2(\mathcal D_Eh)=a_2(h),\qquad
 \rho(\mathcal D_Eh)=\rho(h)\pmod3.\quad}
\tag{D23.P13.1}
\]
More precisely its normalized pair coordinate $f_E$ has exactly
the same value as $f$ in the signed affine harmonic detector of
supplied FP.34. Consequently, using D23.F22 and supplied A.107,
\[
 \boxed{\quad
 \Delta\circ\mathcal D_E
 =a_2+\Psi_3/2
 =d_2
 =a_2+a_2\circ D_L
 =a_2+a_2\circ D_P
 \quad\text{on }\mathcal H.\quad}
\tag{D23.P13.2}
\]
The half-character in this display is reduced modulo two.
This identifies an actual descended defect; it does not assert
that $d_2$ vanishes, or that either descent fixes all full-GT
elements.

\begin{proof}
We keep the exact FP.34 marking. Put
\[
 R=(\mathbb Z/4)[\omega]/(\omega^2+\omega+1),\qquad
 c=(1-\omega)^{-1}=2+3\omega,
\]
and use the affine multiplication $(a,u)(b,v)=(a+ub,uv)$.
The homomorphism $B=C_2\amalg C_3\to R\rtimes\langle-\omega\rangle$
is
\[
 s\longmapsto(c,-1),\quad r\longmapsto(c,\omega),\quad
 x\longmapsto(0,\omega),\quad y\longmapsto(1,\omega).
\tag{D23.P13.3}
\]
Write $f((0,\omega),(1,\omega))=(k,1)$.

For any commutator word $f$ and any $a,b\in R$, the exact
affine identity
\[
 f((a,\omega),(b,\omega))=((b-a)k,1)
\tag{D23.P13.4}
\]
holds. Indeed simultaneous conjugation by the translation
$ac$ changes $(0,\omega),(b-a,\omega)$ to the indicated
pair. A commutator word has linear part one, and conjugation
by a translation leaves a pure translation unchanged.
Finally multiplying the translation module by $b-a$ is an
affine-group homomorphism (no invertibility of $b-a$ is needed),
which gives (D23.P13.4). The argument applies to profinite
words by continuity, since the target is finite.

Use the notation and normalized formula from D23.F22:
\[
 C=f(x,r),\qquad g=f(r,srs),\qquad t=3\Psi_3(f)/2,
 \qquad f_E=y^{-t}(sC^{-1}s)gCx^t.
\tag{D23.P13.5}
\]
Here $C,g\in F$ and $t\in3\widehat{\mathbb Z}$; the latter
uses the already proved evenness of $\Psi_3$.
Since $srs\mapsto(-\omega c,\omega)$, formula (D23.P13.4) gives
\[
 C\longmapsto(ck,1),\qquad
 g\longmapsto(\omega^2ck,1).
\tag{D23.P13.6}
\]
Conjugation by $s$ negates translations, hence
$sC^{-1}s\mapsto(ck,1)$. Therefore
\[
 (sC^{-1}s)gC\longmapsto((2+\omega^2)ck,1)=(k,1),
\]
because $2+\omega^2=1-\omega$ and $(1-\omega)c=1$ in $R$.
Both displayed images of $x,y$ have order three. The normalizing
powers $x^t,y^{-t}$ thus have identity image, and
\[
                    q_{\rm harm}(f_E)=q_{\rm harm}(f).
\tag{D23.P13.7}
\]
This is an exact equality of finite values, not just a comparison
after forgetting a sign or reducing modulo two.

FP.34 gives, for each actual gentle normalized pair,
\[
 1-(1-\omega)k=\omega^\rho(-1)^{a_2}.
\]
The six elements $\omega^j(-1)^e$, $j\in\mathbb F_3$,
$e\in\mathbb F_2$, are distinct in $R$. Applying the formula
to $f_E$ and $f$ proves both claims in (D23.P13.1), independently
of the second congruence proved by the cubic matrices in D23.F22.

D23.F22 gives $a_3\circ\mathcal D_E=\Psi_3/2$. Adding
(D23.P13.1) proves the first two equalities of (D23.P13.2).
The remaining equalities are exactly supplied A.107, equations
(MD.8a)--(MD.8b). This deduction uses genuine endomorphisms
and actual characters on $\mathcal H$ throughout; it makes
no assertion about full-GT descent compatibility.

The exact certificate \path{agents/pentagon/cubic_descent_harmonic.py}
checks the ring, all signed peripheral images, the affine group
order, the translation computation for every $k\in R$, and
the uniqueness of the six harmonic phases. It does not replace
the cited actual-source descent construction in D23.F22 or the
supplied harmonic-frame theorem FP.34.
\end{proof}

\paragraph{Exact source locators and remaining question.}
The signed harmonic representation and recovery formula are
\path{sources/DELTA_ON_GT/agent_five/five_point_results.tex},
FP.34.1--FP.34.3. The pair formula is
\path{work/agents/five_level_six/cubic_descent.tex},
D23.F22.1 and its displayed normalization of $f_E$.
The half-character comparison with $D_L,D_P$ is
\path{sources/DELTA_ON_GT/sections/A.tex}, A.107,
(MD.8a)--(MD.8b).
The identity $\Psi_3/2\circ\mathcal D_E=\Psi_3/2$, or a different
formula for this iterated half-character, is not proved here.

\subsection{The cubic descent for all multipliers}

\begin{statement}{D23.F23}{Proved}{The whole gentle-group descent and its relative quadratic character}
\end{statement}
The endomorphism $\mathcal D_E$ of \textup{D23.F22} extends canonically
to a continuous endomorphism of the full actual gentle group
$\mathcal G=\widehat{GT}_0$, preserving its multiplier.
For $g=(\lambda,f)$ put
\[
 m=(\lambda-1)/2,\quad
 \varepsilon\in\{1,-1\},\quad\varepsilon\equiv\lambda\pmod3,
 \quad b=(1-\varepsilon)/2,
 \qquad C_g=r^{-m}f(x,r).
\]
Define
\[
 \psi_g(r)=r^\lambda,\qquad
 \psi_g(s)=f(r,srs)^{-1}s,\qquad
 \theta_g=\operatorname{Inn}(\tau^b C_g^{-1})\psi_g.
\]
Then $f_{0,g}=s\theta_g(s)$ belongs to $F$ and has abelianization
$(u_g,-u_g)$. The full descent is the explicit normalized automorphism
\[
 \boxed{\quad
 \mathcal D_E(g)=\operatorname{Inn}(x^{u_g})\theta_g,
 \quad f_{\mathcal D_E(g)}=y^{u_g}f_{0,g}x^{-u_g}\in F'.
 \quad}                                                     \tag{D23.F23.1}
\]
This construction is independent of any Galois lift of the multiplier.

Using \textup{D23.P13}, put $\gamma=\Delta\circ\mathcal D_E$.
Its restriction to $\mathcal H=\ker\lambda$ is the duplication
defect $d_2=a_2+\Psi_3/2$. There is a unique continuous character
$\chi:\widehat{\mathbb Z}^{\times}\to\mathbb F_2$ such that
$\gamma|_{G_{\mathbb Q}}=\chi\circ\lambda$.
Consequently
\[
 \widehat d_2=\gamma+\chi\circ\lambda,\qquad
 \widehat e_3=\Delta+\widehat d_2                             \tag{D23.F23.2}
\]
are canonical continuous relative quadratic characters of
$\mathcal G$, restricting respectively to
\[
 a_2+\Psi_3/2,\qquad a_3+\Psi_3/2
\]
on $\mathcal H$. In particular
$\Delta=\widehat d_2+\widehat e_3$ on the full gentle group.
This does not assert nontriviality or independence of these two
characters, and does not determine the whole relative character space.
It also does not assert that the uncorrected $\gamma$ vanishes on
all of $G_{\mathbb Q}$.
By uniqueness of relative extension, $\widehat d_2$ is the
character $d$ already constructed in D23.C1. The new assertion
is its realization through the whole-group cubic descent.

\begin{proof}
As in \textup{D23.F22}, quotient the input free profinite group by
the closed normal subgroup generated by $X^3,Y^3$. The quotient
is $E=C_3\amalg C_3$, with $X=r,Y=srs$. Both the actual input
automorphism and its inverse preserve this normal subgroup, since
$\lambda$ is a unit. They therefore induce an automorphism of $E$.
Write $G=f(r,srs)$. The two-cycle gives $sGs=G^{-1}$, so
$(G^{-1}s)^2=1$. The formulas for $\psi_g$ define an endomorphism
of $B$, whose restriction to $E$ is the induced automorphism and
whose action on $B/E=C_2$ is the identity. It is an automorphism.
The normalized pair-composition formula proves directly that
$g\mapsto\psi_g$ is a homomorphism, including arbitrary multipliers.

For clarity we derive the general peripheral conjugator. Write
$Z=(XY)^{-1}$, $G_0=f(Z,X)$, $H_0=f(Y,Z)$ and
$m=(\lambda-1)/2$. The three-cycle for $(\lambda,f)$ is
\[
 G_0 Z^m H_0 Y^m f X^m=1.
\]
Right composition with the actual element $c=(-1,1)$ gives the
pair $(-\lambda,f)$ and its three-cycle
\[
 G_0 Z^{-m-1}H_0Y^{-m-1}fX^{-m-1}=1.
\]
Eliminating $H_0$ from these equations gives
\[
 G_0Z^{-\lambda}G_0^{-1}
   =X^{m+1}f^{-1}Y^\lambda fX^m
   =X^{-m}\Phi_g(Z^{-1})X^m.
\]
Thus
\[
 \Phi_g(Z)=X^m f(Z,X)Z^\lambda f(Z,X)^{-1}X^{-m}.             \tag{D23.F23.3}
\]
Evaluate this at $(r,srs)$, where $Z=y=(rs)^{-2}$.
Since $x=r^{-1}yr$ and $\psi_g(r)=r^\lambda$, we obtain
\[
 \psi_g(x)=r^{m-\lambda}f(y,r)y^\lambda f(y,r)^{-1}r^{\lambda-m}
          =C_g x^\lambda C_g^{-1}.
\]
In the last equality use $m-\lambda+1=-m$ and
$f(y,r)r=rf(x,r)$.

The element $f(x,r)$ is pure, and $\psi_g$ induces on $S_3$
the automorphism $s\mapsto s$, $r\mapsto r^\varepsilon$.
Consequently $v=C_g^{-1}\psi_g(\tau)C_g$ has the same
permutation as $\tau^\lambda=\tau\bmod F$ and satisfies
$v^2=\tau^{2\lambda}$.
Both $v$ and $\tau^\lambda$ commute with $x^\lambda$.
Since $\lambda$ is a unit, they commute with $x$ as well.
Therefore $v\tau^{-\lambda}$ lies in
$C_F(x)=\overline{\langle x\rangle}$.
The equality of squares and injectivity of multiplication by two
in $\widehat{\mathbb Z}$ show that it is one. We have proved
\[
 \psi_g(\tau)=C_g\tau^\lambda C_g^{-1}.
                                                               \tag{D23.F23.4}
\]

The correction by $\tau^b$ is essential for the permutation
normalization. The two possible cases are
\[
\begin{array}{c|c|c|c|c}
 \lambda\bmod6&\varepsilon&m\bmod3&b&
       \tau^b C_g^{-1}\bmod F\\\hline
 1&1&0&0&1\\
 5&-1&2&1&(sr)r^2=s.
\end{array}
\]
In the first case $\psi_g$ already acts identically on $S_3$.
In the second it fixes $s$ and inverts $r$, and conjugation by
$s$ changes this into the identity on $S_3$. Thus in both cases
$\theta_g$ induces the identity on $S_3$; by (D23.F23.4) it
sends $\tau$ to $\tau^\lambda$.

It follows that $f_{0,g}=s\theta_g(s)$ belongs to $F$.
The order-two relation $\theta_g(s)^2=1$ gives the exact identity
$s f_{0,g}s=f_{0,g}^{-1}$. On $F^{\rm ab}$, conjugation by $s$
exchanges the two coordinates. Hence if
$f_{0,g}^{\rm ab}=(u,v)$, then $(v,u)=(-u,-v)$, or $v=-u$.
This proves the asserted abelian form. Conjugation by $x^{u_g}$
changes the pair coordinate to $y^{u_g}f_{0,g}x^{-u_g}$, whose
abelianization is zero, and preserves the cusp formula.

The normalized automorphism restricts to the actual free-group
automorphism
$x\mapsto x^\lambda$, $y\mapsto f_E^{-1}y^\lambda f_E$.
Its two-cycle follows from the order-two relation just used.
Its three-cycle follows from the order-three relation in $B$:
its image of $r$ is $f_E^{-1}r x^{-m}$, and expansion of its
cube, using the cyclic action of $r$ on $x,y,(xy)^{-1}$, gives
\[
 f_E(z,x)z^m f_E(y,z)y^m f_E(x,y)x^m=1.
\]
Thus (D23.F23.1) is an actual element of $\mathcal G$ of
multiplier $\lambda$.

For the composition law, each (D23.F23.1) is an inner correction
of $\psi_g$. Since $\psi$ is a homomorphism, its normalized
corrections have the same outer composition law. Two such
corrections with the same multiplier both induce the identity
on $S_3$, so an inner element relating them has trivial image
in $S_3$, whose center is trivial. It is therefore in $F$.
Since both automorphisms send $\tau$ to $\tau^\lambda$, that
element centralizes $x$ and is $x^n$. Conjugation by $x^n$
changes the pair coordinate by the abelian vector $(-n,n)$.
Both coordinates belong to $F'$, so $n=0$. This proves uniqueness
of the normalization and hence the homomorphism law.
All operations, the locally constant choice of $b$, and extraction
of $u_g$ from $F^{\rm ab}$ are continuous, proving continuity.
At $\lambda=1$ we have $b=m=0$ and, by \textup{D23.F22},
$u_g=-3\Psi_3(g)/2$. Thus the extension restricts to that of
\textup{D23.F22}.

Now $\gamma=\Delta\mathcal D_E$ is a continuous character on
$\mathcal G$. Statements \textup{D23.F22} and \textup{D23.P13}
give $\gamma|_{\mathcal H}=d_2$.
The arithmetic comparison needed here is precisely
\textup{DELTA\_ON\_GT A.103}, final paragraph after
\textup{(MD.4e)}: on $G_{\mathbb Q}\cap\mathcal H$,
Anderson duplication in the specified intrinsic convention gives
$d_2=0$. Equivalently this is the physical-descent arithmetic
compatibility of \textup{A.62}, combined with the universal
half-character formula \textup{A.107 (MD.8a)}.
We use this supplied arithmetic comparison, and do not replace
it by an unmarked assertion about a Kummer exponent.

Therefore $\gamma|_{G_{\mathbb Q}}$ vanishes on the kernel of
the cyclotomic character. The cyclotomic map
$G_{\mathbb Q}\twoheadrightarrow\widehat{\mathbb Z}^{\times}$
is a continuous surjection of compact profinite groups, hence
its target is the topological quotient by that kernel.
There is consequently a unique continuous character $\chi$
with the stated factorization. The formulas (D23.F23.2) define
continuous characters and show directly that they vanish on
$G_{\mathbb Q}$. Their restrictions to $\mathcal H$ are as
claimed. Their uniqueness as relative extensions of those
restrictions follows from
$\mathcal G=\mathcal H G_{\mathbb Q}$, the same decomposition
used in \textup{LM7\_PUSH E.65}. This also proves the last
displayed identity on all of $\mathcal G$.
\end{proof}

\paragraph{Geometric normalization.}
For the cubic eigenlocus put
$n=(1+\sqrt{-3})/2$, $\bar n=(1-\sqrt{-3})/2$ and
\[
 t=\frac{\lambda-\bar n}{\lambda-n},\qquad b=1-t^3.
\]
Direct substitution in the Legendre expression gives
\[
 \frac{256(1-\lambda+\lambda^2)^3}
      {1728\lambda^2(1-\lambda)^2}
       =\frac{4(b-1)}{b^2}.                                \tag{D23.F23.5}
\]
The pole values $n,\bar n$ correspond to $b=\infty,1$;
the three cusps $\lambda=0,1,\infty$ correspond to $b=0$.
The degree-two map on the right is the harmonic map
$4u(1-u)$ after $u=1/b$. This records the precise relation
between the $C_3\amalg C_3$ and $C_2\amalg C_3$ orbifolds.
It does not by itself identify the two GT actions: the cyclic
distribution filling retains $0,1,\infty$ in the $t$-sphere,
whereas the cubic modular filling removes the two poles and
retains the three roots of unity. The functorial comparison for
these different fillings remains the required additional step.

\subsection{A restriction on identity relative seeds for the harmonic defect}

\begin{statement}{D23.P15}{Proved}{A normal relative seed cannot separate the two order-48 evaluations by nilpotent factorization}
\end{statement}
Let $B=C_2\amalg C_3$, $F=\widehat F(x,y)$, and
$x=(sr)^{-2}$, $y=sxs$ as in the supplied conventions. Put
$R_4=(\mathbb Z/4)[\omega]$, $\omega^2+\omega+1=0$, and use
the following two marked surjections onto groups of order $48$:
\[
 q_K(x)=(0,\omega),\quad q_K(y)=(1,1),\qquad
 q_H(x)=(0,\omega),\quad q_H(y)=(1,\omega)
 \quad\text{in }R_4\rtimes\langle\omega\rangle.
\]
If $U\triangleleft B$ is an open normal subgroup contained in $F$
and $q_K(U)$ is nilpotent, then
\[
 U\subseteq F'F^3,\qquad q_H(U)\subseteq(R_4,+).
 \tag{D23.P15.1}
\]
Consequently both maps kill $R_{\rm nil}(U)$. In particular, an
actual gentle pair $(1,f)$ with
$f\in R_{\rm nil}(U)$ necessarily satisfies
\[
 \Psi_3(f)/2\equiv0\pmod2,
 \qquad \rho(f)\equiv0\pmod3,
 \qquad a_2(f)=0.
 \tag{D23.P15.2}
\]
Thus an identity relative construction which forces the Kummer
target to factor through the relative pronilpotent quotient cannot
produce $a_2=1$ and $\Psi_3=0$. This statement does not exclude a
nonidentity relative automorphism, a different construction, or an
actual gentle element with those character values.

\begin{proof}
Write $V=(R_4,+)$ and $G=V\rtimes C_3$. Since $U$ is normal in
$F$ and $q_K$ is onto, $N=q_K(U)$ is normal in $G$. If its image
in $C_3$ were nontrivial, $N$ would contain an element with linear
part $\omega$ or $\omega^2$. Its commutators with $V$ generate
$(1-\omega)V=V$, because $1-\omega$ is a unit in $R_4$.
Normality would give $V\subseteq N$, and then $N=G$. But $G$
is not nilpotent: its order-three action on $V$ is nontrivial,
whereas distinct Sylow subgroups commute in a finite nilpotent
group. It follows that $N\subseteq V$.

The linear part of $q_K(u)$ is $\omega^{\operatorname{exp}_x(u)}$.
Therefore every $u\in U$ has $x$ exponent divisible by $3$.
Conjugation by $s$ exchanges $x$ and $y$. Since $sUs=U$, the
$y$ exponent of every $u\in U$ is divisible by $3$ as well.
This proves $U\subseteq F'F^3$. The linear part of $q_H(u)$ is
$\omega^{\operatorname{exp}_x(u)+\operatorname{exp}_y(u)}$, so
it is one. This proves (D23.P15.1), and $q_H(U)$ is abelian.

The universal property of the nilpotent residual now gives
$q_K(R_{\rm nil}(U))=q_H(R_{\rm nil}(U))=1$.
The supplied A.89--A.90 identify the half-character modulo two
with the scalar translation of $q_K(f)$ reduced modulo two;
it is therefore zero. The supplied FP.34 harmonic formula is
\[
 1-(1-\omega)\operatorname{trans}(q_H(f))
       =\omega^{\rho(f)}(-1)^{a_2(f)}.
\]
Its left side is one. The six phases
$\omega^j(-1)^e$, $j\in\mathbb Z/3$, $e\in\mathbb Z/2$, are
distinct in $R_4^\times$, so $\rho\equiv0\pmod3$ and $a_2=0$.
This proves (D23.P15.2).
\end{proof}

\paragraph{Source and check scope.}
The two marked quotients and their character formulas are the
supplied \texttt{DELTA\_ON\_GT/sections/A.tex}, statements
A.89--A.90, and \texttt{agent\_five/five\_point\_results.tex},
FP.34. The relative lifting theorem itself is supplied
\texttt{LM7\_PUSH/E\_relative\_cm\_lift\_agent.tex}, E.58;
its missing D.50 proof is not audited here. The argument above
is an independent proof of this restriction on its identity
relative application, with no assertion about other seeds.

\subsection{An exact finite detector for the next cubic-descent character}

\begin{statement}{D23.P16}{Proved}{The iterated half-character and a genuinely invariant finite quotient}
\end{statement}
Write $b=\Psi_3/2\pmod2$ on the actual multiplier-one gentle
group $\mathcal H$. There is an explicit permutation quotient
$q:B\twoheadrightarrow Q\subseteq S_{24}$, of order $3456$,
with the following properties. Put
\[
 E=\langle r,srs\rangle,\qquad H=F'F^3.
\]
Then
\[
 |q(E)|=1728,\quad |q(F)|=576,\quad
 q(H)=q(F)'\simeq C_2^6.
 \tag{D23.P16.1}
\]
For $h=(1,f)\in\mathcal H$, set $C=f(x,r)$. Its image belongs
to $q(H)$. The element
\[
 q(rC^{-1}r^{-1}C)
 \tag{D23.P16.2}
\]
acts on a specified four-point fiber by translation by the scalar
\[
 \boxed{\ b(\mathcal D_E h)+b(h)\ \in\mathbb F_2.\ }
 \tag{D23.P16.3}
\]
The kernel of $q$ is preserved by every actual gentle
automorphism, so an induced finite-automorphism condition in this
quotient is a necessary condition without any unproved
representation-rigidity hypothesis.

Exact enumeration leaves eight $q(C)$ values after the finite
conditions listed below, four for each scalar in (D23.P16.3).
All eight extend to automorphisms of $Q$ and are compatible,
as ordinary common-word values, with both original values of
$b$. Thus these tested conditions do not prove $b\mathcal D_E=b$.
No actual gentle or full-GT lift of the eight finite candidates
is claimed.

\begin{proof}
Start with the marked off-diagonal quotient of A.90:
\[
 q_0:F\longrightarrow \mathbb F_4\rtimes C_3,
 \qquad x\longmapsto(0,\omega),\quad y\longmapsto(1,1).
\]
It acts faithfully on the four elements of $\mathbb F_4$.
Induce this action from $F$ to $B$, using the six representatives
\[
 1,r,r^2,s,sr,sr^2
\]
and the action on the left cosets of $F$. The resulting
24-point action is $q$. Its kernel is precisely
$\operatorname{Core}_B(\ker q_0)$, since the original four-point
action is faithful. The zero fiber is the coset represented by
one; restriction of $q(w)$ to this fiber, for $w\in F$, is
the affine action of $q_0(w)$.

For reproducibility, the six $s$-edges have trivial Schreier
words. The six $r$-edges, in the displayed order, have words
\[
 1,1,1,y^{-1},x^{-1},xy.
\]
The target coset of each edge is determined by the ordinary
$S_3$ action of $s$ and $r$. These data specify $q$ completely.
The exact certificate reconstructs the edges in $C_2*C_3$, checks
$s^2=r^3=1$ and the zero-fiber $x,y$ marking, and enumerates
the finite groups. It gives the derived-series orders
\[
\begin{array}{c|rrrr}
 q(B)&3456&576&64&1\\
 q(E)&1728&192&64&1\\
 q(F)&576&64&1&
\end{array}
\]
and $q(F)^{\rm ab}\simeq C_3^2$, while $q(F)'$ has exponent
two. Therefore $q(H)=q(F)'$ and (D23.P16.1) follows.
These orders and identifications are exact finite computations,
not independent assertions about profinite quotients of the same
orders.

For the character formula, put $g=f(r,srs)$ and
$t=3\Psi_3(f)/2$. The supplied gentle three-cycle, applied to
$(r,srs)$, gives
\[
 g=(sCs)rC^{-1}r^{-1}.
 \tag{D23.P16.4}
\]
Indeed $(r\cdot srs)^{-1}=y$, $rxr^{-1}=y$,
$rf(x,r)r^{-1}=f(y,r)$, and
$sf(x,r)s=f(y,srs)$. The standard three-cycle is
$f(y,r)f(srs,y)f(r,srs)=1$, and the two-cycle gives
(D23.P16.4), with the displayed orientation.

F22 supplies $C^{\rm ab}=-t[x]$, so $C\in H$ because $3\mid t$.
It also supplies the normalized descendant pair
\[
 f_E=y^{-t}(sC^{-1}s)gC x^t
     =y^{-t}(rC^{-1}r^{-1}C)x^t.
 \tag{D23.P16.5}
\]
Since $H$ is normal in $B$, the middle factor maps into $q(H)$.
Its action on the zero fiber is a translation by some
$k\in\mathbb F_4$. On that fiber $q_0(x^t)=1$, because
$3\mid t$, and $q_0(y^{-t})$ is translation by $b(h)$.
Thus $q_0(f_E)$ is translation by $k+b(h)$. The actual
gentleness of $\mathcal D_E h$, together with A.90, says that
this is the scalar $b(\mathcal D_E h)$. It follows that $k$ is
a scalar and proves (D23.P16.3).

To justify kernel invariance, let $\Phi_h$ be the normalized
automorphism of $F$. In $q_0(F)$ every commutator value is a
translation, and translations centralize $q_0(y)$. Hence
\[
 q_0\Phi_h(x)=q_0(x),\qquad
 q_0\Phi_h(y)=q_0(f^{-1}yf)=q_0(y).
\]
Thus $\Phi_h$ preserves $\ker q_0$, with equality because
its inverse has the same property. As it is also an actual
automorphism of $B$, it preserves the core of this subgroup in
$B$. This proves the asserted invariance of $\ker q$.
The same applies to $\mathcal D_E h$. Its raw automorphism
\[
 \psi_f(r)=r,\qquad\psi_f(s)=g^{-1}s
\]
differs from $\mathcal D_E h$ by an inner automorphism of $B$;
it therefore preserves $\ker q$ as well. Testing whether these
two images induce an automorphism of $Q$ is consequently valid.

Here is the exact enumeration and its scope. Enumerate the 64
elements $c\in q(H)$ and set
$g=(scs)rc^{-1}r^{-1}$ in $Q$. Impose that $g^{-1}s$ lie in
the conjugacy class of $s$ under $q(F)$, which has 24 elements.
Exactly eight $c$ survive. Every surviving middle factor
$rc^{-1}r^{-1}c$ has scalar zero-fiber translation, with four
values of each scalar. Each surviving pair satisfies the
two-cycle and both the $m=0$ and $m=-1$ three-cycles at
$(r,srs)$, generates the prescribed pair groups, and has
\[
 \left|\left\langle(s,g^{-1}s),(r,r)\right\rangle
               \subseteq Q\times Q\right|=3456.
\]
Both projections of this graph group are onto $Q$, so it is
the graph of an automorphism.

The original pair conditions in $q(F)'$ are also enumerated:
$(sf)^2=(f^{-1}r)^3=1$, together with scalar zero-fiber
translation, leave exactly two values, one for each $b$.
Both give graph groups of order $3456$. The common-word map
of three marked pairs
\[
 F\longrightarrow q(E)\times q(E)\times q(F),
 \quad X\longmapsto(r,x,x),\quad
 Y\longmapsto(srs,r,y)
\]
has image order $7077888$ and derived-image order $786432$.
Membership tests show that all sixteen tuples consisting of
one of the eight $(g,c)$ and either original $f$ value belong
to this derived image. Thus a single ordinary commutator word
realizes each of these finite tuples. This last sentence is
strictly a finite common-word statement.
\end{proof}

\paragraph{Certificates and further test.}
The source action is reconstructed by
\path{cubic_descent_iterate.py}. The complete finite
enumeration in the proof is independently rerunnable as
\path{cubic_descent_iterate_check.py}, with literal
permutations and results in its JSON file. An additional
Enriquez $N=3$ distribution calculation, simultaneously on
the two substituted pairs and the original pair, leaves every
one of the sixteen tuples even after imposing
$c_3\equiv2b\pmod4$. Its projected commutator-image order is
$3710851743744$; all three compatible residues of $c_3\pmod{12}$
pass for each tuple. This is a further exact finite test, not
a distribution theorem or an actual lift. The corresponding
record is
\path{cubic_descent_iterate_distribution_full_N3.json}.

\paragraph{Source boundary.}
The actual-source inputs are supplied A.89--A.90 and A.100 in
\path{DELTA_ON_GT/sections/A.tex}, and the newly proved
D23.F22 in
\path{work/agents/five_level_six/cubic_descent.tex}.
The scalar correction, the stable finite quotient, and the
enumerated limitation above are new. No equality
$b\mathcal D_E=b$ and no nonzero value on an actual gentle
element has been established here.


\section{Orbital pentagons and the two residual parities}
% DELTA23: author proof and root's independent exact marked-path audit complete.
\subsection{An orbital image of the pentagon with its boundary corrections}

\paragraph{D23.C16 (the five explicitly marked orbital cofaces).}
Use Enriquez's notation
\[
 W_{n,3}=\{(z_1,\ldots,z_n)\in(\mathbb C^\times)^n:
                     z_i^3\ne z_j^3\ (i\ne j)\},
 \qquad \zeta=e^{2\pi i/3}.
\]
There is an algebraic map
\[
 \mathcal O_n:W_{n,3}\longrightarrow M_{0,n+2},\qquad
 (z_1,\ldots,z_n)\longmapsto
       (\zeta z_1,\zeta^2z_1,z_1,z_2,\ldots,z_n).
\tag{D23.C16.1}
\]
All points on the right are distinct. In particular $\mathcal O_3$
induces a homomorphism from the cyclic covering subgroup
$K_{3,3}=\pi_1(W_{3,3})$ to the pure five-point mapping class group.

Fix the following one marking for all five faces. The source base
configurations and connecting paths are in the positive chamber
$0<z_1<z_2<z_3$. Send the fifth target point to infinity by
$M(p)=1/(p-z_3)$, use a projection of angle
$0<\theta<\pi/3$, and straighten the remaining four strands in
reverse label order. Finally conjugate every target word by the
same half twist $a=\sigma_1$. Write
\[
 \begin{gathered}
 a=\sigma_1,\quad b_0=\sigma_2,\quad c=\sigma_3,
 \quad\delta=ab_0,\\
 x=a^2,\quad y=b_0^2,\quad K=c^2,\quad
 Z=(ab_0)^3,\quad V=(b_0c)^3.
 \end{gathered}
\tag{D23.C16.2}
\]
The group is $B_4/\langle z_4\rangle$, where
$z_4=(ab_0c)^4$. Thus $Z$ is the boundary twist $T_{45}$;
$K=T_{34}$ and $V=T_{15}$. In particular
\[
 \delta^3=Z,\qquad [K,x]=[K,V]=1,
       \qquad [Z,x]=[Z,y]=1.
\tag{D23.C16.3}
\]

For the rank-four cyclic source use the actual basis from
Enriquez's Section~8:
\[
 U=x_{01}^3,\qquad Y_j=x_{01}^j x_{12}x_{01}^{-j}\quad(0\le j<3).
\]
The three middle coface maps, followed by $\mathcal O_3$, have
these literal images in the common marking:
\[
\begin{array}{c|cccc}
 &U&Y_0&Y_1&Y_2\\\hline
\widetilde\partial_1&1&Z&1&1\\
\widetilde\partial_2&1&xK^{-1}&V&K(xV)^{-1}\\
\widetilde\partial_3&Z&xZ^{-1}&yZ^{-1}&(xy)^{-1}.
\end{array}
\tag{D23.C16.4}
\]
The exceptional faces are
\[
 \mathcal O_3\partial_0(f)=f(K,Z),\qquad
 \mathcal O_3\partial_4(f)=f(\delta,K).
\tag{D23.C16.5}
\]
The second formula is an extension through the partially labelled
family: rotating the first coordinate through $120$ degrees
permutes the first three target points and gives $\delta$.
Its cube is $Z$, and is not set to one in the five-point group.

Consequently, for every actual full-GT element $h=(1,f)$, put
\[
 b=\Psi_3(h)/2,\qquad C_f=f(x,r)\in F,
 \qquad A^E_f=x^{-3b}C_f^{-1}\in F'.
\]
Here $A^E_f$ is precisely the pure cubic frame of the descent
$\mathcal D_E$ in D23.F22; evaluations of $C_f^{-1}$ or $A^E_f$
at another pair mean evaluation of their marked pure free-group
word. The orbital pentagon identity is
\[
 \boxed{\quad
 f(K,Z)\,A^E_f(x,V)\,K^{-3b}\,f(\delta,K)
                  =A^E_f(x,y)Z^b.
 \quad}
\tag{D23.C16.6}
\]
This is a genuine five-point identity obtained from the original
pentagon, with the two boundary exponents retained.
There is also a specified literal lift to $B_4$, obtained by using
the actual map $(p_1,\ldots,p_5)\mapsto
(1/(p_1-p_5),\ldots,1/(p_4-p_5))$ into the configuration space
of four finite points. In this lift the identity is
\[
 f(K,Z)A^E_f(x,V)K^{-3b}f(\delta,K)
                   =A^E_f(x,y)Z^b z_4^{-b}.
\tag{D23.C16.6L}
\]
In particular a representation with $q(z_4)=-I$ must retain
the scalar $(-I)^b$ on the right. Projective vanishing cannot
be substituted for this literal equation.

\begin{proof}
We first make the base paths explicit. The positive-chamber
connecting paths have no crossings in the chosen target projection.
Indeed put $q=z_1/z_3\in(0,1)$. Then
\[
 \operatorname{Re}(e^{i\theta}M(\zeta^2z_1))
 -\operatorname{Re}(e^{i\theta}M(z_1))
 =\frac{q[3\cos\theta(1+q)-\sqrt3\sin\theta(1-q)]}
        {2z_3(1-q)(1+q+q^2)}>0.
\]
Also $M(z_2)<M(z_1)$ on the real axis, and the imaginary parts
of $M(\zeta z_1),M(\zeta^2z_1)$ give their strict order after the
positive projection. The order is therefore uniformly
$4<3<2<1$ throughout the whole chamber. Thus none of the five
coface identifications acquires a separate hidden conjugator.
The one common conjugation by $a$ changes the observed cluster
rotation $b_0a$ to $ab_0$ and preserves $K$, because $a$ commutes
with $c$.

Use positive lollipop meridians for the base powers
$Z_i=z_i^3$. The generator $x_{01}$ rotates $z_1$ through
$120$ degrees and fixes $z_2$; $x_{12}$ is the lift of a positive
collision loop of $Z_2$ about $Z_1$, along the positive real
connecting arc. Hence $Y_j$ is obtained by rotating the first
coordinate $j$ times, following that collision loop, and rotating
back. This fixes the directions of the sector paths as well as
their endpoints.

The three lifted cofaces of Enriquez's Lemmas~8.2--8.4 are, on
appropriate homotopy-equivalent open sets,
\[
 (z_1,z_2)\longmapsto
 (az_1,z_1,z_2),\quad(z_1,cz_1,z_2),\quad(z_1,z_2,cz_2),
\]
where in this display only $0<a\ll1$ and $c>1$ close to one are
real constants, not braid generators. The collision lollipops
are chosen larger than the indicated splitting distance.

For the first face, the $U$ loop rotates the first four target
points together and is the full four-strand twist, hence trivial
in the sphere quotient. The $Y_0$ loop is the collision of target
points $4,5$, giving $Z$. The other two root-collision divisors
are unmarked regular points in this family, so their meridians
are trivial. This proves the first row. In the ratio
$u=z_2/z_1$, $U=x_{01}^3$ is inertia at $u=\infty$, not at
$u=0$. A meridian at the latter point instead maps to $Z^{-1}$;
confusing those two meridians leaves the value on $f\in F'$
unchanged but gives an incorrect generator table.

For the second face, $U$ again rotates the first four points
and dies. The $Y_0$ loop moves point $5$ around the close pair
$3,4$ without an internal twist. Its value is the twist on
$3,4,5$ divided by $K$, namely $xK^{-1}$. The first rotated
lollipop is the meridian of the remaining orbit point, along
its lower-sector arc. In the common marking its value is
$T_{15}=V$. Both missing peripheral loops $0,\infty$ become
regular, so the three remaining meridians multiply to one;
this gives $Y_2=(xK^{-1}V)^{-1}=K(xV)^{-1}$.

For the third face, $U$ rotates precisely the first three target
points and gives $Z$. The positive-root collision moves the close
pair $4,5$ around point $3$, with no internal twist. Its value is
$xZ^{-1}$. Rotation of the cluster supplies
$Y_j=\delta^j(xZ^{-1})\delta^{-j}$; the braid relations give
$Y_1=yZ^{-1}$ and $Y_2=(xy)^{-1}$. These observations prove
(D23.C16.4). The positive collision arcs in $\partial_0$ give
$K,Z$; the partial cluster rotation and the same collision arc
in $\partial_4$ give $\delta,K$. This proves (D23.C16.5).

There is an independent discrete certificate of the stated braid
identifications. The file
\path{agents/characters/records/orbit_coface_artin_simplifications.json}
records the sixteen crossing words, their one common conjugator,
the simplified words in (D23.C16.4)--(D23.C16.5), and the discarded
power of $z_4$. The faithful Artin action verifies each equality
before quotienting by that specified central power. The numerical
path tracker was used to discover those words; the boundary and
sector-path argument above is the geometric justification.

The path extraction has also been independently certified without
floating-point sampling. In
\path{code/certify_orbital_paths.py}, take source base coordinates
$(z_1,z_2)=(1,4)$, with the extra coordinate $15/2$ when needed.
Use the regular twelve-gon on the unit circle for rotations of
$z_1$: four edges give $x_{01}$, and twelve give $U$.
A collision meridian is a positive diamond of radius $1/20$
times its positive real centre, reached along the positive real
segment and followed by the return segment. For $Y_j$, precede
and follow this lollipop by the prescribed $j$ rotations and
their inverses. These are exactly the positive meridians and
sector arcs specified above. Choose the splitting constants
$a=1/5$ and $c=101/100$, and use the projection factor $1+i/10$.
The diamond about $1$ contains $c$ and $c^{-1}$ in its interior,
but neither $0$ nor the other cube roots of unity; its minimum
distance from $1$ is $1/(20\sqrt2)$. The rotation polygons stay
in the annulus $\cos(\pi/12)\leq|z_1|\leq1$, while their other
source coordinate is $4$. These observations verify that the
source paths and their lifted cofaces stay in the required open
configuration spaces. They are homotopic, with the same sector
markings, to the circular lollipops used to define the cofaces.

All coordinates on a polygon edge are affine functions of its
real parameter with coefficients in $\mathbb Q(\sqrt3,i)$.
After the stated M\"obius normalization, the difference of two
projected real coordinates has a numerator of degree at most
three over $\mathbb Q(\sqrt3)$. Its field norm has degree at
most six over $\mathbb Q$. The certificate isolates every real
root in $(0,1)$ by rational Sturm intervals, separates the two
quadratic-field conjugates by exact signs, and checks the signs
of the derivative and imaginary difference at each retained
root. It also checks the endpoint orders and that no two
crossing intervals overlap. Thus every crossing, its order,
and its sign are proved by rational arithmetic. The resulting
sixteen literal $B_4$ words agree with the recorded words before
any quotient by the centre. The complete root intervals and
polynomials are in \path{records/orbital_paths_exact.json}.

More precisely, in the literal $B_4$ lift the first row of
(D23.C16.4) is $(z_4,Zz_4^{-1},1,1)$ and the second is
\[
       (z_4,xK^{-1}z_4^{-1},Vz_4^{-1},K(xV)^{-1}).
\]
The third row is unchanged. The two generators of the exceptional
zeroth face have images $K,Zz_4^{-1}$, while those of the fourth
face have images $\delta,K$. These are precisely the central
powers retained by the exact Artin certificate.

For the effect on $f$, let its index-three Schreier coefficients
be $(m_0,m_1,m_2)$. D23.F22's integral calculation gives
\[
                    (m_0,m_1,m_2)=(2b,-b,-b),
\]
and the coefficient of $U$ is zero. Filling both fixed punctures
is the map
\[
 \varepsilon_3:(U,Y_0,Y_1,Y_2)\longmapsto
                (1,x,y,(xy)^{-1}),\qquad
 \varepsilon_3(f)=f(r,x)=C_f^{-1}.
\]
Since $K$ commutes with $x,V$, row two of (D23.C16.4) differs
from $C_f^{-1}(x,V)$ by the central factor
$K^{-m_0+m_2}=K^{-3b}$. Since $Z$ commutes with $x,y$,
row three differs from $C_f^{-1}(x,y)$ by
$Z^{-m_0-m_1}=Z^{-b}$. The first row has value $Z^{m_0}=Z^{2b}$.
In particular the half-character correction has been computed
integrally, with no choice of a cube root.

Enriquez's lifted pentagon (51), which is the original full-GT
pentagon inside the covering group, reads
\[
 \partial_0(f)\widetilde\partial_2(f)\partial_4(f)
             =\widetilde\partial_3(f)\widetilde\partial_1(f).
\]
Apply $\mathcal O_3$ and substitute the five formulas. This gives
\[
 f(K,Z)C_f^{-1}(x,V)K^{-3b}f(\delta,K)
                         =C_f^{-1}(x,y)Z^b.
\]
Finally $C_f^{-1}=x^{3b}A^E_f$, and $x$ commutes with both $K$
and $Z$. Cancel the common factor $x^{3b}$ to obtain
(D23.C16.6).
In the literal lift, the first face contributes an additional
$z_4^{-m_0}=z_4^{-2b}$, while the second contributes
$z_4^{-m_0-m_1}=z_4^{-b}$. The central input in the zeroth face
does not change its commutator-word evaluation. Consequently
the same calculation gives (D23.C16.6L). Its total braid exponent
is $-6b$ on each side: the right side has $6b-12b$.
\end{proof}

\paragraph{D23.C16a (the specific cyclotomic action being compared).}
In the orbifold group
\[
 \Gamma=\langle u,v\mid(uv)^2=1\rangle^{\wedge}
           \simeq\widehat{\mathbb Z}\amalg C_2,
\]
Enriquez's Section~13 action, for multiplier one and the canonical
module parameter $g=f$, is
\[
 u\longmapsto u,\qquad
 v\longmapsto f(u,v^2)^{-1}v f(u,v^2).
\]
Filling the two fixed punctures in the order-three cover imposes
$u^3=1$ and gives $B$ with $r=u,s=uv$, so $v^2=x$.
The induced action has $v\mapsto C_fvC_f^{-1}$.
The three-cycle at $(r,srs)$ shows that its action on $s$ is
$f(r,srs)^{-1}s$, exactly the raw action $\psi_f$ of D23.F22.
Thus the descent being compared is a specified cyclotomic orbifold
filling, and not a separately chosen action on an isomorphic group.
The compatibility still sought is with the standard full-GT action
on the five-point periodic frame.

\paragraph{D23.C17 (a precise periodic-factorization residual).}
Let $h=(1,f)$ be an actual full-GT element, and retain the notation
of D23.C16. Let $A_h$ be its original pure cubic frame and put
\[
 D_h=A_h(x,y)^{-1}A^E_f(x,y),\qquad
                g=f(\delta,K).
\]
Then, in the five-point group,
\[
             \boxed{\ (\delta D_hg^{-1}KgD_h^{-1})^3=1.\ }
\tag{D23.C17.1}
\]
Its original signed modular residue is
\[
                  q_3(D_h)=(-I)^{a_3(h)+\Psi_3(h)/2}.
\tag{D23.C17.2}
\]
Therefore excluding the residue $-I$ for this specific correlated
periodic factorization would prove the cubic comparison defect
$e_3=0$ on full GT. This requires less than proving $D_h=1$;
no such exclusion is asserted here.

\begin{proof}
Set $R=f(K,Z)$ and $M=A^E_f(x,V)K^{-3b}$, so $M$ commutes with
$K$. The orbital identity gives
\[
                 (A^E_f)^{-1}R=Z^b g^{-1}M^{-1}.
\]
The standard braid action has
$\Phi(\delta)=A_h\delta A_h^{-1}$ and
$\Phi(K)=RKR^{-1}$. The first identity lifts literally from the
three-strand central quotient: any central discrepancy has braid
exponent zero, while the central generator has exponent six.
Also $(\delta K)^3=z_4=1$ in the sphere quotient. Substitution,
and the fact that $Z$ centralizes $\delta,D_h$, yield
\[
 \Phi(\delta K)
   =A_h Z^b(\delta D_h g^{-1}K gD_h^{-1})Z^{-b}A_h^{-1}.
\]
Cubing proves (D23.C17.1).

The pure frames $A_h,A^E_f$ belong to $F'$. Thus their signed
modulo-three images lie in
$\operatorname{SL}_2(\mathbb F_3)'=Q_8$.
Write the corrected frame as $T_h=(-I)^{a_3(h)}R_0^j$.
Its image in the abelian quotient $C_3$ is that of
$U^{2\rho(h)}$, and $R_0=-UL$ has twice the class of $U$ there;
therefore $j=\rho(h)\bmod3$. D23.F22 gives the same shear for
$\mathcal D_E(h)$ and corrected frame
$(-I)^bR_0^{\rho(h)}$. The shear and order-three factors cancel
in the ratio $q_3(A_h)^{-1}q_3(A^E_f)$, proving (D23.C17.2).
Uniqueness of a periodic frame alone does not make this ratio
trivial: the original and descended actions must still be compared.
\end{proof}

% DELTA23: quadratic orbital identity; root independently certified all13 paths.
\subsection{The harmonic orbital identity and its shared periodic braid}

\paragraph{D23.C18 (the quadratic orbital cofaces and residual).}
Retain the literal four-strand marking
\[
 a=\sigma_1,\quad b_0=\sigma_2,\quad c=\sigma_3,\qquad
 x=a^2,\ y=b_0^2,\ K=c^2,\ Z=(ab_0)^3,\ V=(b_0c)^3,
\]
and put $s=ab_0a$, $r=ab_0$, $z_4=(ab_0c)^4$.
There is a quadratic orbital map
\[
 \mathcal O^H_n:W_{n,2}\longrightarrow M_{0,n+2},\qquad
 (z_1,\ldots,z_n)\longmapsto(-z_1,0,z_1,z_2,\ldots,z_n).
\tag{D23.C18.1}
\]
Use the positive source chamber, send the last target point to
infinity by $p\mapsto1/(p-z_3)$, and straighten the remaining
strands in reverse label order. Unlike D23.C16, there is no
final common conjugation. All initial target points are real,
and their transformed order is uniformly $4<3<2<1$ throughout
the positive chamber. Thus all face connecting paths have the
same marking without additional conjugators.

For the cyclic source basis
\[
 U=x_{01}^2,\qquad Y_0=x_{12},\qquad
 Y_1=x_{01}x_{12}x_{01}^{-1},
\]
the three middle faces have the following literal $B_4$ images:
\[
\begin{array}{c|ccc}
 &U&Y_0&Y_1\\\hline
 \widetilde\partial_1&z_4&Zz_4^{-1}&1\\
 \widetilde\partial_2&z_4&xK^{-1}z_4^{-1}&Vz_4^{-1}\\
 \widetilde\partial_3&Z&xZ^{-1}&yZ^{-1}.
\end{array}
\tag{D23.C18.2}
\]
The exceptional zeroth face sends its two generators to
$K,Zz_4^{-1}$. The fourth face sends them to $s,K$.
Here the partial half rotation of $z_1$ exchanges the two
outer points of the cluster $(-z_1,0,z_1)$ and gives $s$;
it satisfies $s^2=Z$.

Let $h=(1,f)$ be an actual multiplier-one full-GT element and put
\[
 p=\Psi_2(h),\qquad b=\Psi_3(h)/2,\qquad
 J_f=f(s,x)\in F,\qquad B^H_f=x^{-2p}J_f\in F.
\tag{D23.C18.3}
\]
Then $(J_f)^{\rm ab}=(p,-p)$ and
$(B^H_f)^{\rm ab}=(-p,-p)$. The harmonic orbital identity is
\[
 \boxed{\quad
 f(K,Z)B^H_f(x,V)K^{-p}f(s,K)
                  =B^H_f(x,y)Z^p z_4^{-p}.
 \quad}
\tag{D23.C18.4}
\]
It is a literal identity in $\widehat B_4$ and therefore also
holds after imposing $z_4=1$. No assertion that $B^H_f$ is a
frame of an actual gentle descent is needed or made.

Let $\beta_h$ be the original pure harmonic frame and define
\[
 D^H_h=\beta_h(x,y)^{-1}B^H_f(x,y),\qquad g_H=f(s,K).
\]
In the sphere quotient one has
\[
 \boxed{\quad
 (sD^H_hg_H^{-1}c g_HD^{H,-1}_h)^3=1.
 \quad}
\tag{D23.C18.5}
\]
For the original signed modular quotient $q_3$, its residual is
\[
 q_3(D^H_h)\in\langle S\rangle\simeq C_4,
 \qquad
 \epsilon_2(q_3(D^H_h))=a_2(h)+b=d_2(h).
\tag{D23.C18.6}
\]
Thus C17 and C18 give correlated five-point periodic
factorizations with residual parities $e_3$ and $d_2$.
They do not by themselves prove equality of these parities.

\begin{proof}
Use the same coface paths as in D23.C16, now with degree two.
The three lifted cofaces are
$(az_1,z_1,z_2)$, $(z_1,cz_1,z_2)$ and
$(z_1,z_2,cz_2)$, with real splitting constants $a$ small and
$c>1$ close to one. In this sentence the two constants are not
braid generators. The positive-root lollipop has the same
boundary twist description as in C16. The only other root
is $-z_1$; the fixed point $0$ remains marked. Consequently
the first two rows in (D23.C18.2) are the stated boundary
twists, and the third follows from $sxs^{-1}=y$ and $s^2=Z$.
The zeroth face follows the positive collision arcs, and the
fourth is the half rotation followed by the same collision.

There is an independent exact geometric certificate of all
thirteen generator paths. The program
\path{code/certify_harmonic_orbital_paths.py} uses the cubic
certificate's rational Sturm engine, with six edges of the
regular twelve-gon for a half rotation and all twelve for $U$.
It uses the target coordinates $(-z_1,0,z_1,z_2,z_3)$ and the
projection factor $1+i/10$. The positive collision diamonds
and positive-chamber connecting paths have the same elementary
separation properties as in C16. Exact rational isolation of
every crossing polynomial, its derivative sign, and its
imaginary-coordinate sign proves all thirteen literal crossing
words, with no sampled inference. The full certificate is
\path{records/harmonic_orbital_paths_exact.json}.
The separate faithful Artin replay
\path{agents/characters/code/check_harmonic_orbit_coface_artin.py}
checks every simplification in (D23.C18.2), with its exact
central exponent; it also checks all braid identities used below.

In the index-two Schreier abelianization the coefficient of $U$
is zero and those of $Y_0,Y_1$ are $(p,-p)$.
The filling $U\mapsto1,Y_0\mapsto x,Y_1\mapsto y$ sends $f$
to $J_f=f(s,x)$, since $s^2=1$ and $sxs=y$ in the central
three-strand quotient. This gives the asserted abelianization.
The equality lifts to the marked free pure subgroup of $B_3$:
both sides have total braid exponent zero, so their possible
central discrepancy is zero. Equivalently the $U$ coefficient
is zero and $s^2=Z$ is central in the literal lift.

Rows one, two and three therefore give, respectively,
\[
 Z^p z_4^{-p},\qquad J_f(x,V)K^{-p},\qquad J_f(x,y).
\]
The two central exponents in row two cancel because
$p+(-p)=0$, as do the $Z$ exponents in row three.
The commutator-word evaluation in the zeroth face is
$f(K,Zz_4^{-1})=f(K,Z)$.
Apply Enriquez's lifted pentagon (51) in the form
$\partial_0(f)\widetilde\partial_2(f)\partial_4(f)
=\widetilde\partial_3(f)\widetilde\partial_1(f)$.
Then substitute $J_f=x^{2p}B^H_f$ and cancel the common
factor $x^{2p}$, which commutes with $K,Z$ and therefore
with $f(K,Z)$. This proves (D23.C18.4).
Its two sides have the same literal braid exponent $-10p$.

Put $R=f(K,Z)$ and $M_H=B^H_f(x,V)K^{-p}$.
Both $x$ and $V$ commute with $c$, so $M_H$ commutes with $c$.
Equation (D23.C18.4) gives, in the sphere quotient,
\[
       (B^H_f)^{-1}R=Z^p g_H^{-1}M_H^{-1}.
\]
The standard full-GT braid action has
$\Phi(s)=\beta_hs\beta_h^{-1}$ and
$\Phi(c)=RcR^{-1}$. The first equality lifts literally
from the three-strand central quotient by braid exponents;
the second is the standard Drinfeld braid formula, using
the two-cycle to reverse $f(Z,K)$.
Since $(sc)^3=z_4$, and $Z$ centralizes $s,D^H_h$, it follows
that
\[
 \Phi(sc)=\beta_h Z^p
  (sD^H_hg_H^{-1}c g_HD^{H,-1}_h)Z^{-p}\beta_h^{-1}.
\]
Cubing proves (D23.C18.5).

It remains to identify the finite residual. Write
$q_3(x)=U^2$ and $q_3(s)=S$, allowing the central lift sign
of $s$, which does not change a commutator-word evaluation.
In the projective affine realization over $\mathbb F_4$ used
in A.87 and A.106,
\[
 x=(0,\omega),\qquad s=(c_0,1),\qquad
 c_0=(1-\omega)^{-1}.
\]
A.90 and the two-cycle give
$J_f=f(s,x)=(c_0b,1)$ in this affine quotient (the inverse
has the same translation in characteristic two).
The inverse image of the subgroup generated by this projective
$s$ in $\operatorname{SL}_2(\mathbb F_3)$ is
$\langle S\rangle\simeq C_4$. Hence
\[
 W_H:=q_3(B^H_f)^{-1}U^{2p}
     =q_3(J_f)^{-1}\in\langle S\rangle,
       \qquad \epsilon_2(W_H)=b.
\]
The equality uses $U^3=1$:
$q_3(B^H_f)^{-1}U^{2p}=q_3(J_f)^{-1}U^{6p}$.
For the original frame
$W_h=q_3(\beta_h)^{-1}U^{2\rho(h)}\in\langle S\rangle$.
A.102 supplies the full-GT congruence $\rho(h)=p\pmod3$.
Consequently
\[
 q_3(D^H_h)
   =W_hU^{-2\rho(h)-4p}q_3(J_f)
   =W_h W_H^{-1}.
\]
Taking the quotient $C_4\to C_2$ proves (D23.C18.6).
This finite computation does not presume an automorphism of
$\widehat B$ extending the filled index-three subgroup action.
\end{proof}

\paragraph{D23.C19 (the shared periodic conjugacy).}
Let $P=(ab_0c)^2$. There is the literal braid identity
\[
                P(sc)P^{-1}=rK.
\tag{D23.C19.1}
\]
Write
\[
 H_h=sD^H_hg_H^{-1}c g_HD^{H,-1}_h,\qquad
 E_h=rD_hg^{-1}KgD_h^{-1},
\]
with the cubic notation $D_h,g,A_h,b$ of C17.
Then both residual periodic elements satisfy the specific
common conjugacy
\[
 E_h=L_hH_hL_h^{-1},\qquad
 L_h=Z^{-b}A_h^{-1}\Phi(P)\beta_hZ^p.
\tag{D23.C19.2}
\]
Indeed C17 and C18 express $\Phi(rK)$ and $\Phi(sc)$ as
the stated conjugates of $E_h,H_h$; applying $\Phi$ to
(D23.C19.1) gives (D23.C19.2). The braid identity follows
also from conjugation by $ab_0c$ shifting $a$ to $b_0$ and
$b_0$ to $c$: $sc=(ab_0c)a$ and $rK=(ab_0c)c$.
Thus the two residual torsion conditions are coupled by the
same original action. Reading their two frame parities from
this common conjugacy, or from an equivalent finite quotient,
remains the unresolved comparison; arbitrary unrelated
conjugators would discard precisely this extra data.

\paragraph{D23.C20 (why the affine shared conjugacy alone is insufficient).}
The shared conjugacy C19 imposes no parity constraint after applying
$\eta_3$ and retaining only the supplied full-GT covariance
$\eta_3\Phi=\operatorname{Ad}(0,U^{2\rho})\eta_3$.
This is stronger than the separate observation that both cube
tests admit their two possible defect signs.

Indeed suppress $\eta_3$ from the notation and put
\[
 V_0=(0,U^{2\rho}),\quad
 T=A_h^{-1}V_0,\quad W=\beta_h^{-1}V_0,\quad z=Z,
 \quad d_E=rK,\ d_H=sc,\ p_0=P.
\]
Here $T,W,z$ are linear affine matrices, and $z=-I$ commutes
with both $T,W$. C17--C19 and the covariance give exactly
\[
 E_h=\operatorname{Ad}(z^{-b}T)(d_E),\qquad
 H_h=\operatorname{Ad}(z^{-p}W)(d_H),\qquad
 L_h=(z^{-b}T)p_0(z^{-p}W)^{-1}.
\tag{D23.C20.1}
\]
For example $A_h^{-1}V_0=T$ converts the conjugator of
$\Phi(rK)$ directly into the first expression; the harmonic
expression is identical with $\beta_h,W,p$ in place of
$A_h,T,b$. These formulas use no equality between the two
frame parities. Since $p_0d_Hp_0^{-1}=d_E$, substitution proves
$E_h=L_hH_hL_h^{-1}$ identically for every
$T\in C_6$, $W\in C_4$ and every $b,p$.
Thus the affine translation class in this particular conjugacy
cannot recover the missing parity by itself. Additional
same-word constraints on the full orbital factors have not
been discarded by this statement and remain a separate test.

\subsection{The exceptional orbital pair generates a six-dimensional unitary group}

\paragraph{D23.F25.}
Let $K_1,K_2,K_3$ be the literal Krammer matrices over
$\mathbb F_9=\mathbb F_3[i]$, $i^2=-1$, used in D23.P6--P9.
Set $X=K_1K_2$, $Y=K_3^2$, and $Z=(XY)^{-1}$.
There is a nondegenerate Hermitian form $H$ such that
\[
 \langle K_1,K_2,K_3\rangle=U(H),\qquad
 \langle X,Y\rangle=\det^{-1}(\{1,-1\})\subset U(H),
 \qquad \langle X,Y\rangle'=SU(H).
\tag{D23.F25.1}
\]
The respective group orders are
\[
 \begin{aligned}
 |U(H)|&=182699779456696320,\\
 |SU(H)|&=45674944864174080,\\
 |\langle X,Y\rangle|&=91349889728348160.
 \end{aligned}
\tag{D23.F25.2}
\]
Moreover $SU(H)$ is perfect, its centre is $\{I,-I\}$, and
$PSU(H)$ is simple.  The projective orders of $X,Y,Z$ are
respectively $36,6,3$, whereas their literal orders are $36,6,6$.
In particular the equal literal orders of $Y,Z$ do not give equal
conjugation-action orders on the derived group.

\begin{proof}
All matrix entries below use the encoding $a+3b\leftrightarrow a+bi$.
The exact invariant form is
\[
 H=\begin{pmatrix}
 0&1&1&6&6&0\\
 1&0&1&1&5&6\\
 1&1&0&0&1&1\\
 3&1&0&0&1&6\\
 3&8&1&1&0&1\\
 0&3&1&3&1&0
 \end{pmatrix}.
\tag{D23.F25.3}
\]
The certificate checks $H=\overline H^{\,T}$, $\operatorname{rank}H=6$,
and $\overline K_j^{\,T}HK_j=H$ for each $j$.  It also computes
the full invariant Hermitian-form space for $X,Y$, of dimension
one over $\mathbb F_3$.  The determinants are
\[
 \det K_j=i,\qquad \det X=\det Y=-1,\qquad\det Z=1.
\tag{D23.F25.4}
\]

There is a particularly small source of all the required root
elements.  The matrix $T=Y^2=K_3^4$ is
\[
 T=\begin{pmatrix}
 1&0&0&0&0&0\\0&1&0&0&0&0\\0&0&1&0&0&0\\
 0&0&0&1&0&0\\0&0&0&0&1&0\\0&2&3&2&3&1
 \end{pmatrix}.
\tag{D23.F25.5}
\]
Thus $\operatorname{rank}(T-I)=1$, $(T-I)^2=0$, and $T^3=I$.
Its image line is the sixth coordinate line.  An exact orbit
enumeration under $X,Y$ gives $22204$ distinct projective lines,
all isotropic for $H$.  This is the entire isotropic-line set,
not a sample.  To verify the count without a group-order table,
diagonalize a nondegenerate Hermitian form to the identity using
Gram--Schmidt and surjectivity of the norm $\mathbb F_9^\times
\to\mathbb F_3^\times$.  For a coordinate, the numbers of vectors
of norm $0,1,2$ are $1,4,4$.  If $a_n$ is the number of norm-zero
vectors in dimension $n$, then
\[
       a_1=1,\qquad a_n=4\cdot9^{n-1}-3a_{n-1}.
\]
This gives $a_6=177633$ and $(a_6-1)/8=22204$ isotropic lines.
The executable lists the complete orbit, with the first nonzero
coordinate normalized to one, and checks each norm directly.

We supply the generation argument needed to turn this orbit into
a group identification.  Write $h(v,w)=\overline v^{\,T}Hw$.
A rank-one unitary element $I+v\varphi$ with square-zero
$v\varphi$ fixes the hyperplane $\ker\varphi$ pointwise.
Unitarity implies that hyperplane is $v^\perp$; square-zero then
implies $h(v,v)=0$.  Consequently it has the form
\[
       T_{v,c}(w)=w+c\,h(v,w)v,\qquad \overline c=-c.
\tag{D23.F25.6}
\]
For one isotropic line over $\mathbb F_9$, there are exactly two
nonidentity such transformations, inverse to one another.
Conjugation carries their image line by the conjugating isometry.
Thus the orbit calculation proves that $\langle X,Y\rangle$
contains every unitary transvection of the six-dimensional space.

Here is an elementary proof that these transvections generate the
special unitary group in this situation.  On a hyperbolic plane
with form matrix $\left(\begin{smallmatrix}0&1\\1&0\end{smallmatrix}\right)$,
a determinant-one unitary matrix has diagonal entries in
$\mathbb F_3$ and off-diagonal entries in $i\mathbb F_3$.
Conjugating by $\operatorname{diag}(i,1)$ identifies this group
with $SL_2(\mathbb F_3)$.  Its upper and lower elementary matrices
are precisely the two transvection root groups, and elementary
row reduction generates $SL_2(\mathbb F_3)$ from them.
Every nondegenerate Hermitian plane over $\mathbb F_9$ is
hyperbolic: after diagonalization, norm surjectivity gives a
nonzero isotropic vector, and a hyperbolic partner is obtained
by subtracting half its norm times that vector.

The subgroup generated by all transvections is transitive on
vectors of norm one in every dimension $n\ge2$.  If two such
vectors span a nondegenerate plane, its special unitary group
carries one to the other: extend each to an orthonormal basis
of the plane and adjust the second vector by a norm-one scalar
to make determinant one.  This plane group is generated by
ambient transvections, as just proved.  Proportional vectors
also lie in such a plane.  In the remaining case write
$w=\alpha v+u$, where $h(v,v)=h(w,w)=1$, $u\ne0$ is isotropic
in $v^\perp$, and $N(\alpha)=1$.  Choose an isotropic partner
$u'$ in $v^\perp$ with $h(u,u')=1$.  Take $\gamma$ with
$N(\gamma)=2$, and choose $\beta$ so that
$\operatorname{Tr}(\overline\beta\gamma)=1$.
Then $z=\beta u+\gamma u'$ has norm one, is orthogonal to $v$,
and satisfies $N(h(w,z))=2$.  Both planes $\langle v,z\rangle$
and $\langle w,z\rangle$ are nondegenerate, giving the required
two-step move.  The trace choice exists because the nonzero
linear functional $\beta\mapsto\operatorname{Tr}
(\overline\beta\gamma)$ is onto $\mathbb F_3$.

Induction now proves generation of $SU_n$.  Given $g\in SU_n$,
first multiply by a product of transvections to fix a chosen
norm-one vector.  Its restriction to the orthogonal complement
belongs to $SU_{n-1}$ and is generated by its transvections by
induction.  The base $SU_1$ is trivial, and the plane case was
proved above.  This establishes the needed $SU_6$ generation
without invoking a classification of finite subgroups.

We have proved $SU(H)\subseteq\langle X,Y\rangle$.
Equation (D23.F25.4) proves the reverse containment in the
determinant preimage and supplies its nontrivial determinant
coset.  It also supplies all four determinant cosets in the
full braid image.  This proves the first two assertions of
(D23.F25.1).  Norm-one vector counts and stabilizers give
\[
 |U_n(3)|=3^{n(n-1)/2}\prod_{j=1}^{n}(3^j-(-1)^j),
 \qquad |SU_n(3)|=|U_n(3)|/4.
\]
For clarity, the stabilizer of a norm-one vector is $U_{n-1}$,
the group is transitive on those vectors by the argument above,
and their number is $3^{n-1}(3^n-(-1)^n)$, obtained from the
same norm recurrence.  This proves the order formula and
(D23.F25.2).

We also prove the structural facts used below.  In dimension
six, choose two orthogonal hyperbolic pairs, the first containing
the image vector of a given transvection.  An isometry which
scales that vector by $\alpha$ of norm two and its partner by
$\overline\alpha^{-1}$, and scales the second pair by
$\alpha^{-1},\overline\alpha$, has determinant one.
Its conjugation multiplies the transvection parameter by two.
Thus $[g,T]=T$ for the associated order-three transvection.
Every transvection is a commutator in $SU(H)$, so the generation
argument proves that $SU(H)$ is perfect.  Determinants now imply
that the derived groups of both groups in (D23.F25.1) are exactly
$SU(H)$.

To identify the centre and prove projective simplicity, use a
hyperbolic basis $e_1,e_2,e_3,f_1,f_2,f_3$.
An isometry fixing every isotropic line is diagonal on these
six lines.  The lines $e_i+e_j$, $f_i+f_j$ and $e_i+f_j$ for
$i\ne j$ then force all its diagonal entries to be equal.
It is therefore scalar.  A scalar in $SU(H)$ has fourth and
sixth power one, hence is $I$ or $-I$.  A central element
centralizes every transvection and therefore fixes each of
its image lines, proving that this is exactly the centre.

The action of $PSU(H)$ on isotropic lines is faithful and
primitive.  Here are explicit reasons for primitivity.  The
stabilizer of a line $v$ has exactly three orbits: $v$ itself,
the other lines orthogonal to $v$, and the lines not orthogonal
to $v$.  For nonorthogonal pairs, scale representatives to a
hyperbolic pair and complete to a hyperbolic basis.  For two
orthogonal independent lines, complete their span to two
hyperbolic pairs and then to a basis.  These constructions
give the asserted orbit transitivity; any determinant can be
corrected on the unused nondegenerate complement, without
moving the specified lines.  A block containing $v$ must be
a union of these stabilizer orbits.  The putative block of
all lines orthogonal to $v$ is excluded by moving $v$ to a
different orthogonal line $w$: the two perpendicular-line
sets intersect and differ.  The putative block consisting of
$v$ and its nonorthogonal lines is excluded by moving $v=e_1$
to $w=f_1$: these two sets intersect, but the isotropic line
$f_1+e_2$ belongs to the former and not the latter.
Thus only singleton and whole-set blocks remain.

For completeness, the elementary normal-subgroup argument is
as follows.  In a faithful primitive action, every nontrivial
normal subgroup is transitive, since its orbits are blocks.
Let $N$ be a nontrivial normal subgroup of $PSU(H)$ and let
$P$ be a point stabilizer.  Then $PSU(H)=NP$.  The order-three
transvection group at that point is abelian and normal in $P$;
its conjugates generate $PSU(H)$.  Modulo $N$, all these
conjugates have the same image, because every element is a
product from $N$ and $P$.  The quotient is consequently
abelian.  Perfectness forces it to be trivial, so $N=PSU(H)$.
This proves simplicity with no unproved finite-simple-group
identification.

Finally exact multiplication gives the local orders stated in
the assertion: no power $X^j$ with $0<j<36$ is scalar; no power
$Y^j$ with $0<j<6$ is scalar; and $Z^3=-I$ with $Z,Z^2$
nonscalar.  The centre of $U(H)$ is scalar by the same line
argument.  Moreover an element acting trivially by conjugation
on $PSU(H)$ acts trivially on $SU(H)$: its commutators would
otherwise define a homomorphism from the perfect group
$SU(H)$ to its centre.  It must therefore be scalar.
Thus these projective orders are also the conjugation-action
orders on both $SU(H)$ and $PSU(H)$.

The exact executable is
\path{agents/five_level_six/krammer_orbital_unitary.py}; its
JSON record is \path{krammer_orbital_unitary.json}, and its
complete line orbit is \path{krammer_orbital_isotropic_orbit.npz}
in the same directory.  The program checks only the stated
finite matrices and exhaustive line orbit.  Generation,
perfectness, simplicity and the subsequent Goursat deductions
are the mathematical proofs above, not assertions delegated
to a computer algebra group's name.
\end{proof}

\subsection{The six exceptional kernels and the remaining central sign}

\paragraph{D23.F26.}
Use $X,Y,Z=(XY)^{-1}$ and $S=SU(H)$ from D23.F25.  For the
six ordered pairs
\[
 (X,Y),\ (Y,Z),\ (Z,X),\ (Y,X),\ (Z,Y),\ (X,Z),
\tag{D23.F26.1}
\]
the simultaneous commutator-word evaluation is onto $S^6$.
Their six kernels on $F'$ are distinct even after quotienting
each target by its centre.  In particular the cyclic pair
$(Y,Z)$ does not have a swap-stable derived kernel.

This simultaneous map remains onto the direct product with
any prescribed finite word image whose order has prime divisors
only $2,3,7$.  Such finite images can encode all the old Krammer
pair values, their index-three induced evaluations of $C_f$,
and all needed finite residues of $b=\Psi_3/2$.

On the six exceptional values, impose the two-cycle, the
multiplier-one and multiplier-minus-one three-cycles, including
the correctly conjugated reversed cycles.  Whenever a tuple
satisfies these equations, independently multiplying two of
its three cyclic values by $-I$ gives another tuple satisfying
the same equations and all the same conjugation and generation
conditions.  In particular these finite conditions do not fix
the central sign of $f(X,Y)$ while the old data are held fixed.
They do not resolve the quadratic comparison solely through
the central factor $(-I)^b$ in the lifted orbital relation.
This conclusion does not discard the possible force of its
noncentral, projective periodic-product condition.

\begin{proof}
Every pair in (D23.F26.1) generates the same group
$H_0=\langle X,Y\rangle$, since any two of $X,Y,Z$ determine
the third.  Its derived subgroup is $S$ by D23.F25.  Write
$\overline S=S/\{I,-I\}$.  The conjugation-action orders of
the two marked generators on $S$ or $\overline S$ are,
in the six respective cases,
\[
       (36,6),\ (6,3),\ (3,36),\ (6,36),\ (3,6),\ (36,3).
\tag{D23.F26.2}
\]
These are all distinct.  Suppose two kernels of
$F'\to\overline S$ were equal.  The resulting identification
between their two quotient groups would intertwine conjugation
by each of the two free generators: both kernels are normal
in the full free group, and the quotient identification is
induced by the same words.  Thus the two ordered action orders
would agree.  Equation (D23.F26.2) rules this out.  This is a
kernel argument, not just a failure to find a linear intertwiner.
It also explains why the literal orders $6,6$ at $(Y,Z)$ are
misleading: $Z^3=-I$, whereas $Y^3$ is noncentral.

By D23.F25, $\overline S$ is simple.  A subdirect product of
copies of one nonabelian simple group is a product of diagonal
factors.  One can see the required statement inductively from
Goursat's lemma: a new simple factor is either independent or
is identified with a simple quotient of the preceding product;
an epimorphism from a product of simple groups to a simple
group factors through one coordinate, because the images of
the commuting factors are commuting normal subgroups.
Such an identification would equate two source kernels.
Their distinction therefore gives the full image
$\overline S^{\,6}$.

Let $J$ be the simultaneous image in $S^6$ and let
$Z_0=\{I,-I\}^6$.  The last paragraph gives $JZ_0=S^6$.
Because $Z_0$ is central, this makes $J$ normal and $S^6/J$
abelian.  But $S^6$ is perfect by D23.F25, so $J=S^6$.
This proves full independence including the six Schur-centre
coordinates, rather than only projective independence.

Now let $M$ be a finite word target with prime divisors among
$2,3,7$, and put $M_0=\operatorname{im}(F'\to M)$.
Every nontrivial quotient of $S$ is $S$ or $\overline S$:
a normal subgroup either lies in the centre, or surjects to
the simple projective quotient; in the latter case its quotient
in $S$ is abelian and hence trivial by perfectness.
Both nontrivial quotients have order divisible by $13$ and
$61$, as follows directly from D23.F25's order formula.
A nontrivial quotient of $S^6$ has a nontrivial image of some
factor $S$, so its order also has these prime divisors.
It cannot be a quotient of $M_0$.  Goursat's lemma proves
\[
 \operatorname{im}(F'\longrightarrow S^6\times M_0)
                         =S^6\times M_0.
\tag{D23.F26.3}
\]
Finite images of the ordinary free commutator subgroup and its
profinite closure agree, so all of these finite prescriptions
have common ordinary-word representatives.

Here is why the additional old data described in the statement
really fall under this assertion.  The old pair groups have
orders with prime divisors $2,3,7$ by D23.P6 and D23.C9;
the latter's perfect factors are $SU_3(3)$.  Given a finite
quotient $q:F\to M$ of this kind, induce it from the normal
index-three subgroup $F$ to $E=C_3\amalg C_3$.  The resulting
homomorphism takes values in $M\wr C_3$, which has no new
prime divisor.  Under the homomorphism from the parameter free
group to $E$ sending its two generators to $x,r$, the value
$C_f=f(x,r)$ lies in $F$, and its distinguished induced
coordinate is exactly $q(C_f)$.  Thus these $C_f$ values can
be encoded by ordinary finite homomorphisms of the parameter
free group, even if the original kernel of $q$ is not
$E$-stable.  Finally finite index-three Schreier coefficients
are encoded by the abelianization of the index-three subgroup
modulo the required modulus, followed by its finite coset
action.  To retain $b\bmod m$ use $\Psi_3\bmod2m$ and
restrict to its even values.  For the powers in the displayed
orbital matrices the required moduli have prime divisors
only $2,3$.  Finite products of all these targets remain of
the specified prime type.  No invariance of a congruence
kernel has been assumed in this induction argument.

We now verify the central-sign assertion at the level of all
the stated equations.  Denote the six values in order by
$v_0,\ldots,v_5$.  Two-cycle gives $v_{j+3}=v_j^{-1}$ for
$0\leq j\leq2$, and the two cycles at $(X,Y)$ are
\[
 v_2v_1v_0=I,\qquad
 v_2 Z^{-1}v_1Y^{-1}v_0X^{-1}=I.
\tag{D23.F26.4}
\]
For any $\alpha,\beta\in\mathbb F_2$ replace the first three
values by
\[
 v_0^*=(-I)^\alpha v_0,\qquad
 v_1^*=(-I)^\beta v_1,\qquad
 v_2^*=(-I)^{\alpha+\beta}v_2,
\tag{D23.F26.5}
\]
and replace their three reversed values by the inverses of
these new values.  Both products in (D23.F26.4) acquire the
same scalar $(-I)^{2\alpha+2\beta}=I$.
The reversed cycles use
$Yv_4Y^{-1},X^{-1}v_5X,v_3$, so the same verification applies
with the required inner corrections.  This proves preservation
of every cycle in the full outer three-puncture orbit.
Conjugation of a generator by a value $v_j$ is also unchanged
by its scalar multiplier.  Accordingly every finite test
depending only on the induced generator automorphism,
generation, or conjugacy classes of peripheral products is
unchanged.  Taking $\alpha=1$ changes the sign of $v_0$.
Equation (D23.F26.3) realizes the resulting tuple by one common
word while preserving all the encoded old data.

The exact literal orbital identity of D23.C16 contains
\[
 R\,A_E(x,V)\,Y^{-3b}\,f(X,Y)
       =A_E(x,y)\,Z_{123}^{\,b}(-I)^b,
\tag{D23.F26.6}
\]
where $Z_{123}=(K_1K_2)^3$; the letter $Z$ in
(D23.F26.1)--(D23.F26.4) instead denotes $(XY)^{-1}$.
We keep these two matrices explicitly distinct.
The last sign is essential: in the braid lift the two sides
without it have exponent sums $-6b$ and $6b$, so the exact
central correction is $z_4^{-b}$; the Krammer image of $z_4$
is $-I$.  At fixed old data, the scalar change (D23.F26.5)
can select either lift of the required projective value of
$f(X,Y)$ without affecting the other finite conditions just
proved.  Thus once such a projective value has a solution of
the literal cycle conditions, its sign in (D23.F26.6) supplies
no further obstruction on the old data.  This precise assertion
does not assert that its noncentral value always satisfies
the periodic-product condition, nor that one common finite
word is an actual gentle or full-GT parameter.

The exact action orders and scalar powers used here are part
of \path{agents/five_level_six/krammer_orbital_unitary.py} and
its JSON record.  The central-sign verification is an identity
in any group with central involution and is written out in
(D23.F26.4)--(D23.F26.5); it requires no numerical approximation.
\end{proof}

\subsection{What the non-Ferrari affine quotient sees in the cubic residual}

\paragraph{D23.F27.}
In the literal non-Ferrari quotient $\eta_3$ of supplied FP.31,
\[
 a\mapsto(0,U),\quad b_0\mapsto(0,L),\quad
 c\mapsto(e_1,U),\qquad
 U=\begin{pmatrix}1&1\\0&1\end{pmatrix},\quad
 L=\begin{pmatrix}1&0\\-1&1\end{pmatrix},
\]
retain $\delta=ab_0$, $K=c^2$ and the cubic residual $D_h$ of
D23.C17.  For an actual gentle parameter, the exceptional value
$g=f(\delta,K)$ has the necessary form
\[
 \eta_3(g)=(u e_1,N_\rho),\quad u\in\mathbb F_3,\qquad
 N_0=I,\quad N_1=\begin{pmatrix}0&1\\2&0\end{pmatrix},\quad
 N_2=\begin{pmatrix}1&1\\1&2\end{pmatrix}.
\tag{D23.F27.1}
\]
Here $\rho$ is the usual frame residue modulo three.
The nine finite affine translations of the descended cubic
frame allowed by its exact frame equation give all three
values of $u$, for each of the two values of $b=\Psi_3/2\bmod2$.
This is a necessary finite list, not a lifting assertion for
each entry.

For every value in (D23.F27.1), the periodic cube equation
of D23.C17 holds with either possible residual sign
\[
          \eta_3(D_h)=(0,(-I)^e),\qquad e\in\mathbb F_2.
\tag{D23.F27.2}
\]
Thus this non-Ferrari periodicity test does not exclude the
sign $e=1$.  The affine periodic coinvariant remains nonzero
for both signs.  No smaller comparison gate or vanishing of
$e_3$ is inferred from this calculation.

\begin{proof}
We first identify the marked exceptional pair, not just its
abstract finite group.  In the cubic detector of D23.P3, let
$\eta$ be the affine representation over $\mathbb F_9$ with
\[
 \eta(s)=(0,iS),\qquad \eta(r)=(e_1,R),\quad
 S=\begin{pmatrix}0&1\\-1&0\end{pmatrix},\quad R=-UL.
\]
On the index-two subgroup $E=\langle r,srs\rangle$, all
linear parts have entries in $\mathbb F_3$, so real-part
projection of translations is a homomorphism.  Follow it
by the origin change $t=(0,2)^{\mathsf t}$, as in D23.P10
and D23.C14.  The resulting three marked elements are
\[
 x=\left(0,\begin{pmatrix}2&1\\0&2\end{pmatrix}\right),\quad
 r=\left(\binom22,\begin{pmatrix}0&2\\1&2\end{pmatrix}\right),\quad
 v=srs=\left(\binom11,\begin{pmatrix}2&2\\1&0\end{pmatrix}\right).
\]
Conjugation by the linear matrix
\[
                         H=\begin{pmatrix}2&0\\1&2\end{pmatrix}
\tag{D23.F27.3}
\]
sends the pair $(x,v)$ to
\[
 \eta_3(\delta)=\left(0,\begin{pmatrix}0&1\\2&1\end{pmatrix}\right),
 \qquad
 \eta_3(K)=\left(-e_1,\begin{pmatrix}1&2\\0&1\end{pmatrix}\right).
\tag{D23.F27.4}
\]
These are literal equalities, checked on both inputs.
In the original modular group $\tau=sr$ fixes $x=\tau^{-2}$
by conjugation and sends $r$ to $v$.  Consequently
\[
                    f(x,v)=\tau C_f\tau^{-1},\qquad C_f=f(x,r).
\tag{D23.F27.5}
\]
The conjugation in this equation must be evaluated before real
projection: $\tau$ is outside $E$, and that projection is not
being extended to the whole modular group.

Here are the full finite frame constraints and their evaluation.
Let $A_E$ be D23.F22's descended pure cubic frame and put
$C_\rho=U^2M_\rho U^{-2}$, with
$M_\rho=\left(\begin{smallmatrix}1&\rho\\\rho&1+\rho^2\end{smallmatrix}\right)$.
D23.F22 gives its linear part
\[
                  A_0=(-I)^b C_\rho^{-1}.
\]
D23.P3 gives the complete affine value of the descended pair
coordinate $f_E$, using $\rho(D_Eh)=\rho(h)$ and
$a_3(D_Eh)=b$.  Thus the translation $t_E$ of $\eta(A_E)$
satisfies the exact affine equation
\[
       \eta(A_E)\eta(r)\eta(A_E)^{-1}
                         =\eta(f_E)^{-1}\eta(r).
\tag{D23.F27.6}
\]
Solving its two linear translation equations over $\mathbb F_9$
gives the following complete parametrizations, with $z\in\mathbb F_9$.
The last column is the first translation coordinate after
forming $C_f=A_E^{-1}x^{-3b}$, applying (D23.F27.5), real
projection with the stated origin shift, and conjugation by $H$.
Write $\operatorname{Im}(a+bi)=b\in\mathbb F_3$.
\[
\begin{array}{cc|c|c}
 \rho&b&t_E&u\\\hline
 0&0&(z,-z)&\operatorname{Im}z\\
 0&1&(z,1+i-z)&2-\operatorname{Im}z\\
 1&0&(z,2+2i+z)&-\operatorname{Im}z\\
 1&1&(z,2+i+z)&1+\operatorname{Im}z\\
 2&0&(2+2i,z)&1-\operatorname{Im}z\\
 2&1&(1,z)&2+\operatorname{Im}z.
\end{array}
\tag{D23.F27.7}
\]
In every row the second resulting translation coordinate is
zero and the linear part is $N_\rho=C_\rho^{-1}$.
For an explicit check of the powers appearing here, D23.P3's
affine coordinates are
\[
 \eta(\tau)=((0,-i),iU^{-1}),\qquad
 \eta(x)=((1+i,1+2i),-U^2),\qquad
 \eta(x)^3=((-i,1+2i),-I).
\]
Substituting these coordinates in the affine multiplication
law verifies every row of (D23.F27.7).  In particular no central
linear factor has been discarded as an affine translation.
The table proves (D23.F27.1) and explains why its set is
independent of $b$.

Since $D_h$ is evaluated in the linear first-three-strand
subgroup, D23.C17.2 gives exactly (D23.F27.2).  Put
$\epsilon=(-1)^e$.  The translation $u e_1$ in (D23.F27.1)
commutes with $\eta_3(K)$ because $U^2e_1=e_1$.  Direct affine
multiplication therefore gives
\[
 \eta_3(\delta D_hg^{-1}KgD_h^{-1})
             =(\epsilon v_\rho,L_\rho),
\]
where
\[
\begin{array}{c|c|c|c}
 \rho&L_\rho&v_\rho&\ell_\rho\\\hline
 0&\left(\begin{smallmatrix}0&1\\2&2\end{smallmatrix}\right)
       &(0,1)^{\mathsf t}&(2,1)\\
 1&\left(\begin{smallmatrix}1&1\\0&1\end{smallmatrix}\right)
       &(2,2)^{\mathsf t}&(0,2)\\
 2&\left(\begin{smallmatrix}1&0\\2&1\end{smallmatrix}\right)
       &(1,0)^{\mathsf t}&(1,0).
\end{array}
\tag{D23.F27.8}
\]
Each $L_\rho$ is a nonidentity unipotent with
$(L_\rho-I)^2=0$.  In characteristic three this implies
$I+L_\rho+L_\rho^2=0$.  Hence every affine lift has cube one,
for either sign of its translation.  The last column verifies
$\ell_\rho L_\rho=\ell_\rho$ and
$\ell_\rho v_\rho=1$.  Its affine periodic coinvariant is
therefore $\epsilon$, nonzero for either sign.
Conjugation by $(0,-I)$ interchanges these two nonzero classes;
requiring ordinary affine conjugacy does not distinguish them.

There is also a simpler general reason the displayed residual
sign cannot enter this test directly.  In this quotient
$\eta_3(D_h)$ commutes with $\eta_3(\delta)$, so the whole
periodic product is conjugate by $\eta_3(D_h)$ to
$\eta_3(\delta g^{-1}Kg)$.  The explicit calculation above
checks that the possible actual exceptional values do satisfy
the cube equation, rather than merely observing that the
residual sign drops out of it.

The certificate \path{agents/five_level_six/orbital_affine216.py}
rechecks the two exact marked conjugacies, retains the
$\mathbb F_9$ conjugation by $\tau$, enumerates all $81$
translation candidates in each of the six frame equations,
and checks all $108$ remaining signed frame cases.  Its
record is \path{orbital_affine216.json}.  The full joint
orbital identities contain additional correlated evaluations;
this proposition only establishes the stated obstruction for
the periodic residual in this specific quotient.
\end{proof}

\subsection{The harmonic exceptional pair is independent of the cubic pair}

\paragraph{D23.F28.}
In the same literal Krammer representation as D23.F25 put
\[
 S_0=K_1K_2K_1,\qquad Y=K_3^2,\qquad Z_H=(S_0Y)^{-1}.
\]
Then $\langle S_0,Y\rangle=U(H)$, with derived group $SU(H)$.
Its three local elements have literal and projective orders
$24,6,12$, respectively.  The full six-marking harmonic orbit
and the full six-marking cubic orbit of D23.F26 therefore have
twelve distinct kernels on $F'\to PSU(H)$ and simultaneous
image $SU(H)^{12}$ before projectivizing.

Consequently the two exceptional Krammer families have no
common derived word coordinate.  Their two central phase choices
can also be varied independently while preserving all their
two-cycle, positive and negative three-cycle equations and
holding the old finite Krammer data fixed.  This is a finite
word-image assertion, not a proof that their two noncentral
orbital equations can always be solved or that any assignment
is an actual gentle parameter.

\begin{proof}
Both generators preserve D23.F25's Hermitian form.  The same
matrix $T=Y^2$ is a rank-one unitary transvection.  The exact
orbit of its image line under $S_0,Y$ again consists of all
$22204$ isotropic lines.  D23.F25's complete generation proof
therefore gives $SU(H)\subseteq\langle S_0,Y\rangle$.
Since $\det S_0=i^3=-i$ generates the four norm-one scalars,
the determinant image is the full determinant group.  This
proves $\langle S_0,Y\rangle=U(H)$; its derived group is
$SU(H)$ by D23.F25's perfectness proof.

The direct order computation gives
\[
 (\operatorname{ord}S_0,\operatorname{ord}Y,
        \operatorname{ord}Z_H)=(24,6,12),
\]
and no nonidentity power below these respective orders is
scalar.  The six ordered conjugation-action orders are thus
\[
       (24,6),(6,12),(12,24),(6,24),(12,6),(24,12).
\]
Together with D23.F26's six ordered pairs
$(36,6),(6,3),(3,36),(6,36),(3,6),(36,3)$ they are twelve
different ordered pairs.  All six harmonic pairs generate
$U(H)$, since any two of $S_0,Y,Z_H$ determine the third.
Their commutator images are therefore all $SU(H)$.

If two projective derived kernels agreed, their quotient
identification would intertwine the actions of both free
generators.  The ordered action orders would agree, contradicting
the displayed list.  D23.F25 proves $PSU(H)$ simple and
$SU(H)$ perfect.  Applying the subdirect-product argument
of D23.F26 to twelve factors gives $PSU(H)^{12}$ first and
then $SU(H)^{12}$, because any central gluing would yield an
abelian quotient of that perfect product.  Its independence
from every old finite target with prime divisors $2,3,7$
is proved by the same Goursat argument: a nontrivial quotient
of a special unitary factor retains the prime divisors $13,61$.

Finally apply D23.F26.5 separately within the harmonic and
cubic triples.  Each operation multiplies two cyclic values
by $-I$, preserves all the equations, and changes its first
exceptional value's central sign.  The twelve-factor
surjectivity realizes the two changes independently by a
common ordinary commutator word, with the specified old data
unchanged.  The restriction to finite word images is essential
and is part of the assertion.

The exact certificate is
\path{agents/five_level_six/krammer_harmonic_orbital.py},
with record \path{krammer_harmonic_orbital.json}.
It checks the form, complete line orbit and all local scalar
orders.  The group-theoretic deductions use the complete
generation, perfectness, simplicity and Goursat proofs in
D23.F25--F26.
\end{proof}


\section{A coupled Hurwitz quotient preserved by full GT}
\subsection{A genuine Hurwitz quotient retaining both marked sign detectors}

\begin{statement}{D23.P17}{Proved}{A degree-1344 five-point representation with both required marked group quotients}
\end{statement}
There is an explicit permutation representation
\[
 q:B_4/Z(B_4)\longrightarrow S_{1344},\qquad
 (\sigma_1,\sigma_2,\sigma_3)\longmapsto(A,B,C),
\]
with the following two properties, using exactly the standard
pentagon markings
\[
 X_1=A^2,\quad X_2=B^2,\quad X_3=C^2,\quad
 X_4=(ABA)^2,\quad X_5=(CBC)^2.
 \tag{D23.P17.1}
\]
First, $\langle X_1,X_2\rangle$ has a marked quotient onto
$E_0=\mathbb F_3^2\rtimes\operatorname{SL}_2(\mathbb F_3)$,
of order $216$, sending its two generators to
\[
 \xi(x)=\begin{pmatrix}2&1&0\\0&2&0\\0&0&1\end{pmatrix},
 \qquad
 \xi(y)=\begin{pmatrix}2&0&1\\2&2&0\\0&0&1\end{pmatrix}.
 \tag{D23.P17.2}
\]
Second, $\langle X_3,X_4\rangle$ has a marked quotient onto
$S_4$, sending its generators to $(123)$ and $(1234)$.
Consequently this single genuine five-point representation
retains both the exact cubic detector of D23.P10 and the
harmonic word detector of supplied FP.37.

The final quotient maps in both assertions are quotients of
the permutation \emph{image groups}. No equivariant map of
the degree-36 set onto nine points, or of the degree-72 set
onto four points, is asserted or needed.

\begin{proof}
Use $G=E_0$ in its natural affine realization. An element will
be written as $(u,v,a,b,c,d)$, meaning the affine map
\[
 \binom{z_1}{z_2}\longmapsto
 \begin{pmatrix}a&b\\c&d\end{pmatrix}\binom{z_1}{z_2}
 +\binom uv.
\]
Let $\mathcal C$ be the conjugacy class of
$(0,0,2,1,0,2)=\xi(x)$. It has 36 elements, all of order six.
Consider the finite set
\[
 \mathcal N=\mathcal C^4/G
\]
of all four-tuples modulo simultaneous conjugation. There is
no imposed product-one condition. The left Hurwitz action is
\[
 \sigma_i:(g_1,\ldots,g_i,g_{i+1},\ldots,g_4)
 \longmapsto
 (g_1,\ldots,g_i g_{i+1}g_i^{-1},g_i,\ldots,g_4).
 \tag{D23.P17.3}
\]
This action respects simultaneous conjugation and satisfies
the Artin braid relations. Its full twist acts on a tuple by
simultaneous conjugation by its product. It therefore induces
an action of $B_4/Z(B_4)$ on $\mathcal N$.

The complete enumeration gives $|\mathcal N|=7796$ and 17
$B_4$-orbits. Exactly one of these has degree 1344. It contains
the tuple
\[
\begin{split}
 &(0,0,0,1,2,1),\quad (0,0,0,1,2,1),\\
 &(0,0,1,1,2,0),\quad (0,1,1,1,2,0).
\end{split}
 \tag{D23.P17.4}
\]
The displayed tuple has product $(2,2,2,1,0,2)$ of order six,
which is noncentral in $G$. Restricting the Hurwitz action to
this orbit defines $q$.

Here is a direct canonical enumeration of the entire set,
so the assertion of uniqueness is not limited to a selected
product or to already promising orbits. Fix the first factor
to be the smallest element of $\mathcal C$ in the displayed
six-coordinate ordering. Its centralizer in $G$ has order six.
Every simultaneous-conjugacy class of four-tuples has a
representative with this first factor; two such representatives
are equivalent exactly under that six-element centralizer.
Thus enumerate all $36^3$ possible remaining factors and
identify centralizer orbits. Conjugating the new first factor
back to the fixed one computes each Hurwitz generator. The
certificate explicitly checks the three braid relations and
$(ABC)^4=1$ on all 7796 points before extracting the 1344 orbit.

For the first marked quotient, first perform the analogous
enumeration of $\mathcal C^3/G$. It contains a $B_3$-orbit
$\mathcal O$ of degree 36, represented by
\[
 (0,0,0,1,2,1),\quad (0,0,1,1,2,0),\quad
 (0,1,2,0,2,2).
 \tag{D23.P17.5}
\]
Write $a,b$ for its two Hurwitz permutations. The exact group
$P=\langle a^2,b^2\rangle$ has order 1944. Adjoin the two
natural nine-point permutations in (D23.P17.2) on a disjoint
set. The joint group generated by
\[
 (a^2,\xi(x)),\qquad(b^2,\xi(y))
\]
also has order 1944. Both projections are onto, so projection
to $P$ is an isomorphism and the other projection defines
the required marked epimorphism $P\twoheadrightarrow E_0$.

Forgetting the fourth entry defines a $B_3$-equivariant map
from $\mathcal N$ to $\mathcal C^3/G$. Within the 1344 orbit,
the inverse image of $\mathcal O$ is a $B_3$-invariant subset
of size 432, mapping onto $\mathcal O$. Hence restriction to
this subset and its equivariant image give a marked epimorphism
$\langle X_1,X_2\rangle\twoheadrightarrow P$. Composing the
two image-group epimorphisms proves the first assertion. This
argument retains the fourth point in the full $B_4$ action;
only the stated $B_3$ restriction uses forgetting.

For the second assertion, decompose the 1344-point set into
orbits under $\langle X_3,X_4\rangle$. Its ninth component,
when components are numbered from zero in increasing degree
and then discovery order, has degree 72. Its image group $M$
has order $995328$. Adjoining on four new points the marked
permutations $(123),(1234)$ gives a joint group of the same
order $995328$. As above, equality of orders and surjectivity
of both projections prove that the two marked generators of
$M$ map onto the specified generators of $S_4$. Restriction
to this invariant component gives
$\langle X_3,X_4\rangle\twoheadrightarrow M$, completing the
proof.
\end{proof}

\paragraph{Exact certificates.}
The full reconstruction is
\path{hurwitz_cubic216_four.py}; its JSON records all
retaining orbits, the unique degree-1344 orbit, and its literal
representatives. The raw matrices of $G$, class enumeration,
and degree-36 prerequisite are reconstructed by
\path{hurwitz_cubic216_triples.py order6}. The selected
orbit and the actual $432\to36$ set map are stored in
\path{hurwitz_cubic216_four_13.npz}. The short independent
replay \path{hurwitz_cubic216_markings.py} verifies both
image-group order equalities and the forgetting equivariance
directly from the stored permutations. Its JSON explicitly
distinguishes the set map from the two group quotient maps.

\paragraph{Scope.}
This proves a new finite representation and two exact marked
quotients, not a decision of the pentagon equation. The old
explicit P12 escape word fails the literal pentagon in this
representation (support 1304 out of 1344 points) and also
fails the first-pair three-cycle, so it cannot be reused
as a common finite shadow without additional corrections.
No actual gentle or full-GT lift is inferred from any finite
word computation here.

\subsection{Full-GT stability of the new Hurwitz representations}

\begin{statement}{D23.P18}{Proved}{Functorial stability of an isolated finite Hurwitz orbit}
\end{statement}
Let $G$ be a finite group, let $\mathcal C$ be a conjugacy class
in $G$, and let $\mathcal N_n=\mathcal C^n/G$ carry the left
Hurwitz action
\[
 (g_i,g_{i+1})\longmapsto(g_i g_{i+1}g_i^{-1},g_i).
\]
Every multiplier-one element $h\in\widehat{GT}$ induces a
permutation $V_h$ of $\mathcal N_n$ satisfying
\[
 V_h q_n(\beta)V_h^{-1}=q_n(\Phi_h(\beta))
 \qquad(\beta\in\widehat B_n),
 \tag{D23.P18.1}
\]
where $\Phi_h$ is the standard Drinfeld automorphism. Therefore
an orbit whose degree occurs exactly once among all
$B_n$-orbits in $\mathcal N_n$ is preserved. Its permutation
representation has kernel invariant under the multiplier-one
subgroup of full GT, and its transform
under $\Phi_h$ is conjugate to the original representation by
an actual permutation of that orbit.

In particular, the degree-1344 representation of D23.P17 is
stable in this sense. If $F=f(X_1,X_2)$ and $R=f(X_3,X_4)$ in
its permutation image, there exists $V\in S_{1344}$ with
\[
 \boxed{\quad
 VAV^{-1}=A,\qquad
 VBV^{-1}=F^{-1}BF,\qquad
 VCV^{-1}=RCR^{-1}.
 \quad}                                                     \tag{D23.P18.2}
\]
The analogous degree-36 triple representation is also stable.
These are necessary conditions for actual full GT. They are
not asserted for every actual gentle element.

\begin{proof}
We first fix the rightmost-strand convention, including the
inverse in the Hurwitz action. Let $M_{n,1}\subset B_{n+1}$
be the subgroup fixing the last strand. Forgetting that strand
gives the split Artin sequence
\[
 1\longrightarrow F_n\longrightarrow M_{n,1}
   \longrightarrow B_n\longrightarrow1,
 \tag{D23.P18.3}
\]
whose section uses $\sigma_1,\ldots,\sigma_{n-1}$. Its free
kernel has the peripheral basis
\[
 u_i=\sigma_n\cdots\sigma_{i+1}\sigma_i^2
          \sigma_{i+1}^{-1}\cdots\sigma_n^{-1},
 \qquad 1\le i\le n.
\]
In particular $u_n=\sigma_n^2$, and the kernel is its normal
closure in $M_{n,1}$. If $\alpha_\beta$ denotes conjugation by
the section element $\beta$ on $F_n$, the Artin relations give
\[
\begin{split}
 \alpha_{\sigma_i}(u_i)&=u_{i+1},\\
 \alpha_{\sigma_i}(u_{i+1})&=u_{i+1}^{-1}u_i u_{i+1},\\
 \alpha_{\sigma_i}(u_j)&=u_j\quad(j\ne i,i+1).
\end{split}                                                 \tag{D23.P18.4}
\]
Consequently the chosen left Hurwitz action on a homomorphism
$\varphi:F_n\to G$ is
\[
 q_n(\beta)[\varphi]=[\varphi\circ\alpha_\beta^{-1}].
 \tag{D23.P18.5}
\]
This is precisely the tuple action in the statement.

The sequence (D23.P18.3) remains a split exact sequence after
full profinite completion. Here is a direct topology check.
Characteristic finite-index subgroups of the finitely generated
free kernel are cofinal among its finite-index subgroups. For
such a subgroup $K$, the action on $F_n/K$ factors through a
finite quotient of $B_n$. The associated finite semidirect
product is a quotient of $M_{n,1}$ whose restriction to $F_n$
has kernel $K$. Thus the induced topology on $F_n$ is its full
profinite topology. The completed section and kernel have
trivial intersection, and their compact product contains the
dense discrete semidirect product; hence it is all of
$\widehat M_{n,1}$. This proves the asserted completed split
sequence. Since $M_{n,1}$ has finite index in $B_{n+1}$, its
completion is the corresponding open subgroup of
$\widehat B_{n+1}$.

We use the full-profinite Drinfeld braid automorphism theorem,
with its standard formula
\[
 \Phi_h(\sigma_1)=\sigma_1^\lambda,
 \qquad
 \Phi_h(\sigma_i)=f(y_i,\sigma_i^2)^{-1}
        \sigma_i^\lambda f(y_i,\sigma_i^2),
\quad
 y_i=\sigma_{i-1}\cdots\sigma_2\sigma_1^2
            \sigma_2\cdots\sigma_{i-1}.
 \tag{D23.P18.6}
\]
The action induces the identity on the symmetric-group
quotient, because $\lambda$ is odd and the conjugators are
pure. It preserves the section $\widehat B_n$ and restricts
there to its standard Drinfeld action. It therefore preserves
$\widehat M_{n,1}$.

Let $N=\widehat F_n$. Under the rightmost forgetting map,
the element $f(y_n,\sigma_n^2)$ maps to $f(y_n,1)=1$.
Thus this conjugator belongs to $N$, and
$\Phi_h(u_n)$ is $N$-conjugate to $u_n^\lambda$.
Since $N$ is the closed normal closure of $u_n$ in
$\widehat M_{n,1}$, $\Phi_h(N)\subseteq N$. Applying the same
argument to $h^{-1}$ gives equality. Put
$\kappa_h=\Phi_h|N$.

Every $u_i$ is a conjugate of $u_n$ by a braid in the section;
one may take $v_i=\sigma_i^{-1}\cdots\sigma_{n-1}^{-1}$ and
$u_i=v_i u_n v_i^{-1}$. The Artin action of any braid carries
a peripheral generator to an $N$-conjugate of the generator
with the corresponding puncture label. Since $\Phi_h(v_i)$
has the same permutation as $v_i$, it follows that
$\kappa_h(u_i)$ is $N$-conjugate to $u_i^\lambda$.

Define
\[
 V_h[\varphi]=[\varphi\circ\kappa_h^{-1}]
 \quad\text{on }\operatorname{Hom}_{\rm cont}(N,G)/G.
 \tag{D23.P18.7}
\]
This is a bijection, and when $\lambda=1$ it preserves the
condition $\varphi(u_i)\in\mathcal C$ for all $i$. Thus it
restricts to a permutation of the entire finite set
$\mathcal C^n/G$.
The semidirect product gives the exact covariance
\[
 \kappa_h\alpha_\beta\kappa_h^{-1}
       =\alpha_{\Phi_h(\beta)}.
\]
Combining it with (D23.P18.5)--(D23.P18.7) proves (D23.P18.1)
with the displayed orientation. If a different section is
used and a kernel cocycle appears, the covariance differs
only by an inner automorphism of $N$. Precomposition by that
inner automorphism becomes simultaneous conjugation by its
$G$-image, and hence disappears on $\operatorname{Hom}(N,G)/G$.

Because $\Phi_h$ is onto $\widehat B_n$, equation
(D23.P18.1) carries each $B_n$-orbit bijectively onto another
$B_n$-orbit of the same degree. A degree occurring exactly
once must therefore be preserved. Restriction of (D23.P18.1)
to it gives permutation conjugacy and kernel invariance.

For D23.P17 the complete four-tuple orbit degree list is
\[
 1,4,9,9,9,12,12,12,12,72,324,360,405,
 1344,1512,1512,2187.
\]
The degree 1344 therefore occurs exactly once in the entire
$\mathcal C^4/G$. The complete triple list is
\[
 1,1,2,2,4,4,4,4,4,36,40,60,60,
\]
so the degree 36 is likewise isolated.

Finally evaluate (D23.P18.6) on the three generators of
$B_4$. For $i=2$, $y_2=\sigma_1^2$, giving $F^{-1}BF$.
For $i=3$, replace $y_3$ by $z_3=(\sigma_1\sigma_2)^3$:
indeed $z_3=\sigma_1^2y_3$, and $\sigma_1^2$ commutes with
both $y_3$ and $\sigma_3^2$. A commutator word has zero
exponent sum in its first input, so this commuting factor
does not change its value. Thus the conjugator is
$f(X_4,X_3)=f(X_3,X_4)^{-1}=R^{-1}$, by the two-cycle.
This gives exactly (D23.P18.2).
\end{proof}

\paragraph{Additional functorial restrictions.}
The permutation (D23.P18.7) commutes with postcomposition by
every automorphism of $G$ preserving $\mathcal C$. It preserves
the image subgroup of a homomorphism because $\kappa_h$ is
onto. It also preserves the conjugacy class of the peripheral
product when $\lambda=1$: the boundary word $u_1\cdots u_n$
is $z_n^{-1}z_{n+1}$, and the standard action sends each full
twist $z_j$ to $z_j^\lambda$. These restrictions may be imposed
in a simultaneous-conjugacy test; no extra geometric rigidity
of $G$ or of a matrix representation is being assumed.

\paragraph{Primary source and check scope.}
The full-profinite automorphism theorem used in
(D23.P18.6) is Lochak--Schneps, \emph{The
Grothendieck--Teichm\"uller group and automorphisms of braid
groups}, LMS Lecture Notes 200 (1994), \S5, Proposition 10,
author-PDF p.23 (PDF index 22):
\url{https://webusers.imj-prg.fr/~leila.schneps/Fls.pdf}.
Its formula uses the displayed $y_i$ marking. Nakamura,
RIMS-1691, \S5.4, equation (5.4.2), author-PDF p.28,
\url{https://www.kurims.kyoto-u.ac.jp/preprint/file/RIMS1691.pdf},
records the equivalent partial-full-twist formula in Galois
scope; that source is used to check the marking, not to
replace the full-GT theorem by a Galois statement.
The new proof here is the induced finite-Hurwitz stability
argument and its application to the two completely enumerated
isolated orbits.

The independent exact word replay
\path{hurwitz_stability_artin_check.py} verifies
(D23.P18.4), all displayed peripheral transports, the boundary
full-twist relation, and the commuting input replacement in
$B_4$ and $B_5$, using the faithful Artin representation.
Its JSON keeps the literal peripheral braid words. The
published full-GT automorphism theorem is a cited input, not
independently reproved by that finite word check.

\subsection{The unique first-pair value in the degree-36 Hurwitz orbit}

\paragraph{D23.H1.}
Let $A,B\in S_{36}$ be the literal braid generators of the unique
degree-$36$ three-tuple Hurwitz orbit in D23.P17, and set
$X=A^2$, $Y=B^2$. Their pure image has order $1944$ and derived
subgroup of order $648$. Each of the six marked cubic values
$\xi(f)$ in D23.P10 has nine lifts in that derived subgroup.
For each of these six labels $(\rho,a_3)$, exactly one lift
satisfies the two-cycle and both multiplier-one and
multiplier-minus-one three-cycles. Exactly that same lift admits
a permutation $P\in S_{36}$ with
\[
 PAP^{-1}=A,\qquad PBP^{-1}=f(X,Y)^{-1}Bf(X,Y).
\]
Consequently this necessary permutation-conjugacy condition
leaves all six cubic labels. It does determine the complete
first-pair value on this orbit, not merely its order-$216$
projection. No actual gentle lift or full-GT element is constructed
by this finite calculation.

\begin{proof}
The certificate constructs the joint marked source on the
$36$ Hurwitz points and the nine affine points of $E_0$ using
exactly $(X,\xi(x))$ and $(Y,\xi(y))$. Exhaustive closure has
$1944$ elements, and each permutation on the first set determines
its affine coordinate uniquely. This proves that the affine
projection used here is the specified marked quotient; it is
not an abstract group identification. An exact Schreier--Sims
computation followed by enumeration gives $648$ elements in
the derived subgroup. Their projections have precisely nine
preimages for each of P10's six specified affine matrices.

Put $D=AB$, $S=ABA$, and $Z_0=(XY)^{-1}$. The certificate
checks $D^3=S^2=1$ and the literal conjugation identities
\[
 DXD^{-1}=Y,\qquad DYD^{-1}=Z_0,\qquad
 SXS^{-1}=Y,\qquad SYS^{-1}=X.
\]
Thus for a common commutator word with value $F$, its cyclic
values are $F_1=DFD^{-1}$ and $F_2=D^2FD^{-2}$, and its
reversed value is $SFS^{-1}$. The three exact tests are
\[
 SFS^{-1}=F^{-1},\qquad F_2F_1F=1,\qquad
 F_2Z_0^{-1}F_1Y^{-1}FX^{-1}=1.
\]
All $54$ candidates are checked. There is exactly one passing
value above each label.

For completeness, simultaneous permutation conjugacy is checked
by a separate exhaustive method. If $PAP^{-1}=A$ and
$PBP^{-1}=F^{-1}BF$, the image $P(0)$ of one chosen point
determines every other image, since the original braid action
is transitive. For each of its $36$ possible values, propagate
\[
 P(Au)=A(Pu),\qquad P(Bu)=(F^{-1}BF)(Pu)
\]
along the two generators, rejecting any inconsistent image or
repeated target. A surviving assignment is checked on both
entire permutations. This procedure loses no possible
conjugator. It gives exactly the same unique candidate above
each label, independently of imposing the cycle equations.

The full enumeration is reproducible in
\path{code/check_hurwitz36_gt_stability.py}; all $54$ values,
test outcomes, and an explicit conjugator for every passing
value are retained in
\path{records/hurwitz36_gt_stability.json}.
The input generators are
\path{agents/pentagon/hurwitz_cubic216_order6_triple_9.npz}.
The full-GT necessity of preserving this unique Hurwitz orbit
uses the functorial argument of D23.P18. The finite counts and
conjugators proved here do not themselves require that theorem.
\end{proof}

\subsection{A Galois normalization of the remaining mixed state}

\paragraph{D23.H2 (Galois normalization).}
Put $k=\Q^{\mathrm{ab}}$ and $\mathcal H_5=\GT\cap\ker\lambda$.
For every $h\in\mathcal H_5$ there is
$\sigma\in G_k\subset\mathcal H_5$ such that
\[
 \rho(h\sigma)=0\pmod3,\qquad
 a_2(h\sigma)=0,\qquad a_3(h\sigma)=\Delta(h).
\]
Consequently the existence of a full-GT element with $\Delta=1$
is equivalent to the existence of one with
$(\rho\bmod3,a_3,a_2)=(0,1,0)$.
This is a symmetry reduction of the same gate, not an exclusion
of any mixed state.

\begin{proof}
The supplied DELTA\_ON\_GT A.88 and A.102 identify
$\rho\bmod3=\Psi_2\bmod3$ on $\mathcal H_5$ and make this
an additive character. On $G_k$, the arithmetic comparison
in A.88 identifies $\Psi_2$ with the Kummer character of $2$,
up to the common sign determined by the monodromy marking.
Only its kernel and surjectivity modulo three are used here.
The source convention is recorded separately in A.100;
no equality between differently marked integral characters is
needed. The two frame bits are characters on the multiplier-one
gentle group by LM7\_PUSH E.65. Its Galois comparison and
DELTA\_ON\_GT A.12--A.13 give
\[
 a_2(\sigma)=a_3(\sigma)=\kappa_{r_0}(\sigma),
 \qquad r_0=(3/4)^{1/4}.
\]

The extension $k(\sqrt[3]{2})/k$ is cyclic of degree three.
Indeed $\zeta_3\in k$, and $\sqrt[3]{2}\notin k$: its real
cubic field is not normal over $\Q$, whereas every finite
subextension of the abelian extension $k/\Q$ is normal.
The extension $k(r_0)/k$ has degree two, as proved in the
supplied A.12 and A.25. Equivalently,
$k(r_0)=k(\sqrt[4]{3})$ because $\sqrt2\in k$.
These two finite Galois extensions have coprime degrees,
so their intersection is $k$ and their compositum has Galois
group $C_3\times C_2$. The restriction map from $G_k$ to that
group is surjective. Hence
\[
 G_k\longrightarrow\mathbf F_3\times\mathbf F_2,
 \qquad \sigma\longmapsto(\rho(\sigma),a_2(\sigma))
\]
is surjective, regardless of the common sign convention on
the first factor.

Choose $\sigma$ with
$\rho(\sigma)=-\rho(h)$ modulo three and
$a_2(\sigma)=a_2(h)$. Additivity of the three characters gives
the displayed normalization, since
$a_3(h)+a_2(h)=\Delta(h)$ on $\mathcal H_5$.
Multiplication by $\sigma$ takes place within actual full GT.
Thus this argument does not replace the lifting question by
an assertion about ordinary free-group words.
\end{proof}

% DELTA23: exact compact structure inside the new genuine Hurwitz quotient.
\subsection{A compact mixed quotient with two different octahedral markings}

\paragraph{D23.C21 (the mixed degree-seventy-two quotient).}
Use the genuine degree-$1344$ four-strand Hurwitz action constructed
by the pentagon workstream, and its mixed $(X_3,X_4)$ component
number $9$ of degree $72$. Let $H$ be the image of that pair.
Its derived-series orders are
\[
 995328,\quad82944,\quad27648,\quad1728,\quad64,\quad1,
\]
and $H^{\rm ab}\simeq C_2^2\times C_3$.
There is an explicit small quotient
\[
 H/H'''\simeq C_6\times(S_4\times_{S_3}S_4),
 \qquad |H/H'''|=576,
\tag{D23.C21.1}
\]
where the two maps to $S_3$ are the quotients by the Klein
four-groups, with their uniquely compatible marking.
Up to separately specified conjugations in the two $S_4$ factors,
the marked generator pairs are
\[
 ((123),(1234)),\qquad ((123),(14)).
\tag{D23.C21.2}
\]
The first has peripheral orders $(3,4,4)$ and is the existing
harmonic detector. The second has peripheral orders $(3,2,4)$;
its extra phase is not identified here with $\Psi_3/2$.
In particular A.107's affine modulo-three computation does
not transfer to this different finite target solely from an
abstract Nielsen reordering of modular generators.

The two $S_4$ factors together have order $96$, derived orders
$96,48,16,1$, and the same-word image of $F''$ is the full
$V_4\times V_4$. Thus their two Klein-group coordinates are
independent before the actual gentle equations are imposed.

The retained first-pair degree-$36$ image has derived orders
\[
                       1944,648,162,81,1.
\]
Its joint same-word image with $H$ is the full direct product,
of order $1934917632$. This excludes a direct common-quotient
comparison between those two pair images alone. It does not
make the full pentagon automatic in the genuine Hurwitz action.

Finally, each of the degree-$24$ components numbered
$9,11,13,15$ for $(X_4,X_5)$ has group
$\operatorname{SL}_2(\mathbb F_3)$ and exactly the same marked
free-group kernel as the pure pair $(U^2,L_0^2)$.
Their actual $f$ values are therefore transported from the
same original signed modulo-three free-word value.

\begin{proof}
All assertions are verified from the literal permutations
in the supplied Hurwitz arrays by
\path{agents/characters/code/hurwitz_mixed72_structure.py}.
The record
\path{agents/characters/records/hurwitz_mixed72_structure.json}
contains the derived orders, block labels, quotient generators,
and the exact four-letter conjugator below.

Among the minimal invariant block systems of the mixed action
is a system of eighteen blocks of size four. Its image $G$
has order $144$ and derived orders $144,12,4,1$.
Its centre is $C_6$. Its derived subgroup is $A_4$, and adjoining
the second marked generator to this subgroup gives an $S_4$
complement of order $24$ meeting the centre trivially.
Consequently $G=C_6\times S_4$.
The $S_4$ action on the four order-three subgroups of $G'$ is
faithful; in the recorded initial labels its generators are
\[
             [2,0,1,3],\qquad[0,1,3,2].
\]
Conjugation by $[2,1,0,3]$ changes these to
$p=[1,2,0,3]$ and $v=[3,1,2,0]$, exactly the second pair
in (D23.C21.2). Their orders, and the order of their product,
are $3,2,4$.

Joining this eighteen-block action with the already certified
marked harmonic $S_4$ gives order $576$ and derived orders
$576,48,16,1$. Both factors kill $H'''$, and the index of
the joint kernel is exactly $|H|/|H'''|=576$, so this kernel
is $H'''$. Joining just the two four-letter actions gives
order $96$. Their common quotient is $S_3$, and their full
second derived subgroup is $V_4\times V_4$. By Goursat's lemma
the common quotient of $G$ and the first $S_4$ has order
$144\cdot24/576=6$; the corresponding quotient of $S_4$
is $S_3$. The centre $C_6$ of $G$ maps trivially to this
centreless common quotient. The fibre product therefore
splits off that entire central factor. The orders prove
the fibre-product description in (D23.C21.1).

The first/mixed product assertion follows by adjoining their
literal degree-$36$ and degree-$72$ permutations with the
same two free generators. The computed joint order is
$1944\cdot995328$; no assumption about a presentation or an
actual profinite lift is used.
For each stated $(X_4,X_5)$ component, adjoining the eight-point
nonzero-vector action of $(U^2,L_0^2)$ leaves its order equal
to $24$. Both projections are onto, so the joined group is
the graph of a marked isomorphism. This proves equality of
their free-group kernels and the resulting equality of all
same-word evaluations.
\end{proof}

\paragraph{D23.C22 (the complete finite gentle test in the octahedral core).}
Let $p=(123)$ and $v=(14)$, and induce the marked homomorphism
$F\to S_4$, $x\mapsto p,y\mapsto v$, from
$F=\ker(B\to S_3)$ to $B=C_2\amalg C_3$ on twenty-four
letters. Its image has order $2654208$; the free pure image
has derived-series orders
\[
                         442368,110592,4096,1.
\tag{D23.C22.1}
\]
This is the full $S_3$ Nielsen core of the $(3,2,4)$ marking.
The marked $(p,t)$ quotient, $t=(1234)$, already factors
through its pure image: adjoining it leaves order $442368$.

Among all $110592$ possible images of $f\in F'$,
exactly $576$ satisfy the two-cycle, $2304$ satisfy the
multiplier-one three-cycle, and exactly twelve satisfy both.
All twelve satisfy the $m=-1$ three-cycle as well.
Their two marked four-letter values are precisely
\[
 f(p,t)=t^{2\nu}p^a,\qquad
 f(p,v)=((14)(23))^\kappa p^a,
 \quad (a,\nu,\kappa)\in\mathbb F_3\times\mathbb F_2^2,
\tag{D23.C22.2}
\]
each triple occurring once. Every one of these twelve
residues also has an exact finite action conjugator on the
twenty-four-letter representation and both pure torsion
frames in the finite pure image. Thus none is excluded by
these complete finite cycle, covariance, surjectivity or
pure-frame-conjugacy tests. This does not assert the
existence of an actual profinite gentle lift, nor does it
impose the genuine five-point pentagon.

In particular the second phase in (D23.C22.2) defines a
continuous quadratic character $\kappa_{324}$ on the actual
multiplier-one gentle group. Its arithmetic value, relative
normalization, and possible dependence on the already known
characters have not been identified here. It must therefore
not be counted as a new nonzero relative character for B.40.

\begin{proof}
The induction uses the six representatives
$1,s,r,r^2,sr,sr^2$ of $B/F$. The exact Schreier rewrites,
already used in C11, are repeated with the new target
$x\mapsto p,y\mapsto v$ in
\path{agents/characters/code/induced_s4_324.py}.
Its output contains the complete $s,r,x,y$ permutations and
rewriting data, so the claimed source marking is explicit.
The joined degree-twenty-eight free action with the marked
$(p,t)$ quotient has the same pure order, proving the
factorization asserted above.

The exhaustive replay
\path{agents/characters/code/hurwitz_324_gentle_core.py}
enumerates the derived subgroup of this joined pure image.
For each element $f$ it tests
\[
                         (sf)^2=1,\qquad(f^{-1}r)^3=1.
\]
These are exactly the two and three cycles in the $B$
extension, where conjugation by $s$ exchanges $x,y$ and
conjugation by $r$ cyclically permutes $x,y,(xy)^{-1}$.
It separately tests
\[
 f(z,x)z^{-1}f(y,z)y^{-1}f(x,y)x^{-1}=1,
 \qquad z=(xy)^{-1}.
\]
The surviving images have the counts and values in
(D23.C22.2). The implementation uses right permutation
products, as the induction requires. Converting to left
products and conjugating by $(13)$ in the first four-letter
target and by $(23)$ in the second gives exactly the displayed
left-permutation markings; neither bit changes under this
conversion.

For every surviving element, a transitive-action propagation
search exhausts all twenty-four possible images of the first
letter and constructs the unique permutation intertwining
$(s,r)$ with $(sf,f^{-1}r)$. Thus the proposed action is an
actual automorphism of this finite image. Independent
conjugacy-orbit enumeration under the pure generators gives
pure $s$- and $r$-frame classes of sizes $1152$ and $9216$;
all twelve proposed images lie in these classes.
All finite residue arrays and covariance permutations are
retained in
\path{agents/characters/records/hurwitz_324_gentle_candidates.npz}.

For necessity of this finite covariance on actual gentle
elements, the $(3,2,4)$ triangle group is $S_4$.
An actual gentle action preserves the closed normal subgroup
generated by $x^3,y^2,(xy)^4$, by the peripheral identity,
and the same holds for its inverse. Hence it acts on this
$S_4$ quotient. Its action is inner and fixes $p$.
It therefore preserves, up to pure conjugacy, the subgroup
defining the four-point permutation action. Taking its six
conjugates in $B$ gives the stated core and the induced
permutation covariance. The finite tests are thus legitimate
necessary conditions, although their sufficiency is not claimed.

Finally, write $d_v=(14)(23)$. Equation (D23.C22.2) shows
that on this target
$q\Phi_h=\operatorname{Ad}(p^{-a_h})q$, since $d_v$
centralizes $v$. In the convention
$\Phi_{hk}=\Phi_h\circ\Phi_k$, the free-group parameter is
$f_{hk}=f_h\Phi_h(f_k)$. Therefore
\[
 q(f_{hk})
 =d_v^{\kappa_h}p^{a_h}
   p^{-a_h}d_v^{\kappa_k}p^{a_k}p^{a_h}
 =d_v^{\kappa_h+\kappa_k}p^{a_h+a_k}.
\]
Uniqueness of the six-state expression proves additivity of
$\kappa_{324}$; continuity follows from the finite quotient.
No assertion about its restriction to the arithmetic subgroup
is used in this calculation.
\end{proof}


\section{The orthogonal geometry and its actual GT constraints}
% DELTA23: exact fixed-point coordinates, orthogonal group, and outer signs.
\subsection{The orthogonal geometry of the new Hurwitz block}

\paragraph{D23.C23 (an explicit minus-type quadratic geometry).}
Take block zero of the degree-$1344$ Hurwitz action, with its
$112$ labels in the order recorded in
\path{agents/five_level_six/hurwitz_block_stabilizer.npz}.
For a tuple entry $g_i=(v_i,A_i)\in
\mathbb F_3^2\rtimes\operatorname{SL}_2(\mathbb F_3)$, define
its unique fixed point by
\[
                     q_i=(I-A_i)^{-1}v_i.
\]
All these inverses exist. Normalize the ordered linear tuple by
an element $B\in\operatorname{SL}_2(\mathbb F_3)$ to
\[
 A_1=A_2=\begin{pmatrix}0&1\\2&1\end{pmatrix},\qquad
 A_3=A_4=\begin{pmatrix}1&1\\2&0\end{pmatrix}.
\tag{D23.C23.1}
\]
There are exactly two such normalizers, differing by $-I$.
Thus
\[
 [w]=[\,B(q_1-q_4),B(q_2-q_4),B(q_3-q_4)\,]
                  \in\mathbb P^5(\mathbb F_3)
\tag{D23.C23.2}
\]
is well defined. These $112$ projective vectors are distinct
and are exactly the singular points of the form $w^tQw$, where
\[
 Q=\begin{pmatrix}
 0&1&1&1&1&0\\
 1&1&0&1&0&1\\
 1&0&1&1&0&2\\
 1&1&1&0&1&1\\
 1&0&0&1&1&2\\
 0&1&2&1&2&0
 \end{pmatrix}.
\tag{D23.C23.3}
\]
The form has rank six and determinant one. Hence it has minus
type over $\mathbb F_3$ and Witt index two.

Every block-stabilizer permutation has a unique projective
linear realization on these coordinates. In particular each
of the thirteen recorded stabilizer generators lifts to a
matrix $M$ satisfying
\[
                         M^tQM=Q,\qquad\det M=1.
\tag{D23.C23.4}
\]
The complete coordinate table and all thirteen matrices are
in \path{agents/characters/records/hurwitz_orthogonal112.npz}
and the corresponding JSON record. Consequently the independently
computed block image of order $6531840$ is exactly
\[
 H=\operatorname{PSO}_6^-(3),\qquad
 H'=\operatorname{P\Omega}_6^-(3)\simeq
                  \operatorname{PSU}_4(3),
 \qquad |H'|=3265920.
\tag{D23.C23.5}
\]
In particular $H'$ is nonabelian simple. The space of
$H$-invariant symmetric bilinear forms on this six-dimensional
space has dimension one. Thus any nonzero symmetric form
obtained independently from a cohomological construction and
preserved by these same generators is a scalar multiple of $Q$.
This last assertion identifies multipliers without requiring
an arbitrary scalar normalization of that second form.

\begin{proof}
The affine conjugation formula sends $q_i$ to $u+Bq_i$;
subtracting $q_4$ removes $u$. The simultaneous centralizer of
the ordered tuple (D23.C23.1) in $\operatorname{SL}_2(3)$ is
$\{I,-I\}$, verified directly from its two distinct matrices.
Searching the twenty-four matrices of $\operatorname{SL}_2(3)$
therefore gives the stated projective normalization, with no
choice left except the harmless sign. The certificate checks
the existence of precisely these two transporters for every
one of the $112$ tuples.

The program \path{agents/characters/code/hurwitz_orthogonal112.py}
performs only exact arithmetic over $\mathbb F_3$.
The $112$ equations $w^tQw=0$ in the twenty-one coefficients
of a quadratic form have a one-dimensional solution space;
its normalized matrix is (D23.C23.3). Exact row reduction gives
rank six and determinant one. It then enumerates all
$(3^6-1)/2=364$ projective points and checks that its singular
locus is exactly the recorded set of $112$ points. For an
even-dimensional form of dimension $2m$, the split case has
$(-1)^m\det Q$ square. Here $m=3$ and $-\det Q=2$ is a
nonsquare, giving minus type.

For each stabilizer permutation choose six independent
coordinate vectors. Their images determine a linear map up
to six individual signs. Fixing one sign leaves only
$2^5=32$ possibilities. Exactly one of these maps induces the
given permutation on all $112$ points. It is invertible and
satisfies (D23.C23.4). The same procedure for the identity
permutation confirms faithfulness modulo scalar matrices.
Finally the simultaneous equations $M^tQ'M=Q'$ on the
twenty-one symmetric coefficients have dimension one, proving
the invariant-form assertion independently of the singular
locus computation.

The standard orthogonal order formula gives
\[
 |O_6^-(3)|=2\cdot3^6(3^3+1)(3^2-1)(3^4-1)=26127360.
\]
Taking determinant one and then dividing by $\{I,-I\}$
gives $|\operatorname{PSO}_6^-(3)|=6531840$.
The computed inclusion and equality of orders prove the
first group identification. The spinor computation in C24
is onto on this group and trivial on its derived subgroup;
its projective kernel has order $3265920$, the independently
computed order of $H'$. Hence $H'=\operatorname{P\Omega}_6^-(3)$.
The exceptional isomorphism with $\operatorname{PSU}_4(3)$
and simplicity are the standard finite classical-group facts
for this exact parameter, not an inference from order alone.
For explicit primary reference conventions, the GAP Character
Table Library entry \texttt{U4(3)} lists both
\texttt{POmega(-1,6,3)} and \texttt{PSU(4,3)} as its group
constructions; the ATLAS entry gives its order and the
$112$-point action.
\footnote{Primary pages checked in this run:
\url{https://www.math.rwth-aachen.de/~Thomas.Breuer/ctbllib/ctbltoc/data/U4(3).html}
and \url{https://brauer.maths.qmul.ac.uk/Atlas/clas/U43/}.
Cached pages and exact locators are in
\path{agents/characters/records/orthogonal_primary_sources.txt}.}
\end{proof}

\paragraph{D23.C24 (the full normalizer and its two visible outer signs).}
Write $S=H'$. In this statement
\[
 GO^*(Q)=\{M\in GL_6(3):M^tQM=\mu(M)Q,
                                 \ \mu(M)\in\mathbb F_3^\times\}
\]
includes nonsquare similitudes, and $PGO^*(Q)$ is its quotient
by the two scalar matrices. Then
\[
 N_{\mathfrak S_{112}}(H)=N_{\mathfrak S_{112}}(S)
       =PGO^*(Q),\qquad |PGO^*(Q)|=26127360,
\tag{D23.C24.1}
\]
and the maps determinant and multiplier give
\[
 PGO^*(Q)/H\simeq C_2\times C_2.
\tag{D23.C24.2}
\]
Both are well defined projectively, since the dimension is
even and both scalar squares are one. Furthermore
\[
                   PGO^*(Q)/S\simeq D_8,
\]
and its central subgroup of order two is exactly $H/S$.

The spinor norms of the thirteen stabilizer matrices, in
their recorded order, are
\[
                       1,1,1,2,2,2,2,2,2,2,2,2,2.
\tag{D23.C24.3}
\]
Here $2=-1$ is the nonsquare class. The norm of $-I$ is one,
so this character descends to $H$ and has kernel $S$.
Explicit matrices and permutations for the two other outer
generators are in
\path{agents/characters/records/hurwitz_orthogonal_signs.npz}.
One is a reflection with $(\det,\mu)=(2,1)$; the other has
$(\det,\mu)=(1,2)$. Their commutator has determinant and
multiplier one but spinor norm two.

\begin{proof}
The exact point-stabilizer orbit sizes are $1,30,81$ for
both $H$ and $S$. The orbit of length thirty is precisely
orthogonality for (D23.C23.3). A permutation normalizing either
group must preserve that uniquely distinguished orbital graph.
The full automorphism group of this graph is $PGO^*(Q)$;
this is the standard $O_6^-(3)$ graph with parameters
$(112,30,2,10)$. Brouwer's primary graph record identifies
its full automorphism group, including the nonsquare
similitudes, as $U_4(3).D_8$ of order $26127360$.
\footnote{A. E. Brouwer, \emph{The $O_6^-(3)$ graph},
\url{https://aeb.win.tue.nl/graphs/U4_3.html}, Group and
Construction paragraphs, checked in this run. Its star
notation is retained here to avoid confusing isometries with
all similitudes.}
Conversely every similitude normalizes $SO(Q)$ and its
spinor kernel, proving both equalities in (D23.C24.1).

For the exact spinor certificate, reflect in a nonsingular
vector $v$ by
\[
 R_v=I-\frac{2}{v^tQv}\,v(v^tQ).
\]
An orthogonal Gram--Schmidt basis reduces each isometry to
the identity by these reflections. If $Mv-v$ is nonsingular,
reflect in it; otherwise reflect successively in $Mv+v$
and $v$. The latter two vectors are nonsingular because
$v$ is and $Mv-v$ is singular. Multiplying their quadratic
norms gives the spinor norm. The exact replay
\path{agents/characters/code/hurwitz_orthogonal_signs.py}
checks every reflection and the resulting identity matrix,
the determinant parity, (D23.C24.3), and the norm of $-I$.

The same certificate constructs a reflection $R$ and a
nonsquare similitude $M$ with the determinant and multiplier
values stated above. Their squares belong to $S$, while
their commutator belongs to $H\setminus S$.
These two independent outer cosets give (D23.C24.2).
For any similitude $C$ of multiplier $\mu$,
$CR_vC^{-1}=R_{Cv}$ and
$Q(Cv)=\mu Q(v)$. An element of $SO(Q)$ is a product of an
even number of reflections; hence conjugation by $C$ preserves
its spinor norm. It follows that $H/S$ is central in the
outer quotient. The nontrivial commutator of $R,M$, together
with their involutory images, identifies this order-eight
quotient as $D_8$.
\end{proof}

\paragraph{D23.C24a (one common outer type across the twelve blocks).}
Suppose the independently computed full Hurwitz group has
socle $S^{12}$, with its twelve factors supported on the
twelve blocks, and its local transition maps lie in $H$
after the chosen group-transversal identifications.
Any permutation normalizing this full group has local maps
in $PGO^*(Q)$, and all twelve local maps have the same
determinant/multiplier pair in $PGO^*(Q)/H$.
Indeed it permutes the simple factors and their blocks.
Conjugating a group transition from block $i$ to block $j$
shows that the two local maps have equal cosets modulo $H$.
The group is transitive on the blocks, so these cosets are
constant. This is a necessary condition on the actual GT
conjugator supplied by the Hurwitz functoriality theorem.
The identification of this common pair with the original
harmonic and cubic frame bits is a further calculation;
it is not asserted by the orthogonal group identification alone.

\subsection{The complete kernel of the twelve-block Hurwitz action}

\begin{statement}{D23.F30}{Proved}{Twelve orthogonal simple factors and the exact binary quotient}
\end{statement}
Let $\Gamma=\langle A,B,C\rangle\leq\mathfrak S_{1344}$ be
the literal Hurwitz image of D23.P17. Let $Q_{12}$ be its action
on the twelve blocks recorded in
\path{agents/pentagon/hurwitz_cubic216_four_13_block0.npz},
and let $K$ be the kernel of that action. Use group-transversal
identifications of all blocks with block zero, as specified below.
The block-stabilizer image on that block is
\[
 L=\operatorname{PSO}_6^-(3),\qquad |L|=6531840,
 \qquad S=L'=\operatorname{P\Omega}_6^-(3)
              \simeq\operatorname{PSU}_4(3),
 \qquad |S|=3265920.
\]
Here the exact orthogonal identification and standard simplicity
input are D23.C23. Then
\[
 K'=S^{12},\qquad
 K/S^{12}=\mathcal C\subset\mathbb F_2^{12},\qquad
 \dim\mathcal C=9,
\]
where the twelve factors of $S^{12}$ act independently on the
twelve blocks. In the literal block labeling the four triples are
\[
 \mathcal T=\bigl\{
 \{0,9,11\},\{1,6,7\},\{2,4,10\},\{3,5,8\}\bigr\},
\]
and the code is exactly
\[
 \mathcal C=\left\{c\in\mathbb F_2^{12}:
           \sum_{i\in T}c_i\text{ is independent of }T\in\mathcal T
                              \right\}.
\]
More precisely, $K$ is the full preimage of this code in $L^{12}$.
The quotient $\Gamma/S^{12}$ has a faithful permutation
realization on $24$ points and order
\[
                         2^9\cdot216=110592.
\]
The socle and the eventual derived subgroup of $\Gamma$ are
both $S^{12}$. Thus the displayed order-$110592$ group is its
largest solvable quotient.

\begin{proof}
Write $q_A,q_B,q_C\in\mathfrak S_{12}$ for the three block
permutations and put $I_0=\{v:\operatorname{block}(v)=0\}$.
A breadth-first traversal on the twelve blocks chooses
$T_j\in\Gamma$ with $T_j(I_0)=I_j$ and $T_0=1$. If a
generator $g$ carries $i$ to $j$, the exact transition
\[
                    T_j^{-1}gT_i
\]
stabilizes $I_0$. The Schreier lemma therefore says that these
transitions generate the entire block stabilizer, before and
after restriction to $I_0$. The certificate retains the literal
braid words and all the $1344$-point permutations. Their
restrictions give thirteen distinct nonidentity generators,
which reduce to four generators without changing the group.
Exact Schreier--Sims gives order $6531840$, derived order
$3265920$, and a perfect derived group. The stabilizer of one
of the $112$ points has orbit sizes $1,30,81$. D23.C23
identifies these same generators as the full group $L$ above;
that identification is not inferred merely from their order.

Similarly enumerate all $216$ elements of $Q_{12}$ and choose
word lifts $T_q\in\Gamma$ of them. For each generator $g$ and
each $q\in Q_{12}$, the word
\[
                       T_{q_gq}^{-1}gT_q
\]
lies in $K$. These words generate $K$ by the Schreier lemma.
There are $153$ distinct nonidentity full permutations in this
Schreier list. Distinct full permutations are retained even when
their block-zero restrictions agree or are the identity.
Their block-zero image has order $6531840$, the same as $L$, so the
projection of $K$ onto that block is all of $L$. Conjugation
by the $T_i$ proves this for every block. Since $K$ fixes all
blocks, its action identifies it with a subgroup of $L^{12}$.
The action is faithful because the blocks cover all $1344$
points.

We next check the correlations, not just the separate projections.
The stabilizer of zero in $Q_{12}$ has precisely the orbits
\[
 \{0\},\quad\{9,11\},\quad\{1,2,3,4,5,6,7,8,10\}.
\]
There are therefore only two ordered-pair types for distinct
blocks. Normalize the restriction of each of the above
Schreier generators to block $i$ by conjugating with $T_i$.
The joint restrictions to $(0,9)$ and to $(0,1)$ both generate
groups of order
\[
          42664933785600=6531840^2.
\]
Each lies in $L\times L$ and projects onto both factors;
equality of orders proves that each is the full product.
By the two orbit types, the projection of $K$ onto any two
distinct blocks is therefore $L\times L$. Its derived
projection is $S\times S$, of exact order
$10666233446400=3265920^2$.

For clarity, the simple-product inference used next needs
nonabelian simplicity. If a subgroup $J\leq S^m$ projects
onto every factor and every pair, where $S$ is nonabelian
simple, then $J=S^m$. Prove this by induction. Having obtained
the first $m-1$ factors, Goursat's lemma between $S^{m-1}$
and the last factor either gives the full product or identifies
a simple quotient of $S^{m-1}$ with the last $S$. Every normal
subgroup of $S^{m-1}$ is a product of factors: commuting an
element of such a normal subgroup with one factor isolates
that coordinate, and simplicity and trivial center force
each nontrivial projected factor to be contained in the
normal subgroup. Thus a simple quotient selects one factor.
The resulting diagonal relation contradicts the assumed full
projection onto that pair. This proves the induction.
Apply the lemma to $K'\leq S^{12}$. Its single and pair
projections are $S$ and $S\times S$, respectively, so
$K'=S^{12}$.

The coordinate characters $L\to L/S=C_2$ consequently
identify $K/S^{12}$ with a binary linear code. On the $153$
Schreier generators the exact computation tests membership
in the derived subgroup of $L$ in each coordinate. Gaussian
elimination on the resulting $153$ by $12$ matrix gives rank
nine. Every row has the same parity on the four triples
displayed in the statement. Those triples are exactly the
images of $\{0,9,11\}$ under $Q_{12}$. The equality of their
four parities imposes three independent equations, hence
defines a nine-dimensional subspace. The inclusion and the
rank computation prove the displayed equality for $\mathcal C$.
As $K$ already contains $S^{12}$, it is the full preimage of
that code in $L^{12}$.

Here is the faithful quotient explicitly. For $g\in\Gamma$
and a block $i$, put $j=q_g(i)$ and let
\[
  e_g(i)\in\mathbb F_2
       \quad\text{be the image of }T_j^{-1}gT_i\text{ in }L/S.
\]
On the $24$ pairs $(i,e)\in\{0,\ldots,11\}\times\mathbb F_2$
define
\[
                    g(i,e)=(q_g(i),e+e_g(i)).
\]
The transition identity proves multiplicativity. Its kernel
consists precisely of those elements of $K$ whose twelve
coordinates belong to $S$, namely $S^{12}$. Its order is
$|\mathcal C|\,|Q_{12}|=2^9\cdot216$, also independently
checked directly from the three resulting permutations.
The block quotient has derived orders
$216,72,18,9,1$, so this quotient is solvable. Since $S^{12}$
is perfect, it belongs to every derived subgroup of $\Gamma$;
solvability of the quotient proves the eventual-derived claim.

Finally $S^{12}$ is a minimal normal subgroup of $\Gamma$:
normal subgroups of this direct product are products of its
simple factors, and $\Gamma$ acts transitively on them.
Its centralizer in the full permutation group is trivial.
Indeed a commuting permutation must preserve each factor's
support, and the centralizer of $S$ on one block is trivial:
a permutation commuting with a transitive group is determined
by the image of a point, which must be fixed by its point
stabilizer, whose orbit sizes $1,30,81$ leave just that point.
Any different minimal normal subgroup of $\Gamma$ would
centralize $S^{12}$, which is impossible. This proves the
socle assertion.
\end{proof}

\paragraph{Certificates and inputs.}
The complete Schreier construction and exact order checks are
\path{agents/five_level_six/hurwitz_block_stabilizer.py}
and its JSON/NPZ records. The joint restrictions, all local
transports, binary Gaussian elimination, signed quotient,
and the directly verified orbital graph parameters
$(112,30,2,10)$ are in
\path{agents/five_level_six/hurwitz_block_correlations.py}
and its JSON/NPZ records. The classical-group identification
and simplicity are explicit inputs from D23.C23. No abstract
group name or simple-product conclusion is obtained from
coinciding orders alone.

\subsection{The solvable block quotient is the earlier Burau quotient}

\begin{statement}{D23.F31}{Proved}{An exact marked Burau identification}
\end{statement}
The order-$110592$ signed permutation quotient in D23.F30 is
marked-isomorphic to the group generated by the literal
matrices of supplied FP.36.2 over
$R=(\mathbb Z/4)[\omega]/(\omega^2+\omega+1)$, modulo
$\langle\omega I\rangle$:
\[
 A_B=\begin{pmatrix}1-\omega&\omega&0\\1&0&0\\0&0&1\end{pmatrix},
 \quad B_B=\begin{pmatrix}1&0&0\\0&1-\omega&\omega\\0&1&0\end{pmatrix},
 \quad C_B=\begin{pmatrix}1&0&-1\\0&1&-1\\0&0&-\omega\end{pmatrix}.
\]
The isomorphism sends the positive Hurwitz generators
$A,B,C$ to exactly $A_B,B_B,C_B$. The simultaneous flip of
all twelve binary pairs maps to $-I$. The projective quotient
therefore has order $55296$ and is the already supplied
Burau quotient, with its original markings.

\begin{proof}
Encode $a+b\omega\in R$ by $a+4b$ with $0\leq a,b<4$.
The certificate uses the exact multiplication
\[
 (a+b\omega)(c+d\omega)
     =(ac-bd)+(ad+bc-bd)\omega
\]
modulo four. It constructs the three matrices displayed
above and identifies a matrix with its three scalar multiples
by $1,\omega,\omega^2$, choosing a canonical representative.

Starting with the identity, simultaneously traverse the signed
permutation group and these matrix classes using the three
marked positive generators. Whenever a permutation already
occurred, check that its newly computed matrix class is
exactly the stored one. This check is made for every one of
the $3\cdot110592$ directed generator edges. All checks pass.
There are $110592$ permutations and also $110592$ distinct
matrix classes. Thus the marked homomorphism is well defined
and injective as well as surjective. In particular this is an
explicit marked isomorphism, not merely a match of orders or
derived series. The code also finds the global binary flip
and checks that its matrix class is precisely $-I$.

The determinants of all three matrix generators are $-\omega$.
Their determinant values therefore lie in $\mu_6\subset R^\times$,
and this determinant character is unchanged by multiplying
the matrix by $\omega I$. The character
\[
                 \epsilon(g)=\det(g)^3\in\{1,-1\}
\]
takes $-I$ to $-1$. Hence the central scalar subgroup
$\langle-I\rangle$ splits off, and projection to the
projective group restricts to an isomorphism on
$\ker\epsilon$. Its order is $110592/2=55296$.
This proves the projective identification with the exact
matrices and central relation of FP.36.
\end{proof}

\paragraph{Consequences and exact limitation.}
The signed quotient has first-pair image of order $192$.
Adjoining the marked cubic detector of D23.P10 raises the
joint image order to $3456$, so that detector is not a marked
quotient of this first-pair image. Its mixed image has order
$144$, and adjoining the marked $S_4$ generators
$(123),(1234)$ leaves the order $144$, proving retention of
that harmonic word test. These are checked in the same
correlation certificate.

For every commutator word $f$ evaluated at any two elements of
the signed quotient, $\epsilon(f)=1$. Thus its projective
value determines its entire signed value; the extra scalar
$-I$ supplies no new commutator-word constraint. In particular
the signed quotient alone adds no new identity comparing
the cubic frame with the harmonic or half-Kummer character.
This does not assert a general failure of such an identity on
actual gentle elements. The new possible information in the
$1344$-point construction lies in its twelve orthogonal
factors and their actual full-GT normalizer, not in this
already known solvable quotient.

The exhaustive marked comparison is
\path{agents/five_level_six/hurwitz_signed_burau.py}; its
JSON records all counts and the scalar-flip check. The proof
uses the matrices and markings supplied in FP.36, and gives
a new marked identification of the Hurwitz quotient with them.

\subsection{An orientation constraint on the orthogonal Hurwitz blocks}

\begin{statement}{D23.P19}{Proved}{Actual GT acts by isometries on a fixed quadratic block}
\end{statement}
Use the degree-1344 Hurwitz orbit of D23.P17--P18. Its
degree-112 block described in D23.C23 has coordinates
\[
 w=(q_1-q_4,q_2-q_4,q_3-q_4)\in\mathbb F_3^6/\{\pm1\},
 \qquad q_i=(I-A_i)^{-1}v_i,
\]
where the four affine elements are $(v_i,A_i)$ and their
linear parts have been simultaneously transported in
$SL_2(\mathbb F_3)$ to
\[
 A_1=A_2=\begin{pmatrix}0&1\\2&1\end{pmatrix},\qquad
 A_3=A_4=\begin{pmatrix}1&1\\2&0\end{pmatrix}.
 \tag{D23.P19.1}
\]
The canonical symmetric relative cup-product form in these
coordinates is exactly
\[
 Q=\begin{pmatrix}
 0&1&1&1&1&0\\
 1&1&0&1&0&1\\
 1&0&1&1&0&2\\
 1&1&1&0&1&1\\
 1&0&0&1&1&2\\
 0&1&2&1&2&0
 \end{pmatrix}.                                            \tag{D23.P19.2}
\]
In particular, it is the same form as the independently
reconstructed D23.C23 form, with no rescaling.

Let $h\in\widehat{GT}$ have multiplier one, and suppose that
the actual permutation $V_h$ of D23.P18 preserves this block.
Then its restriction is induced by a linear map $T_h$ with
\[
                    T_h^tQT_h=Q.                          \tag{D23.P19.3}
\]
The two possible choices of $T_h$ differ by $-I$ and give the
same projective action. Thus the actual conjugator belongs
to the projective isometry group on each fixed block. The
larger projective similitude group is not an equally valid
substitute for this necessary condition.

This statement does not identify the determinant or the
spinor norm of $T_h$, and does not by itself decide
$\Delta(h)$. It does not assume compatibility with deletion
of a coordinate of an affine tuple.

\begin{proof}
Put $k=\mathbb F_3$ and $M=k^2$, with alternating form
\[
       \langle u,v\rangle=u^tJv,
       \qquad J=\begin{pmatrix}0&1\\2&0\end{pmatrix}.
\]
Let $\Pi$ be the free group on $t_1,\ldots,t_4$, and put
$t_5=(t_1t_2t_3t_4)^{-1}$. Write $\widehat\Pi$ for its full
profinite completion. The linear representation
$\rho:\widehat\Pi\to SL_2(k)$ has
$\rho(t_i)=A_i$ for $1\leq i\leq4$ and
\[
 A_5=(A_1A_2A_3A_4)^{-1}.
\]
All five $I-A_i$ are invertible. For example,
$A_1A_2A_3A_4=\left(\begin{smallmatrix}2&1\\0&2\end{smallmatrix}\right)$,
so this holds at the fifth puncture as well.

\paragraph{Affine lifts are cohomology classes.}
An affine lift of $\rho$ is a map
$g\mapsto(c(g),\rho(g))$ into $M\rtimes SL_2(k)$, where
\[
                 c(gh)=c(g)+\rho(g)c(h).                  \tag{D23.P19.4}
\]
Thus it is a left 1-cocycle. Its four generator values are
arbitrary; the universal property of the free profinite group
extends any such assignment continuously, since the target
semidirect product is finite. Conjugating the affine lift by
a translation changes all its fixed points $q_i$ by the same
vector. Consequently the quotient by translations is
$H^1(\widehat\Pi,M)$, of dimension $8-2=6$, and setting
$q_4=0$ gives precisely the difference coordinates in the
statement. The simultaneous centralizer of (D23.P19.1) in
$SL_2(k)$ is $\{\pm I\}$. Quotienting further by this
centralizer gives the projective nonzero classes and the zero
class. D23.C23 identifies the 112 points of the indicated
Hurwitz block with the singular nonzero projective classes.

For clarity, this description concerns the fiber of the
linear-part representation class. In D23.C23 every point in
the indicated block has that linear-part class. Its other
translates are the other eleven blocks; the ordinary Hurwitz
action commutes with projection to linear parts. Thus fixing
one of these blocks fixes the corresponding linear
representation class up to $SL_2(k)$ conjugacy.

\paragraph{Relative cohomology and its orientation.}
Use the five peripheral subgroups
$C_i=\overline{\langle t_i\rangle}\cong\widehat{\mathbb Z}$
and the relative continuous cochain complex, the shifted
mapping cone of
\[
 C^\bullet(\widehat\Pi,M)\longrightarrow
              \bigoplus_{i=1}^5 C^\bullet(C_i,M).
\]
Because $I-A_i$ is invertible, both
$H^0(C_i,M)$ and $H^1(C_i,M)$ vanish. Indeed these groups are
respectively $\ker(A_i-I)$ and $M/(A_i-I)M$. The long exact
sequence therefore identifies relative $H^1$ with ordinary
$H^1(\widehat\Pi,M)$. This is the usual parabolic situation,
with no local invariant vectors.

For trivial coefficients $k$, the same sequence gives
\[
 H^2(\widehat\Pi,\{C_i\};k)
   =\operatorname{coker}\left(
       k^4\longrightarrow k^5\right)\cong k,
 \quad
 (a_1,\ldots,a_4)\longmapsto
       (a_1,\ldots,a_4,-\textstyle\sum_i a_i).              \tag{D23.P19.5}
\]
Here $H^2(\widehat\Pi,k)=0$: a continuous finite-valued
2-cocycle factors through a finite quotient, and the
corresponding finite extension splits after pullback to the
free profinite group, by its projectivity. The quotient in
(D23.P19.5) is oriented by the sum of the five coordinates.

Suppose an automorphism $\alpha$ of $\widehat\Pi$ carries
each $t_i$ to a conjugate of $t_i^\lambda$, for one
$\lambda\in\widehat{\mathbb Z}^{\times}$. It acts by
$\lambda\bmod3$ on each of the five trivial-coefficient
peripheral $H^1$ groups, since conjugation acts trivially on
them. Consequently its action on (D23.P19.5) is also
$\lambda\bmod3$. This proves the orientation character
directly, without realizing $\alpha$ by a homeomorphism.
The cochain maps for conjugated peripheral subgroups are
defined using their usual conjugation identifications; these
are compatible with restriction and cup product. Changing
one of the conjugating elements gives the corresponding
inner cochain homotopy and does not change the cohomology
map or this orientation character.

Relative cup product, followed by the alternating coefficient
pairing $J$, therefore gives a symmetric bilinear pairing
\[
 H^1(\widehat\Pi,M)\times H^1(\widehat\Pi,M)
       \longrightarrow H^2(\widehat\Pi,\{C_i\};k)\cong k.
 \tag{D23.P19.6}
\]
The symmetry is the product of the degree-one cup sign and
the alternating coefficient sign. Its naturality shows that
pullback by $\alpha^{-1}$ multiplies it by
$\lambda^{-1}\bmod3$, when the coefficient forms are
identified isometrically.

There is no unproved discrete-to-profinite comparison in
this calculation. Restriction to the dense discrete free
group identifies $H^0$ and $H^1$ for these finite modules:
the same arbitrary generator-value description proves the
claim in degree one. It identifies the corresponding
peripheral groups in degrees zero and one. The discrete
free group also has $H^2=0$, so (D23.P19.5) is carried
isomorphically to its discrete counterpart. Restriction of
cochains respects cup product. These observations suffice
for all the comparison maps used in (D23.P19.6); no
comparison in higher degrees is required.

\paragraph{Exact cup-product matrix.}
Here is an explicit bar-chain computation, fixing the sign
and the coordinate identification. For two cocycles $c,d$
put
\[
 c_i=c(t_i),\quad d_i=d(t_i),\quad
 c_i=(I-A_i)q_i,\quad
 p_i=t_1\cdots t_i,\quad P_i=A_1\cdots A_i.
\]
The fifth value is determined by $c(p_5)=0$. In normalized
bar chains the chain
\[
                  Z=\sum_{i=1}^5[p_{i-1}\mid t_i]
\]
has boundary $\sum_i[t_i]$. Thus
$(Z,-\sum_i[t_i])$ is a relative cycle. A relative lift of
$c$ has boundary primitives $-q_i$, because
$\delta(-q_i)(t_i)=(I-A_i)q_i=c_i$. The ordinary cup
cochain is
\[
               (g,h)\longmapsto
                 \langle c(g),\rho(g)d(h)\rangle.
\]
Its relative boundary component is
$\langle-q_i,d_i\rangle$. Evaluation on the displayed
relative cycle consequently gives
\[
 B(c,d)=\sum_{i=1}^5
       \left(\langle c(p_{i-1}),P_{i-1}d_i\rangle
                       +\langle q_i,d_i\rangle\right).
 \tag{D23.P19.7}
\]
This also independently fixes the orientation convention in
the matrix calculation. Evaluating (D23.P19.7) on the eight
fixed-point coordinate vectors $(q_1,q_2,q_3,q_4)$ gives a
symmetric matrix of rank six whose radical is exactly the
two-dimensional diagonal translation subspace. On the
section $q_4=0$ its matrix is exactly (D23.P19.2).
The certificate evaluates all 36 entries, checks the full
eight-dimensional radical, and compares the resulting
matrix literally with D23.C23. In particular the form is
nondegenerate on $H^1$, and no appeal to invariant-form
uniqueness is needed for its identification.

\paragraph{Application to actual GT.}
By D23.P18, the actual permutation is induced by
precomposition with $\kappa_h^{-1}$, where $\kappa_h$ is the
restriction of the Drinfeld automorphism to the free kernel
of the rightmost forgetting map. That proof gives the
required peripheral powers for $t_1,\ldots,t_4$. It also
gives the fifth: the boundary product is
$t_1t_2t_3t_4=z_4^{-1}z_5$, and the standard action sends
each full twist $z_j$ to $z_j^\lambda$. The two nested full
twists commute, so the boundary product, and hence its
inverse $t_5$, has the same power $\lambda$.

If the block is fixed, choose $B\in SL_2(k)$ with
\[
 B\rho(\kappa_h^{-1}(g))B^{-1}=\rho(g)
                  \qquad(g\in\widehat\Pi).
\]
The induced map on affine lifts is
\[
                     c(g)\longmapsto
                       Bc(\kappa_h^{-1}(g)).              \tag{D23.P19.8}
\]
It is linear on cocycles, preserves coboundaries, and is
bijective. Since $B\in SL_2(k)$, it preserves $J$.
The naturality calculation therefore gives
\[
                B(T_hc,T_hd)=\lambda^{-1}B(c,d).
\]
For multiplier one this is (D23.P19.3). The only ambiguity
in $B$ is the simultaneous centralizer $\{\pm I\}$, which
changes $T_h$ by $-I$ and has no effect on its projective
action or on its multiplier. This proves the statement.
\end{proof}

\paragraph{Source and check scope.}
The parabolic cup pairing and the reversal of symmetry type
are standard; a primary reference is Dettweiler--Wewers,
\emph{Variation of parabolic cohomology and Poincar\'e
duality}, arXiv:math/0411119v1, Corollary 2.2 and \S2.3,
printed p.10. Their Theorem 2.5, printed p.12, supplies an
explicit formula in their right-cocycle convention:
\url{https://arxiv.org/pdf/math/0411119}.
Here the needed orientation character, the comparison in the
free profinite case, and formula (D23.P19.7) in the actual
left-cocycle marking were proved directly above.
The exact full-twist power normalization is also stated
explicitly in Lochak--Nakamura--Schneps,
arXiv:2602.12462v1, \S2, Theorem 2.1, printed pp.2--3:
$F(\omega_i)=\omega_i^\lambda$ with
$\omega_i=(\sigma_1\cdots\sigma_{i-1})^i$.
\url{https://arxiv.org/pdf/2602.12462}.

The exact certificate is
\path{hurwitz_cup_pairing.py}; its record is
\path{hurwitz_cup_pairing.json}. It identifies the matrix and
checks all thirteen supplied block-stabilizer matrices
against it. It does not purport to enumerate, approximate,
or construct an unknown full-GT element.

\subsection{Three exact local lifts of each normalized cubic state}

\paragraph{D23.H3 (canonical detector and its local lifts).}
Let $A,B,C\in\mathfrak S_{1344}$ be D23.P17's literal braid
generators, and let $V$ be the actual full-GT conjugator of
D23.P18. Suppose $\rho=0\pmod3$. On each of the four
orthogonal blocks numbered $7,8,10,11$, the restriction of
$V$ is one of exactly three explicitly recorded projective
linear maps, for either value $a_3\in\{0,1\}$.
All these maps have determinant one, quadratic multiplier
one, and spinor norm one. The candidates for the two values
of $a_3$ need not be the same maps.

The degree-$36$ quotient used to impose this constraint is
intrinsic to the finite braid action. In particular the
argument does not assume that deletion of a tuple coordinate
commutes with the full-GT action on tuples.

\begin{proof}
The group $G_3=\langle A,B\rangle$ has the following orbit
degrees, counted with multiplicity:
\[
                 4,8,8,12,80,108,160,160,240,240,324.
\]
The detector subset of D23.P17 is exactly the union of the
unique orbits of degrees $108$ and $324$. The Drinfeld action
preserves the three-strand subgroup and is an automorphism
there. Thus $V$ normalizes $G_3$ and preserves both of these
orbits separately.

On the degree-$108$ orbit the permutation image has order
$11664$ and exactly one minimal nontrivial block system.
It has $36$ blocks of size three and agrees with the
restriction of D23.P17's detector map. On the degree-$324$
orbit the image again has order $11664$. There are exactly
four minimal nontrivial block systems, each with $108$ blocks
of size three. The join of these four equivalence relations
has $36$ classes of size nine, and is again exactly the
detector map. A normalizer permutes the minimal block
systems, so it preserves their join. Both degree-$36$
quotients are therefore preserved intrinsically.

Their labelled braid generators are the same literal
degree-$36$ generators in D23.H1. The induced permutations
must satisfy that finding's simultaneous conjugacy equations.
Its uniqueness assertion identifies both induced permutations
with the same recorded $\theta_{0,a_3}$. Hence, writing $\psi$
for the detector map on its $432$-point domain, we have
\[
                      \psi(Vv)=\theta_{0,a_3}\psi(v).
\tag{D23.H3.1}
\]
This uses only a canonical finite quotient and the actual
braid covariance, not a topological forgetting assertion.

The twelve blocks are permuted by $V$, by D23.F30's socle
description. Alternatively the unique minimal block system
of the full transitive action gives this permutation directly.
The unique degree-$108$ orbit meets precisely blocks
$7,8,10,11$, in $27$ points apiece. Their detector images are
four disjoint sets of nine labels. Both
$\theta_{0,0}$ and $\theta_{0,1}$ preserve each of these four
sets. Equation (D23.H3.1) consequently fixes those four
blocks individually. The normalizer assertion D23.C24
places each local map in the projective similitude group of
the recorded quadratic form. The same conclusion also
follows from the cohomological construction D23.P19, which
already restricts the multiplier to one.

Use the group-transversal identifications of each block with
block zero. In each of these four blocks the full detector
domain has $108$ points; every one of the $36$ labels has
exactly three preimages. Choose the first six linearly
independent projective coordinate vectors in that domain.
They form a basis of the six-dimensional vector space.
By (D23.H3.1), each basis vector has three possible projective
images. Fixing the scalar of the first image leaves two
choices of scalar for each of the other five. This exhausts
all possible projective linear maps satisfying (D23.H3.1)
on the chosen basis: there are at most $3^6 2^5$ candidates.

For each multiplier $\mu\in\{1,2\}$, the certificate first
imposes the pairwise Gram equations on the basis images.
Their diagonal equations vanish because all six basis
vectors and all allowed images are singular. Exact rank
checks discard singular maps. Exactly six projective
similitudes remain at this stage in each of the eight cases
$(a_3,j)$, $a_3=0,1$, $j=7,8,10,11$. Testing (D23.H3.1)
on all $108$ points leaves exactly three in every case.
Every surviving map has $(\det,\mu)=(1,1)$.
The independent reflection certificate computes spinor
norm one for each. No possible projective similitude was
omitted: its six images and their relative scalars occur
in precisely this enumeration.

The complete permutation and matrix lists, together with
the canonical block-system checks, are in
\path{records/hurwitz_h1_orthogonal_lifts.npz} and the JSON
record of the same name. The executable certificate is
\path{code/hurwitz_h1_orthogonal_lifts.py}. Independent
checks of the fixed blocks and all spinor norms are in
\path{agents/characters/records/hurwitz_h3_independent_spinor.json}.
This establishes necessary local candidates, not the
existence of an actual GT element for any of them.
\end{proof}

\subsection{The harmonic braid equation fixes the eleventh local map}

\paragraph{D23.H4 (a unique local involution in the mixed state).}
Retain the actual full-GT hypotheses and block numbering of
D23.H3. If $\rho=0\pmod3$, the restriction of the actual
conjugator $V$ to block $11$ is uniquely determined by $a_3$.
It is the identity for $a_3=0$, and the explicit involution
in \path{records/hurwitz_block11_harmonic.npz} for $a_3=1$.
This conclusion imposes no condition on the harmonic bit
$a_2$ by itself.

\begin{proof}
Put $S=ABA$, $X=C^2$, and $Y=S^2$ in the literal Hurwitz
action. Both $X,Y$, as well as $C$, preserve block $11$.
Their local image $J=\langle X_{11},Y_{11}\rangle$ has
order $72$ and derived subgroup of order eight. Join this
marked pair to the eight-point nonzero-vector action of
\[
 U^2=\begin{pmatrix}1&2\\0&1\end{pmatrix},\qquad
 L^2=\begin{pmatrix}1&0\\1&1\end{pmatrix}
          \quad\text{over }\mathbb F_3.
\]
The joined group still has order $72$, and its derived
subgroup still has order eight. Thus the projection from
$J$ to the specified $SL_2(3)$ is well defined and onto;
it is injective on $J'$, since $SL_2(3)'$ also has order
eight.

The word $f$ has zero exponent sums, so its value
$R_{11}=f(X_{11},Y_{11})$ lies in $J'$. Its image under
the marked projection is $f(U^2,L^2)$. By the linear
part of D23.P10's exact actual-gentle formula, this is
$M_\rho$, and hence is the identity when $\rho=0$.
For precision P10 uses the signed generators $-U^2,-L^2$;
replacing both by their negatives leaves the value of a
zero-exponent-sum word unchanged. Injectivity on $J'$ now
gives $R_{11}=1$.

The third braid covariance equation of D23.P18 is
$VCV^{-1}=RCR^{-1}$. On block $11$ it therefore says
that $V_{11}$ commutes with $C_{11}$. Among the three
candidates of D23.H3 exactly one commutes with $C_{11}$
for each of the two values of $a_3$. In the fixed recorded
ordering these are candidate zero for $a_3=0$ and candidate
two for $a_3=1$. Their orders are respectively one and two.
All group orders, the marked graph assertion, and all six
commutation tests are checked exactly by
\path{code/hurwitz_block11_harmonic.py}.
\end{proof}

\subsection{Global coupling of the orthogonal blocks}

\begin{statement}{D23.P20}{Proved}{Innerness of the normalized Hurwitz action}
\end{statement}
Let $\Gamma$ be the degree-1344 permutation image of
D23.P17, and let $N$ be the kernel of its degree-12 block
action. Use the group structure proved in D23.F30:
\[
 H=PSO^-_6(3),\quad S=P\Omega^-_6(3),\quad
 N'=S^{12},\quad N/S^{12}=\mathcal C\subset\mathbb F_2^{12}.
\]
Here $\mathcal C$ consists exactly of the binary vectors whose
sums on the four parallel triples
\[
 \{0,9,11\},\quad\{1,6,7\},\quad
 \{2,4,10\},\quad\{3,5,8\}                             \tag{D23.P20.1}
\]
are equal. Thus $\dim\mathcal C=9$.

For an actual multiplier-one full-GT element with
$\rho\equiv0\pmod3$, the permutation $V_h$ of D23.P18
belongs to $N$. In particular, its automorphism of
$\Gamma$ is inner. This conclusion applies to the normalized
mixed state $(\rho,a_3,a_2)=(0,1,0)$ obtained in D23.H2.

On the four blocks $7,8,10,11$, the three local possibilities
per block in D23.H3 must have indices
\[
                  (i,i,i,j),\qquad 0\leq i,j<3.            \tag{D23.P20.2}
\]
Thus nine assignments remain on those four blocks, for each
of the two cubic bits. This is only a necessary finite
constraint; no actual mixed full-GT element is constructed.

\begin{proof}
We first recall why the actual permutation fixes the twelve
blocks. The first-pair block quotient is the marked
$SL_2(\mathbb F_3)$ quotient and has $f$-value $M_\rho$.
The mixed block commutator quotient is the $C_3$ quotient
of the marked $A_4=S_4'$ value; the supplied FP.37 formula
$f(p,t)=t^{2a_2}p^\rho$ has trivial image in it when
$\rho=0$. These exact marked quotient identifications are
recorded in
\path{hurwitz_cubic216_block_character_projections.json}.
Thus for $\rho=0$ all three transformed braid generators
have their original degree-12 images. The actual induced
permutation centralizes that degree-12 action. Its
centralizer is trivial: a commuting permutation is determined
by the image of one point, and propagation through the
three generators rejects eleven of the twelve possible
images. The certificate replays this complete calculation.
Hence $V_h$ fixes every block individually.

Choose the twelve braid transversals used in D23.F30, so
that every transition map of an element of $\Gamma$ is an
element of $H$ in these local coordinates. Write $V_j$ for
the local permutation of $V_h$ on block $j$. D23.P19
gives its linear orthogonal realization. The exact exhaustive
lift calculation D23.H3 shows that on each of blocks
$7,8,10,11$, every possibility has determinant one and
multiplier one. Thus at least one $V_j$ belongs to $H$.

For any other block $k$, choose $g\in\Gamma$ carrying
block $j$ to $k$. Because $V_h$ normalizes $\Gamma$ by
D23.P18 and fixes all blocks, the element
$g'=V_hgV_h^{-1}$ has the same source and target blocks.
Let $L,L'\in H$ be their normalized transition maps. Then
\[
                         V_k L=L'V_j,
\]
so $V_k=L'V_jL^{-1}\in H$. This propagates membership in
$H$ to all twelve blocks. No classification of the entire
permutation normalizer of $H$ is needed for this step.

The quotient $H/S$ has order two. Let
$v=(v_0,\ldots,v_{11})\in\mathbb F_2^{12}$ be the vector
of its values on the $V_j$. Since $N$ contains $S^{12}$
and has binary image $\mathcal C$, it is exactly the
preimage of $\mathcal C$ in $H^{12}$. Normalization of
$\Gamma$ implies
\[
                       v+qv\in\mathcal C
                 \quad\hbox{for all }q\in\Gamma/N.       \tag{D23.P20.3}
\]
Indeed $V_hgV_h^{-1}g^{-1}\in N$ has precisely this binary
vector, where $q$ is the block permutation of $g$.

We now prove that (D23.P20.3) forces $v\in\mathcal C$.
Let $\pi:\mathbb F_2^{12}\to\mathbb F_2^4$ take the four
triple sums in (D23.P20.1). Then
\[
 \mathcal C=\pi^{-1}\langle(1,1,1,1)\rangle.
\]
On the four triples the three braid generators induce the
permutations
\[
      (1\,2\,3),\qquad(0\,1\,3),\qquad(1\,2\,3).
\]
They have order three and together act transitively. If the
class of $w=\pi(v)$ modulo the constant vector is fixed by
one such permutation $a$, write
$aw+w=\epsilon(1,1,1,1)$. Adding its three translates gives
$0=3\epsilon(1,1,1,1)=\epsilon(1,1,1,1)$ in
characteristic two. Hence $aw=w$. Applying this to the
three generators makes $w$ a constant vector, by
transitivity. Thus $v\in\mathcal C$, as claimed.
Since $V_h\in H^{12}$, this proves $V_h\in N$.

Finally D23.P18 gives $V_hA=AV_h$ for the first braid
generator $A$. Its block permutation cycles
$7\mapsto10\mapsto8\mapsto7$ and fixes $11$.
The exact local transitions show that this equation
identifies the three H3 candidate indices on the cycled
blocks. All three possibilities on block 11 commute with
its local $A$ transition. Checking the $3^4$ possibilities
therefore gives exactly (D23.P20.2), independently for
$a_3=0$ and $a_3=1$.
\end{proof}

\paragraph{Exact certificate and scope.}
\path{hurwitz_global_innerness.py} replays the degree-12
centralizer, the four-triple action, all 4096 binary sign
vectors, and both local $3^4$ coupling calculations. Its
record is \path{hurwitz_global_innerness.json}. The binary
enumeration agrees with the short proof: exactly the 512
code vectors normalize, and exactly 32 remain after the
four H3 signs are set to zero. This last count uses the
separately checked H3 spinor values; innerness itself uses
only their determinant-one, multiplier-one membership in
$H$.

The group structure input is
\path{five_level_six/hurwitz_block_structure.tex}, D23.F30;
the local exhaustive lift input is D23.H3 with
\path{work/code/hurwitz_h1_orthogonal_lifts.py} and its
JSON/NPZ records. The necessary actual-source interpretation
comes from D23.P18--P19. Innerness of this finite quotient
does not establish innerness of the full profinite GT
action, nor the vanishing of $\Delta$.

\subsection{The normalized actual conjugator belongs to the simple socle}

\begin{statement}{D23.F32}{Proved}{Eliminating the remaining binary block coordinates}
\end{statement}
Let $h=(1,f)$ be an actual full-GT element with
\[
                         \rho(h)=0\pmod3,\qquad a_2(h)=0.
\]
For either value of $a_3(h)$, the actual degree-$1344$
Hurwitz conjugator $V_h$ of D23.P18 belongs to
\[
                 S^{12}=\operatorname{P\Omega}_6^-(3)^{12}
                           \subset\Gamma.
\]
In particular this applies to the normalized hypothetical
mixed state $(\rho,a_3,a_2)=(0,1,0)$ of D23.H2. The statement
is a necessary condition, not a contradiction for that state.

\begin{proof}
First determine the exact signed solvable values, retaining
the central ambiguity until it is resolved. By D23.F31, the
signed $24$-point quotient is marked-isomorphic to the
matrices FP.36.2 modulo $\langle\omega I\rangle$.
Let $\overline F$ and $\overline R$ be the projective Burau
values of $f(X_1,X_2)$ and $f(X_3,X_4)$.
The supplied FP.37 uses exactly these positive half twists,
$X_1=A^2$, $X_2=B^2$, $X_3=C^2$, and
$X_4=(ABA)^2$. Its actual-full-GT consequences are
\[
 a_f=\rho\pmod3,\qquad \nu_4=a_2,
\]
and its exact finite formulas give
\[
 \overline F=F^0_{\rho,a_2}Z^c,
 \qquad
 \overline R=(e_0t_0^2)^{\nu_4}p_0^{a_f},
 \qquad c\in\mathbb F_2.
\]
Here $Z=(ABA)^4$ in the first-pair image, and the literal
table of FP.37 has $F^0_{0,0}=I$. For the present normalized
state, $a_f=\nu_4=0$, hence $\overline R=I$ and
$\overline F=Z^c$. The same source's exact pentagon formula
(FP.37.6) is
\[
                     I=\mathcal P(f)=w_0^c
\]
in this case. Its displayed $w_0$ is a nonidentity involution
(its matrix has nonzero off-diagonal entries), so $c=0$.
Thus both projective values are the identity. The determinant
cube character of D23.F31 is one on every commutator
evaluation and is minus one on the remaining central scalar
$-I$. The signed values are consequently
\[
                             F_{24}=R_{24}=I.
\]
This uses the actual pentagon in the supplied Burau quotient;
the conclusions are not imposed on arbitrary gentle residues.

It follows from D23.P18 that the induced conjugation of the
signed quotient fixes each of its three marked generators.
In particular the induced permutation of the twelve blocks
centralizes $Q_{12}$. That transitive group's point-stabilizer
orbits have sizes $1,2,9$, as checked in D23.F30. Its
permutation centralizer is therefore trivial: the image of
one point under a commuting permutation must be a fixed
point of its stabilizer, and there is exactly one. Thus
$V_h$ fixes all twelve blocks.

Use the group-transversal identifications of D23.F30, and
write $V_i$ for the local permutation on block $i$.
D23.H3 says that, for either normalized cubic value, the
four local maps on blocks $7,8,10,11$ belong to $S$.
It suffices below to use just one such block.
For any generator $g$ taking block $i$ to block $j$, let
$L(g,i)=T_j^{-1}gT_i\in L$ be its local transition.
The actual covariance gives
\[
            V_jL(g,i)=L(\Phi_h(g),i)V_i,
\]
where both transition maps belong to $L$ because
$g,\Phi_h(g)\in\Gamma$ and the blocks are fixed by $V_h$.
Starting at a block with $V_i\in S\subset L$ and traversing
the transitive block action proves $V_j\in L$ for every $j$.
This propagation does not assume in advance that $V_h$
belongs to the finite braid image.

Let $v_i\in L/S=\mathbb F_2$ be the twelve local cosets.
Passing the last displayed equation to $L/S$, the equality
of the signed generator and its transformed signed value
gives
\[
                          v_{q_g(i)}=v_i
                \qquad(g=A,B,C).
\]
Indeed the two transition characters are identical because
$F_{24}=R_{24}=I$ and $\Phi_h(A)=A$ in that quotient.
Transitivity of $Q_{12}$ makes all $v_i$ equal. One is zero
by D23.H3, so every $v_i$ is zero. Thus every local $V_i$
belongs to $S$. D23.F30 identifies the independent product
of these twelve local groups with an actual subgroup
$S^{12}\subset\Gamma$. Therefore $V_h\in S^{12}$, as
claimed.
\end{proof}

\paragraph{Exact dependency scope.}
The only supplied arithmetic normalization used to reach this
case is D23.H2. The finite Burau input is supplied FP.37,
especially its explicit first-pair table and equations
(FP.37.4)--(FP.37.6), in the original positive marking.
D23.F31 removes the possible scalar lift ambiguity. The
remaining proof uses the actual covariance D23.P18, the
checked block structure D23.F30, and the exhaustive local
lift theorem D23.H3. It does not infer an actual GT element
from finite solutions and makes no half-Kummer comparison.

\subsection{The unique action on four orthogonal blocks}

\begin{statement}{D23.P21}{Proved}{A central phase and its actual two-cycle normalization}
\end{statement}
Restrict the first two braid generators of D23.P17 to the
union of orthogonal blocks $7,8,10,11$, of total degree
448. The resulting $B_3$ image has order 11664 and maps
onto the degree-36 detector image of order 3888. Its kernel
is the central subgroup of order three generated by
\[
                         z=(AB)^3.
\]
The pure-pair image $P=\langle A^2,B^2\rangle$ has order
5832, and $P'$ has order 1944. The same central kernel is
contained in $P'$.

Each of the six actual first-pair values determined by
D23.H1 has exactly three lifts in $P'$. Exactly one lift
satisfies the two-cycle. Consequently the entire first-pair
value on these 448 points is determined by $(\rho,a_3)$.

For $\rho=0$, use the H3 candidate indices and block order
$(7,8,10,11)$. The unmodified first generator has already
reduced the possibilities to $(i,i,i,j)$ in D23.P20.
Covariance of the second generator further gives
\[
\begin{array}{c|c}
a_3&\text{allowed pairs }(i,j)\\ \hline
0&(0,0),(1,1),(2,2)\\
1&(0,2),(1,1),(2,0).
\end{array}                                                \tag{D23.P21.1}
\]
Combining this with the independently computed block-11
condition in D23.H4 fixes the action on all four blocks:
\[
\begin{array}{c|c}
a_3& (i_7,i_8,i_{10},i_{11})\\ \hline
0&(0,0,0,0)\\
1&(0,0,0,2).
\end{array}                                                \tag{D23.P21.2}
\]
In particular the normalized mixed actual full-GT state,
if it exists, has the unique indicated nontrivial action
on these four blocks. Existence is not asserted.

\begin{proof}
The detector map on the 432 indicated points is an exact
$B_3$-equivariant map to the degree-36 action. The union
of the four orthogonal blocks is invariant under $A,B$;
the remaining sixteen points cause no ambiguity in the
map between the permutation image groups. Exact
stabilizer-chain calculations give the four orders
\[
 |\langle A,B\rangle|=11664,\quad
 |\langle A_{36},B_{36}\rangle|=3888,\quad
 |P|=5832,\quad |P'|=1944.
\]
The images of $P$ and $P'$ on the detector have orders
1944 and 648 by D23.H1. Hence each of the three maps
has kernel of order three. Direct evaluation shows that
$z=(AB)^3$ is nonidentity, has order three, projects to
the identity on the detector, commutes with $A,B$, and
belongs to $P'$. It therefore generates every one of
these kernels. The certificate verifies all of these
facts, including membership in $P'$.

Put $s=ABA$. The braid relation gives
\[
        sA^2s^{-1}=B^2,\qquad sB^2s^{-1}=A^2.
\]
For every actual commutator associator $f$ satisfying
the two-cycle, its value $F=f(A^2,B^2)$ consequently
satisfies
\[
                         sFs^{-1}=F^{-1}.                 \tag{D23.P21.3}
\]
If one lift $F$ of a prescribed detector value satisfies
this equation, then replacing it by $Fz^k$ changes the
two sides by $z^k$ and $z^{-k}$, respectively. Since
$z$ is central of order three, at most one of the three
lifts can satisfy the equation. Enumerating the 1944
elements of $P'$ finds exactly three lifts of each of
the six H1 values, and exactly one of them satisfies
(D23.P21.3). All six selected values also satisfy the
cyclic norm condition under $AB$; this is a check of
consistency, not a sufficiency claim for actual GT.

Now form each of the nine permutations from D23.P20
on the 448-point set and test the exact covariance
\[
                     VBV^{-1}=F^{-1}BF,
\]
using the selected $F$ for its $(\rho,a_3)$ value.
This complete nine-element calculation yields exactly
(D23.P21.1). It uses the literal permutations on all
448 points, not only their orthogonal matrices or their
degree-36 quotients.

Finally D23.H4 proves that at $\rho=0$ the mixed value
$R=f(X_3,X_4)$ is the identity on block 11. Thus $V_{11}$
must commute with the local third braid generator. Its
unique H3 candidate is index zero for $a_3=0$, and
index two for $a_3=1$. Substitution in (D23.P21.1)
gives (D23.P21.2).
\end{proof}

\paragraph{Certificate and exact input locators.}
The full finite certificate is
\path{hurwitz_first_four_block_lifts.py}, with JSON and
NPZ records of the same stem. The NPZ contains the
448-point ordering, its two generators, both
$\rho=0$ first-pair values, and every permutation
passing the displayed covariance test.

The actual detector values come from D23.H1,
\path{work/records/hurwitz36_gt_stability.json}.
The exhaustive orthogonal lift lists come from D23.H3,
\path{work/notes/hurwitz_h1_orthogonal_lifts.tex} and its
code/records. The final block-11 condition is D23.H4,
\path{work/notes/hurwitz_block11_harmonic.tex}, with
\path{work/code/hurwitz_block11_harmonic.py} and its
JSON/NPZ records. These are necessary actual-source
constraints, not constructions of an associator from
the finite values.

\subsection{Solvability of all mixed Hurwitz images and its exact consequence}

\begin{statement}{D23.F33}{Proved}{Mixed derived length six and independence of first-pair perfect coordinates}
\end{statement}
In the degree-$1344$ representation of D23.P17, put
\[
 M_{34}=\langle X_3,X_4\rangle,\quad
 M_{45}=\langle X_4,X_5\rangle,\quad
 M_{51}=\langle X_5,X_1\rangle.
\]
Each of these finite groups has derived length at most six.
The same bound holds for the joint image of the free source
under any finite collection of evaluations in these groups,
including their complete peripheral Nielsen closures.

Let $P$ be any finite image of the same free source obtained
from one or several first-pair evaluations, and let
$P^{(\infty)}$ be its eventual derived subgroup. In the joint
image of the ordinary commutator subgroup $F_2'$ in
$P'\times M$, where $M$ is any such joint mixed image,
one has
\[
                    P^{(\infty)}\times\{1\}\subseteq
                 \operatorname{im}(F_2'\to P'\times M).
\]
Thus the perfect coordinates of the first-pair image can
be varied independently while all those mixed word values
are fixed. This is a statement about common finite word
images. It does not preserve the full gentle identities,
the actual GT automorphism, or its canonical conjugator.

\begin{proof}
Compute the orbits of each literal generator pair on the
full set of $1344$ points, by an exhaustive traversal.
For $(X_3,X_4)$ there are sixteen components; for
$(X_4,X_5)$ there are thirty. They partition the complete
set, and the restriction of each pair to each component is
compared, entry by entry, with its stored permutation record.
Thus restriction embeds the full pair image in the product
of the listed component images; no component is discarded
because its order is inconvenient.

Here are all the distinct component-derived series. The last
two columns give multiplicities in the two decompositions.
Every series ends at the trivial group.
\[
\begin{array}{l|r|r}
 \text{orders of successive derived groups}&M_{34}&M_{45}\\\hline
 2,1&1&1\\
 12,4,1&1&4\\
 24,8,2,1&0&4\\
 36,4,1&1&0\\
 72,8,2,1&4&0\\
 331776,82944,27648,1728,64,1&0&6\\
 419904,104976,11664,729,1&1&1\\
 995328,82944,27648,1728,64,1&2&0\\
 1259712,314928,11664,729,1&0&2\\
 2654208,663552,221184,13824,512,8,1&0&12\\
 3779136,314928,11664,729,1&2&0\\
 7962624,663552,221184,13824,512,8,1&4&0
\end{array}
\]
The exact certificate obtains each derived subgroup from
the commutators and their normal closure in the preceding
permutation group. Schreier--Sims gives its order, and a
membership-checked reduction of the generator list changes
neither that subgroup nor the next step. The largest number
of nontrivial derived steps is six. A subgroup of a finite
product of groups of derived length at most six has the same
bound, proving it for both complete pair images.

Let $D=ABCABA$ be the reversing half twist in the same
positive marking. Its literal permutation satisfies
\[
                  DX_3D^{-1}=X_1,\qquad
                  DX_4D^{-1}=X_5.
\]
Thus $M_{51}$ is conjugate to $M_{34}$, with its two
generators reversed, and has the same derived-length bound.
Precomposing any of these evaluation maps with a Nielsen
automorphism does not change its target group. Any finite
joint image is again a subgroup of a product of solvable
groups of derived length at most six. This proves the
first assertion, including all such common-word closures.

For the second assertion, let $L=F_2'$ and write
$a:L\twoheadrightarrow P'$ and $b:L\to M$ for the two
evaluation maps. Replace $M$ by $b(L)$ if necessary, so
that $b$ is onto. The subgroup
\[
                         N=a(\ker b)
\]
is normal in $P'$. There is a well-defined surjection
$M\twoheadrightarrow P'/N$: choose a lift $w\in L$ of
an element of $M$ and send it to $a(w)N$. Different lifts
differ by $\ker b$, and the map is a homomorphism and is
onto. Consequently $P'/N$ is solvable. The perfect group
$P^{(\infty)}$ has trivial image in any solvable group:
its image is perfect and lies in every term of that group's
derived series, whose last term is trivial. Therefore
$P^{(\infty)}\subseteq N$. For each $s$ in this subgroup
there is a single ordinary commutator word $w\in\ker b$
with $a(w)=s$. This gives precisely the claimed inclusion.
All target groups are finite, so using the closed profinite
commutator subgroup instead gives exactly the same image.
\end{proof}

\paragraph{Exact records and scope.}
The complete forty-six component reconstruction and all
derived series are retained in
\path{agents/five_level_six/hurwitz_mixed_large_components.py}
and its JSON record. The code starts from the original
$1344$-point generators, checks all component restrictions,
and checks the displayed reversing conjugacy. No simplicity
or abstract recognition of the large first-pair factors is
needed for this proof. In particular it does not claim an
isomorphism with a product of alternating groups from a
coincidence of orders. It only localizes possible additional
common-word correlations to the solvable quotient of the
first-pair image. The actual GT normalizer and the exact
five-factor pentagon remain additional constraints.

\subsection{Sixteen actual-source candidates in the retained mixed quotient}

\begin{statement}{D23.C25}{Proved}{The normalized degree-72 covariance gate}
\end{statement}
Use the genuine degree-1344 Hurwitz action of D23.P17--P18, with
left permutation composition, and put
\[
 X=C^2,\qquad Y=(ABA)^2,\qquad R=f(X,Y),\qquad J=\langle X,Y\rangle.
\]
Suppose an actual multiplier-one full-GT element has
$(\rho,a_3,a_2)=(0,\epsilon,0)$, where $\epsilon\in\{0,1\}$.
Let $V$ be its actual permutation supplied by P18.  The two
degree-72 $J$-orbits numbered 9 and 10 are each preserved by $V$.
On either orbit, exactly sixteen values of $R$ satisfy the following
necessary conditions together:
\begin{enumerate}
\item $R$ belongs to the derived subgroup of the marked mixed image;
\item its marked $(p,t)=((123),(1234))$ quotient is the identity;
\item there is a permutation $T$ fixing the 24 points in orthogonal
      blocks $7,8,10$ and satisfying
      $TXT^{-1}=RXR^{-1}$ and $TYT^{-1}=Y$.
\end{enumerate}
For each such $R$, the permutation $T$ is unique.  The sixteen
values form an elementary abelian group of order sixteen.  Their
four distinct $T$ values have orders $1,2,2,2$, with four values
of $R$ over each $T$.

The two marked degree-72 representations have the same kernel.
Their sixteen surviving values therefore match in a graph with
sixteen pairs.  In the saved lexicographic ordering of their four
$T$ values, the only possible pairs are
\[
                   (0,0),\quad(1,3),\quad(2,1),\quad(3,2).
\]
All these conclusions are identical for $\epsilon=0$ and
$\epsilon=1$.  They impose necessary finite restrictions; they do
not construct an actual associator or prove the existence of the
normalized unequal-bit state.

\begin{proof}
The reversed standard free-group automorphism is
\[
            x\longmapsto f(x,y)x f(x,y)^{-1},
            \qquad y\longmapsto y.
\]
Indeed it is obtained from the usual gentle automorphism by
interchanging the two free generators and using the two-cycle.
It is an automorphism of the free profinite group.  Its images in
any finite quotient therefore still generate that quotient.
P18 and the fixed three-strand full twist give
\[
                VXV^{-1}=RXR^{-1},\qquad VYV^{-1}=Y.
\]
Consequently $VJV^{-1}=J$, so $V$ permutes the $J$-orbits.

The exact orbit enumeration has two orbits of degree 72.  Each
meets each of blocks $1,2,3,4,5,6,7,8,10$ in eight points and
meets no other block.  P21 gives the actual maps on blocks
$7,8,10,11$.  Substitution of those literal permutations, for
either cubic bit, fixes every one of the 24 indicated points in
each degree-72 orbit.  A normalizer of $J$ sends a $J$-orbit onto
a $J$-orbit; since these orbits contain fixed points, each is
preserved individually.

On each component, C21 gives the mixed image order 995328 and
derived order 82944, and the exact marked $S_4$ quotient.
The supplied FP.37.2 formula for full GT is
$f(p,t)=t^{2a_2}p^\rho$.  Thus the normalized state has identity
image in that quotient.  The joined permutation representation
on $72+4$ points has the same group order.  Enumeration of its
derived subgroup finds exactly 6912 values with last four
coordinates fixed.

For each of these 6912 values $R$, solve the two displayed
intertwining equations by propagation.  Begin with the prescribed
images of the 24 known points.  Whenever $T(a)$ is known, the
equations force the images of $X(a)$ and $Y(a)$, respectively,
to be $(RXR^{-1})(T(a))$ and $Y(T(a))$.  A conflict of either
images or inverse images rejects the candidate.  Transitivity
of $\langle X,Y\rangle$ ensures that propagation reaches every
point.  Therefore this procedure is both necessary and
sufficient for the stated finite intertwiner problem, and it
has at most one solution.  Its exact result is sixteen accepted
values and four distinct intertwiners.  Direct multiplication
checks that the sixteen values generate a group of order
sixteen and that all fifteen nonidentity elements are
involutions.  The subgroup is therefore $C_2^4$.

Finally join the two degree-72 representations on 144 points.
Its image has order 995328, equal to both projections; hence
both projections are isomorphisms and the marked kernels
coincide.  Its derived image has order 82944.  Membership
testing of the $16^2$ pairs in this derived image gives
exactly sixteen pairs, bijective under both projections.
Reading their already unique intertwiners gives the four pairs
in the statement.  This is a same-word calculation in finite
quotients and needs no assertion that finite candidates lift
to actual gentle or full-GT elements.
\end{proof}

\paragraph{Exact certificates and scope.}
The exhaustive source-marked calculation is
\path{agents/characters/code/hurwitz_mixed72_partial.py},
with JSON and NPZ records of the same stem.  Its NPZ records the
global point labels, both restricted generators, the 24 marked
fixed points, all sixteen $R$ values, their intertwiners, and
the four lexicographically ordered distinct intertwiners.
The graph comparison is
\path{agents/characters/code/hurwitz_mixed72_joint.py}, with
\path{agents/characters/records/hurwitz_mixed72_joint_probe.json}.
The actual-source inputs are P18, P21, H4, and the supplied
FP.37.2 marked $S_4$ formula.  The finite comparison is stronger
than the compact order-576 quotient of C21, but it still leaves
both cubic states possible.

\paragraph{Further local tests: an explicitly recorded failure.}
For each of the 72 global permutations retained after C25 in
the subsequent global calculation, all other mixed-pair orbit
restrictions were tested independently.  Every restriction
normalizes its marked image and sends $X$ to a conjugate under
that image's derived subgroup.  Explicit local conjugators
were found for every case; no rejection was inferred from a
search cutoff.  The exact records are
\path{agents/characters/records/hurwitz_global_mixed_local.json}
and the accompanying NPZ, generated by the code of the same
stem.  Thus these separate local tests do not reduce the 72
global candidates.  Common-word compatibility across the
components and the full equation involving $C$ remain
additional constraints.

\subsection{A complete list of global necessary conjugators}

\begin{statement}{D23.P22}{Proved}{At most 144 global conjugators before the mixed tests}
\end{statement}
For an actual multiplier-one full-GT element in the normalized
state
\[
                     \rho=0,\qquad a_2=0,
\]
and either fixed value of $a_3$, its actual permutation
$V_h$ on the degree-1344 orbit belongs to a completely
enumerated list of 144 permutations. Every permutation in
that list fixes each of the twelve orthogonal blocks,
belongs to $S^{12}$, commutes with the literal first braid
generator $A$ and the nested full twist
$Z=(ABA)^2=(AB)^3$, and has the unique four-block
restriction determined in D23.P21.

The list is exhaustive for precisely these conditions.
It consists of 18 choices on block zero and eight choices
on block one, with all other blocks determined by transport.
The two degree-72 mixed component tests of D23.C25,
including their common-word graph, leave 72 permutations
for each value of $a_3$. These remaining permutations are
necessary candidates only. In particular they are not
asserted to arise from actual GT elements or to satisfy all
remaining finite evaluations.

\begin{proof}
Write $Z=(AB)^3$. In $B_3$ this is the central full twist,
so $A$ and $Z$ commute. The standard Drinfeld action fixes
$A$ and $Z$ when $\lambda=1$. The actual covariance
D23.P18 consequently implies
\[
                         V_hA=AV_h,\qquad V_hZ=ZV_h.
 \tag{D23.P22.1}
\]
By D23.F32, the normalized actual $V_h$ belongs to the
independent simple product $S^{12}$.

The degree-12 action of $\langle A,Z\rangle$ has orbits
\[
       \{0,9\},\qquad\{1,2,3,4,5,6\},\qquad
       \{7,8,10\},\qquad\{11\}.                          \tag{D23.P22.2}
\]
On the first orbit, $A$ fixes both blocks and $Z$
interchanges them. On the second orbit, the degree-three
action of $A$ and the degree-two action of $Z$ generate a
regular $C_3\times C_2$ action. This is an exact check of
the two displayed block permutations. The other two orbits
are already fixed by D23.P21's complete local lists.

We use the following elementary transport fact. Let an
abelian group generated by commuting permutations $a,z$
act transitively on a set of blocks. A permutation $v$
fixing each block and commuting with $a,z$ is determined
by its restriction to one block. That restriction must
centralize the image of the stabilizer of that block.
Conversely every such local centralizer element extends
uniquely: transport it by any word taking the base block
to the desired block. Two choices differ by a stabilizer
element, so this is well defined. The extension commutes
with both generators by construction.

For the orbit $\{0,9\}$, the stabilizer is the image of
$\langle A,Z^2\rangle$. For the regular six-block orbit,
it is the image of $\langle A^3,Z^2\rangle$. In the
group-transversal coordinates of D23.F30 the corresponding
local centralizers in $S$ have exact orders
\[
 \begin{split}
 \left|C_S\bigl(A_0,(Z^2)_0\bigr)\right|&=18,\\
 \left|C_S\bigl((A^3)_1,(Z^2)_1\bigr)\right|&=8.
 \end{split}                                               \tag{D23.P22.3}
\]
The first local pair has element orders $(3,2)$ and the
second has orders $(4,2)$. The seed certificate constructs
the actual group $S$ as the derived subgroup of the
112-point block-stabilizer image, computes both complete
centralizers by stabilizer chains, enumerates every element,
and verifies its membership in $S$ and commutation with
both local operators.

Transport as above gives a bijection between the $18\cdot8$
seed choices and all possible extensions of the known
four-block action satisfying (D23.P22.1). The global
certificate constructs every extension as a literal
permutation on all 1344 points and checks both equations
in (D23.P22.1) on every point. It also checks the block
labels and distinguishes all 144 permutations. The
transporters belong to $\Gamma$, which normalizes the
independent factors of $S^{12}$, so every extension belongs
to $S^{12}$. This proves exhaustiveness for the stated
conditions, including the equal-sign control case $a_3=0$.

For completeness, the further mixed-component restriction
has the following actual-source interpretation. Put
$X_3=C^2$, $X_4=Z$, and $R=f(X_3,X_4)$. The same
covariance gives
\[
     V_hX_3V_h^{-1}=RX_3R^{-1},\qquad
     V_hX_4V_h^{-1}=X_4.                                  \tag{D23.P22.4}
\]
The corresponding transformation of the free pair is an
automorphism for actual gentle $f$ (it is the swapped
standard action, using the two-cycle). Hence $V_h$
normalizes the image $\langle X_3,X_4\rangle$ and permutes
its orbits. Each of the two degree-72 orbits contains 24
points in blocks $7,8,10$, on which the already determined
$V_h$ is the identity. They are therefore preserved
individually, rather than exchanged.

Each degree-72 image has order 995328, derived order
82944, and the exact marked $S_4$ quotient. At
$\rho=a_2=0$, the required marked $S_4$ value is the
identity. D23.C25 enumerates the 6912 elements of the
derived group with this marked value, and propagates
(D23.P22.4) from the known 24 points. Exactly sixteen
$R$ values and four induced permutations remain in each
component. The common-word joint image has the same order
995328 as either component, so the two values are tied by
its exact graph. In the stored four-element permutation
lists, the permitted index pairs are
\[
                       (0,0),(1,3),(2,1),(3,2).
\]
Restricting each of the 144 global permutations to both
components and imposing this graph leaves exactly 72.
Only the pairs $(0,0)$ and $(1,3)$ occur among them.
Equivalently, all eighteen block-zero seeds remain, while
four of the eight block-one seeds remain. The count is
the same for the two cubic-bit control cases. This is a
necessary restriction; other mixed and first-pair conditions
have not been omitted from the final problem by calling
these 72 permutations solutions.
\end{proof}

\paragraph{Certificates and exact scope.}
The canonical seed computation is
\path{hurwitz_seed_centralizers.py}, with JSON and NPZ
records. It independently rebuilds both centralizers and
keeps the previously recorded numbering of their elements
so that downstream indices remain stable.

The global construction and mixed restriction are
\path{hurwitz_global_conjugators.py}, with JSON and NPZ
records. The arrays \texttt{all\_V\_0}, \texttt{all\_V\_1}
contain the complete lists of 144 permutations
(each row has 1344 entries); the arrays
\texttt{mixed72\_V\_0}, \texttt{mixed72\_V\_1} contain
the 72 remaining permutations in each control case.
Their corresponding label arrays keep both seed indices
and both component-permutation indices.

The mixed inputs are the exact D23.C25 records
\path{characters/records/hurwitz_mixed72_partial.npz} and
\path{characters/records/hurwitz_mixed72_joint_probe.json},
with their code. The actual-source dependencies are
D23.P18, D23.P21, D23.F32 and the supplied two-cycle;
the group structure is D23.F30. The theorem establishes
a complete finite list that every normalized actual
conjugator must enter. It establishes neither existence
of a mixed full-GT state nor its impossibility.

\subsection{The full first-pair covariance test}

\begin{statement}{D23.F34}{Proved}{Thirty-six first-pair survivors and an explicit cycle-compatible value}
\end{statement}
Let $A,B,C$ be the literal degree-$1344$ permutations of D23.P17,
and put
\[
 G_3=\langle A,B\rangle,\qquad P=\langle A^2,B^2\rangle,
 \qquad r=AB,\quad s=ABA,\quad z=r^3=s^2.
\]
The exact orders are
\[
 |G_3|=3134566563840000,\qquad
 |P|=1567283281920000,\qquad
 |P'|=261213880320000.
\]
Moreover $P/P'$ is cyclic of order six, generated by the image of
$B^2$.

Consider the two lists of $72$ global candidate permutations in
\path{pentagon/hurwitz_global_conjugators.npz}, with keys
\texttt{mixed72\_V\_0} and \texttt{mixed72\_V\_1}.
The index is zero based. These are the already constructed normalized
candidates with $(\rho,a_2,a_3)=(0,0,b)$, where $b=0,1$;
this statement verifies the following additional conditions on the
literal lists, without claiming that their members lift to actual GT.
Each candidate commutes with $A$ and $z$. Exactly $36$ candidates in
each list normalize the full group $G_3$.

The accepted index set is the same in both lists:
\[
\begin{gathered}
 0,1,4,5,10,11,12,13,18,19,22,23\\
 26,27,30,31,34,35,38,39,40,41,46,47\\
 50,51,52,53,56,57,60,61,64,65,68,69
\end{gathered}
\]

For every one of those candidates $V$, the explicit permutation
\[
                   F_V=rVr^{-1}V^{-1}                    \tag{D23.F34.1}
\]
lies in $P'$, restricts on the ordered $448$-point subset to the
unique normalized value of D23.P21, and satisfies
\[
 \begin{split}
 V B V^{-1}&=F_V^{-1}BF_V,\\
 sF_Vs^{-1}&=F_V^{-1},\\
 (r^2F_Vr^{-2})(rF_Vr^{-1})F_V&=1.
 \end{split}                                             \tag{D23.F34.2}
\]
Thus first-pair subgroup normalization discards half of the listed
candidates, while membership in the commutator image, exact D23.P21
restriction, and the two displayed cycle conditions discard none of
the remaining candidates. These tests alone do not separate the two
cubic bits.

\begin{proof}
All groups here are finite permutation groups on the original
$1344$ points. Construct $G_3$ and $P$ from the indicated generators,
construct $P'$ as the normal closure of their commutators, and compute
orders by exact stabilizer chains. This gives the three displayed
orders. Their ratio gives $|P/P'|=6$. Testing the six powers of
$B^2$ for membership in $P'$ shows that its image has order exactly
six, proving the cyclic assertion.

For each of the $144$ input permutations, the certificate checks
$VA=AV$ and $Vz=zV$ literally. Because $V$ fixes $A$ by conjugation,
\[
             VG_3V^{-1}\subseteq G_3
       \quad\Longleftrightarrow\quad VBV^{-1}\in G_3.
\]
The two finite groups in this inclusion have equal order; hence the
inclusion is equivalent to equality. An exact membership test in the
full $G_3$, rather than its separate orbit projections, accepts
exactly $36$ indices for each bit. All accepted and rejected indices
and their original four seed labels are retained in the JSON record.
For an actual full-GT permutation of D23.P18, normalization of $G_3$
is necessary because the standard GT action restricts to an
automorphism of the first $B_3$.

The covariance and cycle identities have a direct proof, independent
of the membership computations. Since $rAr^{-1}=B$ and $V$ commutes
with $A$, the element
\[
                       F_V V=rVr^{-1}
\]
commutes with $B$. This is the first equation in (D23.F34.2).
Also $s=rA$ and $AV=VA$, so
\[
                    F_V=sVs^{-1}V^{-1}.
\]
As $s^2=z$ commutes with $V$, conjugating this expression by $s$
gives $sF_Vs^{-1}=V(sVs^{-1})^{-1}=F_V^{-1}$.
For the third equation let $\alpha=\operatorname{Ad}(r)$.
Then $F_V=\alpha(V)V^{-1}$ and $\alpha^3(V)=V$, whence
\[
 \alpha^2(F_V)\alpha(F_V)F_V
 =\alpha^3(V)\alpha^2(V)^{-1}
  \alpha^2(V)\alpha(V)^{-1}\alpha(V)V^{-1}=1.
\]
These are precisely the full-permutation equations subsequently
replayed by the certificate.

For each accepted $V$, compute (D23.F34.1), test its membership in
the previously constructed $P'$, and restrict it to the literal
ordered subset stored with D23.P21. All $72$ accepted values lie in
$P'$ and agree entry by entry with the corresponding
\texttt{first\_f\_0} or \texttt{first\_f\_1} array. This proves the
remaining assertions. No search over a small family of possible
corrections is being used to infer nonexistence: the successful values
are explicitly constructed and saved.
\end{proof}

\paragraph{Reproducibility and limitation.}
The self-contained replay is
\path{agents/five_level_six/hurwitz_first_pair_covariance.py};
its JSON contains all input indices and membership outcomes, and
its NPZ contains every accepted explicit $F_V$. The group calculations
start from the original generators, not from cached group orders.
The input list is
\path{agents/pentagon/hurwitz_global_conjugators.npz};
the exact $448$-point comparison uses
\path{agents/pentagon/hurwitz_first_four_block_lifts.npz}.
Additional common-word conditions with the mixed evaluations,
the full pentagon, and existence of an actual profinite GT element
are not inferred from this certificate.

\subsection{Full mixed commutator covariance still permits every candidate}

\begin{statement}{D23.C26}{Proved}{Exactly 432 mixed commutator values per candidate}
\end{statement}
Let $M=\langle C^2,Z\rangle$ in the degree-1344 Hurwitz image,
where $Z=(ABA)^2$, and let $K$ be the kernel of its signed
degree-24 quotient.  For every one of the 72 candidate
permutations $V$ of P22 with $(\rho,a_3,a_2)=(0,1,0)$,
there are exactly 432 elements $R\in M'\cap K$ such that
\[
                       RCR^{-1}=VCV^{-1}.                 \tag{D23.C26.1}
\]
Thus the entire mixed-word commutator condition, its normalized
signed value, and the full third-generator covariance do not
exclude any of these candidates.  The assertion concerns
ordinary commutator-word values in the finite image.  It does
not assert the gentle cycle relations, compatibility with the
first-pair value of the same word, or existence of an actual
GT element.

\begin{proof}
The exact Schreier kernel construction in the root certificate
\path{work/code/hurwitz_mixed_signed_kernel.py} supplies a
generating set for $K$.  It traverses all 144 elements of the
signed mixed quotient while retaining literal lifts in $M$;
the resulting Schreier generators therefore generate the full
kernel, not merely a subgroup of it.

Split the twelve orthogonal blocks into the three disjoint
$C$-invariant sets
\[
       I_{09}=\{0,9\},\qquad
       I_9=\{1,2,3,4,5,6,7,8,10\},\qquad I_{11}=\{11\}.
\]
The restrictions of $K$ to these sets have orders
\[
       |K_{09}|=104976=324^2,\quad
       |K_9|=110592,\quad |K_{11}|=24.
\]
The exact full kernel order is
\[
 |K|=278628139008=104976\cdot110592\cdot24.
\]
Restriction embeds $K$ into the product of its three
projections, because their supports exhaust all 1344 points.
Equality of these finite orders proves the literal product
decomposition
\[
                       K=K_{09}\times K_9\times K_{11}.
                                                               \tag{D23.C26.2}
\]
In particular there is no further same-word relation between
these three coordinates inside $K$.

For each $V$, exact enumeration gives 54 conjugators satisfying
(D23.C26.1) in the first coordinate, eight in the second, and
six in the third.  On $I_9$, the 72 candidates yield four
distinct targets; all 110592 elements of $K_9$ are tested and
each target has eight conjugators.  On $I_{09}$ there are six
distinct targets, with 54 conjugators each.  On block 11, H4
says that $V$ commutes with $C$; enumeration of the group of
order 24 gives a centralizer of order six.  The first-coordinate
certificate is the root's
\path{work/records/hurwitz_mixed_fixed_pair_witnesses.npz};
the nine-block certificate is specified below.

By the product decomposition every candidate therefore has
$54\cdot8\cdot6=2592$ conjugators in $K$.  Each nonempty
conjugator set is a right coset of
\[
                    D=C_K(C),\qquad |D|=2592.              \tag{D23.C26.3}
\]
The group $D$ is constructed explicitly from the three local
centralizers: if $r_0,r$ solve the same conjugacy equation,
then $r_0^{-1}r$ centralizes $C$.  This gives complete local
centralizer lists and hence the complete product centralizer.

The independent full mixed image calculation gives
\[
       |M|=40122452017152,\qquad |M'|=557256278016,
       \qquad |M/M'|=72.
\]
The image of $M'$ in the signed quotient has order 12.
Consequently
\[
               [K:K\cap M']=6.                            \tag{D23.C26.4}
\]
This index is retained; one cannot replace $K\cap M'$ by $K$.

Form the finite subgroup $L=\langle M',D\rangle=M'D$, using
the complete derived-group stabilizer chain and the explicitly
constructed centralizer.  Its exact order is
\[
        |L|=3343537668096=6|M'|.
\]
Thus $D$ maps onto all six classes of $K/(K\cap M')$.
Every right coset $r_0D$ in $K$ therefore meets $M'$ in exactly
\[
                   |D\cap M'|=2592/6=432
\]
elements.  Apply this to the conjugator coset for each $V$.
The certificate also computes $D\cap M'$ explicitly, checks
its order 432, chooses a six-element coset transversal in $D$,
and corrects one conjugator for every $V$ into $M'$.
All 72 corrected representatives are checked directly both
for membership in $M'$ and for the full equation on all
1344 points.  This proves the assertion without inferring
any actual profinite lift from the finite solutions.
\end{proof}

\paragraph{Reproducible data.}
The kernel-product and complete nine-block computation is
\path{agents/characters/code/hurwitz_mixed_kernel_product.py},
with records
\path{agents/characters/records/hurwitz_mixed_full_kernel_order.json}
and \path{agents/characters/records/hurwitz_mixed_nine_cycle_probe.npz}.
The final exact derived-fiber calculation is
\path{agents/characters/code/hurwitz_mixed_derived_fibers.py},
with JSON and NPZ records of the same stem.  In the latter,
\texttt{derived\_representatives} is indexed by P22's 72
normalized mixed candidates, and
\texttt{derived\_centralizer\_generators} generates the
order-432 right factor of every fiber.  The full mixed group
and its derived stabilizer chains are the independently
computed inputs in
\path{agents/five_level_six/hurwitz_full_mixed_groups_state.pkl}.

\subsection{The exact maximal solvable quotient of the first Hurwitz pair}

\begin{statement}{D23.F35}{Proved}{Perfect residual and a marked solvable quotient of order $93312$}
\end{statement}
In the notation of D23.F34, the full pure-pair image
$P=\langle A^2,B^2\rangle$ has the following exact derived series:
\[
\begin{array}{c|r}
 i&|P^{(i)}|\\\hline
0&1567283281920000\\
1&261213880320000\\
2&16325867520000\\
3&4081466880000\\
4&50388480000\\
5&16796160000\\
6&16796160000.
\end{array}                                               \tag{D23.F35.1}
\]
In particular its perfect residual is
\[
 N=P^{(5)}=P^{(6)},\qquad |N|=16796160000=360^4.
\]
This order identity is not an assertion identifying $N$ with a
product of alternating groups.

Restrict the first pair to the $448$-point subset of D23.P21, and
simultaneously take its signed degree-$24$ quotient from D23.F30.
Their marked joint permutation image $Q$ has degree $472$, order
$93312$, and derived series
\[
               93312,15552,972,243,3,1.                  \tag{D23.F35.2}
\]
The natural marked surjection $P\twoheadrightarrow Q$ has kernel
exactly $N$. Therefore
\[
                Q\cong P/P^{(\infty)},\qquad
                |P'/P^{(\infty)}|=15552.
\]
Thus this concrete joint image is the maximal solvable quotient of
the full first-pair image, with its original two free generators
retained.

\begin{proof}
Start with the literal degree-$1344$ generators of D23.P17 and form
$P$ from their squares. At each stage, form the normal subgroup
generated by all commutators of the current generators and compute
its exact order with a stabilizer chain. This gives precisely the
orders in (D23.F35.1). When a generator list is unnecessarily long,
the certificate replaces it by a shorter list of explicit products
in those generators. The replacement subgroup is contained in the
previous one by construction, and equality of their computed finite
orders proves equality of the groups. The random choice of products
is only an implementation aid: no probabilistic conclusion about
generation or group recognition is used.

The sixth derived subgroup is contained in the fifth. Their equal
finite orders prove equality, so the fifth is perfect and the series
has stabilized. This proves the assertion about $N$, including that
it is the intersection of the whole derived series.

For each pure generator, concatenate its restriction to the ordered
$448$-point invariant subset with its signed action on $24$ points.
Both are actual quotients of $P$: the former is restriction to an
invariant subset, and the latter is the signed quotient proved in
D23.F30. Hence the resulting $472$-point group is a marked quotient
$Q$ of $P$, without any presentation or order-based identification.
The literal two arrays are checked against the separately
constructed root record
\path{work/records/hurwitz_first_solvable_joint.npz}.
Computing all its successive derived groups gives exactly
(D23.F35.2), and in particular proves that $Q$ is solvable.

A homomorphism sends every term of a derived series into the
corresponding term of the target series. Since $Q^{(5)}=1$, the
kernel of $P\to Q$ contains $P^{(5)}=N$. On the other hand,
\[
        |\ker(P\to Q)|=\frac{1567283281920000}{93312}
                      =16796160000=|N|.
\]
Thus the kernel equals $N$. Finally the image of the commutator
subgroup under a surjection is the target commutator subgroup,
so $P'/N\cong Q'$ has order $15552$. Every solvable quotient of
$P$ kills its perfect residual; consequently $Q=P/N$ has the stated
maximal property.
\end{proof}

\paragraph{Exact replay and scope.}
The complete certificate is
\path{agents/five_level_six/hurwitz_first_perfect_residual.py},
with JSON and NPZ records of the same stem. It starts from the
original $1344$-point generators, constructs every derived group,
checks stabilization, constructs the marked $472$-point quotient,
and checks its complete derived series. No cached group order is
accepted as an input. The NPZ retains generating lists at every
step and the literal two quotient generators. This statement
identifies the maximal solvable quotient only; common-word lifting
and actual gentle or full-GT realization are separate assertions.

\subsection{The exact remaining first--mixed common-word test}

\paragraph{D23.H7 (proved; whole-proof check completed).}
Let $P=\langle A^2,B^2\rangle$ and
$M=\langle C^2,(ABA)^2\rangle$ in the full degree-$1344$
action. The marked joint quotient of $P$ given by its
degree-$448$ restriction and the signed degree-$24$
action has order $93312$, with derived order $15552$.
Write it as $Q$. By D23.F35 it is the maximal solvable
quotient of $P$.

The same-word image $J$ in $Q\times M$, under the marked
generators $(A^2,C^2)$ and $(B^2,(ABA)^2)$, has
\[
 \begin{split}
 |J|&=77998046721343488,\\
 |J'|&=1083306204463104.
 \end{split}
\]
The common quotient in Goursat's lemma therefore has order
$48$. The common quotient for the two derived projections
has order eight.

For each of the $36$ normalized mixed-state permutations
retained in D23.F34, there is an ordinary word
$w\in F_2'$ whose two full evaluations are
\[
         w(A^2,B^2)=F_V,\qquad
         w(C^2,(ABA)^2)=R_V,
\]
where $F_V$ is the explicit value in D23.F34 and
$R_V$ satisfies the full degree-$1344$ equation
\[
                    VC V^{-1}=R_V C R_V^{-1}.
\]
Among the $432$ mixed commutator conjugators for a fixed
$V$, exactly $216$ have the required common-word image
in $Q'\times M'$.

Only these two evaluations are asserted. The word is not
asserted to satisfy the free two-cycle or three-cycle,
the other three pentagon evaluations, or the actual GT
conditions. In particular a value at the swapped mixed
pair is not silently replaced by the inverse of $R_V$.

\begin{proof}
The quotient $Q$ is the literal joint permutation image
on $448+24$ points. It is a quotient of $P$ because
both component maps are restrictions or quotients of
the given braid representation. Exact stabilizer-chain
calculations give order $93312$ and its complete
solvable derived series. The independent calculation
D23.F35 gives
\[
             |P^{(\infty)}|=16796160000,
 \qquad |P|/|P^{(\infty)}|=93312.
\]
Since every map to a solvable group kills the perfect
residual, the equality of orders identifies the kernel
of $P\to Q$ with $P^{(\infty)}$.

Join the two generators of $Q$ to the literal generators
of $M$ on a disjoint $1344$-point set. This constructs
$J$ on $1816$ points, directly from the two free
generators, and not from a presumed presentation.
Its order and derived order are as displayed. The
individual orders are
\[
 |M|=40122452017152,
 \qquad |M'|=557256278016.
\]
Both projections of $J$ are onto, as are the projections
of $J'$ onto $Q'$ and $M'$. Thus the two common quotient
orders are exactly
\[
          \frac{|Q|\,|M|}{|J|}=48,
 \qquad \frac{|Q'|\,|M'|}{|J'|}=8.
\]

For the normalized mixed state, the prescribed first
value in $Q$ is the pair
\[
                 q_*=(F_{448,1},I_{24}),
\]
where $F_{448,1}$ is D23.P21's unique value. It is the
same for all $36$ candidates, because their first
restrictions agree and their signed action is trivial.
The certificate additionally checks $(q_*,1)\in J'$.
This is an exact finite statement of independence of
this particular cubic value from the entire single
mixed-pair image $M$.

For each candidate $V$, the full mixed-conjugacy
calculation gives its $432$ values in the coset
\[
                      r_V C_*,\qquad
                  C_*=C_{K\cap M'}(C),\quad |C_*|=432,
\]
where $K$ is the kernel of the signed mixed quotient.
The representatives and the exact centralizer
generators are retained in
\path{agents/characters/records/hurwitz_mixed_derived_fibers.npz}.
Let
\[
 N_* =\{c\in C_*:(1,c)\in J'\}.
\]
This is a normal subgroup: it is the intersection with
$C_*$ of the kernel of the common quotient map on $M'$.
Exhaustive exact membership testing gives
$[C_*:N_*]=2$. Testing the two coset representatives
against $(q_*,r_Vc)\in J'$ accepts exactly one coset
for each of the $36$ candidates. Every accepted coset
has size $432/2=216$, proving the stated count.
The computation saves a full $1344$-point witness $R_V$
for each candidate.

Membership in $J'$ is exactly membership in the image
of the ordinary commutator subgroup under the two
quotient evaluations. Choose a word $w_0\in F_2'$
with values $(q_*,R_V)$. Its first full value $P_0$ may
differ from the specified $F_V$, but their quotient
values in $Q$ agree. Hence $F_VP_0^{-1}$ belongs to
$P^{(\infty)}$. By D23.F33 there is a single ordinary
commutator word $u$ with full values
\[
                 (F_VP_0^{-1},1)\in P\times M,
\]
because $M$ is solvable. The word $w=uw_0$ has
the two exact full values asserted in the statement.
This adjustment preserves the mixed value but carries
no claim about identities of the free word itself.
\end{proof}

\paragraph{Exact records.}
The marked join and its derived group are constructed by
\path{code/hurwitz_solvable_commonword.py}. The common
quotient and all mixed-fiber membership tests are in
\path{code/hurwitz_solvable_fiber_filter.py}; their JSON
and NPZ records retain the generator permutations,
all $36$ indices, both centralizer coset representatives,
and the chosen witnesses. The maximal solvable quotient
input is D23.F35, with its complete full-$P$ derived
series; no abstract alternating-group recognition is
needed.

\subsection{Uniqueness of the full first-pair value}

\begin{statement}{D23.F36}{Proved}{The explicit first value is forced by the two-cycle}
\end{statement}
Use the full first-pair groups $G_3=\langle A,B\rangle$,
$P=\langle A^2,B^2\rangle$, and
$N=P^{(\infty)}$ from D23.F34--F35. The perfect residual has
trivial center:
\[
                             Z(N)=1.                     \tag{D23.F36.1}
\]
Let $V$ be any of the $36$ retained candidates for either normalized
cubic bit, and let $F_0=F_V$ be the explicit value in D23.F34.
Then $F_0$ is the unique $F\in P'$ satisfying all three conditions
\[
 \begin{split}
 F\bmod N&=F_0\bmod N,\\
 VBV^{-1}&=F^{-1}BF,\\
 sFs^{-1}&=F^{-1},\qquad s=ABA.
 \end{split}                                             \tag{D23.F36.2}
\]
The first condition is exactly equality of the combined
$448$-point and signed $24$-point values, by D23.F35.
Thus for an actual normalized full-GT element with that canonical
permutation $V$, its complete first value is necessarily $F_V$.
There is no remaining first-value fiber that can evade a subsequent
finite pentagon test by choosing another commutator value.

\begin{proof}
The center assertion is an exact finite computation on the perfect
residual already constructed in D23.F35. Start from its saved four
degree-$1344$ generators, rebuild the group and its stabilizer chain,
and check its order $16796160000$. Compute its centralizer in itself
by simultaneous commutation with all four generators. The resulting
subgroup has order one. Since commuting with a generating set is
equivalent to commuting with every group element, this is exactly
$Z(N)$, not a center test on a smaller quotient. The certificate
replays this calculation without assuming an abstract
identification of $N$.

Suppose that $F$ satisfies (D23.F36.2), and set
\[
                              c=FF_0^{-1}.
\]
The common quotient value gives $c\in N$. Both $F$ and $F_0$
conjugate $VBV^{-1}$ to $B$, so
\[
                             cB=Bc.                       \tag{D23.F36.3}
\]
Both also satisfy the two-cycle. Substituting $F=cF_0$ into that
identity and using $sF_0s^{-1}=F_0^{-1}$ yields
\[
                         scs^{-1}=F_0^{-1}c^{-1}F_0.
                                                               \tag{D23.F36.4}
\]
The braid relation gives $sBs^{-1}=A$. Conjugating
(D23.F36.3) by $s$ and using (D23.F36.4) shows that
$F_0^{-1}c^{-1}F_0$ commutes with $A$. Equivalently,
$c$ commutes with $F_0AF_0^{-1}$.

These last two commuting relations imply that $c$ centralizes
all of $G_3$. Indeed, D23.F34 proves $F_0\in P'\subseteq G_3$,
$VA=AV$, and $VG_3V^{-1}=G_3$, and its covariance gives
$F_0^{-1}BF_0=VBV^{-1}$. Therefore
\[
 \begin{split}
 \langle B,F_0AF_0^{-1}\rangle
 &=F_0\langle F_0^{-1}BF_0,A\rangle F_0^{-1}\\
 &=F_0\langle VBV^{-1},VAV^{-1}\rangle F_0^{-1}\\
 &=F_0VG_3V^{-1}F_0^{-1}=G_3.
 \end{split}
\]
Consequently $c\in N\cap Z(G_3)\subseteq Z(N)$.
Equation (D23.F36.1) forces $c=1$, and hence $F=F_0$.
This proves uniqueness over the entire stated fiber; no assumption
about the size of the separate centralizer $C_N(B)$ is required.

Finally the actual first associator value belongs to $P'$ and
satisfies the two-cycle. Its $448$-point value is fixed by D23.P21,
and its signed $24$-point value is the identity in the normalized
state by D23.F32. Its covariance is D23.P18. These are precisely
(D23.F36.2), so the actual-source conclusion follows. The
three-cycle, already checked for $F_0$ in D23.F34, is not needed
for the uniqueness argument.
\end{proof}

\paragraph{Certificate and exact scope.}
The replay is
\path{agents/five_level_six/hurwitz_first_value_uniqueness.py},
with its JSON record. It rebuilds $N$ from the
\texttt{derived\_5} generators in the complete D23.F35 record,
checks its order and trivial center, checks the literal swapping
conjugation by $s$, and replays the covariance and duality identities
for all $72$ accepted candidate instances. The group-theoretic
uniqueness proof applies to every possible $F$, not just those
explicitly saved by a previous enumeration. It does not by itself
assert that any surviving $V$ comes from an actual full-GT element.

\subsection{The exact remaining pentagon value}

\begin{statement}{D23.P25}{Proved}{The required value at the fourth--fifth pair is a commutator}
\end{statement}
Use the positive braid marking
\[
 X_1=A^2,\quad X_2=B^2,\quad X_3=C^2,\quad
 X_4=(ABA)^2,\quad X_5=(CBC)^2,
\]
and put $D=ABCABA$. Let $f$ satisfy the two-cycle, and
write $F=f(X_1,X_2)$ and $R=f(X_3,X_4)$. In the
five-factor pentagon order of supplied FP.37.3, the
pentagon equation is equivalent to
\[
 \boxed{\quad f(X_4,X_5)=[D,FR]
            =DFRD^{-1}R^{-1}F^{-1}.\quad}                 \tag{D23.P25.1}
\]
Here $[a,b]=aba^{-1}b^{-1}$. The formula keeps the actual
inverse of $D$ and remains valid in the raw braid action,
where $D^2$ is the nontrivial full-twist deck generator.

\begin{proof}
The reversing half twist satisfies
\[
 DX_1D^{-1}=X_3,\quad DX_2D^{-1}=X_2,\quad
 DX_4D^{-1}=X_5.
\]
It also sends $X_3$ back to $X_1$, since $D^2$ is
central. Equivariance of evaluation and the two-cycle give
\[
 f(X_5,X_1)=DR^{-1}D^{-1},\qquad
 f(X_2,X_3)=DF^{-1}D^{-1}.
\]
Thus, putting $H=f(X_4,X_5)$, the literal pentagon is
\[
 \mathcal P(f)
   =R(DR^{-1}D^{-1})(DF^{-1}D^{-1})HF
   =RDR^{-1}F^{-1}D^{-1}HF.
\]
Solving $\mathcal P(f)=1$ for $H$ gives
$H=DFRD^{-1}R^{-1}F^{-1}$, exactly (D23.P25.1).
No involutivity assumption on $D$ was used.
\end{proof}

\paragraph{Exact replay and remaining test.}
\path{hurwitz_required_f45.py} verifies all displayed
conjugacies in both the 1344-point and 8064-point actions,
and verifies $D^2$ equals the literal full twist. For
all 36 recorded D23.H7 witnesses it computes the
required value in (D23.P25.1) and checks the original
five-factor product entry by entry. The raw required
value projects to the base required value in every case.
The NPZ stores all these required values.

These identities do not prove that the required value
belongs to $\langle X_4,X_5\rangle'$, or that a common
word realizes $F,R$ and that value while satisfying
the gentle identities. They give the exact remaining
value to impose in those tests.

\subsection{A necessary twelve-candidate reduction from the remaining pair}

\paragraph{D23.H8 (proved; whole-proof check completed).}
Normalize a hypothetical mixed actual full-GT element as in
D23.H2 to $(\rho,a_3,a_2)=(0,1,0)$. Use the zero-based
candidate indices in D23.P22 and D23.F34. Its complete
degree-$1344$ conjugating permutation $V$ must have one of
the following twelve indices:
\[
       0,1,18,19,22,23,26,27,52,53,60,61.
\]
For each of these indices, exactly three values of the full
mixed commutator $R=f(X_3,X_4)$ survive the necessary
fourth--fifth pair membership condition. The complete list is
specified in this table. The second column is the index in
the stored $216$-element centralizer subgroup, with the
marking defined below.
\[
\begin{array}{c|l@{\qquad}c|l}
 V\text{ index}&\text{three indices}&V\text{ index}&\text{three indices}\\\hline
 0&108,128,136&1&109,129,137\\
 18&113,121,141&19&112,120,140\\
 22&41,49,69&23&40,48,68\\
 26&45,53,61&27&44,52,60\\
 52&116,124,132&53&117,125,133\\
 60&80,88,96&61&81,89,97
\end{array}
\]
All other twenty-four candidates of D23.F34 are excluded by
this necessary condition. The statement does not assert
existence of an actual element for any surviving pair.

\begin{proof}
Let $F=f(X_1,X_2)$. D23.F36 is essential here: it proves
\emph{uniqueness} of the full first value, and hence
$F=F_V$ for an actual normalized element. An enumeration
for one convenient first value alone would not justify
the present exclusion.

The full mixed covariance calculation D23.C26 gives
$432$ possible values of $R$ for each of the $36$
candidates, in a right coset $r_VC_*$ with
\[
                    C_*=C_{K\cap M'}(C),\qquad |C_*|=432.
\]
Use the marked common-word group $J'$ of D23.H7 and put
\[
 C_{216}=\{c\in C_*:(1,c)\in J'\}.
\]
D23.H7 proves that this is an index-two subgroup of $C_*$.
For each candidate, its stored common-word representative
$R_V$ then gives the \emph{complete} remaining fiber
\[
                         R_V C_{216}.
\]
This is necessary because the actual first and mixed values
are evaluations of the same element of $\widehat F_2'$.
In a finite target, the image of that profinite commutator
subgroup is the ordinary derived image: the ordinary free
group is dense, and finite images are discrete. Thus the
exact finite $J'$ membership test applies to the actual
profinite value without assuming an ordinary-word lift in
the free group itself.

Construct $C_{216}$ by enumerating all $432$ elements of
$C_*$ and testing the literal permutation
$(1_{472},c)$ in the exact $1816$-point group $J'$.
There are exactly $216$ accepted elements; their generated
group has order $216$. Save this complete ordered list.
The order here is merely a certificate index, with no
arithmetic meaning.

Set $D=ABCABA$, $X_4=(ABA)^2$, and $X_5=(CBC)^2$.
The full marked pair group $M_{45}=\langle X_4,X_5\rangle$
has exact orders
\[
 |M_{45}|=13374150672384,\qquad
 |M_{45}'|=557256278016.
\]
These are obtained directly from the two degree-$1344$
permutations, followed by their normal commutator subgroup
and exact stabilizer chains.

For each of the $36\cdot216=7776$ possible pairs $(V,R)$,
compute the necessarily prescribed value
\[
                        H=[D,F_VR].
\]
By D23.P25, an actual full-GT element must satisfy
$H=f(X_4,X_5)$, and hence $H\in M_{45}'$.
Exact membership in the full group $M_{45}'$ accepts
precisely the $36$ pairs in the displayed table, three
for each of twelve candidates and zero for the other
twenty-four. The original five-factor pentagon product
is checked directly for every accepted pair, using the
actual inverse of $D$.

Every stage has exhausted a complete necessary finite
fiber: the $36$ choices of $V$, the unique full first
value from D23.F36, and all $216$ common-word mixed
values. Consequently the rejected candidates cannot be
rescued by another first value or another value inside
the same mixed commutator fiber. Further same-word
conditions involving $M_{45}$, the other gentle
substitutions, and the raw cover are still additional
conditions on the survivors.
\end{proof}

\paragraph{Exact certificates.}
The full pair group is constructed by
\path{code/hurwitz_pentagon45_probe.py}; the exhaustive
fiber test is
\path{code/hurwitz_pentagon45_fibers.py}.
The resulting JSON and NPZ records retain every accepted
$R$, its prescribed $H$, and its centralizer index.
The ordered centralizer is
\path{agents/characters/records/hurwitz_commonword_centralizer216.npz},
and the base representatives are the $R$ witnesses in
\path{records/hurwitz_solvable_fiber_filter.npz}.
The proof uses the complete groups, not separate orbit
membership as a substitute for joint membership.

\subsection{Retaining the raw product and the sixfold deck action}

\paragraph{D23.H6 (raw product-fixed Hurwitz cover).}
There is a transitive degree-$8064$ Hurwitz action of $B_4$
lying over the degree-$1344$ quotient of D23.P17. Its fibers
have size six, with cyclic deck group generated by the
literal full twist $z_4=(ABC)^4$. The actual multiplier-one
full-GT action preserves this raw orbit, commutes with its
deck group, and satisfies the same exact three-generator
covariance equations as D23.P18, now in the raw action.

This statement does not identify the kernel of the raw
permutation image over the degree-$1344$ image with a
single central $C_6$. That kernel may have additional
fiber translations. The global full twist supplies its
specified constant cyclic subgroup.

\begin{proof}
Let $E_0=\operatorname{ASL}_2(\mathbb F_3)$ and let
$\mathcal C$ be the size-$36$ order-six conjugacy class
used in D23.P17. Take its first recorded four-tuple and
let $p$ be its ordered product. The exact affine
multiplication table gives
\[
             |p|=6,\qquad C_{E_0}(p)=\langle p\rangle.
\]
For each of the $1344$ recorded simultaneous-conjugacy
classes choose its recorded tuple $t$, of product $p_t$,
and choose $g$ with $gp_tg^{-1}=p$. All choices form a
coset of the displayed centralizer. Conjugating $t$ by
these six choices gives six distinct ordered tuples.
The certificate checks distinctness within every fiber
and across all $1344$ fibers. Consequently the complete
preimage with product exactly $p$ consists of $8064$
raw tuples.

Hurwitz moves act on raw adjacent entries by
$(u,v)\mapsto(uvu^{-1},u)$. They preserve the ordered
product exactly. Applying all three moves to all $8064$
tuples therefore gives the stated raw permutation action.
The certificate verifies the Artin relations and its
projection to the original three generators. On raw
tuples the full twist is simultaneous conjugation by
their product $p$; in the recorded fiber numbering this
is precisely the shift by one in $\mathbb Z/6$. It is
transitive on every fiber. Since the base action is
transitive, the raw action is transitive as well.

For the actual-source assertion use the exact standard
split mixed braid sequence in D23.P18, before taking
simultaneous-conjugacy classes of representations.
Its free kernel is $\widehat F_4$, with peripheral
generators $u_1,\ldots,u_4$. Let $\kappa_h$ be the
restriction of the Drinfeld automorphism. The same
semidirect-product identity gives
\[
       \kappa_h\alpha_\beta\kappa_h^{-1}
                        =\alpha_{\Phi_h(\beta)}
\]
as actual automorphisms of this free kernel, for the
specified standard section. Hence precomposition by
$\kappa_h^{-1}$ gives exact Hurwitz covariance already
on raw representations. Each $u_i$ is sent to a conjugate
of itself when the multiplier is one. Moreover D23.P19
shows that $u_1u_2u_3u_4=z_4^{-1}z_5$ is fixed exactly.
Thus precomposition preserves both the four chosen
conjugacy classes and the exact product $p$.

After passage to simultaneous-conjugacy classes this is
the actual action of D23.P18, which preserves the unique
degree-$1344$ orbit. Its full product-$p$ preimage is the
raw orbit just constructed, so that orbit is preserved.
Precomposition commutes with every postcomposition by
an inner automorphism of $E_0$, in particular with
conjugation by $p$. This proves the deck commutation
and the stated raw covariance. Inverses give bijectivity.
\end{proof}

\paragraph{An exact initial lift test.}
Fix the cyclic fiber numbering $(i,k)$, $k\in\mathbb Z/6$,
from the construction. A braid generator acts as
\[
              (i,k)\longmapsto(q(i),k+e_q(i)).
\]
A deck-equivariant lift of a base permutation $V$ has
the form $(i,k)\mapsto(V(i),k+t_i)$. Its commutation
with that braid permutation is equivalent to
\[
             t_{q(i)}-t_i=e_q(V(i))-e_q(i)\pmod6.
\]
Propagating these equations for $A$ and the nested full
twist $(AB)^3$ on each connected base orbit is an
exhaustive test: choose one value $t_i=0$, propagate,
and reject precisely a conflicting cycle. Every other
solution differs by a constant on that connected orbit.
All $72$ base candidates of D23.P22 pass for either
cubic bit. Thus these commutation equations alone
give no further exclusion. The canonical linear action
on cocycles and the remaining raw covariance equations
are additional tests, not consequences of this initial
potential calculation.

The full construction, permutation arrays, fiber phases,
and all passing potentials are retained in
\path{code/hurwitz_product_fixed_cover.py} and
\path{records/hurwitz_product_fixed_cover.npz}, with its
JSON record. The certificate proves the finite claims;
the actual-source stability is the preceding argument.

\subsection{Canonical linear lifts of the product-fixed cover}

\begin{statement}{D23.P23}{Proved}{The stronger raw linear lift test remains nonempty}
\end{statement}
In the product-fixed degree-8064 cover of D23.H6, the
preimage of each projective 112-point block splits into
three exact linear-monodromy framings, of 224 points each.
Each framing is the set of nonzero singular vectors in a
six-dimensional $\mathbb F_3$ cocycle space. Its two-to-one
map to the projective block identifies opposite vectors.

The actual multiplier-one full-GT action is linear between
these framed cocycle spaces. For a fixed projective
permutation on a block, every possible linear lift commuting
with the cyclic deck group differs from any one such lift
by a constant deck shift in $C_6$ on that entire block.
Thus there are exactly six possible local lifts.

For each of the 72 global necessary candidates in D23.P22,
and for both cubic-bit control cases, the six-choice local
lifts admit a global assembly commuting with the raw
generators $A$ and $z_3=(AB)^3$. In fact there are
$6^4=1296$ such assemblies per candidate. This stronger
linearity condition therefore gives no exclusion at this
stage. Full raw covariance and common-word conditions
remain additional requirements.

\begin{proof}
Write the fixed affine product as $p=(v,P)$. Since $I-P$
is invertible, $p$ has the unique fixed point
\[
                         q_\infty=(I-P)^{-1}v.
\]
Conjugating all affine maps by the translation taking
$q_\infty$ to zero changes the product to $(0,P)$.
Fix an exact ordered four-tuple of linear parts
$\rho(t_i)=A_i$, with $A_1A_2A_3A_4=P$.
The translations of its affine lifts are left cocycles
$c$ for this representation. The exact product condition
is the linear equation
\[
                       c(t_1t_2t_3t_4)=0.                 \tag{D23.P23.1}
\]
The space of cocycles has dimension eight. The map in
(D23.P23.1) is onto: on a coboundary it takes the value
$(I-P)q$, and $I-P$ is invertible. Its kernel therefore
has dimension six and intersects the translation
coboundaries trivially. Projection identifies that kernel
linearly with the six-dimensional cohomology quotient in
D23.P19. In particular the fixed-point differences
$(q_1-q_4,q_2-q_4,q_3-q_4)$ are linear coordinates on it.
Returning from the translated origin leaves these
difference coordinates unchanged.

The linear centralizer of $P$ in $SL_2(\mathbb F_3)$ is
cyclic of order six. Its element of order two is $-I$;
the simultaneous centralizer of the four linear parts is
exactly $\{\pm I\}$. Thus the linear-part class underlying
one projective block has three exact framings with product
$P$. The raw construction in D23.H6 has precisely six
points above every projective point. Partitioning those
672 points by their exact four linear matrices gives
three sets of 224 points. In each set the difference
coordinates are 224 distinct nonzero vectors, in opposite
pairs, and their projectivizations are exactly the 112
singular points already identified in D23.C23/P19.
This also gives a direct finite verification of the
framing description, without deducing it from cardinality
alone.

Now let $\kappa_h$ be the actual free-kernel automorphism
of D23.P18. It fixes the boundary product exactly at
multiplier one. Precomposition consequently carries
(D23.P23.1) into the analogous kernel for the transformed
linear representation. It acts linearly on cocycles.
Before the translation of the fixed product point, the
canonical origin is the affine lift all of whose fixed
points are $q_\infty$; precomposition preserves that
origin for the transformed representation. Thus there
is no hidden affine translation in the six-dimensional
map. The transformed exact linear representation depends
only on the initial linear representation, not on its
cocycle. An actual map therefore carries an entire
framing linearly to one entire framing.

The deck generator is simultaneous conjugation by $p$.
It cyclically permutes the three framings; its cube fixes
each framing and negates its difference-coordinate vectors.
On one chosen framing a prescribed projective map has
exactly two linear representatives, differing by $-I$.
The target framing has three choices. Deck equivariance
extends either linear choice from the chosen framing
uniquely to the other two. Consistency under the deck cube
holds because a linear map takes $-w$ to the negative of
its value on $w$. Hence these are exactly six lifts, and
the deck group acts simply transitively on them.

Here is the exact exhaustive calculation used in the
certificate. For each block choose the lexicographically
first exact linear framing and list its 224 difference
vectors. Choose six linearly independent vectors. For a
given projective permutation, each of their images has
two vector representatives. Fix the sign of the first
image and enumerate the remaining $2^5$ signs. Exactly
one linear matrix carries all 224 vectors to the prescribed
projective images. This gives one canonical lift, which
is then extended by deck equivariance as above.

In the raw numbering $(i,k)$ of D23.H6, this lift has
the form
\[
                  (i,k)\longmapsto(V(i),k+p_i).
\]
Every other canonical linear lift is obtained by replacing
$p_i$ by $p_i+d_{b(i)}$, where $b(i)$ is its one of twelve
projective blocks and $d_j\in\mathbb Z/6$ is constant
on that block. If a raw braid acts by
$(i,k)\mapsto(q(i),k+e(i))$, commutation is equivalent to
\[
 d_{b(q(i))}-d_{b(i)}
        =e(V(i))-e(i)-p_{q(i)}+p_i\pmod6.                 \tag{D23.P23.2}
\]
The certificate first checks that the right side is
constant on every source block, for each of $A,z_3$.
It then solves all resulting equations on the twelve
blocks by exact propagation and rejects an inconsistent
cycle. Every one of the 72 candidates passes, in each
cubic-bit control case. The resulting literal raw
permutation is checked against $A,z_3$ and the deck
generator on all 8064 points.

The block graph for $A,z_3$ has the four connected
components displayed in D23.P22. Once a solution to
(D23.P23.2) exists, all solutions differ by an arbitrary
constant in $\mathbb Z/6$ on each component. Therefore
there are exactly $6^4$ assemblies per candidate. The
proof and computation have imposed linearity, deck
equivariance, and these two commutations; they have not
silently imposed the remaining raw braid covariance.
\end{proof}

\paragraph{Certificate and scope.}
\path{hurwitz_raw_linear_lifts.py} reconstructs all thirty-six
linear framings from the literal raw affine tuples in
D23.H6. It verifies the coordinates, opposite-vector
deck cube, unique projective matrix lifts, block-constant
phase equations and final exact raw permutations.
The JSON record reports every passing candidate; the
NPZ retains a particular phase solution and its block
shifts for each. The proof of actual raw stability is
D23.H6, and the linearity argument above uses its exact
standard-section normalization, not a hypothetical
coordinate-forgetting compatibility.

\subsection{All recorded mixed states survive full raw braid covariance}

\begin{statement}{D23.P24}{Proved}{Exact raw simultaneous covariance for the 36 recorded witnesses}
\end{statement}
For every one of the 36 normalized mixed candidate
permutations retained in D23.F34 and D23.H7, take the
recorded first value $F$ and the recorded common-word
mixed value $R$ of D23.H7. Their unique lifts to the
product-fixed raw degree-8064 representation admit a
canonical linear lift $\widetilde V$ satisfying
\[
\begin{split}
 \widetilde V\widetilde A\widetilde V^{-1}
     &=\widetilde A,\\
 \widetilde V\widetilde B\widetilde V^{-1}
     &=\widetilde F^{-1}\widetilde B\widetilde F,\\
 \widetilde V\widetilde C\widetilde V^{-1}
     &=\widetilde R\widetilde C\widetilde R^{-1}.
\end{split}                                                \tag{D23.P24.1}
\]
It also commutes with the raw nested full twist and the
cyclic deck group. There are exactly six such lifts among
the canonical linear lifts of P23, differing by the global
deck group, for each recorded $(V,F,R)$.

Thus this full raw covariance requirement does not reject
any of these 36 recorded mixed witnesses. This statement
does not assert the full gentle identities at every
substituted pair, the remaining pentagon-pair common-word
conditions, or a lift to an actual profinite GT element.

\begin{proof}
The raw first-three-strand image has order
\[
                     3134566563840000,
\]
the same as its degree-1344 image in D23.F34. The natural
projection is onto, so equality of orders makes it an
isomorphism of the marked groups. In particular the
pure-pair subgroups and their derived subgroups also map
isomorphically. The analogous raw mixed-pair image
$\langle\widetilde C^2,\widetilde Z_3\rangle$ has order
\[
                        40122452017152,
\]
again the same as its base image. Its projection is
likewise a marked isomorphism. The first order is computed
in \path{hurwitz_raw_first_group.py}; the second is the
independent root record
\path{work/records/hurwitz_raw_mixed_group.json}.

Consequently each of the recorded $F,R$ has a unique raw
lift, denoted $\widetilde F,\widetilde R$. Any common
ordinary word realizing the base pair realizes these unique
raw values as well. No new choice of associator value is
introduced at this step. The certificate computes the lifts
by evaluating the strong-generator straight-line programs
of the base marked groups in the raw generators. Every
computed strong-generator lift is checked to project
literally to its base permutation; memoization avoids
expanding possibly long words. This is an exact group
calculation, not a search over possible fiber phases of
$F$ or $R$.

Start with the particular canonical linear raw lift of
$V$ from D23.P23. Write its fiber phase as $t_i$. All
remaining lifts commuting with $A$ and $Z_3$ are obtained
by adding constants $d_0,d_1,d_2,d_3\in\mathbb Z/6$
on the four block orbits
\[
 \{0,9\},\quad\{1,2,3,4,5,6\},\quad
 \{7,8,10\},\quad\{11\}.
\]
Let $o(i)$ denote the corresponding orbit number of the
base point $i$. If a raw generator is
$(i,k)\mapsto(q(i),k+e(i))$ and its required transform
is $(j,k)\mapsto(q'(j),k+e'(j))$, where
$q'=VqV^{-1}$, the covariance equation is exactly
\[
 d_{o(q(i))}-d_{o(i)}
     = e'(V(i))-e(i)-t_{q(i)}+t_i \pmod6.                 \tag{D23.P24.2}
\]
For $B$ use the required transform
$\widetilde F^{-1}\widetilde B\widetilde F$; for $C$
use $\widetilde R\widetilde C\widetilde R^{-1}$.

The certificate imposes (D23.P24.2) at every one of the
1344 base points. The $B$ equations are consistent for
all 36 candidates and leave two independent constants.
Adding the $C$ equations for each of the recorded H7
witnesses is also consistent in all 36 cases. The graph
is then connected, leaving just one constant in
$\mathbb Z/6$. This constant is the global deck shift,
which commutes with every raw braid generator. Hence
there are exactly six solutions per recorded witness.

Finally the resulting $\widetilde V$ is reconstructed
as an actual permutation of all 8064 points, and all
three equations (D23.P24.1), commutation with $Z_3$, and
deck equivariance are verified entry by entry. The
construction remains within the canonical linear lifts
because only block-constant deck shifts have been added.
This proves the stated existence and count, with precisely
the finite scope claimed.
\end{proof}

\paragraph{Records and remaining requirement.}
\path{hurwitz_raw_simultaneous_covariance.py} is the
complete replay. Its JSON lists all 36 tests and the
two-to-one reduction in the number of free phase
components. Its NPZ retains one actual raw permutation
$\widetilde V$, and the exact $\widetilde F,\widetilde R$,
for every retained base candidate.

The first values are D23.F34's
\path{five_level_six/hurwitz_first_pair_covariance.npz};
the mixed witnesses are D23.H7's
\path{work/records/hurwitz_solvable_fiber_filter.npz}.
That common-word input concerns the first pair and the
full mixed pair $(X_3,X_4)$. It does not by itself
specify $f(X_4,X_5)$ or all remaining pentagon factors.
Those requirements and the complete gentle closure are
still additional constraints after the result here.

\subsection{The solvable quotient obtained from the raw orthogonal blocks}

\begin{statement}{D23.F37}{Proved}{The raw block quotient restores the cubic Burau scalar}
\end{statement}
In the raw degree-$8064$ action of D23.H6, lift the twelve
projective blocks of D23.F30. Each lifted block has $672$ points.
The permutation image $H$ of the stabilizer of one block has
\[
 |H|=39191040,\qquad |H'|=|H''|=6531840,\qquad H/H'\cong C_6.
\]
The induced action on twelve copies of the regular coset set
$H/H'$ is a degree-$72$ quotient $Q_{72}$ of the raw braid image.
Its derived series has orders
\[
              331776,55296,13824,6912,768,1.
\]
It is exactly, with the original three braid generators marked,
the literal matrix group generated by FP.36.2 over
$R=\mathbb Z/4[\omega]/(\omega^2+\omega+1)$, before quotienting
by scalar matrices. The natural map to D23.F31's signed
$24$-point quotient has central kernel
\[
                      \langle\omega I\rangle\cong C_3,
\]
generated by the four-strand full twist $(ABC)^4$.
This kernel is contained in the derived subgroup: the two derived
orders are $55296$ and $18432$, respectively. In particular it
would be incorrect to describe this extension as adding only an
abelian character of the whole braid image.

Nevertheless every one of the five marked evaluation pair images
maps isomorphically to its signed degree-$24$ counterpart:
\[
\begin{array}{c|rrrrr}
\text{pair}&12&23&34&45&51\\\hline
\text{image order}&192&192&144&48&144.
\end{array}
\]
Thus no value at any one of these five pairs gains an independent
scalar choice in the raw solvable quotient. The marked cubic
$216$ quotient still does not factor through its first-pair image,
whereas the marked harmonic $S_4$ quotient does factor through
its $34$ image. This concerns the induced solvable quotient only,
not all information in the raw degree-$8064$ action.

\begin{proof}
Write $q$ for the block permutation of a raw braid permutation
$g$, and choose twelve raw braid transversals $T_i$ from block
zero to block $i$. The certificate obtains them by an exhaustive
traversal of the twelve-block action. The Schreier elements
\[
                       T_{q(i)}^{-1}gT_i
\]
for the three generators and all twelve blocks generate the
stabilizer. Restricting these elements to the original $672$-point
block gives thirteen distinct nonidentity generators of its full
image $H$. The exact permutation-group order and two successive
derived orders are as displayed. Equality $H'=H''$ proves
perfectness of $H'$. The quotient has order six and is abelian,
so it is cyclic.

Let $L(g,i)$ denote the restriction of
$T_{q(i)}^{-1}gT_i$ to the base block. On pairs $(i,hH')$ define
\[
                  g(i,hH')=(q(i),L(g,i)hH').
\]
The identity
$L(gg',i)=L(g,q'(i))L(g',i)$ proves the action law.
Since $H'$ is normal in $H$, this is well defined. Enumerating
the six cosets by exact membership in $H'$ therefore constructs
an actual quotient of the raw image, not an independently chosen
finite representation. The three resulting degree-$72$
permutations generate a group of the stated order and complete
derived series.

For the matrix identification, encode $a+b\omega\in R$ by
$a+4b$, $0\leq a,b<4$, and use the literal FP.36.2 matrices
\[
 A=\begin{pmatrix}1-\omega&\omega&0\\1&0&0\\0&0&1\end{pmatrix},
 \quad B=\begin{pmatrix}1&0&0\\0&1-\omega&\omega\\0&1&0\end{pmatrix},
 \quad C=\begin{pmatrix}1&0&-1\\0&1&-1\\0&0&-\omega\end{pmatrix}.
\]
Traverse the full degree-$72$ permutation group while multiplying
these three matrices along the identical generator edges. There
are $331776$ distinct permutation elements. Every one of the
$995328$ generator edges gives the same matrix whenever a
permutation is reached again, and all $331776$ resulting
matrices are distinct. This gives a well-defined bijective marked
homomorphism onto the literal matrix group. It is an exact
relation check, not recognition by matching group orders.
The full twist maps to $\omega I$.

D23.F31 identifies the signed quotient with the same matrix group
modulo $\langle\omega I\rangle$. Since both identifications use
the same three generators, their diagram commutes. Hence the
projection has precisely the displayed central kernel. The
certificate independently verifies this projection by joining
the two marked permutation triples: the joint image has order
$331776$, equal to $|Q_{72}|$. It also computes the two derived
orders. Their quotient is three, which proves that the entire
kernel lies in $Q_{72}'$.

Finally form each of the five literal pairs
$X_1=A^2$, $X_2=B^2$, $X_3=C^2$, $X_4=(ABA)^2$,
$X_5=(CBC)^2$ in both quotients. For every pair the two individual
orders and the order of their marked joint image are equal to
the corresponding entry in the table. Both projections are onto,
so equality proves that each pair projection is an isomorphism.
In particular the same word at that pair has a unique raw
solvable value above its signed value.

For the final marked quotient tests, join the $12$ pair with the
explicit nine-point affine $216$ detector of D23.P17. Its joint
order is $3456$, strictly larger than $192$, so that marked map
does not factor through the first-pair image. Joining the $34$
pair with the literal $S_4$ pair $((123),(1234))$ leaves order
$144$, so its marked map does factor. These are tests of the
specified two-generator maps, not statements inferred only from
the abstract orders of potential target groups.
\end{proof}

\paragraph{Exact certificates and limitation.}
The construction and derived groups are in
\path{agents/five_level_six/hurwitz_raw_block_solvable.py};
the marked quotient and all five pair checks are in
\path{hurwitz_raw_solvable_markings.py}; the complete matrix
isomorphism replay is \path{hurwitz_raw_burau.py}.
Each has its JSON record, and the block construction also retains
all transversals, cosets, local transitions and generators in NPZ.
No abstract identification of $H$ from its order is required,
and no assertion that this is the maximal solvable quotient of
the entire raw braid image is made. Raw perfect-kernel and
global common-word constraints remain available beyond this
particular quotient.

\subsection{The raw pentagon leaves six normalized candidates}

\begin{statement}{D23.P26}{Proved}{An exact six-candidate necessary reduction}
\end{statement}
Normalize a hypothetical actual full-GT element with nonzero
$\Delta$ by D23.H2 to $(\rho,a_3,a_2)=(0,1,0)$.
Use the zero-based indices in the array
\texttt{mixed72\_V\_1} of D23.P22. Its degree-$1344$
conjugating permutation must have one of the six indices
\[
                         1,19,23,27,53,61.
\]
For each such candidate exactly three of D23.H8's mixed
values survive the literal degree-$8064$ pentagon. The full
retained list, indexed in D23.H8's ordered $216$-element
centralizer subgroup, is
\[
\begin{array}{c|l}
V\text{ index}&\text{mixed centralizer indices}\\\hline
1&109,129,137\\
19&112,120,140\\
23&40,48,68\\
27&44,52,60\\
53&117,125,133\\
61&81,89,97
\end{array}
\]
Thus there are eighteen remaining pairs $(V,R)$ in this
necessary test. No assertion that any pair is realized by
an actual gentle or full-GT element is made.

\begin{proof}
We first justify why a full-GT element must satisfy the
\emph{literal} pentagon in the raw braid action, although
that action does not factor through $B_4/\langle z_4\rangle$.
Write $a,b,c$ for the positive braid generators and put
\[
\begin{aligned}
X_1&=a^2,&X_2&=b^2,&X_3&=c^2,\\
X_4&=(aba)^2,&X_5&=(cbc)^2,&z_4&=(abc)^4.
\end{aligned}
\]
Supplied FP.30.1 and FP.30.3 identify their images with
the five standard cyclic boundary twists and identify the
full-GT five-cycle in the order
\[
 \mathcal P(f)=f(X_3,X_4)f(X_5,X_1)f(X_2,X_3)
                    f(X_4,X_5)f(X_1,X_2).
                                                        \tag{D23.P26.1}
\]
The image of this product in
$\widehat{B_4/\langle z_4\rangle}$ is one.
The kernel of the natural map from $\widehat B_4$ is
the closed normal subgroup generated by $z_4$, hence its
procyclic central closure. It is canonically
$\widehat{\mathbb Z}$: the continuous exponent-sum map
$e:\widehat B_4\to\widehat{\mathbb Z}$ takes $z_4$ to
$12$, so the map $t\mapsto z_4^t$ is injective, as
multiplication by $12$ is injective on
$\widehat{\mathbb Z}=\prod_p\mathbb Z_p$. The same
central-kernel and exponent-sum argument is recorded in
supplied FP.26, paragraph \emph{Lifting the five-point
frame}.

Consequently $\mathcal P(f)=z_4^t$ for some
$t\in\widehat{\mathbb Z}$. Each factor has exponent
sum zero, since $f\in\widehat F_2'$ and an evaluation
followed by $e$ is an abelian homomorphism. Applying $e$
therefore gives $12t=0$, whence $t=0$. This proves
\[
                   \mathcal P(f)=1\quad\hbox{in }\widehat B_4.
                                                        \tag{D23.P26.2}
\]
In particular its degree-$8064$ raw image must be the
identity. This conclusion is stronger than identity
modulo the global deck group and requires no claim that
the raw action descends to the five-point quotient.

Next fix the complete first value. By D23.F36 an actual
normalized element with a candidate $V$ has
$F=f(X_1,X_2)=F_V$, uniquely, not merely modulo a
solvable quotient. The proof of F36 uses the trivial
center of the full perfect residual and the two-cycle,
so no unenumerated first-value fiber remains. D23.H7
and the exhaustive D23.H8 calculation then leave exactly
the twelve candidates and their three mixed values listed
in H8. These are all the possibilities after the necessary
base pentagon-pair membership test.

Let tildes denote the raw permutations. The projection of
the raw first-three-strand group to its degree-$1344$
image is a marked isomorphism by D23.P24, as is that of
$\langle\widetilde X_3,\widetilde X_4\rangle$.
For the remaining pair the direct stabilizer-chain
calculations give
\[
\begin{split}
|\langle\widetilde X_4,\widetilde X_5\rangle|
  &=|\langle X_4,X_5\rangle|=13374150672384,\\
|\langle\widetilde X_4,\widetilde X_5\rangle'|
  &=|\langle X_4,X_5\rangle'|=557256278016.
\end{split}                                             \tag{D23.P26.3}
\]
The projection is onto because it sends the two marked
generators to their base counterparts. Equality of finite
orders makes it an isomorphism, including on the derived
subgroups. Thus each possible base value $F,R,H$ has a
\emph{unique} marked raw lift. In particular an actual
evaluation of the same profinite $f$ must have exactly
those raw values; there is no separate scalar or deck
choice for an associator value.

Use $D=abcaba$, retaining its actual inverse. By D23.P25,
the base value of $H=f(X_4,X_5)$ is necessarily
$[D,F_VR]$, and H8 has checked its membership in the
full derived pair group. Write $\widetilde H$ for its
unique lift from that marked group. Equation (D23.P26.2)
is equivalent to the further literal identity
\[
 \widetilde H=
 [\widetilde D,\widetilde F_V\widetilde R].
                                                        \tag{D23.P26.4}
\]
The certificate computes both sides for every one of the
$12\cdot3=36$ complete H8 pairs. It evaluates the
strong-generator straight-line programs in the raw marked
generators, with every lifted permutation checked against
the base projection. It also directly multiplies the
original five factors in (D23.P26.1), so (D23.P26.4)
is not checked merely by a rearranged symbolic convention.

Exactly all three mixed values pass for each odd index
$1,19,23,27,53,61$. All three fail for each even partner
$0,18,22,26,52,60$. In every failing case the five-factor
product has order two, projects to the base identity,
and has fiber phases $0$ and $3$, both occurring.
It is therefore a nonconstant fiber permutation, not a
global power of the order-six deck transformation.
Both descriptions are recorded as additional exact checks;
nonidentity alone is enough to violate (D23.P26.2).
The identity associator gives the identity raw product as
an equal-sign control. The retained indices and mixed
values are precisely the displayed table.

The completeness assertion now follows from F36, H8,
and the unique raw lifts: any actual normalized mixed
element must appear among those thirty-six pairs, and
its raw pentagon removes exactly the rejected eighteen.
The eighteen survivors remain subject to all common-word
conditions involving the third pair and all other actual
gentle or full-GT conditions.
\end{proof}

\paragraph{Certificates and source locators.}
The pair-group orders and marked projection are rebuilt
by \path{agents/pentagon/hurwitz_raw_pair45.py}; its JSON
and NPZ retain both marked generator pairs. The complete
raw test is
\path{agents/pentagon/hurwitz_raw_pentagon45.py}, with
the JSON table and NPZ arrays of required values, unique
marked lifts, and all thirty-six literal pentagon products.
The strong-generator evaluator is
\path{agents/pentagon/hurwitz_marked_raw_lift.py}.
The base exhaustive fiber is
\path{code/hurwitz_pentagon45_fibers.py} and its records;
the first-value uniqueness input is
\path{agents/five_level_six/hurwitz_first_value_uniqueness.tex}.
The supplied universal marking and central-kernel input are
\path{sources/DELTA_ON_GT/agent_five/five_point_results.tex},
FP.30.1, FP.30.3, and FP.26's paragraph
\emph{Lifting the five-point frame}. The exact raw orbit
and its actual multiplier-one action are D23.H6.

\subsection{The eighteen raw-pentagon survivors also satisfy raw covariance}

\begin{statement}{D23.P27}{Proved}{No further exclusion from simultaneous raw braid covariance}
\end{statement}
Every one of the eighteen pairs $(V,R)$ retained by
D23.P26 has a canonical linear raw lift satisfying
\[
\begin{split}
 \widetilde V\widetilde A\widetilde V^{-1}&=\widetilde A,\\
 \widetilde V\widetilde B\widetilde V^{-1}
   &=\widetilde F_V^{-1}\widetilde B\widetilde F_V,\\
 \widetilde V\widetilde C\widetilde V^{-1}
   &=\widetilde R\widetilde C\widetilde R^{-1}.
\end{split}
\]
It commutes with the raw nested full twist and the deck
group. For each pair exactly six lifts satisfy these
conditions, differing by a global deck transformation.

\begin{proof}
The first and mixed values have unique marked raw lifts
by D23.P24. Use those lifts, with the exact P26 values
of $R$, not the earlier H7 witness for the same $V$.
Start from the canonical linear lift of D23.P23.
Every possible remaining linear lift commuting with
$A$ and $Z_3$ adds a constant in $\mathbb Z/6$ on each
of the four block orbits
\[
\{0,9\},\quad \{1,2,3,4,5,6\},\quad
\{7,8,10\},\quad \{11\}.
\]
The exact covariance equations are D23.P24.2. The
certificate imposes both the $B$ and $C$ equations
at every base point for each of the eighteen pairs.
It records each equation together with its reversed
equation and solves the resulting four-vertex potential
system over $\mathbb Z/6$. In all eighteen cases the
system is consistent and its graph is connected.
There is therefore one free constant, giving exactly
six solutions. This constant is a global deck shift.

The certificate reconstructs a complete permutation of
the 8064 raw points from each solution and checks all
three displayed equations, commutation with $Z_3$, and
commutation with the deck generator literally on every
point. The shifts remain block-constant, so every lift
is still linear in P23's exact framing model. This
proves the statement for the complete P26 list.
\end{proof}

\paragraph{Scope and records.}
The complete replay is
\path{agents/pentagon/hurwitz_raw_survivor_covariance.py}.
Its JSON lists all eighteen cases; its NPZ retains the
raw $V,F,R,H$ for each pair. This result closes a gap
between the earlier covariance witnesses and the new
pentagon survivors, but it is a negative test result:
it removes no remaining candidate and does not assert
that any is an actual profinite GT element. The full
third-pair common-word condition remains separate.

\subsection{The complete six-coset gentle test in the mixed degree-72 image}

\begin{statement}{D23.C27}{Proved}{Eight cycle-compatible residues, already satisfied by C26}
\end{statement}
Let $x,y$ be the component-9 generators of C25 and put
$z=(xy)^{-1}$. Among C25's sixteen mixed residues, exactly
the following eight pass the full induced six-coset
commutator-word and gentle two/three-cycle test, including
the $m=-1$ three-cycle:
\[
                   0,2,4,5,6,7,9,12.                      \tag{D23.C27.1}
\]
Indices refer to the fixed C25 NPZ ordering.  However, all
432 full mixed commutator values for every candidate $V$
in C26 already project into this list.  Thus C27 does not
remove any of those full covariance values, including any
subsequent subset obtained from the first/mixed ordinary-word
condition.  No profinite gentle lifting assertion is made.

\begin{proof}
The joined marked representations $(x,y)$ and $(x,z)$ have
image order 2985984, while their joined derived image has
order 82944, equal to the derived order in either component.
The two derived projections are therefore isomorphisms.
Consequently the value $f(x,z)$ of any commutator word is
uniquely determined by $a=f(x,y)$; write this value as
$\theta(a)$.  The certificate obtains it from the literal
joined derived graph, for every one of the sixteen $a$.

If $f$ satisfies the two-cycle and three-cycle, its six
values on
\[
       (x,y),(x,z),(y,z),(y,x),(z,y),(z,x)
\]
must respectively be
\[
 a,\quad t,\quad ta^{-1},\quad a^{-1},\quad
 (ta^{-1})^{-1},\quad t^{-1},\qquad t=\theta(a).
                                                               \tag{D23.C27.2}
\]
This follows from duality and
$f(z,x)f(y,z)f(x,y)=1$.

Induce the marked degree-72 representation of the free
subgroup to $B=C_2*C_3$ along its six cosets.  The exact
Schreier representatives and rewrite words are those
certified in C11--C22.  The present certificate converts
that right-action construction explicitly to left
permutations and verifies the source marking
\[
          s^2=r^3=1,\quad X=(sr)^{-2},\quad Y=sXs,
          \quad (X,Y)|_{\text{first block}}=(x,y).
\]
The induced free image has order 20542695432781824 and
derived order 570630428688384.  In each of the six
72-point blocks the marked pair is conjugate to exactly
one pair in the preceding displayed list.  The literal
intertwiners are found exhaustively and verified on both
generators, so (D23.C27.2) determines a unique candidate
$F$ on all 432 points.

For each of the sixteen $a$, test membership of $F$ in
the induced derived image and the full matrix equations
\[
             (sF)^2=1,\qquad(F^{-1}r)^3=1.
\]
The latter are the normalized gentle two/three-cycle
relations in the induced $B$ quotient.  In addition test
the $m=-1$ equation
\[
 F(Z,X)Z^{-1}F(Y,Z)Y^{-1}F(X,Y)X^{-1}=1,
 \qquad Z=(XY)^{-1},
\]
where the other two evaluations are obtained by the
explicitly verified cyclic conjugation by $r$ or $r^{-1}$,
according to the recorded marking.  These are equations
on all 432 points, including the reflected cosets; the
test does not assume that an equation in the first
72-point block propagates to the others.
Exactly the eight indices (D23.C27.1) pass.  Each rejected
index fails both the derived-image membership test and
the three-cycle; the eight accepted indices also pass
the negative three-cycle.

Finally C26 gives every full mixed fiber as
$R_0 C_0$, with $|C_0|=432$ and explicit generators.
Enumerate this finite group and project it to component 9.
Its image has order two.  For every one of the 72
representatives $R_0$, both resulting component-9 values
lie in (D23.C27.1), each occurring 216 times.  Direct
evaluation of these complete fibers proves the last
assertion.  In particular any further ordinary-word
restriction on a fiber remains within the accepted list.
\end{proof}

\paragraph{Exact certificate.}
\path{agents/characters/code/hurwitz_mixed72_s3_core.py}
rebuilds the reflection graph, the induced $B$ action,
the six marked intertwiners, all norm and derived-image
tests, and all 72 fiber projections.  Its records are
\path{agents/characters/records/hurwitz_mixed72_s3_core.json},
the NPZ of the same stem, and
\path{agents/characters/records/hurwitz_mixed72_s3_fibers.json}.
The finite source inputs are C25, C26 and the exact
six-coset Schreier data of C22.  This computation imposes
necessary gentle identities and does not replace the
unresolved actual profinite lifting problem by a finite
membership claim.

\subsection{The full mixed six-orientation word and cycle test}

\begin{statement}{D23.C28}{Proved}{All eighteen remaining values pass the full mixed core test}
\end{statement}
Put $x=C^2$, $y=(ABA)^2$, $z=(xy)^{-1}$ and
$M=\langle x,y\rangle$ in the degree-1344 Hurwitz image.
The assignment
\[
                 \theta(x)=x,\qquad\theta(y)=z
                                                               \tag{D23.C28.1}
\]
defines an automorphism of the finite group $M$.
All eighteen raw-pentagon survivors indexed by
\[
        i\in\{1,19,23,27,53,61\},\qquad j\in\{0,1,2\}
\]
have compatible values of one commutator word on all six
peripheral orientations of this full mixed representation.
Those values satisfy the gentle two-cycle, three-cycle and
$m=-1$ three-cycle in the entire induced $C_2*C_3$ quotient.
This is a finite word-and-identity assertion.  It does not
assert that these values lift to an actual gentle element,
nor does it include a finite automorphism-generation or
pure torsion-frame conjugacy claim.

\begin{proof}
The joined marked representations $(x,y)$ and $(x,z)$ have
image order $40122452017152$, equal to $|M|$.
Both projections are onto, since $x,z$ also generate $M$.
Consequently the first projection is an isomorphism and
the second defines the automorphism (D23.C28.1).  In particular
\[
             \theta^2=\operatorname{Ad}(x),\qquad
             \kappa=\operatorname{Ad}(z)\theta
                  \text{ interchanges }y,z.               \tag{D23.C28.2}
\]
The first identity follows directly on $x,y$; the second
uses $\theta(z)=z^{-1}yz$.

For a proposed mixed value $a=f(x,y)$, put
\[
                  t=\theta(a),\quad b=ta^{-1},\quad c=t^{-1}.
                                                               \tag{D23.C28.3}
\]
Duality and the three-cycle require $b=f(y,z)$ and
$c=f(z,x)$.  Conversely these three cyclic values determine
the six oriented values once the reflection graphs are
imposed.  Indeed, as marked representations of the same
free group,
\[
 q_{xz}=\theta q_{xy},\qquad
 q_{zx}=\theta q_{yx},\qquad
 q_{zy}=\kappa q_{yz}.                                    \tag{D23.C28.4}
\]
Thus the six-kernel intersection is already the intersection
of the three cyclic kernels.  For every retained candidate
the certificate additionally checks
$\kappa(b)=b^{-1}$, so the value determined by the third
reflection graph is exactly the dual value used in the norm
calculation below.  This consistency is checked explicitly;
it is not inferred merely from membership of $(a,b,c)$.

The $M$-orbits numbered $4,8,12,13$ have total degree 492.
Restriction to their union has image order $|M|$, proving
that it is faithful.  In this smaller faithful action form
the cyclic-three image
\[
 Q=\big\langle (x,y,z),(y,z,x)\big\rangle
                      \subset M^3.
\]
Its exact order is
\[
                 |Q|=389356498710072424831060715520393216.
\]
The derived-membership test has a useful exact simplification.
Because $x,y$ have orders $6,12$ and $|M/M'|=72$, the marked
abelianization is $C_6\times C_{12}$.  A free word with
exponents $(u,v)$ has, in the first two cyclic coordinates,
the abelian values
\[
       (u\bmod6,v\bmod12),\qquad
       (-v\bmod6,u-v\bmod12).
\]
These recover both $u,v$ modulo 12.  The two generators of
$Q$ have order 12, so this proves
\[
                 Q/Q'\simeq C_{12}^2.                    \tag{D23.C28.5}
\]
In particular, a triple in $(M')^3$ belongs to $Q'$ exactly
when it belongs to $Q$.  All triples (D23.C28.3) are in
$(M')^3$.  Exact stabilizer-chain membership finds all
eighteen triples in $Q$, hence in $Q'$.  The resulting
derived order is
\[
                |Q'|=2703864574375502950215699413336064.
\]
An independent direct derived calculation and a separate
Schreier calculation over the formal source quotient
$C_{12}^2$ both gave this same order; the proof of
(D23.C28.5) avoids requiring either large calculation for
future replays.

For the full norm equations, induce the original marked
degree-1344 action from the free subgroup to $B=C_2*C_3$.
With the C22 Schreier sections, the six blocks of the
commutator-word value are
\[
       a,\quad a^{-1},\quad b,\quad t^{-1},\quad
       x^{-1}tx,\quad yb^{-1}y^{-1}.                      \tag{D23.C28.6}
\]
The last two inner corrections are verified against both
marked generators of their induced blocks.  Write $F$ for
the resulting permutation on 8064 points, and $s,r$ for
the induced torsion generators, so that
$X=(sr)^{-2}$ and $Y=sXs$.
For all $36\cdot216=7776$ values in the H7 finite fibers,
the exact equations
\[
                  (sF)^2=(F^{-1}r)^3=1
\]
hold.  The full $m=-1$ equation is checked on these same
8064 points, using the verified cyclic conjugation by
$r$ to evaluate the other orientations.  It also holds
for every value.  In particular it holds for the eighteen
remaining values considered here.  Together with the
derived-image membership and reflection consistency above,
this proves the stated finite test, with no actual
profinite lifting inference.
\end{proof}

\paragraph{Exact certificates.}
The reflection graph and its exact straight-line evaluator
are \path{agents/characters/code/hurwitz_mixed_reflection_map.py}.
This evaluates on the 29 strong generators, verifies their
reconstruction, and checks (D23.C28.2).  The faithful cyclic
image is rebuilt by
\path{agents/characters/code/hurwitz_full_mixed_s3_group.py}.
The full norm and final membership scripts are
\path{agents/characters/code/hurwitz_full_mixed_gentle_norms.py}
and \path{agents/characters/code/hurwitz_full_mixed_s3_membership.py}.
Their matching JSON/NPZ records retain every tested index
and the passing cyclic tuples.  The eighteen-element list
comes from the literal raw-pentagon record
\path{agents/pentagon/hurwitz_raw_pentagon45.npz}; its
construction and the uniqueness of the first value are
separate inputs to the main necessary reduction.

\subsection{The second mixed pair is already in the same cyclic system}

\begin{statement}{D23.C29}{Proved}{An exact braid identity identifies the mixed45 cycle value}
\end{statement}
In $B_4$, let $A,B,C$ be the positive Artin generators and put
\[
 s=ABA,\quad x=C^2,\quad y=s^2,\quad
 X_5=(CBC)^2,\quad z=(xy)^{-1},\quad z_4=(ABC)^4.
\]
Then
\[
              s^{-1}X_5s=A^2z\,z_4,\qquad
              [A^2,x]=[A^2,y]=1.                          \tag{D23.C29.1}
\]
Consequently, for every profinite commutator word $f$,
\[
                   s^{-1}f(y,X_5)s=f(y,z).                \tag{D23.C29.2}
\]
In the Hurwitz finite image, if $a=f(x,y)$ is an actual
gentle value, the C28 reflection therefore gives
\[
                   f(X_4,X_5)=s\theta(a)a^{-1}s^{-1}.
                                                               \tag{D23.C29.3}
\]
All eighteen raw-pentagon survivors satisfy this equality
with their required $H_{45}$ values.

\begin{proof}
Here is a derivation directly from the Artin relations.
Write $a=A$, $b=B$, $c=C$, $s=aba$ and $t=bcb=cbc$.
The relations $ac=ca$, $aba=bab$ and $bcb=cbc$ give
\[
 b^{-2}t^2=cb^2c,\qquad
 c^2scb^2c=(abc)s(abc).
\]
Consequently
\[
 z_4=c^2s b^{-2}t^2s
     =c^2s^2a^{-2}(s^{-1}t^2s),
\]
and hence $s^{-1}t^2s=a^2s^{-2}c^{-2}z_4=A^2zz_4$.
The literal braid identity (D23.C29.1) is also verified by the
faithful Artin action on the free group of rank four;
the certificate records the two explicit words
\[
 (ABA)^{-1}(CBC)^2(ABA),\qquad
 A^2\bigl(C^2(ABA)^2\bigr)^{-1}(ABC)^4
\]
and checks equality of their four reduced free-generator
images.  The two commutations are checked by the same
exact word method (and also follow respectively from
$AC=CA$ and centrality of $(ABA)^2$ in the first $B_3$).
The full twist $z_4$ is central in $B_4$.

Conjugation by $s$ fixes $y$.  The element $A^2z_4$
centralizes both $y$ and $z$, so replacing the second
input $z$ by $(A^2z_4)z$ changes a word value only by
the central multiplier raised to the word's second
abelian exponent.  This exponent is zero for a commutator
word.  The assertion holds first for discrete words and
then for the closed commutator subgroup by continuity,
proving (D23.C29.2).

The gentle three-cycle and duality give
$f(y,z)=f(x,z)f(x,y)^{-1}$.  C28 identifies
$f(x,z)=\theta(a)$ in the marked finite image, proving
(D23.C29.3).  The certificate evaluates its right side
for each of the eighteen survivors and compares it with
the required $H_{45}$ permutation on all 1344 points;
all comparisons agree.

More generally, the six-orientation commutator-word
evaluations of the mixed45 pair are the cyclic
mixed34 evaluations after conjugation and central
input modifications.  Such central factors cancel
in the two/three-cycle identities, including the
$m=-1$ equation because the three peripheral
multipliers have product one.  Thus a separate
mixed45 norm computation adds no constraint beyond
the corresponding full cyclic system.  Compatibility
of that system with a first-pair value of the same
word remains an additional condition.
\end{proof}

\paragraph{Certificate.}
\path{agents/characters/code/hurwitz_mixed45_cycle_link.py}
replays the literal Artin identities and all eighteen
finite comparisons.  Its JSON/NPZ records have the same
stem.  The literal $z_4$ in (D23.C29.1) is retained;
it is not omitted before passing to commutator words.

\subsection{Common words for the complete mixed cyclic system}

\paragraph{D23.H9 / D23.F38 (proved; two whole-proof checks).}
Every one of the eighteen pairs retained by D23.P26 has
one ordinary commutator word with its prescribed complete
first-pair value $F_V$ and all six mixed peripheral
evaluations from D23.C28. In particular this word has
the prescribed values at all five pentagon pairs, and
their degree-$8064$ raw five-factor product is the identity.
The assertion concerns these finite evaluations and
finite cycle tests. It does not assert that the word
satisfies the free profinite gentle equations or belongs
to the actual full-GT group.

\begin{proof}
Let $P=\langle A^2,B^2\rangle$, $N=P^{(\infty)}$, and
let $Q=P/N$ be the marked quotient of order $93312$
from D23.F35. Let $M$ be the mixed pair group and let
$E\subset M^3$ be its faithful cyclic-three image in
D23.C28. The latter has a faithful permutation action
on $1476$ points, constructed from three copies of
the same $492$-point restriction of $M$.

On the disjoint union of the $472$-point action of $Q$
and this action of $E$, form the two literal joint
generator permutations. They generate a marked group
$J\subset Q\times E$ of degree $1948$ with exact order
\[
       |J|=756909033492380793871582030971644411904.
\]
Both joint generators have order $12$.
For each retained $R$, put
\[
       t=\theta(R),\qquad b=tR^{-1},\qquad c=t^{-1}.
\]
The required target tuple is
\[
                 (q_*,R,b,c),\qquad
                 q_*=(F_{448,1},I_{24})\in Q'.
\]
The exact complete stabilizer-chain test accepts all
eighteen target tuples in $J$.

Here membership is in fact derived membership. D23.C28
proves that $E/E'\cong C_{12}^2$, with the two original
free generators as its basis. Since the joint
generators of $J$ also have order $12$, its
abelianization is a quotient of $C_{12}^2$. The
surjection $J\to E$ gives the reverse quotient, so
\[
                      J/J'\cong C_{12}^2.
\]
The tuple projects into $E'$, and hence its image in
this abelianization is zero. Thus every accepted tuple
lies in $J'$.

The certificate gives an additional direct check of
this last step. Recover a straight-line expression for
each accepted tuple in the two original joint
generators, by the exact stabilizer-chain straight-line
programs. Both exponent sums are $0$ modulo $12$ for
each of the eighteen expressions. If their integer
exponent sums are $12a,12b$, multiplication of the
free expression by $x^{-12a}y^{-12b}$ changes neither
finite value and makes both exponent sums exactly
zero. It therefore gives an ordinary word in $F_2'$.
This does not assume that a word with exponent sums
merely zero modulo $12$ is itself a commutator word.

Choose such a word $w_0$. Its first full value $P_0$
has the same quotient value in $Q$ as $F_V$, so
$F_VP_0^{-1}\in N$. The complete mixed cyclic image
$E$ is solvable by D23.F33. Its order does not affect
this fact: it is a subgroup of a finite product of
the solvable mixed group $M$.
D23.F33 therefore supplies a further ordinary
commutator word $u$ whose two values in $P\times E$
are $(F_VP_0^{-1},1)$. Then $w=uw_0$ has exact
full first value $F_V$ and all the prescribed mixed
cyclic values. The three reflection graphs and the
checked identity $\kappa(b)=b^{-1}$ in D23.C28
give its other three mixed orientations as well.

D23.C29 identifies its fourth--fifth value as
$sbs^{-1}$, since the other input factors are central
and the free word is a commutator word. This equals
the prescribed $H_{45}$ on the eighteen retained
pairs. The reversing half twist identifies its
fifth--first value with $DR^{-1}D^{-1}$ using the
now imposed mixed duality. The second--third value
is $DF_V^{-1}D^{-1}$: the first pair is conjugated
to its reversal by $ABA$, and the duality of $F_V$
is the exact identity of D23.F34. Thus all five
values are those used in D23.P26.

Finally each of the five pair images has a unique
marked raw lift, by D23.P24 and D23.P26. Evaluation
of the same word $w$ must give precisely these lifts.
Their literal five-factor product is the identity
for every retained pair by D23.P26. The full mixed
finite norm equations hold by D23.C28. None of these
finite conclusions changes the free word into an
actual gentle or full-GT element.
\end{proof}

\paragraph{Certificates and independent comparisons.}
The complete joint construction and eighteen membership
and exponent tests are
\path{agents/five_level_six/hurwitz_first_full_cyclic_core.py},
with its JSON, NPZ and exact group-state record.
The smaller $Q\times M_{34}\times M_{45}$ test was
also performed independently by
\path{code/hurwitz_three_pair_commonword.py}
and \path{agents/five_level_six/hurwitz_three_pair_fast.py}.
Both obtain the same full order
$10866225302113486007489789952$ and derived order
$75459897931343652829790208$, and accept all thirty-six
base pairs of D23.H8. That weaker test is not used as
a substitute for the full mixed cyclic-system test above.
The independently written parallel proof is retained as
\path{agents/five_level_six/hurwitz_full_core_commonword.tex}.
D23.F38 names that same finding, not an additional reduction.

\subsection{The abelian restrictions on the actual pure torsion frames}

\paragraph{D23.H10 (proved; independent whole-proof check).}
Let $h=(1,f)$ be an actual multiplier-one gentle element,
with the positive marking
\[
 B=C_2\amalg C_3=\langle s,r\rangle,\qquad
 F=\langle x,y\rangle,\quad x=(sr)^{-2},\quad y=sxs.
\]
Let its actual pure frames be
\[
       sf=\beta s\beta^{-1},\qquad
       f^{-1}r=A r A^{-1},\qquad A,\beta\in\widehat F.
\]
Then
\[
             A\in\widehat F',\qquad
       \beta^{\mathrm{ab}}=(-\rho,-\rho),\qquad
                      \rho\in2\widehat{\mathbb Z}.
\]
In particular, in any finite marked quotient whose pure
image has abelianization $C_{12}^2$ on $x,y$, a normalized
actual element with $\rho=0\pmod3$ must admit a conjugator
from the \emph{derived pure image} between $r$ and
$f^{-1}r$. Its harmonic conjugator between $s$ and $sf$
must have abelian coordinate either $(0,0)$ or $(6,6)$.
An arbitrary pure conjugator is insufficient for either
conclusion. These are separate conjugator requirements;
the same conjugator is not required to serve both torsion
generators.

\begin{proof}
The existence and uniqueness of the actual pure frames,
and the diagonal abelian form of $\beta$, are the supplied
LM7\_PUSH E.65.1. Rearranging the cubic frame equation gives
\[
                    f=(rAr^{-1})A^{-1}.
\]
The action $R$ of conjugation by $r$ on
$F^{\mathrm{ab}}=\widehat{\mathbb Z}^2$ is the standard
order-three peripheral rotation. Its characteristic
polynomial is $t^2+t+1$, so $\det(R-I)=3$.
Since $f\in\widehat F'$, abelianizing the last display gives
\[
                       (R-I)A^{\mathrm{ab}}=0.
\]
Multiply by the integral adjugate. Multiplication by $3$
is injective on $\widehat{\mathbb Z}^2=\prod_p\mathbb Z_p^2$,
so $A^{\mathrm{ab}}=0$. This proves $A\in\widehat F'$.
One must perform this argument in the torsion-free
profinite integral module: its reduction modulo $12$ alone
has nonzero order-three fixed vectors.

For clarity, the parity restriction on $\rho$ can be read
directly from E.65.11. In the free class-two quotient,
write
\[
          \beta^{-1}=x^\rho y^\rho[x,y]^d.
\]
The multiplier-one three-cycle kills the central
commutator coefficient of $f$, since three times that
coefficient is zero in $\widehat{\mathbb Z}$. The
interchange $\theta(x)=y$, $\theta(y)=x$ and
$f=\theta(\beta^{-1})^{-1}\beta^{-1}$ then give
\[
                            \rho^2+2d=0.
\]
Reduction in $\mathbb Z_2$ implies that $\rho$ is even;
at odd primes multiplication by $2$ is a unit, giving the
stated inclusion $\rho\in2\widehat{\mathbb Z}$.

A continuous surjection to a finite group sends the
closed commutator subgroup onto the finite derived
subgroup. Thus the cubic conjugator in the finite pure
image lies in its derived subgroup. Its harmonic
abelian coordinates are the reduction of
$(-\rho,-\rho)$. If additionally $\rho=0\pmod3$,
the only possibilities modulo $12$ are $\rho=0,6$,
giving precisely the two displayed harmonic cosets.
\end{proof}

\paragraph{Source and scope.}
The exact supplied inputs are LM7\_PUSH E.65.1 and the
paragraph \emph{The integral sign that fixes the lift},
especially E.65.11, in
\path{sources/LM7_PUSH/E_cm_character_agent.tex}.
This argument does not identify the full character
$\rho$ with a Kummer character. Supplied A.102 gives only
the congruence modulo three needed for D23.H2; no stronger
arithmetic normalization is inserted here.

\subsection{Finite source surjectivity and pure frames still retain every candidate}

\begin{statement}{D23.P28}{Proved}{The full mixed finite source tests exclude no remaining $V$}
\end{statement}
Let $E$ be the finite image of $B=C_2*C_3$ obtained by
inducing D23.C28's faithful degree-$492$ mixed-pair image
from its index-six free subgroup. Its degree is $2952$ and
its order is
\[
             |E|=2336138992260434548986364293122359296.
\]
For every one of D23.P26's eighteen pairs, let $f_E$ be
the complete six-coset value supplied by the gentle
identities and C28's exact common-word membership.
If $s,r$ are the marked generators of orders two and
three, then
\[
                       \langle sf_E,f_E^{-1}r\rangle=E.
                                                        \tag{D23.P28.1}
\]
In addition, for each of the six remaining $V$ there
is a retained $R$ for which both new torsion generators
are conjugate to the old ones by a \emph{single element
of the derived pure image}. The explicit labels, using the
zero-based position of $R$ in each H8 fiber, are
\[
                   (1,0),(19,0),(23,1),(27,2),(53,0),(61,0).
                                                        \tag{D23.P28.2}
\]
Thus imposing finite source surjectivity and these derived-pure
torsion-frame conditions cannot exclude any of the six
remaining candidates. This is not a profinite realization
statement.

\begin{proof}
The induced representation has marked free generators
$X=(sr)^{-2}$ and $Y=sXs$, and its first degree-$492$
fiber is the original faithful mixed-pair representation.
This uses the exact Schreier cocycles in C28's induced
source construction. Write $x,y$ for that original pair,
$z=(xy)^{-1}$, and $\theta$ for its verified automorphism
$\theta(x)=x$, $\theta(y)=z$. If $a$ is a retained mixed
value and $t=\theta(a)$, its six literal blocks in this
induced marking are
\[
 a,\quad a^{-1},\quad ta^{-1},\quad t^{-1},\quad
 x^{-1}tx,\quad y(at^{-1})y^{-1}.
                                                        \tag{D23.P28.3}
\]
D23.C28 proves membership in the full derived free-image
group and the required norm identities. The certificate
reconstructs (D23.P28.3), including all conjugations,
and rechecks $(sf_E)^2=(f_E^{-1}r)^3=1$.

For each of the eighteen residues it constructs an
explicit permutation $T$ of the full $2952$-point set
such that
\[
                  TsT^{-1}=sf_E,\qquad
                  TrT^{-1}=f_E^{-1}r.                    \tag{D23.P28.4}
\]
The construction is an exact labeled-graph traversal.
The old action has four orbits of sizes
$144,216,1296,1296$. On each orbit, try the possible
image of its first point, propagate the two generator
equations, and reject a choice as soon as the assignments
are inconsistent or cease to be injective. A successful
choice on each orbit gives a permutation, and both
identities in (D23.P28.4) are checked entry by entry on
the entire set. All eighteen residues have such a $T$.

Since $f_E\in E$, the new generated group is a subgroup
of $E$. Equation (D23.P28.4) identifies it with the
conjugate group $TET^{-1}$, so it has the same finite
order and equals $E$. This proves (D23.P28.1).
No invariance of the original finite quotient kernel
under an actual gentle automorphism was assumed; the
displayed intertwiners are computed conclusions.

For the six labels (D23.P28.2), exact membership in the
full group $E$ additionally verifies $T\in E$. Its
permutation fixes each of the six induced fibers. The
action of $E$ on those fibers is the original quotient
$B/F_2\cong S_3$, so its kernel is exactly the image
of $F_2$. Hence $T$ belongs to that pure image.
Moreover restrict the pure image to the three cyclic
fibers numbered $0,2,3$. The two marked restricted
generators are literally C28's degree-$1476$ generators.
This is a surjection onto that cyclic core, and the
two orders are equal: both are $|E|/6$.
Thus the restriction is a marked isomorphism. Exact
membership of the six restricted $T$ values in C28's
full derived group proves that all six $T$ belong
to the derived pure image. The two identities
(D23.P28.4) therefore give the stronger torsion
conjugacies with a derived-pure conjugator for both.
The certificate makes no claim that the other
twelve intertwiners are the only possible ones or that
those twelve residues fail the pure-frame condition.
One retained residue for each $V$ already proves the
stated inability of these tests to exclude any $V$.

The distinction between a pure and a derived-pure
conjugator is relevant to the actual source. If
$\Phi(r)=f^{-1}r=ArA^{-1}$ with $A\in\widehat F_2$,
then $f=(rAr^{-1})A^{-1}$. On
$\widehat F_2^{\mathrm{ab}}\cong\widehat{\mathbb Z}^2$
this gives $(R-I)A^{\mathrm{ab}}=0$, where the order-three
rotation matrix has $\det(R-I)=3$. Multiplication by
$3$ is injective on $\widehat{\mathbb Z}^2$, so
$A^{\mathrm{ab}}=0$. Thus an actual cubic frame is
derived-pure, although its finite counterpart could
otherwise have a nonzero three-torsion abelian class.
The six verified $T$ satisfy this stronger requirement.
Their zero abelian class also belongs to the allowed
harmonic classes $(0,0)$ or $(6,6)$ when only the known
conditions $\rho\in2\widehat{\mathbb Z}$ and
$\rho\equiv0\pmod3$ have been imposed.
\end{proof}

\paragraph{Certificates and scope.}
The replay is
\path{agents/pentagon/hurwitz_mixed_source_intertwiners.py};
the option \texttt{--rebuild} recomputes the full source
stabilizer chain from its two original generators.
The JSON records all eighteen intertwiners and the six
verified derived-pure ones, while the NPZ retains $s,r,X,Y$,
every complete $f_E$, and every literal $T$.
The source order was also independently recomputed for
the first new generating pair; subsequent redundant
order calculations were stopped once the exact
intertwiner proof covered all eighteen.
The large mixed S3 common-word group and the reflection
map are the separately recorded D23.C28 inputs. These
finite conclusions do not identify an actual associator
or prove $\Delta$ vanishes.

\subsection{An actual-source normalization from eighteen pairs to four}

\begin{statement}{D23.C30}{Proved}{Cubing reduces the existential mixed gate to four pairs}
\end{statement}
If a normalized actual full-GT element with
$(\rho,a_3,a_2)=(0,1,0)$ exists, then one exists whose
degree-1344 Hurwitz data $(V,R)$ is one of the four pairs
\[
                        (1,0),\quad(19,0),\quad(23,1),\quad(27,2)
                                                               \tag{D23.C30.1}
\]
in the fixed eighteen-pair numbering.  In each of these
four pairs, both $V$ and $U=R^{-1}V$ are involutions.
This is an existential normalization by replacing the
element by its cube; it does not reject the other
fourteen possibilities for an element before normalization.

\begin{proof}
The multiplier-one normalized free action associated to
$h=(1,f_h)$ is
\[
            \varphi_h(x)=x,\qquad
            \varphi_h(y)=f_h^{-1}y f_h.
\]
Its composition law is
\[
              f_{hk}=f_h\varphi_h(f_k).                   \tag{D23.C30.2}
\]
This follows directly by composing the two automorphisms
on $y$; uniqueness of the commutator normalization removes
the possible power of $y$.  It is the multiplier-one
specialization, in the present action convention, of the
standard GT composition formula of Lochak--Schneps,
\emph{Inventiones Mathematicae} \textbf{127} (1997),
571--600, equation (2) on printed page 572.

For the mixed finite pair, actual covariance gives
$VxV^{-1}=RxR^{-1}$ and $VyV^{-1}=y$, where $R=f_h(x,y)$.
Thus $U=R^{-1}V$ implements $\varphi_h$ on the finite
mixed image.  Evaluation of (D23.C30.2), or induction on
powers of $h$, yields
\[
                    R_{h^n}=V^nU^{-n}.                    \tag{D23.C30.3}
\]
The actual Hurwitz action is a group action, so its
permutation for $h^n$ is $V^n$.  These statements apply
to powers of a single actual element and do not assume
that different surviving finite candidates can be
realized simultaneously.

Exact enumeration of the eighteen necessary pairs gives
orders dividing six for both $V$ and $U$.  Applying
(D23.C30.3) at $n=3$ gives exactly the four pairs
(D23.C30.1).  All four have $V^2=U^2=1$.
For example the three pairs above each of base indices
$1,19,23,27$ cube respectively to $(1,0),(19,0),(23,1),(27,2)$;
all six pairs above $53,61$ cube to $(1,0)$.
The certificate also checks every odd power in the
finite cyclic group generated by each pair, and each
remains in the original eighteen-element list.

Finally $\rho\bmod3$, $a_2$ and $a_3$ are characters on
the multiplier-one group.  Cubing preserves
$\rho=0\bmod3$, $a_2=0$, and $a_3=1$.
Consequently $h^3$ is still a normalized mixed element,
and its finite data belongs to (D23.C30.1).  This proves
the smaller existential gate.
\end{proof}

\paragraph{Certificate.}
\path{agents/characters/code/hurwitz_candidate_power_closure.py}
computes all orders, every odd power, and the complete
cube map.  Its JSON/NPZ records have the same stem.
The input eighteen-pair list is the main actual-source
necessary reduction, including the uniqueness of the
first value.  The power normalization uses no new
lifting assertion.

\subsection{Three existential normal forms after complex conjugation}

\begin{statement}{D23.C31}{Proved}{Complex conjugation identifies two cube-normalized cases}
\end{statement}
If a normalized actual mixed full-GT element exists, then
one exists whose degree-1344 data $(V,R)$ has one of the
three labels
\[
                           (1,0),\qquad(19,0),\qquad(27,2).
                                                               \tag{D23.C31.1}
\]
This uses cubing as in C30 and, if necessary, conjugation
by the actual element $c=(-1,1)$ of complex conjugation.
It is an existential normal form, not an assertion that
any of the three finite pairs lifts to an actual element.

\begin{proof}
First realize the complex-conjugation symmetry on the
chosen finite Nielsen orbit.  In
$G=\mathbb F_3^2\rtimes SL_2(\mathbb F_3)$, let
$d=\operatorname{diag}(1,-1)$.  The map
\[
            \eta(v,A)=(dv,dAd^{-1})
\]
is an involutive automorphism of $G$.  The exact
36-element order-six class used in P17 satisfies
$\eta(\mathcal C^{-1})=\mathcal C$, verified on every
class element in the certificate.

The standard GT element $c=(-1,1)$ sends each Artin
generator to its inverse.  The arithmetic identification
of this element is Lochak--Schneps,
\emph{Inventiones Mathematicae} \textbf{127} (1997),
571--600, Theorem 1 on printed page 573; the braid action
is the convention used in supplied FP.34 and D23.P18.
In the free-kernel
construction of P18 it therefore takes the four local
conjugacy classes to their inverse classes.  Postcompose
the target representation by $\eta$ to return to the
original class family.  The resulting bijection of the
finite Nielsen set intertwines the three braid
generators with their inverses.  The unique orbit of
degree 1344 in P17 must be preserved.  Hence it gives
a permutation $W$ with
\[
           WAW^{-1}=A^{-1},\quad WBW^{-1}=B^{-1},
           \quad WCW^{-1}=C^{-1}.                         \tag{D23.C31.2}
\]
The transitive labeled-graph calculation exhausts all
1344 possible images of one point and finds exactly
one such $W$; it also verifies $W^2=1$.

Postcomposition by a target-group automorphism commutes
with every precomposition defining a GT action.
Consequently the inserted $\eta$ cancels when one
conjugates a multiplier-one GT action by this symmetry:
the action of $chc^{-1}$ is $WV_hW^{-1}$.
On the normalized free-group action,
\[
              f_{chc^{-1}}(x,y)=f_h(x^{-1},y^{-1}).
\]
Since (D23.C31.2) inverts both $x=C^2$ and
$y=(ABA)^2$, its mixed value is $WR_hW^{-1}$.

The normalized characters remain in the required state.
On the canonical degree-36 detector the exact induced
$W$ sends the six first-pair labels by
\[
                       (\rho,a_3)\longmapsto(-\rho,a_3).
\]
This is checked on all six values.  For the marked
$S_4$ pair $p=(123)$, $t=(1234)$, conjugation by $(13)$
inverts both $p$ and $t$.  Thus the supplied value
$t^{2a_2}p^\rho$ goes to $t^{2a_2}p^{-\rho}$, so $a_2$
is unchanged as well.  In particular
$(\rho,a_3,a_2)=(0,1,0)$ is preserved.

Finally simultaneous conjugation of all eighteen
recorded pairs by $W$ preserves that list.  On the four
cube-normalized pairs of C30 it fixes $(1,0)$ and
$(27,2)$ and interchanges $(19,0)$ with $(23,1)$.
Conjugating the latter case if necessary gives exactly
the three normal forms (D23.C31.1).
\end{proof}

\paragraph{Certificate and scope.}
\path{agents/characters/code/hurwitz_complex_conjugation.py}
checks the inverse-class correction, reconstructs the
unique inverse-generator intertwiner, computes its
degree-36 detector action, verifies the marked $S_4$
inversion, and records its action on every candidate.
The JSON/NPZ records are
\path{agents/characters/records/hurwitz_complex_conjugation_probe.json}
and the NPZ of the same stem.  The actual-source inputs
are the standard $c=(-1,1)$ action, P18 functoriality,
and P17's unique degree-1344 orbit.  No arithmetic
lifting or extension of a finite candidate is assumed.

\subsection{Two existential normal forms from actual conjugate closure}

\begin{statement}{D23.C33}{Proved}{An odd product reduces the mixed existence gate to two pairs}
\end{statement}
If a normalized actual full-GT element with
$(\rho,a_3,a_2)=(0,1,0)$ exists, then one exists whose
degree-1344 data $(V,R)$ has one of the two labels
\[
                              (1,0),\qquad(27,2).          \tag{D23.C33.1}
\]
Both have $V^2=(R^{-1}V)^2=1$, and their finite pairs
are fixed by the complex-conjugation symmetry of C31.
This is an existential normalization.  In fact the
complete odd-word closure of each of the original
eighteen finite pairs, together with its complex conjugate,
stays within the eighteen-pair list; no original individual
case is rejected by this closure test.

\begin{proof}
For an actual multiplier-one element $h$, write
$R_h=f_h(x,y)$ and $U_h=R_h^{-1}V_h$ in the marked
mixed finite image.  The actual Hurwitz action gives
$V_{hk}=V_hV_k$.  The normalized free action is implemented
on this image by $U_h$, so the composition law
$f_{hk}=f_h\varphi_h(f_k)$ from C30 gives
\[
 R_{hk}=R_hU_hR_kU_h^{-1}
        =V_hR_kU_h^{-1}.
\]
Therefore
\[
                    U_{hk}=R_{hk}^{-1}V_hV_k=U_hU_k.
                                                               \tag{D23.C33.2}
\]
Thus the paired action $h\mapsto(V_h,U_h)$ is a
homomorphism on the actual normalized group.  This is
the point needed to pass from powers of one element
to words in two elements.

Let $c=(-1,1)$ and let $W$ be C31's explicit realization
of complex conjugation.  Its naturality gives
\[
 (V_{chc^{-1}},U_{chc^{-1}})
       =(WV_hW^{-1},WU_hW^{-1}).                          \tag{D23.C33.3}
\]
The three character values of $h$ and $chc^{-1}$ are
the same normalized values $(0,1,0)$.  Consequently any
word of odd total length in these two elements and
their inverses still has these values: the quadratic
values add modulo two and $\rho$ remains zero modulo
three.

After C30 and C31, it suffices to consider the labels
$(1,0),(19,0),(27,2)$.  For $(19,0)$, its complex
conjugate has label $(23,1)$.  Exact multiplication
of the two paired actions gives
\[
 \bigl(V_hV_{chc^{-1}}V_h,\,
       U_hU_{chc^{-1}}U_h\bigr)
                    =(V_{27,2},U_{27,2}).                 \tag{D23.C33.4}
\]
The actual element $h(chc^{-1})h$ therefore has
label $(27,2)$ and retains the required character
values.  This replaces the last remaining third
case and proves (D23.C33.1).

For completeness, the certificate enumerates the
entire parity-graded subgroup generated by each of
the eighteen pairs and its conjugate, including
both inverse generators.  It starts at the identity
with even parity and closes the finite set without
a word-length cutoff.  Every odd state belongs
to the original eighteen-pair list.  The graded
group for either of the two final pairs has order
two.  The groups for $(19,0)$ and $(23,1)$ have
order six, with odd states precisely
$(19,0),(23,1),(27,2)$; the remaining cases have
graded orders six or eighteen.  These calculations
confirm the stated scope: a smaller existence
gate, but no exclusion of an element before the
allowed source operations.
\end{proof}

\paragraph{Exact certificate.}
\path{agents/characters/code/hurwitz_candidate_conjugate_closure.py}
checks the full finite closures, records explicit words
for every odd state, and verifies (D23.C33.4).
The JSON and NPZ files of the same stem retain
all states and their parity.  The only source inputs
are the actual GT composition, C31's realized
complex conjugation, and the previously proved
eighteen-case necessary reduction.  No new finite
or profinite lifting assertion is used.

\subsection{A literal spherical five-puncture quotient and its order-five cochains}

\begin{statement}{D23.P29}{Proved}{The $6720$-point extension and all eighteen prescribed cochains}
\end{statement}
The degree-$1344$ Hurwitz action of D23.P17 extends, by
adjoining four sectors, to a literal degree-$6720$ action
of the full five-puncture sphere mapping class group.
The added $5376$ points are exactly the unique orbit of
that degree in the colored $C_+^3C_-$ construction, where
$C_+$ is the original order-six class in $E_0$ and
$C_-=C_+^{-1}$. The original $1344$ points and the added
$5376$ points are invariant under the first three braid
generators, and their restrictions agree literally with
the saved earlier actions.

Write $\sigma_1,\ldots,\sigma_4$ for the four literal
half-twist permutations and
\[
                    r_5=(\sigma_4\sigma_3\sigma_2\sigma_1)^2.
\]
It has order five and conjugates the five marked twists
$X_i$ to $X_{i+3}$, with cyclic subscripts. For every
one of D23.P26's eighteen pairs, the prescribed full
value $f_{12}$ has an explicit permutation $K$ such that
\[
                    f_{12}=r_5K^{-1}r_5^{-1}K.
                                                        \tag{D23.P29.1}
\]
This $K$ fixes all sixty transported orthogonal blocks,
and its restriction to every block lies in the local
simple group $S=P\Omega^-_6(3)$ of D23.F30.

The statement does \emph{not} yet assert that the
independently specified sixty local components assemble
to an element of the actual global pure braid image.
That membership is a separate condition on (D23.P29.1).

\begin{proof}
Use the literal group $E_0=\mathbb F_3^2\rtimes
SL_2(\mathbb F_3)$ of D23.P10. Consider length-five
tuples in this group, with product one, four entries
in $C_+$ and one entry in $C_-$, modulo simultaneous
conjugation. The left Hurwitz half twist replaces
adjacent entries $(u,v)$ by $(uvu^{-1},u)$. This
preserves their product and exchanges their class
positions. The usual Hurwitz identities give the Artin
relations, and product one modulo simultaneous
conjugation gives the two sphere relations.

For the concrete orbit, append the inverse product to
each of the original $1344$ four-tuples and each of
the colored $5376$ four-tuples. Every resulting tuple
has exactly the required class pattern. Canonicalize
simultaneous conjugacy by transporting its first entry
to the least element of its class and then taking the
least tuple under that element's six-element centralizer.
The two lists yield $6720$ distinct canonical tuples.
Each of the four half twists permutes the complete list,
and the original two restrictions agree entry by entry.
The five positions of the exceptional $C_-$ entry give
five sectors of size $1344$. On these sectors the four
generators act as the adjacent transpositions, so this
color quotient is the literal $S_5$ quotient.

The replay directly verifies every adjacent braid
relation and every distant commutation, together with
\[
\begin{split}
 (\sigma_1\sigma_2\sigma_3\sigma_4)^5&=1,\\
 \sigma_1\sigma_2\sigma_3\sigma_4^2
                   \sigma_3\sigma_2\sigma_1&=1.
\end{split}
\]
These are the two sphere relations in the primary
presentation of Lochak--Schneps (1997), printed page
571, with the first relation written in reversed
generator indexing; the change $\sigma_i\leftrightarrow
\sigma_{5-i}$ gives the equivalent presentation.
The full twist is therefore trivial, and the
finite action factors through the required sphere
mapping class group and its profinite completion.
The standard ten pure braid generators are evaluated
with their complete conjugating words. Since the color
quotient is $S_5$, their image is exactly the kernel of
this quotient in the full finite image.

Set
\[
 X_1=\sigma_1^2,\quad X_2=\sigma_2^2,\quad
 X_3=\sigma_3^2,\quad X_4=(\sigma_1\sigma_2\sigma_1)^2,
 \quad X_5=(\sigma_3\sigma_2\sigma_3)^2.
\]
The replay checks $X_4=\sigma_4^2$, $r_5^5=1$, and
$r_5X_ir_5^{-1}=X_{i+3}$ literally. This is precisely
the rotation convention of supplied FP.2 and the primary
Lochak--Schneps introduction, printed page 572.

Let $\Omega_0$ be the original sector and
$\Omega_k=r_5^k\Omega_0$. These five ordered sectors
partition the whole orbit. On $\Omega_0$, put
\[
 E_0=F_V,\quad E_1=DF_V^{-1}D^{-1},\quad E_2=R,
 \quad E_3=H,\quad E_4=DR^{-1}D^{-1},
\]
where $D=ABCABA$ is D23.P25's reversing half twist.
These are exactly the five consecutive evaluations
$f(X_{j+1},X_{j+2})$. Therefore an actual associator's
full value at $(X_1,X_2)$ is forced on every sector by
\[
 f_{12}(r_5^ku)=r_5^k E_{-3k}(u),
                 \qquad u\in\Omega_0,                 \tag{D23.P29.2}
\]
where the subscript is reduced modulo five.
No new associator-value choice is introduced here.

For each of the eighteen retained pairs, the certificate
constructs (D23.P29.2) and verifies
$w^5=1$ for $w=f_{12}^{-1}r_5$. This is the five-cycle
norm, already implied by the base pentagon with the
specified ordering, now checked on all $6720$ points.
To construct a primitive, put $K(u)=u$ on $\Omega_0$
and propagate
\[
                         K(wu)=r_5K(u).                 \tag{D23.P29.3}
\]
The permutation $w$ cyclically traverses the five
sectors. Thus four propagation steps define $K$ on
the whole set, while $w^5=1$ makes the return to
$\Omega_0$ consistent. The result is a bijection and
satisfies $Kw=r_5K$, which is equivalent to
(D23.P29.1). All identities are checked pointwise.

Transport the original twelve orthogonal blocks by
$r_5^k$ to obtain sixty blocks of size $112$. Their
invariance under the full four-generator action is
verified directly. Every computed $K$ fixes all sixty
blocks. On sector $k$, pull it back through $r_5^k$;
on original block $j$, then pull it back through the
verified D23.F30 braid transversal from block zero.
This gives a literal permutation of the original
$112$ points. Rebuild the local group from F30's
thirteen generators and take its derived subgroup
$S$, of order $3265920$. Exact membership in $S$
holds for all $18\cdot60=1080$ such restrictions.
This proves the final local assertion, without replacing
global pure-image membership by local membership.
\end{proof}

\paragraph{Why the global membership question matters.}
Lochak--Schneps (1997), Theorem 2(III'), printed page
573, states for multiplier one that the primitive of
the order-five cocycle belongs to the profinite pure
five-point group. In the present notation this implies
that $f_{12}^{-1}r_5$ is conjugate to $r_5$ by an element
of the actual finite pure image. A cochain in the
larger product of sixty local permutation groups alone
does not establish that property. Conversely, proving
that the actual pure image contains $S^{60}$ would
make every explicitly constructed $K$ legitimate and
would show that this finite order-five test rejects
none of the eighteen pairs.

\paragraph{Records and exact locators.}
The construction and all sphere/marking checks are
\path{agents/pentagon/hurwitz_spherical_five_orbit.py}
and its JSON/NPZ. The latter contains the four full
generators, ten standard pure generators, colors,
sixty-block map, rotation and ordered sectors.
The eighteen full values and cochains are in
\path{agents/pentagon/hurwitz_order5_cochains.py} and
its JSON/NPZ. The primary source is
\url{https://webusers.imj-prg.fr/~pierre.lochak/textes/inventiones.pdf},
printed pages 571--573; the exact statement and its
rotation convention were checked in the opened primary
text and the page-573 scan. Supplied FP.2 and FP.30
record the same source input in the project's marking.

\subsection{The full spherical order-five test has no further obstruction}

\begin{statement}{D23.F39}{Proved}{Sixty independent simple factors and a trivial finite order-five cohomology set}
\end{statement}
Use the literal degree-$6720$ spherical Nielsen action of D23.P29.
Let $\Gamma$ be the image of $M(0,5)$ and $P$ the image of its
pure subgroup $K(0,5)$. Its sixty blocks have size $112$.
With the same local group as in D23.C23, put
\[
 L=\operatorname{PSO}_6^-(3),\qquad
 S=L'=\operatorname{P\Omega}_6^-(3),\qquad |S|=3265920.
\]
Then $\Gamma$ contains the full independent product
\[
                    N=S^{60}\subseteq P'.
\]
The quotient $P/N$ is solvable and has order prime to five.
For $R=(\sigma_4\sigma_3\sigma_2\sigma_1)^2$ in this finite
image and $\alpha=\operatorname{Ad}(R)$, one has
\[
                         H^1(C_5,P)=1.
\]
In particular every finite order-five cocycle is a pure
coboundary. All eighteen prescribed cocycles of P29 have
explicit primitives in $N$; these primitives can each be lifted
to an ordinary commutator word in $K(0,5)$. Thus even requiring
zero integral abelianization of a finite primitive does not
exclude any of the eighteen candidates in this representation.

\begin{proof}
We first prove the full simple-product assertion, rather than
infer independence from separate local images. Label the blocks
$(k,j)$ for $0\leq k<5$, $0\leq j<12$, by transporting the original
twelve blocks with $R^k$. They are numbered $12k+j$.
A breadth-first traversal of the four positive half-twists
chooses literal permutations $T_b$ carrying block zero to block
$b$, with $T_0=1$. Normalize local coordinates by these transports.
For every half-twist $g$ and block $b$, the transition
\[
                    T_{q_g(b)}^{-1}gT_b
\]
restricts to an element of $L$ on block zero. All $240$ such
memberships are checked exactly. Their multiplication rule
therefore embeds $\Gamma$ into $L\wr Q$, where $Q$ is its
sixty-block image. The block-zero stabilizer image is all of
$L$: it is contained in $L$ by this calculation and contains
the old block-stabilizer image from the first three half-twists,
which is $L$ by F30 and C23. The action on the sixty blocks is
transitive, so the same assertion holds on every block.

Exact block calculations give
\[
 |Q|=18807399383040,\qquad
 |Q_P|=156728328192=2^{15}3^{14},
\]
where $Q_P$ is generated by the ten correctly marked pure
braids $A_{ij}$. Their color action is trivial; the four
half-twists give the literal $S_5$ action on the five positions
of the exceptional conjugacy class. Consequently the image
of the pure group is exactly the kernel of this color action.
The quotient $|Q|/|Q_P|=120$ also checks its full permutation
image. The stabilizer of block zero in $Q$ has orbits
\[
 \{0\},\quad \{9,11\},\quad
 \{1,2,3,4,5,6,7,8,10\},\quad \{12,\ldots,59\}.
\]
The three orbits of distinct ordered pairs have representatives
$(0,9)$, $(0,1)$, and $(0,12)$.

Here is one explicit word that separates all three types.
For $g_i=\sigma_i$, use the signed-letter word
\[
 w=(1,-3,-2,-4,-3,1,3,-4,1,-3,2,-4,2,-1,1,3),
\]
where a negative letter denotes the inverse and the displayed
order is ordinary left-to-right group multiplication. Its
block permutation has order twenty. Hence $w^{20}$ belongs to
the block kernel, whose normalized coordinates lie in $L$.
As $L/S=C_2$, the element $u=w^{40}$ belongs to $S^{60}$.
The certificate checks this membership on every block and
finds the local orders
\[
       \operatorname{ord}(u_0)=7,\qquad
       u_1=u_9=u_{12}=1.
\]

Let $J$ be the normal closure of $u$ in $\Gamma$. The product
$S^{60}$ is normal in the ambient wreath product, so
$J\subseteq S^{60}$. It is also contained in $P$, since $u$
fixes all blocks and $P$ is normal. Its block-zero projection
is a nontrivial normal subgroup of $L$ contained in the
nonabelian simple group $S$, and thus is $S$. Transitivity
gives surjectivity onto all sixty copies of $S$.
For any of the three representative pairs, the projection of
$J$ is a subdirect subgroup of $S\times S$. By Goursat's lemma
it is either the full product or the graph of an automorphism.
The displayed unequal local orders rule out a graph.
The three pair orbits therefore imply full projection onto
every distinct pair. The simple-product lemma proved in F30
now gives $J=S^{60}$. This proves $N\subseteq P$; perfectness
of $N$ then gives $N\subseteq P'$.

The kernel of $P/N\to Q_P$ embeds into $(L/S)^{60}=C_2^{60}$.
Moreover each of the five sector restrictions of each pure
block generator belongs to the old order-$216$ block group;
these fifty memberships are checked exactly. That old group
is solvable, with derived orders $216,72,18,9,1$ by F30.
Thus $Q_P$ is a subgroup of a product of five solvable groups,
so is solvable. The preceding extension proves that $P/N$ is
solvable and its order has no factor five.

We give the cohomology argument explicitly. If a finite
solvable group $B$ has order prime to five and $\alpha$ is an
automorphism of order dividing five, every $C_5$ cocycle in
$B$ is a coboundary. For an abelian $B$, use additive notation.
If $\sum_{i=0}^4\alpha^i(c)=0$, then with
$T=\sum_{i=0}^4 i\alpha^i(c)$ one has
$(1-\alpha)T=-5c$. Multiplication by five is invertible on
$B$, so $c=(1-\alpha)(-T/5)$. For a solvable $B$, induct on
its derived length, using a nontrivial abelian last derived
subgroup. First kill the cocycle in the quotient, lift the
chosen conjugator, and then kill the resulting cocycle in
that abelian subgroup by the preceding argument. This proves
the assertion without an unproved coprime-conjugacy input.

Now take any cocycle $f\in P$, with
$\alpha^4(f)\alpha^3(f)\alpha^2(f)\alpha(f)f=1$.
Apply the preceding argument in $P/N$, and choose $p\in P$
lifting a primitive there. The cohomologous cocycle
\[
                          f'=\alpha(p)fp^{-1}
\]
belongs to $N$. The automorphism $\alpha$ permutes the sixty
simple factors in twelve free cycles of length five, because
$R$ carries sector $k$ to sector $k+1$. Also $R^5=1$ literally.
Choose local coordinates by transporting along these cycles;
this removes any intermediate inner coordinate changes.
For one cycle, writing the cocycle as $(c_0,\ldots,c_4)$, its
norm condition is $c_1c_2c_3c_4c_0=1$.
Choose $k_0=1$ and recursively $k_j=k_{j-1}c_j$ for
$j=1,\ldots,4$. The norm condition supplies the closing
equation $k_0=k_4c_0$. Therefore
$c_j=k_{j-1}^{-1}k_j$, exactly the equation
$f'=\alpha(k)^{-1}k$ on that cycle. Do this on all twelve
cycles. It follows that $f=\alpha(kp)^{-1}kp$ in $P$.
This proves $H^1(C_5,P)=1$.

Finally P29 supplies eighteen literal primitives $K$ with
$K(f^{-1}R)K^{-1}=R$. Re-expressing all $60$ local coordinates
in the independently chosen transports above verifies that
each belongs to $S$. Hence these entire permutations lie in
$N\subseteq P'$. The image of the ordinary commutator subgroup
of a group under a finite quotient is the commutator subgroup
of the image: lift each factor of a finite product of
commutators. Thus each primitive has an ordinary commutator
lift in $K(0,5)$, and that lift has zero abelianization.
This is a finite compatibility statement, not a lift of an
eighteen-pair candidate to an actual full-GT element.
\end{proof}

\paragraph{Exact source normalization and abelian constraint.}
Lochak--Schneps, \emph{Inventiones Mathematicae} 127 (1997),
Theorem 2, printed p.~573, gives, at multiplier one,
$f(X_1,X_2)=\alpha(k)^{-1}k$ with $k\in\widehat{K(0,5)}$.
Their p.~572 fixes $R=(\sigma_4\sigma_3\sigma_2\sigma_1)^2$
and $\alpha(X_i)=X_{i+3}$, exactly P29's checked marking.
Equivalently $k(f^{-1}R)k^{-1}=R$; the inverse and side of
the conjugator are thereby fixed. The five-generator
presentation on printed p.~578 has all relators of exponent
sum zero in each generator, so its abelianization is
$\mathbb Z^5$, and the profinite abelianization is
$\widehat{\mathbb Z}^5$. Since $f$ is a commutator word,
$k^{\rm ab}$ is fixed by the cyclic permutation of the five
basis vectors. It therefore has the exact form
$t(1,1,1,1,1)$, with $t\in\widehat{\mathbb Z}$.
Unlike H10, this does not force $t=0$; the eighteen finite
primitives above admit even that stronger normalization.

\paragraph{Certificate and external inputs.}
\path{agents/five_level_six/hurwitz_spherical_socle.py}
reconstructs all sixty transports, checks all $240$ transitions,
both block orders, the three pair orbits, all fifty sector
memberships, the explicit word's order witnesses, and the
eighteen primitives on all sixty blocks. The JSON and NPZ
have the same stem. D23.C23 supplies the exact identification
and nonabelian simplicity of $S$; D23.F30 supplies the old local
group and its solvable block quotient. D23.P29 supplies the
literal spherical representation and the eighteen cocycles.
No unverified full degree-$6720$ group order or maximal
solvable quotient identification is used.

\subsection{The same candidates extend to the full spherical action}

\begin{statement}{D23.P30}{Proved}{Full four-generator spherical covariance leaves all eighteen pairs}
\end{statement}
For every one of D23.P26's eighteen pairs, there is a
unique permutation $V_5$ of the degree-$6720$ orbit
whose restriction to the original sector is the known
$V$, and which satisfies the four literal covariance
equations
\[
\begin{split}
 V_5\sigma_1V_5^{-1}&=\sigma_1,\\
 V_5\sigma_2V_5^{-1}&=F^{-1}\sigma_2F,\\
 V_5\sigma_3V_5^{-1}&=R\sigma_3R^{-1},\\
 V_5\sigma_4V_5^{-1}&=\sigma_4,
\end{split}
\]
where $F$ is the complete $f(X_1,X_2)$ value of P29
and $R=r_5^{-1}Fr_5$ is the complete $f(X_3,X_4)$
value. Every such $V_5$ lies in the independent product
$S^{60}$ proved in D23.F39, and hence in the derived
pure finite image. This supplies no further exclusion
of the eighteen pairs.

\begin{proof}
The rotation marking $r_5X_ir_5^{-1}=X_{i+3}$ shows
that conjugation by $r_5^{-1}$ sends $(X_1,X_2)$ to
$(X_3,X_4)$. Thus the displayed expression for $R$
has the exact intended evaluation marking.
The usual Drinfeld formula gives the first three
displayed images, as in P18. For the last generator
its two conjugator arguments are
\[
 y_4=\sigma_3\sigma_2\sigma_1^2\sigma_2\sigma_3,
                    \qquad \sigma_4^2.
\]
The spherical relation gives $y_4=\sigma_4^{-2}$,
also checked literally in the certificate. A commutator
word evaluated at these commuting elements is one,
so the fourth image is $\sigma_4$ at multiplier one.

Let $g_1,\ldots,g_4$ be the old generators and
$g'_1,\ldots,g'_4$ the four displayed new images.
Start with the prescribed values of $V_5$ at all
$1344$ original-sector points. Propagate the equations
\[
                         V_5(g_i u)=g'_iV_5(u)
\]
along the four labeled generator edges. The original
action is transitive on the whole orbit: the first
three generators are transitive on the original
sector, and $r_5$ cycles all five sectors. Hence a
consistent propagation determines every value and
is unique. Reject a conflict or a noninjective
assignment immediately.

For all eighteen pairs the propagation is consistent,
covers the entire $6720$-point set, and yields a
permutation. The certificate then verifies all four
covariance equations pointwise. This proves the
assertion by an explicit finite construction. In
particular these four-generator images generate a
permutation group conjugate to the original image.
No finite-kernel stability was assumed to infer this
conclusion; the intertwining permutations were
constructed and checked directly.

Finally put $N=S^{60}$. By F39 and P29 the complete
$F$ has a primitive in $N$, so $F\in N$ because
$r_5$ normalizes $N$. The complete $R$ is a conjugate
of $F$, and also belongs to $N$. Hence the four old
and new generator images have the same sixty-block
permutation and the same local $L/S$ transitions in
F39's coordinates. The graph construction preserves
each block: this is true on the prescribed original
sector, and each propagation step preserves this
property because $g_i$ and $g'_i$ have the same block
permutation. Transitivity therefore proves that
$V_5$ fixes all sixty blocks, without assuming
beforehand that an arbitrary normalizer preserves
this block system.
Its old-sector components already lie in $S$ by F32.
The covariance equation propagates membership in $S$
from one block to each adjacent block, since the old
and new local transitions have the same class modulo
$S$. Transitivity reaches every block, proving
$V_5\in S^{60}=N\subseteq P'$. Thus the finite
automorphisms here are implemented even by derived-pure
inner conjugation.
\end{proof}

\paragraph{Scope and replay.}
The complete certificate is
\path{agents/pentagon/hurwitz_full_spherical_covariance.py},
with its JSON and NPZ of the eighteen full permutations.
This result concerns the exact prescribed finite
values and does not produce a profinite associator.
The pure order-five primitives are addressed separately
by P29 and F39; the actual-source common-word membership
is the separately proved H9/F38 input.

\subsection{A literal two-generator restriction of the tuple-kernel action}

\begin{statement}{D23.F40}{Proved}{The last-coordinate action is an evaluation of the same associator}
\end{statement}
Use the rightmost-strand free kernel $N=\widehat F_n$ in
$\widehat B_{n+1}$ and the peripheral basis $u_1,\ldots,u_n$
of P18. Put
\[
 D=u_1\cdots u_n=z_n^{-1}z_{n+1},\qquad u_n=\sigma_n^2,
 \qquad z_m=(\sigma_1\cdots\sigma_{m-1})^m.
\]
For an actual full-GT element $h=(\lambda,f)$, the restriction
$\kappa_h=\Phi_h|N$ satisfies
\[
 \kappa_h(D)=D^\lambda,
 \qquad
 \kappa_h(u_n)=f(D^{-1},u_n)^{-1}u_n^\lambda f(D^{-1},u_n).
                                                        \tag{D23.F40.1}
\]
Thus, for $n\geq2$, its restriction to the literal free
subgroup $\langle D^{-1},u_n\rangle$ is exactly the usual
two-generator GT action with the same $f$.

For any finite target group $G$, the map on representation
classes
\[
 \pi:\operatorname{Hom}(F_n,G)/G\longrightarrow
             \operatorname{Hom}(F_2,G)/G,
 \qquad
 [g_1,\ldots,g_n]\longmapsto
                  [(g_1\cdots g_n)^{-1},g_n]
                                                        \tag{D23.F40.2}
\]
therefore intertwines the actual tuple action with the
standard GT action on two-generator representations.
At multiplier one, on raw tuples of fixed product $p$,
the action of $h^{-1}$ has last coordinate
\[
           g_n\longmapsto
             f(p^{-1},g_n)^{-1}g_nf(p^{-1},g_n).
                                                        \tag{D23.F40.3}
\]
In particular the actual raw permutation preserves the
partition by the last entry. This condition follows from an
explicit invariant subgroup, without assuming compatibility
with an arbitrary operation of deleting branch points.

\begin{proof}
The full-profinite braid-action formula, with the marking of
Harbater--Schneps, Section 0.1, equation (1), is
\[
 \Phi_h(\sigma_n)=f(y_n,\sigma_n^2)^{-1}
                        \sigma_n^\lambda f(y_n,\sigma_n^2),
 \qquad y_n=z_{n-1}^{-1}z_n.
\]
For $n=1$ the assertion reduces to the single peripheral
power map; assume $n\geq2$. The full twists transform by
$\Phi_h(z_m)=z_m^\lambda$, as in the standard braid action
used in P18 and H6. Since $z_n$ commutes with $z_{n+1}$,
this gives $\kappa_h(D)=D^\lambda$.

The central element $z_{n+1}$ and the older twist
$z_{n-1}$ both commute with $D$ and $u_n$. For $D$ this
also follows directly from its expression
$z_n^{-1}z_{n+1}$; for $u_n=\sigma_n^2$, the assertion
for $z_{n-1}$ follows from the distant Artin relations.
Consequently
\[
          y_n=cD^{-1},\qquad
          c=z_{n-1}^{-1}z_{n+1},\qquad [c,D]=[c,u_n]=1.
\]
For every profinite commutator word $f$, multiplying its first
input by a central element of the two-input subgroup has no
effect: the possible extra power is the first exponent sum
of $f$, which is zero. This statement follows first for
ordinary words and then by continuity. Hence
\[
                       f(y_n,u_n)=f(D^{-1},u_n).
\]
Squaring the displayed half-twist formula gives the second
identity in (D23.F40.1). In particular its first argument is
$D^{-1}$ and its second is $u_n$; interchanging them changes
the conjugator to its inverse by the two-cycle identity and
must not be done silently.

The elements $D,u_2,\ldots,u_n$ form a free basis, obtained
from the peripheral basis by replacing $u_1$ with
$u_1\cdots u_n$. Thus $D^{-1},u_n$ generate a free factor
when $n\geq2$. Let $\iota:F_2\hookrightarrow F_n$ send
$x$ to $D^{-1}$ and $y$ to $u_n$. Equation (D23.F40.1)
says exactly
\[
                   \kappa_h\iota=\iota\varphi_h,
 \qquad
 \varphi_h(x)=x^\lambda,\quad
 \varphi_h(y)=f^{-1}y^\lambda f.
\]
P18 defines the actual left tuple action by
$V_h[\psi]=[\psi\kappa_h^{-1}]$. Restriction through
$\iota$ therefore yields
$\pi V_h=\mathcal V_h\pi$, where
$\mathcal V_h[\eta]=[\eta\varphi_h^{-1}]$ is the ordinary
left GT action on two-generator representation classes.
This proves (D23.F40.2).

At multiplier one, $D$ is fixed exactly. Evaluating
$\kappa_h(u_n)$ on a raw representation with product $p$
gives (D23.F40.3). This is the action $V_h^{-1}$, because
that permutation precomposes with $\kappa_h$, not its
inverse. Its last coordinate depends only on $p$ and $g_n$.
The same argument for $h^{-1}$ proves preservation in the
opposite direction and hence preservation of the full
partition by either actual permutation.
\end{proof}

\paragraph{Finite scope and verification.}
For the eighteen P26 survivors, each of the six raw lifts
in P27 preserves the last-entry partition. This passes an
explicit necessary condition, but does not imply that the
induced last-coordinate permutations can all be produced by
evaluations of one common commutator word $f$ in $E_0$.
That stronger requirement uses the complete marked-kernel
correlations of the pairs $(p^{-1},g_n)$.
The exact braid identities and the independent complete
$18\cdot6$ partition replay are recorded in
\path{agents/five_level_six/hurwitz_last_coordinate.py}
and its JSON/NPZ. The primary braid formula is
\url{https://webusers.imj-prg.fr/~leila.schneps/HS.pdf},
printed p.~3, Section 0.1, equation (1); P18 fixes the
rightmost meridians and H6 fixes the raw product convention.

\subsection{Central target extensions do not enlarge a cyclic-centralizer class action}

\begin{statement}{D23.P31}{Proved}{A limitation of central extensions of the cubic target}
\end{statement}
Let $G$ be a finite group, $g\in G$, and assume
$C_G(g)=\langle g\rangle$. For any central extension
\[
              1\longrightarrow Z\longrightarrow
              \widetilde G\xrightarrow{q}G\longrightarrow1
\]
and any lift $\widetilde g$ of $g$, its conjugacy class
$\widetilde C$ maps bijectively to the conjugacy class
$C$ of $g$. The bijection is equivariant for conjugation
and for every Hurwitz braid action on tuple powers.
Consequently, for every $n$, the braid actions on
$\widetilde C^n$ and $C^n$ are isomorphic, and so are
the actions after quotienting by simultaneous target
conjugation. Every individual braid orbit is preserved
with its exact marked permutation action.

In particular, replacing the cubic target
$E_0=\mathbb F_3^2\rtimes SL_2(\mathbb F_3)$ by any
central extension, while using one selected lift of
the original order-six conjugacy class, cannot enlarge
any of the Hurwitz permutation images considered here.
This includes a central extension of order $648$.

\begin{proof}
The inclusion
$C_{\widetilde G}(\widetilde g)\subseteq
q^{-1}(C_G(g))$ is immediate. Conversely, if
$q(h)=g^k$, then $h=\widetilde g^kz$ for some
$z\in Z$. The centrality of $Z$ implies that $h$
commutes with $\widetilde g$. Therefore
\[
 C_{\widetilde G}(\widetilde g)
                  =q^{-1}(\langle g\rangle).
\]
The induced map of coset spaces
\[
 \widetilde G/C_{\widetilde G}(\widetilde g)
                    \longrightarrow G/C_G(g)
\]
is a bijection, and these coset spaces parameterize
the respective conjugacy classes. This proves the
equivariant class bijection. Notice that no assumption
on the order of the chosen lift is needed.

Apply this bijection coordinatewise. The braid action
uses only $(u,v)\mapsto(uvu^{-1},u)$, which commutes
with the homomorphism $q$. Thus the product bijection
is braid equivariant. Central elements act trivially
by conjugation on $\widetilde C$, so its target action
factors through $G$ and agrees exactly with the action
on $C$. Passing to simultaneous conjugacy orbits
therefore remains a bijection. The same reasoning
shows that an individual braid orbit, including one
chosen by a fixed-product condition, has the same
marked permutation action as its image.

For the original cubic target, $|E_0|=216$, the chosen
element has order six, and its conjugacy class has
size $36$ by D23.P17's complete construction. Hence
its centralizer has order $216/36=6$ and contains
the cyclic subgroup of order six generated by it.
They are equal, so the general assertion applies.
\end{proof}

\paragraph{Exact limitation.}
This statement is about the specified selected
conjugacy-class Hurwitz construction. It does not say
that all evaluations in a central extension carry no
new information, nor does it cover noncentral target
extensions or a union of different lifted classes.
It rules out obtaining a new braid permutation image
merely by centrally extending this cyclic-centralizer
target and choosing one lifted class.

\subsection{A same-associator last-meridian test leaves one existential form}

\begin{statement}{D23.C34}{Proved}{The mixed existence gate reduces to the single pair $(1,0)$}
\end{statement}
Let $h=(1,f)$ be an actual full-GT element with
$(\rho\bmod3,a_3,a_2)=(0,1,0)$.
In the eighteen-pair numbering of P26, its degree-1344
Hurwitz data must have one of the nine labels
\[
                   (i,j),\qquad i\in\{1,53,61\},
                   \quad j\in\{0,1,2\}.                  \tag{D23.C34.1}
\]
Among P27's six raw lifts of any such pair, only the
recorded phase-zero lift remains.  For the fixed-product
raw tuples, all the values $f(p^{-1},g_4)$ are the identity.
Consequently, if a normalized mixed actual element exists,
one exists with the single degree-1344 label
\[
                                  (1,0).                 \tag{D23.C34.2}
\]
The last assertion uses cubing as in C30.  Neither the
nine surviving residues nor the single existential
normal form is asserted to lift to an actual element.

\begin{proof}
\textbf{The actual-source formula and its orientation.}\quad
Use F40's rightmost-strand free kernel in $B_5$, with
peripheral basis $u_1,\ldots,u_4$ and
$D=u_1u_2u_3u_4=z_4^{-1}z_5$.
Its full-profinite calculation, from
Harbater--Schneps, Section 0.1, equation (1), gives
\[
 \kappa_h(D)=D,\qquad
 \kappa_h(u_4)=f(D^{-1},u_4)^{-1}u_4f(D^{-1},u_4).
                                                               \tag{D23.C34.3}
\]
In particular this is an evaluation of the same
associator that occurs in the braid covariance.
The raw left action is precomposition by
$\kappa_h^{-1}$, so it is the inverse of the stored
raw permutation $V$ whose last-coordinate action
is the right side of (D23.C34.3).

Let $E_0=\mathbb F_3^2\rtimes SL_2(\mathbb F_3)$.
H6's raw set consists of 8064 literal four-tuples
$(g_1,g_2,g_3,g_4)$ in $E_0$, all with the same
product $p$.  Its last coordinate ranges over the
36-element class $C_+$.  For each of the eighteen
P26 pairs and each of P27's six deck phases, the
certificate reads
\[
   b\longmapsto b'=\operatorname{last}(V^{-1}T),
                 \qquad b=\operatorname{last}(T).
\]
It verifies on all 8064 tuples that this map is
well defined.  Equation (D23.C34.3) requires
\[
                b'=f(a,b)^{-1}bf(a,b),\qquad a=p^{-1}.
                                                               \tag{D23.C34.4}
\]

\textbf{Exact marked kernels and complete fibers.}\quad
For a finite marked pair $(a,b)$, order the subgroup
$\langle a,b\rangle$ by deterministic breadth-first
enumeration starting at the identity and using the
two ordered positive generators.  Record the two
right-multiplication permutations in that ordering.
Two such records are identical if and only if the
two maps from the abstract free group have the
same kernel.  Indeed a marked isomorphism preserves
the enumeration; conversely equal records give a
rooted labeled Cayley-graph isomorphism, hence a
group isomorphism preserving the two generators.
This criterion also supplies the literal isomorphism
needed to transport every possible word value.

The finite inventory checks every one of the
$216^2=46656$ pairs in $E_0^2$.
Conjugation by the explicitly enumerated affine
group of order 432 is used only to accelerate
this enumeration; equality of ambient automorphism
orbits is not substituted for the marked-kernel
criterion.  There are 177 such ambient pair orbits
and 146 distinct marked kernels.  Of these, 64
are onto $E_0$.  Exactly eight kernels occur among
the pairs $(p^{-1},b)$ needed here.  In the fixed
record numbering they are
\[
\begin{array}{c|rrrrrrrr}
\text{kernel}&37&42&49&58&61&66&73&78\\ \hline
|\langle a,b\rangle|&24&6&216&54&216&216&216&18 .
\end{array}                                                \tag{D23.C34.5}
\]
The numbers of original degree-1344 tuples in
these eight classes are respectively
$80,16,240,324,240,168,168,108$.

For each kernel choose its recorded canonical pair
$(a_k,b_k)$ and write $H_k=\langle a_k,b_k\rangle$.
An actual commutator word has value in $H_k'$.
The certificate constructs $H_k'$ as the normal
closure of $[a_k,b_k]$ and exhausts every element.
Using the marked isomorphisms above, each such
element is transported to every occurring pair
$(p^{-1},b)$ with that kernel and tested against
(D23.C34.4).  Thus all occurrences of one kernel
are constrained simultaneously, not independently.
For every one of the $18\cdot6=108$ raw actions,
the complete eight fiber sizes, in the order
(D23.C34.5), are
\[
                         2,1,2,1,2,2,2,1.
\]
There are therefore exactly 32 assignments before
the gentle identities.  All $108\cdot32=3456$
assignments are tested below.  No prior common-word
test on different kernels is needed for this count
or for the ensuing exclusion.

\textbf{All six peripheral orientations.}\quad
On marked kernels the two transformations
$(a,b)\mapsto(b,a)$ and
$(a,b)\mapsto(b,(ab)^{-1})$ generate the peripheral
six-orientation action.  The inventory checks their
exact kernel permutations, including the reflected
identifications.  The eight kernels in (D23.C34.5)
belong to the following seven orbits:
\[
\begin{gathered}
\{1,33,37\},\quad\{39,42,105\},\quad
\{10,25,45,49,61,64\},\quad\{51,58,143\},\\
\{20,66,69\},\quad\{30,68,73\},\quad\{75,78,145\}.
\end{gathered}                                               \tag{D23.C34.6}
\]
Here the integers label marked kernels, not elements
of $E_0$.

For each orbit choose a prescribed base value
$g_0=f(a,b)$ and put $z=(ab)^{-1}$.
Exhaust $g_1=f(b,z)$ over the complete derived
subgroup of $\langle a,b\rangle$.
The positive three-cycle forces
$g_2=f(z,a)=(g_1g_0)^{-1}$, and duality forces the
values on the reverse pairs.  Thus the six values on
\[
 (a,b),(b,z),(z,a),(b,a),(z,b),(a,z)
\]
must be
\[
               g_0,g_1,g_2,g_0^{-1},g_1^{-1},g_2^{-1}.
                                                               \tag{D23.C34.7}
\]
Transport these values back to canonical marked
representatives.  Discard them if two occurrences
of one kernel disagree or if a prescribed fiber
value is not respected.

For every remaining assignment, evaluate the
two-cycle and both three-cycles at the canonical
pair of \emph{each} kernel in (D23.C34.6):
\[
\begin{split}
 f(y,x)f(x,y)&=1,\\
 f(z,x)f(y,z)f(x,y)&=1,\\
 f(z,x)z^{-1}f(y,z)y^{-1}f(x,y)x^{-1}&=1,
                \qquad z=(xy)^{-1}.
\end{split}                                               \tag{D23.C34.8}
\]
This last step includes all reflected kernels and
their inner conjugation corrections through the
literal marked isomorphisms.  It does not presume
that one orientation's equation automatically
propagates to the others.
The negative three-cycle is necessary for an
actual multiplier-one element: the actual product
$(1,f)c=(-1,f)$, where $c=(-1,1)$, has the same
word $f$ and has $m=(-1-1)/2=-1$.

This enumeration is complete for the displayed
necessary identities.  Any actual $f$ supplies
one of the exhausted $g_1$ values, all of its other
values are forced by (D23.C34.7), and its same-kernel
values and all equations (D23.C34.8) must agree.
Conversely the certificate claims only a finite
extension satisfying these equations, not an
actual gentle lift.

\textbf{Outcome and normalization.}\quad
Exactly nine of the 3456 input assignments admit
such extensions.  Their raw labels are
\[
                 (i,j,\text{phase})=(i,j,0),
           \quad i\in\{1,53,61\},\quad j\in\{0,1,2\}.
\]
In each, all eight canonical values are the
identity (element index 6 in the fixed $E_0$ table).
The numbers of extensions on the seven separate
orbits are $4,1,4,3,4,4,3$; correlations between
different orbits may reduce them further and are
not asserted here.
Thus every actual element satisfies
(D23.C34.1), and (D23.C34.4) is the identity
last-coordinate action.
By the complete cube table in C30, every one
of the nine pairs (D23.C34.1) cubes to $(1,0)$.
Cubing preserves $\rho=0$, $a_3=1$ and $a_2=0$.
This proves the existential assertion (D23.C34.2).
\end{proof}

\textbf{Reproducible certificate and status.}\quad
The three scripts are
\path{agents/characters/code/affine216_marked_pair_kernels.py},
\path{agents/characters/code/affine216_last_meridian_fibers.py},
and
\path{agents/characters/code/affine216_last_gentle_extensions.py}.
Run the last with \texttt{--all-fibers} to replay all
3456 inputs independently of the separate nine-coordinate
common-word test.  Its complete JSON/NPZ records have
stem
\path{agents/characters/records/affine216_last_gentle_extensions_all}.
The earlier unqualified stem retains the additional
common-word-filtered replay; it yields the same nine
cases but is not used by this proof.
F40 supplies the actual-source formula, P26--P27 the
complete initial raw list, and C30 the final source
normalization.  The author and the root reviewer have
completed whole-proof and code audits, and the root's
fresh complete 3456-assignment replay passed.
The exact completed scopes and dated review events
are recorded in the check ledger.


\part{Proof drafts and active calculations awaiting their final check}
In this part a proof draft is not yet counted as a completed
finding.  Exact computation records are retained even when the
corresponding geometric or source comparison is still open.
\section{Current five-point calculation}
The full $S_3$ Nielsen closure and all the stated generation and
reflection tests leave all $216$ Krammer assignments. D23.P14
also gives a common ordinary word defeating the combined older
Hurwitz and Krammer equations. No such finite word is asserted
to be an actual gentle or full-GT element.

The two orbital maps now give exact identities in the literal
four-strand braid group, with their full-twist corrections retained.
D23.C17--C19 exhibit residual parities $e_3$ and $d_2$ and their
specified common periodic conjugacy. The known affine covariance
and the individual cube conditions do not compare those parities.
The exceptional Krammer pairs each have derived image $SU_6(3)$,
but their twelve marked derived kernels are independent.
The remaining issue is the correlated noncentral orbital equations,
not a discarded scalar or an unrecorded base path.

D23.P17--P18 construct a degree-$1344$ Hurwitz quotient retaining
both exact marked detectors and prove its preservation by multiplier-one elements of actual
full GT. D23.H1 determines a unique first-pair value for each of
the six cubic labels in its degree-$36$ quotient; all six survive.
The twelve blocks now have an exact six-dimensional orthogonal
geometry. The canonical relative cup form is preserved by actual
multiplier-one GT, and the normalized hypothetical mixed state
has its conjugating permutation in the product of twelve simple
orthogonal groups. D23.H3--H4 and P21 fix its action uniquely on
four blocks using intrinsic finite quotient partitions. The
complete finite list initially had thirty-six necessary conjugators.
D23.F36 fixes the full first value uniquely. D23.H8 then imposes
the complete third-pair derived membership, and D23.P26 imposes
the literal raw pentagon. These are actual necessary restrictions
and leave six conjugators with three mixed values each.
D23.C30--C31 use powers and actual complex conjugation to reduce
mixed existence to three pairs. D23.C33 then uses an odd product
with a complex conjugate to leave two, $(1,0)$ and $(27,2)$,
without asserting their existence in full GT.

D23.H9/F38 proves that every one of the eighteen raw survivors
has a single ordinary commutator word realizing all five values
and every mixed orientation. D23.P27 also gives full raw braid
covariance. D23.H10 and P28 distinguish the necessary derived-pure
cubic frame from an arbitrary finite conjugator and show that
each of the two final existential normal forms passes that test.
These finite statements do not impose the free gentle identities
or construct an actual full-GT lift.

D23.P29 extends the action to a degree-$6720$ spherical orbit.
D23.F39 proves sixty independent simple factors and trivial
finite order-five cohomology; all eighteen pure primitives even
admit commutator lifts. D23.P30 gives full spherical covariance.
This particular extension therefore supplies no further exclusion.

D23.F40 links the actual free-kernel action on target-group
tuples to evaluations of the same associator in $E_0$. D23.C34
exhausts its eight marked last-meridian kernels and all finite
gentle cycle extensions. Only nine raw pairs survive, and all
cube to the single existential form $(1,0)$. Full common-word
correlations across the additional target kernels are active work.
Additional colored Hurwitz covers and the noncentral orbital
comparisons remain available. Those working calculations are not
promoted to a theorem before their proof and checks. No arbitrary
coordinate-deletion compatibility is assumed.

\section{Literature and historical locators}
\small
\begin{enumerate}
\item J. Voight, \emph{Quaternion algebras}, stable version
1.0.5 (2024): (12.4.10), 12.4.12, (12.4.14), Proposition
14.2.1, Main Theorem 14.6.1, Corollary 14.6.2.
\url{https://jvoight.github.io/quat-book-v1.0.5.pdf}.
\item T. Kubota, \emph{Topological covering of SL(2) over a
local field}, J. Math. Soc. Japan 19 (1967), 114--121:
equation (2), p.~114; product formula (4), p.~115.
\url{https://doi.org/10.2969/jmsj/01910114}.
\item H. Hida, \emph{Adjoint L-value as a period integral and
the mass formula of Siegel--Shimura}, 20 September 2023:
Section 1.3, p.~7, the $s_p$ and $\kappa_p$ displays before
(1.5).  The splitting is attributed there to Gelbart,
\emph{Weil's Representation and the Spectrum of the
Metaplectic Group}, LNM 530 (1976), Proposition 2.8.
\url{https://www.math.ucla.edu/~hida/PdFinal.pdf}.
\item W. Herfort and L. Ribes, \emph{Torsion elements and
centralizers in free products of profinite groups}, J. reine
angew. Math. 358 (1985), 155--161: Theorem B and Section 4,
Theorem B', p.~159. \url{https://doi.org/10.1515/crll.1985.358.155}.
\item B. Collas, \emph{Action of a Grothendieck--Teichm\"uller
group on torsion elements of full profinite mapping class groups
in genus zero}, J. Th\'eor. Nombres Bordeaux 24 (2012),
605--622: Theorem B, p.~608.
\url{https://www.numdam.org/article/JTNB_2012__24_3_605_0.pdf}.
\item I. Marin, \emph{Characters of the Grothendieck--Teichm\"uller
group through rigidity of the Burau representation}, Ann. Fac.
Sci. Toulouse 17 (2008), 765--780: Section 4, Lemmas 4.1--4.3;
Section 5 for the general rigidity method.  D23.T1 proves the
precise integral specialization needed here.
\url{https://www.numdam.org/article/AFST_2008_6_17_4_765_0.pdf}.
\item S. Bigelow, \emph{The Lawrence--Krammer representation},
arXiv:math/0204057v1 (2002), Section 4.2, printed p.~13:
the corrected six-case polynomial Krammer matrix formula.
\url{https://arxiv.org/pdf/math/0204057}.
\item J. S. Milne, \emph{Lectures on Etale Cohomology}, v2.21
(22 March 2013), Theorem 3.4.
\url{https://www.jmilne.org/math/CourseNotes/LEC.pdf}.
\item J. S. Milne, \emph{Modular Functions and Modular Forms},
v1.10, Theorem 12.27, p.~124, and the surrounding ray-class
construction. \url{https://www.jmilne.org/math/CourseNotes/MF110.pdf}.
\item H. Nakamura, \emph{On arithmetic monodromy representations
of Eisenstein type in fundamental groups of once punctured
elliptic curves}, RIMS-1691 (2010), Section 5.4, equation
(5.4.2), p.~28.
\url{https://www.kurims.kyoto-u.ac.jp/preprint/file/RIMS1691.pdf}.
\item R. Hain and M. Matsumoto, \emph{Weighted completion of
Galois groups and Galois actions on the fundamental group of
$\mathbf P^1-\{0,1,\infty\}$}, Compositio Math. 139 (2003),
119--167: Section 8.5, paragraph before Theorem 8.7.
\url{https://arxiv.org/abs/math/0006158}.
\item J.-P. Serre and J. Tate, \emph{Good reduction of abelian
varieties}, Ann. Math. 88 (1968), 492--517: Theorem 1, p.~493.
\url{https://doi.org/10.2307/1970722}.
\item A. Grothendieck, Montpellier archive cote 145,
\emph{Printemps 81 (vieille r\'edaction)}: inspected one-based
PDF pages 7, 11, 12, 22, 28; numbered archive scans
6, 10, 11, 21, 27.  The PDF includes a cover, so the two
numbering conventions differ by one.
\url{https://grothendieck.umontpellier.fr/145.pdf}.
The exact historical interpretation and its limits are in
\path{IMPLICATIONS.md}.
\item B. Enriquez, \emph{Quasi-reflection algebras and cyclotomic
associators}, arXiv:math/0408035v2: Section~8, Theorem~8.1,
Lemmas~8.2--8.5 and equations~(50)--(51), printed pp.~30--32;
Section~13, equation~(63), printed pp.~48--49.
\url{https://arxiv.org/pdf/math/0408035}.
\item P. Lochak, H. Nakamura and L. Schneps,
\emph{Eigenloci of 5 point configurations}: Theorem~2,
equation~(HF1), author PDF p.~4; Section~6,
equations~(6.3)--(6.4), author PDF pp.~30--31.
The theorem's Galois scope is retained in D23.F22.
\item P. Lochak and L. Schneps, \emph{The Grothendieck--Teichm\"uller
group and automorphisms of braid groups}, LMS Lecture Notes 200
(1994), Section~5, Proposition~10, author PDF p.~23, equation~(5).
\url{https://webusers.imj-prg.fr/~leila.schneps/Fls.pdf}.
\item D. Harbater and L. Schneps, \emph{Fundamental groups of
moduli and the Grothendieck--Teichm\"uller group}: Section~0.1,
equation~(1), author PDF p.~3; Section~1.1, p.~9;
Section~3, Lemma~1, p.~30.
\url{https://webusers.imj-prg.fr/~leila.schneps/HS.pdf}.
\item P. Lochak and L. Schneps, \emph{A cohomological interpretation
of the Grothendieck--Teichm\"uller group}, Invent. Math. 127
(1997), 571--600: equation~(2), p.~572; Theorems~1--2,
p.~573; Proposition~4, p.~574.
\url{https://webusers.imj-prg.fr/~pierre.lochak/textes/inventiones.pdf}.
\item L. Alp\"oge, \emph{Remark}, supplied four-page PDF:
the prompt's steps 3 and 5 and the Gauss-sum discussion on
p.~3, footnote 3.  Accepted as stipulated; no theorem from
its Lean archive is used here, and those Lean files were not
in the supplied ZIP.  Thus no formal-file locator is invented.
\end{enumerate}
\end{document}
